% Current integrated paper source: Paper4A-FiniteQuantumMechanics-v4_215.tex (v3_017).
% Prior source: Paper4A-FiniteQuantumMechanics-v4_214.tex. New cross-DAG mathematics is scoped, not physical-occurrence certified.
% Paper-series v2.381: complete algebraic exceptional torsion atlas and all-cocycle twist-stability; original claims retained.
% Paper-series v2.380: unified bibliography, exact exceptional double/length/gauge mathematics in I-B; previous theorem labels retained.
% Paper-series v2.379: full equivariant Dirac/Hodge/eta/torsion atlas in Paper I-B; paper-root references synchronized.
% Paper-series v2.378: 2026 exceptional literature consolidation and reference synchronization. The mathematical additions are located in Paper I-B.
% Paper-series v2.377: reference-only rebind to current Paper I-B v1.1; original mathematics retained except new Paper I-B comparison chapter.
% Framework v2.370: derived exceptional-root/Hurwitz comparison in I-A/I-C; other papers reference-only synchronized.
% Framework v2.368: preserve prior II-B convention/unital fixes and verified II-D corrections; current references synchronized.
% Framework v2.358: current 30-root source reference synchronization; selectively updated III-E, III-F, VI-B, VI-C.
% Framework v2.356: bibliography/source-registry synchronization; mathematical body unchanged.
% Framework v2.355: cross-paper native finite-history/ancilla descent and process firewall.
% Framework v2.352: integrated mathematical/operational closure; physical occurrence boundaries preserved.
% Framework v2.347: exact operational test-resolution tiers, one-body response iff criterion, terminating Majorana hierarchy, odd-transfer witness, cutoff cumulants, basis-free regular-CCR covariance. Prior domains retained.
% Framework v2.341: Coherent coset--Hodge and parity cover integration. Previous v2.340 retained.
% Framework v2.340: Exact rank-two 600-cell incidence and typed operator/observable closure.
% Framework v2.320: signed Petersen quantum sectors, mirror-transfer CP no-go, conditional interference and source-visible response; exact finite certificates. Preserves inherited original-source theorem bodies.
% Framework v2.275: exact positive dense-twin word cubature, finite repeatability and eighth-moment classification, nested sampling frames, joint action and POVM-jet exports; preserves v2.274.
% Framework v2.274: two-scalar positive convolution certificates, first-leaked exceptional coset detection, branch-symmetric twin mixing, and finite endpoint moments; all v2.273 branches preserved.
% Framework v2.273: restored missing Problem 3 theorems from v2.268; exact exceptional frames, twin geometry, convolution repeatability and projector/occurrence firewalls; merged on v2.272 without overwriting independent branches.
% Framework v2.272: merged v2.270 covariant native-readout mathematics with v2.271 Problem 4 finite-action Hopf hypercharge/holonomy corrections; no universal bulk-gauge uniqueness claim.
% Independent branches: v2.270 readout/guard exports and v2.270 K16 action-consumed charge closure; all prior scoped assumptions retained.
% Framework v2.267: exact Paper II-A converse/observable certificates, III-F action tests, V-B right orbit and V-C native gauge scope; prior v2.266 preserved.
% Framework v2.266: exact Paper II-C outcome-count/readout converse results and scoped downstream imports; 30-root references synchronized.
% Framework v2.265: II-C coordinate-stage constitution and exact first-symbol host inheritance; prior mathematical domains retained.
% Framework v2.264: II-B reverse and action-bridge integration; preserves prior mathematical domains; active 30-root references synchronized.
% Framework v2.263: Paper I-D adds bridge-controlled root-carry parity and source-relative old-cap cancellation; other root mathematical bodies retained; active 30-root references synchronized.
% Framework v2.262: merges Paper II-F three-gap closure from v2.260 into v2.261, preserving newer II-B/II-D repairs; 30-root references synchronized.
% Framework v2.261: merged Paper II-B GP2-level/geometric repairs; preserves v2.260 II-D repairs; active references synchronized.
% Framework v2.260: Paper II-D signed RS, finite-symbol typing, and conditional curvature/Jordan repairs; references synchronized.
% Framework v2.259: conditional native SU(2) atlas and explicit later discharge; v2.258 mathematical bodies retained elsewhere.
% Framework v2.258: preserves parallel II-A/dynamics update and adds complete readout/frequency closures.
% Framework v2.257: exact minimal complete position readout and arbitrary-preparation frequency law.
% Framework v2.256: exact native inertia and sharp bounded ordered-source closure; preserves concurrent v2.255 work.
% Framework v2.255: preserves parallel detector integration and adds exact finite readout closures.
% Framework v2.254: seven-outcome intrinsic readout and sharp stationary event-frequency criterion.
% Framework v2.253: typed raw coordinate defect in I-A and synchronized VI-A import; active references synchronized.
% Framework v2.252: optional shell-frame source-return calculation in III-G; active references synchronized.
% Framework v2.251: forward common-action execution, joint source jets and scoped coupled control.
% Framework v2.251: exact full-heat variations, fermion source jets and explicit physical completion frontier.
% Framework v2.251: certified interacting source jets, pole stability and scoped completion proofs.
% Framework v2.251: interacting finite-source evaluation and coupled local vacuum classification.
% Framework v2.251: exact finite source execution and native causal common-parent bridge.
% Framework v2.251: constructive parent and scoped source-incidence audit; see release report.
% Framework v2.251: common-action two-jet, corrected source/response closure and synchronized references.
% Framework v2.244: scoped rooted source-order Schur diagnostics and synchronized paper references.
% Framework v2.242: signed rooted phase lift, complete transverse return and exact family cubic.
% Framework v2.241: proved rooted conormal geometry, conditional cap realization and scoped action tests.
% Framework v2.239: coherent 30-paper reference rebinding; selective new proofs on established owners.
% Framework v2.238: Paper IV restructured; inherited mathematics retained.
\documentclass[12pt,a4paper]{article}
\usepackage{mathtools}
\usepackage{float}
\usepackage{caption}
\usepackage{tikz}
\usepackage{tikz-cd}
\usetikzlibrary{arrows,positioning,calc,decorations.pathmorphing}
\usepackage{microtype}
\usepackage{xltabular}
% Paper-series v2.381: new I-B torsion/modular results, bibliography and current source roots.
% Paper-series v2.380: new exceptional bibliography and synchronized paper references.
% Paper-series v2.379 current 31-paper macro/aux registry.
% Paper-series v2.378: current 31-paper AUX/bibliography registry updated for Paper I-B v1.2 literature additions.
% Framework v2.371: current-source reference registry refreshed; old theorem labels retained.
% Framework v2.370: current 30-paper source/AUX root rebinding; native mathematical macros unchanged.
% Framework v2.368: synchronized merged 30-root registry; mathematical notation unchanged.
% Framework v2.358: active 30-paper source registry synchronized; existing mathematical macros unchanged.
% Framework v2.357: selective new source-owned certificates; active reference roots updated for five revised papers.
% Framework v2.356: active 30-root registry synchronized to the 30 papers in this release.
% Framework v2.275: exact positive dense-twin word cubature, finite repeatability and eighth-moment classification, nested sampling frames, joint action and POVM-jet exports; preserves v2.274.
% Framework v2.274: filename and root-source synchronization.
% Framework v2.273: restored missing Problem 3 theorems from v2.268; exact exceptional frames, twin geometry, convolution repeatability and projector/occurrence firewalls; merged on v2.272 without overwriting independent branches.
% Framework v2.272: merged v2.270 covariant native-readout mathematics with v2.271 Problem 4 finite-action Hopf hypercharge/holonomy corrections; no universal bulk-gauge uniqueness claim.
% Independent branches: v2.270 readout/guard exports and v2.270 K16 action-consumed charge closure; all prior scoped assumptions retained.
% Framework v2.269: synchronized 30 active AUX roots; mathematics and bibliography preserved from Macros-v2_161.
% Framework v2.266: Paper II-C converse-derived outcome resources; II-E, III-C, IV-A/B, V-A and VI-B downstream mathematical closures; 30-root registry synchronized.
% Framework v2.265: II-C coordinate-stage constitution and exact first-symbol host inheritance; prior mathematical domains retained.
% Framework v2.264: paper II-B reverse/audit and IV-A observable-minimal exchange; synchronized external references to 30 active roots.
% Framework v2.263: Paper I-C adds bridge-controlled root-carry parity and source-relative old-cap cancellation; other root mathematical bodies retained; active 30-root references synchronized.
% Framework v2.262: integrates Paper II-F radius, Casimir and readout results; active 30-root references synchronized.
% Framework v2.261: merged Paper II-B GP2-level/geometric repairs; preserves v2.260 II-D repairs; active references synchronized.
% Framework v2.260: Paper II-D signed RS, finite-symbol typing, and conditional curvature/Jordan repairs; references synchronized.
% Framework v2.259: conditional native SU(2) atlas and explicit later discharge; v2.258 mathematical bodies retained elsewhere.
% Framework v2.258: preserves parallel II-A/dynamics update and adds complete readout/frequency closures.
% Framework v2.257: exact minimal complete position readout and arbitrary-preparation frequency law.
% Framework v2.256: exact native inertia and sharp bounded ordered-source closure; preserves concurrent v2.255 work.
% Framework v2.255: preserves parallel detector integration and adds exact finite readout closures.
% Framework v2.254: Pasterski bibliography and synchronized 30-paper external-reference registry.
% Framework v2.253: typed raw coordinate defect in I-A and synchronized VI-A import; active references synchronized.
% Framework v2.252: synchronized source registry; optional III-G shell-frame calculation.
% Framework v2.251: forward common-action execution, joint source jets and scoped coupled control.
% Framework v2.251: exact full-heat variations, fermion source jets and explicit physical completion frontier.
% Framework v2.251: certified interacting source jets, pole stability and scoped completion proofs.
% Framework v2.251: synchronized active source registry; mathematical notation unchanged.
% Framework v2.251: synchronized active source registry; mathematical notation unchanged.
% Framework v2.251: synchronized active source registry; mathematical notation unchanged.
% Framework v2.251: synchronized source registry; no mathematical notation changes.
% Framework v2.244: synchronized 30-paper AUX registry and rooted phase-source action-order tests.
% Framework v2.242: signed rooted phase lift, complete transverse return and exact family cubic.
% Framework v2.241: proved rooted conormal geometry, conditional cap realization and scoped action tests.
% Framework v2.239: all 30 active source roots synchronized; no mathematical macro changes.
% Framework v2.238: Paper IV A/B/C split and 30-root reference registry.
% Framework v2.237: Paper VII-D technical split and 28-root reference registry.
% Framework v2.236: Paper VI-C technical-volume reference added; all current root bindings advanced.
% Framework v2.235: 26-root AUX and bibliography bindings synchronized; mathematical macros unchanged.
% BEGIN INLINE SOURCE: Paper00-TeXBuildCompatibility-v1_0.tex
% Rendering-only repairs for inherited long appendices and bitmap text fonts.
% Integer counters through 26 retain the original labels; later values use AA,
% AB, ... . This file changes no action coefficient, state or theorem premise.
\RequirePackage{alphalph}
\makeatletter
\let\BHPrimitiveAlph\@Alph
\newalphalph{\BHExtendedAlph}{\BHPrimitiveAlph}{26}
\let\@Alph\BHExtendedAlph
\AtBeginDocument{%
  \@ifpackageloaded{microtype}{\microtypesetup{expansion=false}}{}%
}
\makeatother

% END INLINE SOURCE: Paper00-TeXBuildCompatibility-v1_0.tex

% Framework v2.164: current reference/bibliography profile; mathematical definitions retained.
% Framework v2.162: finite metric, functional Hilbert and scoped Ward integration.
% Framework v2.161: inherits Macros-v2_57.tex exactly except the current external-reference binding.
% Macros v2.57 / framework v2.156: current reference registry only; mathematical definitions unchanged.
% Framework v2.143 / Macros v2.55: current external-reference bindings; mathematical notation unchanged.
% Framework v2.110 / Macros v2.54: current external-reference registry and bibliography; mathematical notation inherited.
% Framework v2.67 / Macros-v2_53: adds Born-action-cell and joint-EH readout symbols; absolute unit realization is metrological; MATCH_{J/scal}^{EH} remains the sole G24 normalization bridge.
% Framework v2.35 / Macros-v2_52: grouped paper numbering, active AUX registry, and Born-v6_20 bibliography binding synchronized; notation unchanged.
% Framework v2.26: active-source AUX and bibliography bindings synchronized; notation unchanged.
% Framework v2.09: reviewed early-reciprocity integration; scoped physical imports.
% Framework v2.08: Born--Euler mathematics and conditional complete-action potential criterion.
% Framework v2.07: active AUX registry includes Paper I-B v1.2; shared bibliography Born-v6_18.
% Macros-v2_46 / framework v2.04: mathematical atlas and downstream realization correction.
% Macros-v2_46 / framework v2.03: scoped structure integration and current dependency bindings.
% Macros-v2_46 / framework v2.01: shared source bindings; prior revision comments are historical.
% v2.42 / framework v1.99: Synchronized RG/AS owner and bibliography bindings.
% v2.41 / framework v1.99: active manuscript AUX bindings and geometric-closure bibliography synchronized.
% Release-v1.80 centralized source synchronization for scalar-circle/comparison-band integration.
% Release v1.79: Hurwitz hierarchy and kinetic-normalized action response integration; physical mass matching remains conditional.
% Release v1.78: approved M8 closure-batch integration; physical and analytic domains retained.
% Current v1.78 registry: all active TeX roots are synchronized below.
% Earlier release comments are historical scope records, not current status.
% Release-v1.77 synchronization: native same-owner Noether closure, generator--flow/history typing, analytic-certification firewalls, scalaron/K16, LF/QCD, quantum-Ward, and RAR boundary-owner transport are registered; shared mathematical notation is unchanged.
% Release-v1.76 synchronization: registry points to the v1.76 active paper roots.
% Release-v1.76 synchronization: registers the local A+Phi+g Noether/constraint closure and dynamics-frontier propagation; shared mathematical notation is unchanged.
% Release-v1.75 synchronization: registry points to the v1.75 active paper roots.
% Release-v1.75 synchronization: integrates the joint-Dirac Einstein bulk real form and the local odd-scalar heat-coefficient action on explicit analytic first-jet/interior domains; v1.74 Yang--Mills closure is retained. Fermionic reality, pure-metric curvature-squared/exact-nonlocal terms, noncanonical boundaries, physical vacuum selection, and actual solution-history tokens remain open.
% Release-v1.74 synchronization: registry points to the v1.74 active paper roots and retains the v1.73 theorem labels.
% Release-v1.74 synchronization: canonical Macros-v2_34.tex; generated local nonlinear Yang--Mills action, constrained connection-field phase space, source-free Gauss momentum map, and Legendre closure are integrated on their proved domains. The remaining non-gauge ACT_cl sectors and actual solution-history selection remain open.
% Release-v1.71 synchronization: closes the former JOINT_x/JOINT_5
% connection-occurrence gap by canonical structure-group extension of the
% nonflat Born connection, faithful associated representation, and the local
% cotangent--Clifford operator constructor.  Actual fluctuation members,
% actions, vacua, solutions, calibrations, and analytic QFT domains remain
% separately typed.  The active-AUX registry is advanced once.
% Release-v1.71 canonical macro and active-AUX registry.
% ======================================================================
% Macros-v2_39.tex -- canonical publication macro file
% Release v1.71 retains the v1.70 null-ray/QCD boundaries and additionally
% removes the duplicate structural connection actuality premise.  It does not
% claim global bundle topology, select a gauge/Higgs member, or promote the
% finite construction to an interacting continuum theory.
% Release v1.68 constructs the finite light-front total-fraction operators and
% parton number/momentum operator measures, proves their simplex and Hausdorff-limit
% theorems, and separates positive cut measures from scheme-renormalized
% light-ray distributions.  It also separates finite real Schur poles, stable
% norm-resolvent/Riesz pole-residue limits, and continued resonance
% pole/residue limits.  The Paper-VII-C scattering package is claim-indexed:
% local QFT and spectral claims import no S-matrix or infrared premise, while
% asymptotic amplitudes, LSZ, S-matrix poles, completeness, and cross sections
% import only their named massive or massless/long-range clauses.
% Release v1.67 factorizes BND48_H into a same-action Hessian/Schur
% descendant and still-open typed boundary-trace, critical-lift, and
% refinement-limit conditions.  It proves the finite-collar Dirichlet-to-
% Neumann law sqrt(y)tanh(L sqrt(y)), its nonuniform critical limit, and the
% finite-return ultra-infrared turnover.  The Paper-I-A Hardy constant mode is
% an exact geometric near-hit, not yet the physical metric boundary mode.
% The current local f(R) Hessian remains massive and therefore still refutes
% a RAR prediction on the minimal action branch.
% Release v1.66 derives the transverse depth-six normalizer bridge from the
% already generated C6 Reynolds image and Paper-I-A binary-dihedral normalizer.
% The bridge is the unique least 2D3-stable scalar-containing realization of
% the faithful quotient 2D3/C6=C2, has trace defect six and hence K=48, but is
% transverse to the nested active carrier.  Exact K=44 and K=52 contraction,
% semigroup, and energy-midpoint identities are registered as structural
% readouts rather than projector events.  The active E48 carrier gap remains
% unchanged.  A same-action boundary-symbol/Calderon intertwiner is isolated
% as the necessary and still unproved owner bridge to the static DtN sort;
% representation equivalence alone neither derives RAR nor creates a K=48
% action junction, pole, occupation, or cosmological transition.
% Release v1.65 derives the native E_48-centered carrier mass gap from the
% GP1--GP3-generated active spectrum, distinguishes its carrier-free interval
% from a particle located at E_48, and withdraws the unsupported promotion to
% a positive continuum Yang--Mills mass gap.  The latter remains a valid
% conditional transfer theorem on an enlarged continuum branch, but it is not
% an active framework prediction or an unfinished native-gap proof.
% Release v1.64 synchronizes Derivation Batches 14--22.  It registers the
% GP1--GP3 heat-profile orbit, same-Hessian electroweak response and pole
% ownership, K16/K48 action-response corrections, perturbation state and
% non-Gaussianity typing, finite-vacuum/gap/QCD limit repairs, factorized
% position-area ownership, the actual static metric response, and the ordered
% Wilson/string-tension limits.  Strict continuum convergence, positive
% physical gap, surviving higher cumulants, and positive tension remain open.
% Release v1.61 closes TIMEWIDTH16 as an independent direct-production datum.
% JUNC_{2|7} continues to own the finite intrinsic microscopic profile/width,
% while the actual same-action background/production history determines its
% physical-time pullback, support, and duration.  The current Paper-I-A--VII-C AUX
% registry is advanced once; shared mathematical notation is unchanged.
% Release v1.60 derives the final-octet primitive-remnant depletion, canonical
% Hopf-current representative, exact Hopf descent to dOmega/(4pi), and
% closed-FLRW a^{-3} neutral-current shape.  It retypes K=48 as the onset
% of the final one-unit geometric transfer and records cDM/RAR as distinct
% physical projections of that common neutral parent while preserving
% occurrence, Hilbert-owner, stress-matching, and static-kernel premises.
% Release v1.59 transports the v1.58 exceptional-shell/K16 closure through
% Papers III-F and VI-B--VII-C.  It replaces Paper VI-B's obsolete mixed-Hessian
% output typing by the third-variation pair vertex and Bogoliubov occupation
% support, derives the post-seed Higgs-pair feedback vertex and its exact
% alpha--delta kernel factors, propagates the unique K16/K48 action-junction
% split, identifies the QCD delta count with the upstream 2O shell atom only
% as a common-scalar numerical compatibility, and narrows THERM16 to dynamical
% approach/beta selection because the finite Gibbs/KMS target is already
% unique on a supplied simple observable algebra once its Hamiltonian and
% inverse temperature are supplied; direct-sum central weights remain typed.
% Release v1.58 registers the exceptional nearest-shell closure, the derived
% transverse K=16 action event, the unique interior action junction, the
% type-correct ordered pair vertex, the thin-layer squeeze law, and the
% radiation-like/matter-like kinematic split.  REG_0, transfer magnitude,
% physical layer width, and thermalization remain explicitly conditional.
% Release v1.57 registers the active-depth-six obstruction, the native
% K=48 carrier desert, the typed 7|9 owner reduction, and the dynamical
% reheating handoff across Papers III-B--VI-B.
% Release v1.56 registers the least-support n_2 inflation dressing of the
% persistent n_7 Laurent response, the analytic global-convexity and positive
% endpoint-measure proofs, the derived inflation core and K=16 coefficient
% freeze, and the no-particle dark-sector split into static RAR plus the
% conditional homogeneous effective-curvature response at reflected K=48.
% There is no n_6 action, scalaron-DM mass, or second Einstein--Hilbert term.
% Release v1.55 registers the protected-octet resolved/residual reflection
% q_sp+Z_W=8, its j<->8-j and K<->64-K scale covariance, and the derived
% rank-one-exit quadratic grade C_1^2/U=C_2 in Papers III-C, III-E, and VIII.
% The reflection is a typed trace/scale theorem, not a time reversal or a
% physical event-occurrence theorem; reheating transfer remains conditional.
% Release v1.54 registers the exact positive-part Laurent indices, the
% active-domain/terminal-endpoint action split, the no-second-EH firewall, and
% the q_DM=6/K_DM=48 scalaron readout typing across Papers III-D/E/F, VI,
% VIII, and IX.
% Release v1.53 registers the common nonnegative Laurent obstacle projection
% for n_1^+ and n_7^+, discharges the three former exit/freeze assumptions,
% and advances the Paper-III-D/E/F, VIII, and IX consumers.  Shared notation
% and bibliography ownership are unchanged.
% Release v1.51 registers the PQTI derivation batch and its updated
% Paper-III-A/C/D/E/F, VIII, and IX consumers.  Shared notation and
% bibliography ownership are unchanged.
% Release v1.50 advances the Paper-III-B AUX target to v1.40 after the
% rigorous audit repair; shared notation and bibliography ownership are unchanged.
% Release v1.49 advances the Paper-II-E/F AUX target to v3.14 after the
% minimal audit repair; shared notation and bibliography ownership are unchanged.
% Release v1.46 synchronizes Paper II-A's minimal frozen constructor-origin
% theorem and its five dependent AUX targets.  Shared notation is unchanged.
% Release v1.45 updates the Paper-II-B AUX target to v2.23 after the
% GP3-rigidity and appendix-placement repair; shared notation is unchanged.
% Release v1.44 advances the centralized Paper-I-A--VII-C AUX registry after the
% Paper-I-A v5.16 protected-completion scope repair.  Every current
% paper imports this one file; no paper-level opt-in or \externaldocument list
% is permitted.  The bibliography remains centralized here through Born-v6_13.
% Shared bibliography, theorem infrastructure, colors, notation alignment,
% and current Paper-I-A--VII-C external-reference registration.
% This file centralizes common imports, bibliography sorting, colors, theorem
% infrastructure, and the notation shared by the complete series.
% v2.9 (release v1.43) made the Paper-I-A--VII-C AUX registry automatic and
% centralized the bibliography and all external-reference ownership.
% v2.8 release-registry refresh (v1.33) registers the rooted generated-parent
% cascade and the POL/REG/free-law/11+4/Born-group closure revisions.  No
% shared mathematical notation or Paper-II-A fundamental principle is changed.
% v2.8 release-registry refresh (v1.32) registers the universal generated-
% continuation revisions in Papers II-A, III-A, III-C, IV, and VIII.  No
% shared mathematical notation or Paper-II-A fundamental principle is changed.
% v2.8 release-registry refresh (v1.30) registers the CRP4/projective-state
% source revisions and their direct registry-only consumers.  No shared
% mathematical notation or Paper-II-A axiom is changed.
% v2.7 release-registry refresh (v1.27) updates only the current external
% Paper-I-A--VII-C basenames after the placement/assumption synchronization; no
% shared mathematical notation or Paper-II-A axiom is changed.
% v2.7 synchronizes the geometric defect-repair revision: Paper-II-A v3.42,
% Paper-III-C v1.26, Paper-III-D v1.27, Paper-III-E v1.32, and Paper-III-F
% v1.34.  It registers the common trace-defect/spectral-filtration theorem,
% its owner-typed response pullback, and the reduced GRED and TERM interfaces.
% No shared mathematical notation is changed.
% v2.6 synchronizes Paper-III-E v1.31's cone-type repair, minimal-response
% classification, and explicit classical action-role boundary, together with
% Paper-III-F v1.32's downstream response-atlas synchronization.  No shared
% mathematical notation is changed.
% v2.5 synchronizes the Paper-III-D v1.26 faithful-trace-duality proof,
% exact ordered-owner joint locus, and typed alternative-presentation
% bookkeeping repair, together with its Paper-III-E and closure-frontier
% references.  No shared mathematical notation is changed.
% v2.4 synchronizes the Paper-III-C v1.25 orbit/process/owner repair, its
% downstream Papers III-B--F, V, VI, VIII, and IX revisions, and Born-v6_13.
% v2.3 synchronizes the Paper-II-A v3.41 witness/support/principal inheritance
% layer, the Paper-III-A v1.24 joint-realizability/readout firewall, all dependent
% host-rate revisions, and the current Paper-I-A--VII-B source registry.
% v2.2 synchronizes the generated active-realization, determinant-depth,
% dimension-weighted hierarchy, and typed cosmological-event interfaces.
% v2.1 synchronizes the exceptional-support interface revisions: Paper I-A
% v5.9, Paper II-A v3.30, Paper II-E/F v3.7, and Paper IV v4.16.
% v2.0 introduces the stage-safe G0 occurrence notation, synchronizes the
% active Paper-I-A--IV source registry (current rigorous Paper-III closure:
% III-A v1.17, III-B v1.22, III-C--E v1.17, III-F v1.14, IV v4.14), and makes
% external-reference self-skip depend on logical paper role rather than a
% versioned job name.
% ======================================================================

% ======================================================================
% macros.tex — Shared macros for the Born Geometry series
% ======================================================================

% ======================================================================
% 1. PACKAGES, BIBLIOGRAPHY, COLORS, AND UNIT DEFINITIONS
% ======================================================================
% Core packages shared across the complete Paper-I-A--VII-C series.  Loading them here makes
% the macro file a genuine series preamble while remaining compatible with existing
% paper preambles that load the same packages before \input.
\usepackage{amsmath}
\usepackage{amssymb}
\usepackage{amsthm}
\usepackage{mathrsfs}
\usepackage{array}
\usepackage{booktabs}
\usepackage{xcolor}
\makeatletter
\@ifpackageloaded{geometry}{}{\usepackage{geometry}\geometry{margin=1in}}
\makeatother
\usepackage{mdframed}

% Shared support packages required by the macros below and by most papers.
\usepackage{ragged2e}
\usepackage{longtable}
\usepackage{tabularx}
\usepackage{enumitem}
\usepackage[most]{tcolorbox}

% Series-wide numeric bibliography convention using the portable BibTeX path.
% All Papers I-A--VII-C use the ordinary \cite interface.
\usepackage[numbers,sort&compress]{natbib}
\providecommand{\printbibliography}{%
  \bibliographystyle{plainnat}%
  \bibliography{Born-v6_59}%
}

% Shared colors.  \providecolor avoids clobbering a paper-specific override.
\providecolor{postulatecolor}{RGB}{230,242,255}
\providecolor{compare}{RGB}{230,245,230}
\providecolor{definitioncolor}{RGB}{230,242,255}

\providecolor{openriskcolor}{RGB}{255,235,235}
\providecolor{bridgecolor}{RGB}{240,240,255}
\providecolor{gold}{RGB}{255,215,0}

% ======================================================================
% 2. CUSTOM ENVIRONMENTS
% ======================================================================
% Common table column types used in all papers.  Guarded so that a paper may
% define its own versions before loading this file.
\makeatletter
\@ifundefined{NC@find@L}{%
  \newcolumntype{L}[1]{>{\RaggedRight\arraybackslash}p{#1}}%
}{}
\@ifundefined{NC@find@C}{%
  \newcolumntype{C}[1]{>{\Centering\arraybackslash}p{#1}}%
}{}
\@ifundefined{NC@find@Y}{%
  \newcolumntype{Y}{>{\RaggedRight\arraybackslash}X}%
}{}
\makeatother

% ======================================================================
% 3. THEOREM STYLES AND ENVIRONMENTS
% ======================================================================
% Custom underlined theorem style (useful for remarks or notes)
\newtheoremstyle{underlined}
  {3pt}   % Space above
  {3pt}   % Space below
  {}      % Body font (empty means normal/roman text, typical for remarks)
  {}      % Indent amount
  {}      % Theorem head font (can be \itshape, \bfseries, etc.)
  {.}     % Punctuation after theorem head
  { }     % Space after theorem head
  {\underline{\thmname{#1}\thmnumber{ #2}}\thmnote{ (#3)}} % Theorem head formatting

% Standard plain theorem style environments.  Each environment is guarded so
% the shared macro file can be loaded after a paper-specific theorem setup.
\theoremstyle{plain}
\makeatletter
\@ifundefined{theorem}{\newtheorem{theorem}{Theorem}[section]}{}
\@ifundefined{lemma}{\newtheorem{lemma}{Lemma}[section]}{}
\@ifundefined{proposition}{\newtheorem{proposition}{Proposition}[section]}{}
\@ifundefined{corollary}{\newtheorem{corollary}{Corollary}[section]}{}
\makeatother

% Definition style environments (typically non-italicized body)
\theoremstyle{definition}
\makeatletter
\@ifundefined{definition}{\newtheorem{definition}{Definition}[section]}{}
\@ifundefined{postulate}{\newtheorem{postulate}{Postulate}[section]}{}
\@ifundefined{closure}{\newtheorem{closure}{Closure}[section]}{}
\makeatother

% Applying the custom underlined style
\theoremstyle{underlined}
\makeatletter
\@ifundefined{remark}{\newtheorem{remark}{Remark}[section]}{}
\@ifundefined{convention}{\newtheorem{convention}{Convention}[section]}{}
\@ifundefined{principle}{\newtheorem{principle}{Principle}[section]}{}
\makeatother

% Framed proof environment.
% Every \begin{proof} ... \end{proof} is automatically framed/shaded.
% This avoids mdframed/trivlist nesting, which can drop the top border after
% theorem-like environments.
\makeatletter
\renewenvironment{proof}[1][\proofname]{%
  \par\addvspace{6pt}%
  \noindent
  \begin{tcolorbox}[%
    enhanced,
    breakable,
    colback=black!4,
    colframe=black!35,
    boxrule=0.35pt,
    arc=1pt,
    left=6pt,
    right=6pt,
    top=4pt,
    bottom=4pt,
    boxsep=0pt,
    before skip=0pt,
    after skip=6pt,
    parbox=false
  ]%
  \pushQED{\qed}%
  \normalfont\itshape #1\@addpunct{.}\enspace
  \normalfont\ignorespaces
}{%
  \unskip\nobreak\hfill\popQED
  \end{tcolorbox}%
  \par\addvspace{2pt}%
}
\makeatother

\mdfdefinestyle{postulatebox}{
  backgroundcolor=yellow!10,
  linecolor=black,
  linewidth=0.4pt,
  skipabove=\baselineskip,
  skipbelow=\baselineskip,
  nobreak=true
}

% ======================================================================
% PROOF-READABILITY HELPERS
% ======================================================================

\newcommand{\ProofPlan}[1]{%
  \par\smallskip
  \noindent\textbf{Proof plan. }#1%
  \par\smallskip
}

\newlist{proofsteps}{enumerate}{2}
\setlist[proofsteps,1]{
  label=\textbf{(\alph*)},
  ref=\alph*,
  leftmargin=2.2em,
  itemsep=0.45\baselineskip,
  topsep=0.35\baselineskip
}
\setlist[proofsteps,2]{
  label=\textbf{(\roman*)},
  ref=\roman*,
  leftmargin=2.0em,
  itemsep=0.25\baselineskip,
  topsep=0.2\baselineskip
}


% ======================================================================
% 4. NUMBER SYSTEMS AND ALPHABETS
% ======================================================================
\providecommand{\RR}{\mathbb{R}} % Real numbers
\providecommand{\PP}{\mathbb{P}} % Projective space / Probability
\providecommand{\CC}{\mathbb{C}} % Complex numbers
\providecommand{\ZZ}{\mathbb{Z}} % Integers
\providecommand{\NN}{\mathbb{N}} % Natural numbers
\providecommand{\HH}{\mathbb{H}} % Quaternions / Hyperbolic space
\providecommand{\LFRayGeom}{\mathsf{LF\!\mbox{-}RAY\!\mbox{-}GEOM}}
\providecommand{\LFRep}{\mathsf{LF\!\mbox{-}REP}}
\providecommand{\LFState}{\mathsf{LF\!\mbox{-}STATE}}
\providecommand{\LFLim}{\mathsf{LF\!\mbox{-}LIM}}
\providecommand{\LFScat}{\mathsf{LF\!\mbox{-}SCAT}}
% Number systems / projective spaces
\providecommand{\CP}{\mathbb{CP}}
\providecommand{\RP}{\mathbb{RP}}

% ======================================================================
% 5. GROUPS AND LIE THEORY
% ======================================================================
\providecommand{\SO}{\ensuremath{\mathrm{SO}}}     % Special Orthogonal group
\providecommand{\SU}{\ensuremath{\mathrm{SU}}}     % Special Unitary group
\providecommand{\Spin}{\ensuremath{\mathrm{Spin}}} % Spin group
\providecommand{\CS}{\ensuremath{\mathrm{CS}}}     % Chern-Simons (capitalized)
\providecommand{\cs}{\ensuremath{\mathrm{cs}}}     % Chern-Simons (lowercase)

% ======================================================================
% 6. COMMON MANIFOLDS AND TOPOLOGY
% ======================================================================
\providecommand{\Sone}{S^1}                 % Standard circle

% ======================================================================
% 7. MATH OPERATORS (Declared via \DeclareMathOperator)
% ======================================================================
\DeclareMathOperator{\End}{End}   % Endomorphisms
\DeclareMathOperator{\Aut}{Aut}   % Automorphisms
\DeclareMathOperator{\spec}{spec} % Lowercase spectrum
\DeclareMathOperator{\tr}{tr}     % Lowercase trace (if you already have \Tr)
\DeclareMathOperator{\supp}{supp} % Support
\DeclareMathOperator{\TrA}{Tr_{A}} % Partial trace over subsystem A
\DeclareMathOperator{\TrB}{Tr_{B}} % Partial trace over subsystem B
% Lie/algebra operators
\DeclareMathOperator{\Ad}{Ad}
\DeclareMathOperator{\Sym}{Sym}
\DeclareMathOperator{\su}{su}
% ======================================================================
% 8. STANDARD OPERATOR-LIKE SYMBOLS (Safe definitions via \providecommand)
% ======================================================================
\providecommand{\vol}{\operatorname{vol}}   % Volume form
\providecommand{\Vol}{\operatorname{Vol}}   % Total volume
\providecommand{\Ric}{\operatorname{Ric}}   % Ricci tensor
\providecommand{\Rm}{\operatorname{Rm}}     % Riemann tensor
\providecommand{\Tr}{\operatorname{Tr}}     % Uppercase trace
\providecommand{\Spec}{\operatorname{Spec}} % Uppercase spectrum
\providecommand{\diag}{\operatorname{diag}} % Diagonal matrix
\providecommand{\rank}{\operatorname{rank}} % Rank
\providecommand{\scal}{\mathcal R} % Scalar curvature
\providecommand{\Cl}{\operatorname{Cl}}     % Clifford algebra
\providecommand{\diam}{\operatorname{diam}} % Diameter

% ======================================================================
% 9. FUNCTIONS, TRIGONOMETRY, AND CALCULUS
% ======================================================================

% Codifferential naming conventions
% (Avoid \ddag: commonly defined by physics.sty and symbol packages)
\providecommand{\codiff}{\mathrm{d}^{\dagger}} % Codifferential operator

% Quantum mechanics bra-ket notation (needed for basis tables)
\providecommand{\ket}[1]{\lvert #1 \rangle}
\providecommand{\bra}[1]{\langle #1 \rvert}
\providecommand{\braket}[2]{\langle #1 \mid #2 \rangle}

% ======================================================================
% 10. PHYSICAL SCALES AND CONSTANTS (Paper 1c / Born Geometry)
% ======================================================================
\providecommand{\Rp}{\ensuremath{R_{\mathrm P}}}     % matched common EH curvature scale (not defined from observed G)
\providecommand{\hBorn}{\ensuremath{h_{\mathrm B}}} % physical representative 2\pi\hbar of one primitive Born scalar action cell
\providecommand{\UEHjoint}{\ensuremath{U_{J}}} % joint-Dirac EH curvature-coordinate scale 8 kappa_J
\providecommand{\Gread}{\ensuremath{G_{\rm scal}}} % scalar-action coordinate c_op^3/(h_B U); physical matching explicit
\providecommand{\GreadJoint}{\ensuremath{G_{J,\mathrm{raw}}}} % raw ordinary-trace coordinate = G_scal/96
\providecommand{\Rds}{R_{\mathrm{dS}}}    % de Sitter Curvature
\providecommand{\Mpl}{M_{\mathrm{Pl}}}    % Planck mass
\providecommand{\lp}{\ell_{\mathrm{Pl}}}  % Planck length
\providecommand{\Pmax}{P_{\mathrm{max}}}  % Maximum momentum
\providecommand{\Mplbar}{\bar{M}_{\mathrm{Pl}}} % Reduced Planck mass
\providecommand{\eff}{\mathrm{eff}}       % Subscript for effective quantities
\providecommand{\Cosmo}{\ensuremath{\Lambda}}                % Cosmological constant (use \Cosmo in source)
\providecommand{\LCDM}{\ensuremath{\Cosmo\mathrm{CDM}}}  % \LambdaCDM (typesets with \Cosmo)
\providecommand{\Lcut}{\ensuremath{\Lambda_{\mathrm{cut}}}}  % Cutoff scale (brace when used as script)
\providecommand{\Rcs}{\ensuremath{R_{\mathrm{cs}}}}      % Configuration-space radius
\providecommand{\CosmoEff}{\ensuremath{\Cosmo_{\eff}}}   % Effective cosmological constant

% Disambiguation: \Rmom is the momentum-space (S^3) radius. 
% Reserve plain R for the 4D curvature scalar in f(R).
\providecommand{\Rmom}{\ensuremath{R_{\mathrm{mom}}}} 
\providecommand{\Punit}{P_{0}} % Momentum unit set by \Rmom (= \hbar/\Rmom)
\providecommand{\Rvec}{\ensuremath{R_{\mathrm{vec}}}} % invariant vector-space radius in the reduced round sector
% Specific numeric values and expressions
\providecommand{\Vlogcap}{L_{\mathrm{Hopf}}^{(1)}}
\providecommand{\Hopf}{\operatorname{Hopf}}
% Scalar phase (S^1) sector + universal cover notation (Paper I-A)
\providecommand{\tper}{T_\star}      % Time period/scale
\providecommand{\Eper}{E_\star}      % Energy period/scale
\providecommand{\ttilde}{\tilde t}   % Universal cover time
\providecommand{\Etilde}{\tilde E}   % Universal cover energy

% ======================================================================
% 10A. SERIES-WIDE LOGICAL / NOTATIONAL DISAMBIGUATIONS
% ======================================================================
% Use these symbols to avoid collisions among: (i) the Paper--III
% no-insertion/geometric-closure class, (ii) Paper--I-A trimer diagonals,
% (iii) Paper--V-A/V-C family-space diagonals, (iv) scalar-index running,
% and (v) protected shell-clock labels.
\providecommand{\Rprim}{\ensuremath{\mathsf R_{\mathrm{prim}}^{\Sigma}}} % primitive finite structural readouts
\providecommand{\Rloop}{\ensuremath{\mathsf R_{\circ}^{\Sigma}}} % closed-loop readout role
\providecommand{\Rbdry}{\ensuremath{\mathsf R_{\partial}^{\Sigma}}} % boundary-conditioned readout role
\providecommand{\Rgen}{\ensuremath{\mathsf R_{\mathrm{fin}}^{\Sigma}}} % finite generated closure of primitive roles
\providecommand{\PhysOcc}{\operatorname{PhysOcc}}
\providecommand{\Actual}{\operatorname{Actual}}
\providecommand{\GeomLic}{\operatorname{GeomLic}}
\providecommand{\InvDist}{\operatorname{InvDist}}
\providecommand{\Dcl}{\ensuremath{\mathsf{D}_1}}      % unmarked positive Dirac-even closed spectral-readout form class
\providecommand{\Dheat}{\ensuremath{\mathsf D_{\mathrm{heat}}}} % one-parameter heat-semigroup subclass
\providecommand{\DtriI}{\ensuremath{D^{\mathrm{tri}}_1}} % Paper--I-A cyclic-trimer diagonal
\providecommand{\DtriII}{\ensuremath{D^{\mathrm{tri}}_2}}
\providecommand{\DfamI}{\ensuremath{D^{\mathrm{fam}}_1}} % Family-space diagonal selector
\providecommand{\DfamII}{\ensuremath{D^{\mathrm{fam}}_2}}
\providecommand{\alpharun}{\ensuremath{\alpha_{\mathrm{run}}}} % scalar spectral-index running
\providecommand{\Ksh}{\ensuremath{K_{\mathrm{sh}}}} % protected shell-clock coordinate
\providecommand{\ddef}{\ensuremath{\delta_{\mathrm{def}}}} % universal defect scalar after Paper--III realization

% ======================================================================
% 10B. SERIES-WIDE LOGICAL STATUS TAGS
% ======================================================================
% Upright small-caps wrapper for audit/status labels.  It prevents LaTeX from
% trying to use unavailable small-caps italic font shapes inside theorem bodies
% such as TU/lmr/m/scit.
\providecommand{\AuditLabel}[1]{}
\renewcommand{\AuditLabel}[1]{{\normalfont\scshape #1}}

% These tags are bookkeeping labels for consistency audits. They do not
% create new mathematical hypotheses; they record how a datum is being used
% in the proof chain.
\providecommand{\StatusPrimitive}{}
\renewcommand{\StatusPrimitive}{\AuditLabel{Primitive}}
\providecommand{\StatusDerived}{}
\renewcommand{\StatusDerived}{\AuditLabel{Derived}}
\providecommand{\StatusImported}{}
\renewcommand{\StatusImported}{\AuditLabel{Imported}}
\providecommand{\StatusConditional}{}
\renewcommand{\StatusConditional}{\AuditLabel{Conditional}}
\providecommand{\StatusBridge}{}
\renewcommand{\StatusBridge}{\AuditLabel{Bridge-derived}}
\providecommand{\StatusComparison}{}
\renewcommand{\StatusComparison}{\AuditLabel{Comparison-only}}

% ======================================================================
% 10C. SERIES-WIDE CLAIM-WORDING AUDIT LABELS
% ======================================================================
% These are textual audit labels used to keep modal language consistent
% across the series.  They introduce no new mathematical assumptions.
\providecommand{\ClaimForced}{}
\renewcommand{\ClaimForced}{\AuditLabel{forced}}
\providecommand{\ClaimUnique}{}
\renewcommand{\ClaimUnique}{\AuditLabel{unique}}
\providecommand{\ClaimCanonical}{}
\renewcommand{\ClaimCanonical}{\AuditLabel{canonical}}
\providecommand{\ClaimMinimal}{}
\renewcommand{\ClaimMinimal}{\AuditLabel{minimal}}
\providecommand{\ClaimClosed}{}
\renewcommand{\ClaimClosed}{\AuditLabel{closed}}
\providecommand{\ClaimPrediction}{}
\renewcommand{\ClaimPrediction}{\AuditLabel{prediction}}

% ======================================================================
% 10D. SERIES-WIDE THEOREM-CHAIN AUDIT LABELS
% ======================================================================
% These labels mark the direction of load-bearing theorem chains. They are
% audit annotations only; they introduce no new mathematical assumptions.
\providecommand{\ChainInput}{}
\renewcommand{\ChainInput}{\AuditLabel{Chain input}}
\providecommand{\ChainOutput}{}
\renewcommand{\ChainOutput}{\AuditLabel{Chain output}}
\providecommand{\ChainGuard}{}
\renewcommand{\ChainGuard}{\AuditLabel{Chain guard}}
\providecommand{\ChainNoBackflow}{}
\renewcommand{\ChainNoBackflow}{\AuditLabel{No backflow}}

% ======================================================================
% 10E. SERIES-WIDE LOCAL THEOREM-AUDIT LABELS
% ======================================================================
% These labels mark a statement-level dependency audit around individual
% theorem/definition nodes. They are bookkeeping labels only.
\providecommand{\LocalAudit}{}
\renewcommand{\LocalAudit}{\AuditLabel{Local audit}}
\providecommand{\UsesOnly}{}
\renewcommand{\UsesOnly}{\AuditLabel{Uses only}}
\providecommand{\Exports}{}
\renewcommand{\Exports}{\AuditLabel{Exports}}
\providecommand{\DoesNotUse}{}
\renewcommand{\DoesNotUse}{\AuditLabel{Does not use}}

% ======================================================================
% 10F. SERIES-WIDE NUMERICAL-ROW AUDIT LABELS
% ======================================================================
% These labels classify individual rows in numerical, benchmark, status,
% and comparison tables.  They are audit annotations only and do not change
% any theorem, value, formula, or comparison number.
\providecommand{\RowAudit}{}
\renewcommand{\RowAudit}{\AuditLabel{Numerical-row audit}}
\providecommand{\RowTheorem}{}
\renewcommand{\RowTheorem}{\AuditLabel{Theorem-fixed row}}
\providecommand{\RowBridge}{}
\renewcommand{\RowBridge}{\AuditLabel{Bridge-fixed row}}
\providecommand{\RowRepresentative}{}
\renewcommand{\RowRepresentative}{\AuditLabel{Representative row}}
\providecommand{\RowComparison}{}
\renewcommand{\RowComparison}{\AuditLabel{Comparison row}}
\providecommand{\RowCalibrated}{}
\renewcommand{\RowCalibrated}{\AuditLabel{Calibrated row}}
\providecommand{\RowEstimate}{}
\renewcommand{\RowEstimate}{\AuditLabel{Estimate row}}
\providecommand{\RowOpen}{}
\renewcommand{\RowOpen}{\AuditLabel{Open row}}
\providecommand{\RowNoSelector}{}
\renewcommand{\RowNoSelector}{\AuditLabel{No selector}}

% ======================================================================
% 10G. TABLE HELPERS
% ======================================================================
% Table helpers used by Paper V-C.  The unused readout-label preset has been
% removed; only active table helpers remain.
\providecommand{\gpsection}[1]{%
  \addlinespace[0.55em]
  \multicolumn{3}{@{}l}{\textsc{#1}}\\[-0.25em]
  \midrule
}
\providecommand{\mathcell}[1]{%
  \(\begin{gathered}#1\end{gathered}\)%
}

\newcommand{\DerBy}[1]{%
  \begin{tabular}[t]{@{}l@{}}#1\end{tabular}%
}

\newcommand{\TableNotes}[1]{%
  \par\medskip
  \begin{minipage}{0.94\linewidth}
  \footnotesize\emph{Notes.} #1
  \end{minipage}\par
}
% ======================================================================
% 10H. PAPER VI-A / QCD NOTATION
% ======================================================================
% Shared symbols used by Paper VI-A and downstream QCD/Yang--Mills handoffs.
\providecommand{\Pqcd}{\mathfrak P_{\rm VIII}^{\rm QCD}}
\providecommand{\Iqcd}{\mathfrak I_{\rm QCD}^{(7)}}
\providecommand{\muQCD}{\mu_{\rm QCD}^{\rm int}}
\providecommand{\LambdaQCD}{\Lambda_{\rm QCD}}
\providecommand{\LambdaBH}{\Lambda_{\rm BH}^{\rm QCD}}
\providecommand{\HQCD}{\mathcal H_{\rm QCD}^{\rm BH}}
\providecommand{\HQCDfin}{H_{\rm QCD}^{\rm fin}}
\providecommand{\Hraw}{\mathcal H_{\rm raw}}
\providecommand{\HphysQCD}{\mathcal H_{\rm phys}^{\rm QCD}}
\providecommand{\Vcol}{V_C}
\providecommand{\Ps}{P_{\rm sing}}
\providecommand{\PiOct}{\Pi_8}
\providecommand{\MSbar}{\overline{\rm MS}}
\providecommand{\YM}{\mathrm{YM}}
\providecommand{\Ktri}{K_{\triangle}}
\providecommand{\KYM}{K_{\YM}^{\rm BH}}
\providecommand{\aBH}{a_{\rm BH}}
\providecommand{\epsBH}{\varepsilon_{\rm BH}}
\newcommand{\RowConstructive}{\textnormal{\textsc{Constructive}}}
\newcommand{\RowObligation}{\textnormal{\textsc{Bridge obligation}}}
\newcommand{\RowTarget}{\textnormal{\textsc{Target row}}}
\newcommand{\RowValidation}{\textnormal{\textsc{Validation row}}}
\providecommand{\slashed}[1]{\not{#1}}

% ======================================================================
% 10I. PAPER VI-B / CONTINUOUS-QFT NOTATION
% ======================================================================
% Shared symbols used by Paper VI-B and any downstream continuous-QFT or
% Yang--Mills handoff paper.  All definitions are guarded so individual
% papers may override them locally before loading this file if needed.
\providecommand{\A}{\mathfrak A}
\providecommand{\Hphys}{\mathcal H_{\rm phys}}
\providecommand{\QCD}{\mathrm{QCD}}
\providecommand{\BH}{\mathrm{BH}}
\providecommand{\loc}{\mathrm{loc}}

\providecommand{\RBH}{\mathcal R_{\rm BH}}
\providecommand{\DG}{\mathfrak D_G}
\providecommand{\PW}{\mathcal{P\!W}} % Peter--Weyl spectral-symbol subspace
\providecommand{\Toe}{\operatorname{Toe}} % Toeplitz compression/representative
\providecommand{\Hol}{\operatorname{Hol}} % Holonomy

% ======================================================================
% 11. LATE-LOADED PACKAGES (Must be at the end of the preamble)
% ======================================================================
% Load late packages only if a paper has not already loaded them, avoiding
% option clashes in paper-specific preambles.
\makeatletter
\@ifpackageloaded{hyperref}{}{\usepackage{hyperref}}
\makeatother
% --- Series-wide disambiguations (Pass 3) ---
\providecommand{\diracdelta}{\delta_{\rm D}} % Dirac delta distribution (delta reserved for hierarchy)
\providecommand{\Hg}{H_{\mathrm g}} % Hodge harmonic projector (keep H for Hubble)
\providecommand{\Gr}{G_{\mathrm r}} % Green operator (keep G for Newton)
\providecommand{\FD}{\mathrm{D}} % functional/variational derivative symbol (delta reserved)
\providecommand{\Id}{\mathbf{1}} % identity/Kronecker tensor (avoid \delta_{ij})


% BEGIN INLINE SOURCE: Paper00-ActiveExternalReferences-v2_130.tex
% Framework v2.191: current reference binding; existing definitions retained.
% Framework v2.190: current22-root registry with native clock-evolution proofs.
% Framework v2.189: exact current22-root external registry, including retained v2.188 hardening owners.
% Framework v2.187: current registry synchronization only; source release v2.186 preserved.
% Framework v2.186: authoritative current 22-paper registry.
% Framework v2.182: all current 22 logical paper AUX bindings.
% Framework v2.170: current I-B v1.43 external AUX binding.
% Framework v2.167: exact current 22-root registry after CP--BL2 integration.
% Framework v2.166: exact current 22-root compact-release AUX registry.
% Framework v2.165: authoritative current root bindings.
% Framework v2.164: exact current 22 manuscript AUX bindings.
% Framework v2.162: finite metric, functional Hilbert and scoped Ward integration.
% Framework v2.161: current 22 manuscript AUX bindings; source only.
% Framework v2.156 / registry v2.113: synchronized current 22 manuscript roots; source-only registry.
% Framework v2.143 / registry v2.111: active manuscript AUX basenames; no generated AUX or PDF included.
% ======================================================================
% Central external-reference imports, numbers only, no PDF jumping
% ======================================================================

% xr-hyper is loaded centrally so that papers may retain their existing
% \usepackage{xr-hyper} lines without redefining \externaldocument after this
% registry is installed.  If plain xr was loaded earlier, retain it.
\makeatletter
\@ifpackageloaded{xr-hyper}{}{%
  \@ifpackageloaded{xr}{}{\usepackage{xr-hyper}}%
}
\makeatother

% Make external-document registration idempotent.  The registry identifies a
% logical paper by its Paper<number><optional letter> prefix.  Prefix-qualified
% imports remain independent.  Current paper roots contain no local
% \externaldocument calls; this guard also makes accidental duplicates harmless.
\ExplSyntaxOn
\prop_new:N \g_bh_external_refs_prop
\cs_new_eq:NN \BHOriginalExternalDocument \externaldocument

\cs_new_protected:Npn \bh_external_document:nnnn #1#2#3#4
 {
  \tl_set:Nn \l_tmpa_tl {#3}
  \regex_replace_once:nnN
   { \A (?: xref/ )? (Paper[0-9]+[A-F]?) .* \Z }
   { \1 }
   \l_tmpa_tl
  \tl_set:Nx \l_tmpb_tl { #1 | \tl_use:N \l_tmpa_tl }
  \prop_if_in:NVTF \g_bh_external_refs_prop \l_tmpb_tl
   {
    \PackageInfo{BornHopf}
     {Skipping~duplicate~external~refs~for~#3}
   }
   {
    \prop_gput:NVn \g_bh_external_refs_prop \l_tmpb_tl {#3}
    \IfNoValueTF{#4}
     { \BHOriginalExternalDocument[#1][#2]{#3} }
     { \BHOriginalExternalDocument[#1][#2]{#3}[#4] }
   }
 }

\RenewDocumentCommand{\externaldocument}{O{} O{} m o}
 { \bh_external_document:nnnn {#1}{#2}{#3}{#4} }

\NewDocumentCommand{\BHExternalRefs}{m}
 {
  % Do not import the aux file of the logical paper currently being compiled,
  % even when the job and registry use different version suffixes.
  \tl_set:Nx \l_tmpa_tl { \c_sys_jobname_str }
  \regex_replace_once:nnN
   { \A (Paper[0-9]+[A-F]?) .* \Z }
   { \1 }
   \l_tmpa_tl
  \tl_set:Nn \l_tmpb_tl { #1 }
  \regex_replace_once:nnN
   { \A (Paper[0-9]+[A-F]?) .* \Z }
   { \1 }
   \l_tmpb_tl
  \str_if_eq:eeTF { \tl_use:N \l_tmpa_tl } { \tl_use:N \l_tmpb_tl }
   {
    \PackageInfo{BornHopf}
     {Skipping~external~refs~for~#1~because~this~is~the~current~document}
   }
   {
    \IfFileExists{xref/#1.aux}
     { \externaldocument[][nocite]{xref/#1}[#1.pdf] }
     {
      \IfFileExists{#1.aux}
       { \externaldocument[][nocite]{#1} }
       {
        \PackageInfo{BornHopf}
         {External~aux~file~#1.aux~not~found;~no~labels~imported}
       }
     }
   }
 }
\ExplSyntaxOff

% External references are automatic.  This file is the sole owner of both the
% bibliography and the current Paper-I-A--VII-D AUX registry (30 roots including I-C).  Missing auxiliaries
% are informational on the first pass and are populated by the source-only
% two-pass series build.  Basenames deliberately omit temporary download/copy
% suffixes such as ``(1)'' or ``-1''.
  % Paper I-A: intrinsic S3 and exceptional group introduction.
  \BHExternalRefs{Paper1A-Geometry-v5_175}

  % Paper I-B: complete exceptional geometry and finite spectral structures.
  \BHExternalRefs{Paper1B-ExceptionalGeometry-v1_5}
  % Paper I-C: reciprocal product M8 geometry.
  \BHExternalRefs{Paper1D-TechnicalMaterial-v1_65}
  % Paper I-D: terminal technical proofs and calculations.
  \BHExternalRefs{Paper1C-M8Geometry-v1_137}

\ifdefined\BHMathOnly
  % Papers I-A and I-B import only the mathematical atlas labels.
\else
  % Paper II: five physical-mathematical main papers plus one technical proof volume.
  \BHExternalRefs{Paper2A-GeometricConstitution-v3_221}
  \BHExternalRefs{Paper2B-MomentumCarrier-v2_144}
  \BHExternalRefs{Paper2C-ReciprocalLocalization-v2_169}
  \BHExternalRefs{Paper2D-SymmetrySpectra-v2_141}
  % The former II-F export labels remain in II-E; current II-F is the technical volume.
  \BHExternalRefs{Paper2E-KinematicExports-v3_138}
  % Paper II-F: complete technical proof and calculation owner for II-A--II-E.
  \BHExternalRefs{Paper2F-TechnicalMaterial-v1_67}

  % Paper III: Dynamics and Intrinsic Geometric Constitution (six components).
  \BHExternalRefs{Paper3A-GeneratedDynamics-v1_176}
  \BHExternalRefs{Paper3B-SpectralAction-v1_187}
  \BHExternalRefs{Paper3C-ProtectedHierarchy-v1_187}
  \BHExternalRefs{Paper3D-RunningResponse-v1_165}
  \BHExternalRefs{Paper3E-GravityAction-v1_197}
  \BHExternalRefs{Paper3F-CompleteDynamics-v1_263}
  % Paper III-G: complete preserved technical proof owner.
  \BHExternalRefs{Paper3G-TechnicalMaterial-v1_53}

  % Paper IV-A: generated finite quantum mechanics.
  \BHExternalRefs{Paper4A-FiniteQuantumMechanics-v4_214}
  % Paper IV-B: complete finite quantum processes and physical response.
  \BHExternalRefs{Paper4B-QuantumResponse-v1_49}
  % Paper IV-C: extended proofs and finite technical calculations.
  \BHExternalRefs{Paper4C-TechnicalMaterial-v1_49}

  % Paper V-A: The Standard Model.
  \BHExternalRefs{Paper5A-FiniteMatterGeometry-v3_209}

  % Paper V-B: Finite Spectrum and Flavor Closure.
  \BHExternalRefs{Paper5B-TerminalSpectrum-v5_199}

  % Paper V-C: Finite Interaction Boundaries of the Standard Model.
  \BHExternalRefs{Paper5C-InteractionResponse-v2_177}
  % Paper V-D: preserved technical proof and calculation owner.
  \BHExternalRefs{Paper5D-TechnicalMaterial-v1_53}

  % Paper VI-A: Gravity and Closed Cosmology.
  \BHExternalRefs{Paper6A-GravityCosmology-v2_260}

  % Paper VI-B: Geometric Dark Sector and Cosmological Readouts.
  \BHExternalRefs{Paper6B-CosmologyReadouts-v1_206}

  % Paper VI-C: complete technical proof and comparison owner.
  \BHExternalRefs{Paper6C-TechnicalMaterial-v1_50}

  % Paper VII-A: Finite QCD Kernel and Nonperturbative Closure.
\BHExternalRefs{Paper7A-QCDKernel-v1_172}

  % Paper VII-B: Hadronic, Thermal, and Nuclear QCD Readouts.
  \BHExternalRefs{Paper7B-QCDReadouts-v1_157}

  % Paper VII-C: Continuous Quantum Field Theory.
  \BHExternalRefs{Paper7C-ContinuousQFT-v2_202}

  % Paper VII-D: complete technical material for VII-A through VII-C.
  \BHExternalRefs{Paper7D-TechnicalMaterial-v1_53}

\fi

% END INLINE SOURCE: Paper00-ActiveExternalReferences-v2_130.tex


\hypersetup{hidelinks}
\setlength{\emergencystretch}{3em}
\title{\textbf{Born--Hopf Geometry in a Closed Universe IV-A: Finite Quantum Mechanics} \\[0.5em]\large Generated Test Algebra, States, Effects, Reciprocity, and Quantum Constitution}
\author{Peter Kugelmann\thanks{\texttt{pkugelmann@proton.me}}}
\date{Version 4.217 --- 20 November 2026 \\ Framework v3.067}
% Cross-DAG current-series import for the new source-owned mathematical bridge.
\BHExternalRefs{Paper1B-ExceptionalGeometry-v1_10}

\begin{document}
\maketitle
\begin{abstract}
Paper IV-A establishes finite quantum mechanics on the constituted level-one test algebra $\mathcal A_{\rm kin}\cong\mathbb C\oplus M_2\oplus M_2\oplus M_4\oplus M_3\oplus M_3$, with physical realization assumptions explicitly assigned to the geometric prefix. States, effects and Born evaluation follow from the finite algebra; tensor gluing, CP/Kraus representation and dynamical generators follow on their exact occurrence and composition domains. Finite reciprocal coordinates, operational measurement and tomography remain distinct from physically actual preparation and detector histories. Complete quantum responses and their source-visible poles are in Paper IV-B, and the preserved technical derivations are in Paper IV-C.
\end{abstract}
\setcounter{tocdepth}{2}
\tableofcontents
\clearpage
\noindent\textit{Original-source reference convention.} In unchanged inherited passages, Part~I refers to the finite quantum-core material now owned by Paper~IV-A. Former appendix references may resolve to the technical companion Paper~IV-C; their theorem labels and hypotheses remain unchanged.

\section{Framework prefix, geometric primacy and constituted carrier}
\label{p4a:sec:v238-a1}
\noindent\textit{Source and claim boundary.} The finite quantum system is derived from the declared Papers I--III geometric prefix, fixed compact momentum carrier and exact kinematic test algebra. Physical occurrence and optional event/field envelopes are tracked separately.

\paragraph{Framework Prefix and Finite Quantum Closure}
\label{p4:part:closed-qm}

\paragraph{Orientation.}
This opening part fixes the framework-prefix convention for Paper~IV and
proves the closed finite state/effect/operator calculus over
\(\mathfrak F_{\le3}^{\rm core}\) on the explicitly constituted roles of
Section~\ref{p4:subsec:constitutive-role-boundary}.  It is the native finite-algebra core of
the paper: the later parts are representation, interface, resolution-change,
or downstream handoff layers built from this closure.



\subsection{Framework-Prefix Status of Paper~IV}
\label{p4:sec:introduction}



\paragraph{Framework prefix and theorem scope}
\label{p4:subsec:framework-prefix-status}

Paper~IV is the finite-quantum layer over the framework prefix
\[
\begin{aligned}
\mathfrak F_{\le3}^{\rm core}
&:=
\mathrm{G0}+\mathfrak F_{\le2}
+\text{Paper-III reciprocal/scalar closure}\\
&\quad+\text{derived/co-realized }m_p=1
+\text{the GP1 operational cone}.
\end{aligned}
\]
where \(\mathfrak F_{\le1}\) is the Paper-I-A reference layer and
\(\mathfrak F_{\le2}\) is the compact kinematic constitution supplied by
Paper~II.  Paper~II supplies the realized vector A-carrier, the generated and
co-realized compact momentum cell \(m_p=1\), the scalar A-branch realizability
classification, and the admissible underdetermined B-profile target.
Paper~III realizes the scalar branch, derives the Born-reciprocal factors and
scalar cover over that fixed cell, and derives the timed propagation role from
the clocked-Hopf composite and its least constructor closure.  Its action,
hierarchy, analytic-continuation, locality, and continuum domains are imported
only where their own tags are displayed.

The exceptional-core role \(\mathsf{Supp}_{2T}\), the symbol \(\Dcl\),
\((\mathrm{REG}_\sigma)\), \(\mathsf{SF}_{\Lambda_F}\), \(({\rm QLOC})\),
\((\mathrm{DLOC}_H)\), and \(({\rm WICK}_\bullet)\) are not folded into this
unconditional prefix.  Whenever one is used, it remains an explicit
theorem-domain tag.  By contrast, on the generated reciprocal-localization
branch the historical \((\mathrm{GL}_{xp})\) label is a localization
certification, and on the occurring minimal local-action branch
Paper~III-F derives \((\mathrm{VDOM}_x^{\min})\); the legacy
\(({\rm POL})\) label adds no selector there.  Here
\(({\rm QLOC})\) is Paper~III-F's static compressed-observable
separated-support estimate, whereas \((\mathrm{DLOC}_H)\) denotes the
distinct Paper-IV propagation estimate for a selected Hamiltonian.  Neither
is implied by the GP1 cone alone.  The legacy label \(({\rm HT})\) denotes
Paper~III's derived common-rate
hierarchy certification and is not a separate hypothesis.  The current
fixed-shell realization family is
\[
 \mathfrak A_{\rm REG}
 =\{\mathrm{REG}_\sigma:\sigma_{\rm fin}\in\mathbb R\},
 \qquad
 \delta_\sigma=\frac1{6\sqrt2}+\sigma_{\rm fin},
\]
and \(\mathrm{REG}_0\) denotes the source member
\(\sigma_{\rm fin}=0\).  On the generated \(2O\)/Hurwitz interface,
Paper~III-C/III-F derive and anchor precisely this member, so no second REG
selector is available on that branch; unanchored REG branches retain the
full affine family.  Likewise the historical Paper-III-F labels
\((\mathrm{GL}_{xp})\) and \(({\rm POL})\) certify generated localization
and the derived minimal variational domain on the minimal branch.  Enlarged
Palatini, metric-affine, bimetric, or varying-momentum-representation theories
remain different parents and may carry additional domain data.  Finally \((\mathrm{CAL}_t):T_t>0\) denotes the
primitive-turn duration calibration.  The symbol
\(({\rm CAUS})\), by contrast, is the
Paper-III name for the derived GP1 operational cone, not an independent
hypothesis.

\begin{remark}[Profile and scale ledger]
\label{p4:rem:profile-scale-ledger}
The profiles used in Paper~IV are not silently identified:
\begin{center}
\small
\renewcommand{\arraystretch}{1.15}
\begin{tabularx}{\linewidth}{@{}>{\raggedright\arraybackslash}p{0.28\linewidth}>
{\raggedright\arraybackslash}p{0.31\linewidth}>
{\raggedright\arraybackslash}X@{}}
\toprule
\textbf{Object} & \textbf{Role} & \textbf{Independent data} \\
\midrule
\(F_g(D^2/\Lambda_g^2)\) & Generated gravitational heat readout &
\(F_g(u)=e^{-u}\) on that branch; \(\Lambda_g\) remains calibrated \\
\(q(D^2/\Lambda^2)\) & Static coarse localizer & \(q,\Lambda\) \\
\(\mathbf1_{[0,\Lambda_F^2]}(\Delta)\) & Physical sharp test frontier & \(\Lambda_F\) \\
\(C_{\mathcal K}(\tau)\) & Euclidean covariance & \(C_{\mathcal K}(0),S_\tau\) \\
\(e^{-itH_N}\) & Selected dynamical propagation & \(H_N\), DLOC branch \\
\bottomrule
\end{tabularx}
\end{center}
Equalities among these profiles or scales require an explicit same-process or
matching map.  Recovery, spectral flatness, equal-time locality,
dynamical locality, and convergence have separate error terms.
\end{remark}

For each result \(T\), let \(\mathcal R_T\) collect any residual empirical or
existence certificates displayed with it.  The legacy names
\[
 \{\mathsf{PROB},\mathsf{SUB},
   \mathsf{PROC},\mathsf{DYN},\mathsf{FENV},\mathsf{ENV}\}
\]
refer to the precise interfaces of
Section~\ref{p4:subsec:constitutive-role-boundary}; they are not independent
assumptions of matrix form, state normalization, tensor form, complete
positivity, or inner implementation.  Let \(\mathcal I_T\) collect all other
displayed interface/domain tags.  The standing sequent is
\[
\boxed{\mathfrak F_{\le3}^{\rm core}+\mathcal R_T+\mathcal I_T
\ \vdash\ T.}
\]
Only \(\mathcal R_T=\mathcal I_T=\varnothing\) makes \(T\) a prefix-only
native theorem.  A representation theorem may have empty \(\mathcal R_T\)
even though applying it to an actual preparation, composite occurrence,
process, or time flow requires the corresponding antecedent.
The words ``theorem'', ``forced'', ``minimal'', ``unique'', and ``excluded''
are used only relative to this prefix and the additional domain explicitly
stated in the relevant result.  No enriched finite-support datum or
conditional Paper-III action, hierarchy, continuation, locality, or
continuum datum is silently promoted into the core.



\paragraph{Native quantum problem}
\label{p4:subsec:native-quantum-problem}

The prefix \(\mathfrak F_{\le3}^{\rm core}\) supplies the finite
Born--reciprocal carrier, scalar cover, GP1 speed ceiling and Lorentzian cone.
Paper~II-C supplies the exact generated level-one test algebra
\(\mathcal A_{\rm kin}\) and its noncommutative effects.  Paper~III-B also
records \(\End(\mathcal H_m)\) as a mathematically available presentation
algebra, but explicitly withholds quantum-event meaning.  Paper~IV therefore
starts from the generated test algebra, not from every carrier endomorphism.

Paper~II-A already fixes \(S(\mathcal A)\) as the full positive-normalized
state space of a finite algebra.  Effect additivity, block densities, and the
Born trace pairing are consequently theorems; only empirical interpretation
and physical state selection remain interfaces.  GP2 relation-free
composition closes the tensor algebra after two occurrences, absence of
cross-relations, and faithful product-operation realization are certified.
Exact factor-operation interchange first derives commutation, so the tensor
law is not used in its own premise.  An affine state map extends positively
and trace preservingly, and positivity of its factor-local action on one
same-type product forces complete positivity.  Every noncentral self-adjoint element generates a nontrivial
kinematic automorphism flow, while any supplied/selected actual continuous
automorphism flow is blockwise inner.
Representation, apparatus, continuum, sector-specific, and actual dynamical
structures are introduced only on their stated theorem domains.



\paragraph{Constitutive role boundary}
\label{p4:subsec:constitutive-role-boundary}

The following are the residual root or token interfaces of the finite quantum
layer.  Duplicate role-level actualization has been removed.  They introduce no
numerical constant, preferred basis, sector projector, Hamiltonian, apparatus,
or continuum datum.

They all specialize Paper~II-A's stage-safe G0 geometric-actualization
certificate.  Whenever a physical quantum event, preparation, composite,
process, evolution, field bridge, or carrier-resolving test extension is
asserted to occur, its certificate must exhibit a contingent token of a
geometrically licensed, claim-relative invariant distinction with
result-independent provenance.  The labels below specify only the witness
roles: none asserts occurrence or supplies a separate actualization principle.
No instance reopens frozen upstream geometry; a new role or test extension
enters only through a declared downstream dependency frontier.  Analytic,
regularity, continuation, and reconstruction hypotheses remain separately
displayed.  In particular, \(\mathsf{ENV}\)'s full tomography is optional and
is not implied by claim-relative G0.
\begin{enumerate}[label=(\roman*), leftmargin=2.3em]
\item \(\mathsf{QEV}\) is retired as an interface.  When \(\mathcal A\) is
      already generated in the physical test slot, Paper~II-A's finite
      state--test theorem makes its self-adjoint, projection, effect, and
      event roles generated physical content.  A particular performed test is
      an ordinary G0 token and needs no new label.
\item \(\mathsf{PROB}\) is retained only as compatibility shorthand for a G0
      preparation, outcome/readout, finite-frequency record, or probability
      calibration token.  The normalized evaluation \(\omega(E)\), its
      probability-role typing, positivity, normalization, finite additivity,
      and block density form are generated before such a token is selected.
\item \(\mathsf{SUB}\): compatibility name for the G0 occurrence certificate
      that two constituted role occurrences exist, all factor tests persist,
      and no prior cross-relation is present.  GP2 then derives their tensor
      product; \(\mathsf{SUB}\) is not a separate composition law.
      Constructor-linked descendants of one already occurring witness are
      jointly generated by
      Proposition~\ref{p4:prop:common-witness-cooccurrence} and require no
      second occurrence certificate, although every inherited cross-relation
      remains live.
\item \(\mathsf{PROC}\): compatibility name for the G0 process certificate
      asserting the existence or selection of an actual affine state map and
      its factor-local action on an actually usable relation-free product, or
      for a generated system--environment dilation.  The factor-local joint
      action is the type-correct map
      \(\operatorname{id}\otimes\Phi\); its positivity is therefore process
      well-typedness rather than an independent same-copy axiom.  Linearity,
      trace preservation, complete positivity, and Kraus form are derived.
\item \(\mathsf{DYN}\): the G0 trajectory/process certificate that an actual homomorphism from the already oriented
      scalar-succession cover to a closed automorphism flow, or an actual open
      semigroup, has been generated or supplied on the named theorem domain.
      The parameter orientation is imported; selection of the actual flow,
      physical scale, and actual trajectory remains here.  On stationary
      free-scalar blocks, Theorem~\ref{p4:thm:pure-scalar-factor-representation}
      derives the pure scalar-factor representation and therefore removes its
      mathematical generator/scale ambiguity, but it does not actualize that
      representation as the physical DYN member or an actual-history token.
      For a fixed parametrized closed flow, its finite inner derivation and self-adjoint
      generator class modulo \(Z(\mathcal A)_{\rm sa}\) are derived; choosing a
      particular central representative is additional energy-zero/energy-
      observable data.  The finite GKLS representation on its stated open
      domain is likewise derived.
      Accordingly the generic actual-dynamics realization family is
      \[
      \begin{aligned}
        \mathfrak A_{\rm DYN}:=\{\alpha_t:\;&\alpha_t\text{ is an admissible
        actual closed/open evolution role}\\
        &\text{on its stated domain}\}.
      \end{aligned}
      \]
      choosing/actualizing a member is distinct from the derived inner/GKLS
      representation once that member occurs.
\item \(\mathsf{FENV}\): the G0 bridge certificate that an actual central-odd field occurrence has been
      generated and the total nonzero coupling graph is connected.  This
      generates the graded field envelope; it does not declare odd fields to
      be observables.
\item \(\mathsf{ENV}\): compatibility name for an optional carrier-resolving
      test root when no such family is already generated.  If constituted
      physical tests have a connected sector-coupling graph, Paper~II-C's
      graph-closure theorem forces their generated test algebra to be the full
      matrix algebra, which is then automatically carrier resolving; no ENV
      root is needed.  Equivalently, an already generated carrier-resolving
      test family also discharges ENV.  On the level-one central-grading
      domain, an independently adjoined ENV root is a parity-breaking
      extension distinct from \(\mathsf{FENV}\); a field-typed odd bridge by
      itself does not satisfy this test-typing premise.
\end{enumerate}

\begin{proposition}[Structural closure alone cannot entail an occurrence]
\label{p4:prop:structure-does-not-entail-occurrence}
Let \(\mathfrak M\) be any model of the Paper-IV structural language, and
let \(R\) be a nonempty generated role in \(\mathfrak M\).  Extend the
language by an occurrence predicate \(\operatorname{Occ}_R\) but add no
constructor or law connecting that predicate to the structural reduct.  Then
\(\mathfrak M\) has two expansions with the identical structural reduct:
one with \(\operatorname{Occ}_R=\varnothing\), and one with
\(\operatorname{Occ}_R=\{x\}\) for any chosen \(x\in R\).  Consequently no
sentence of the structural language alone entails that a preparation,
performed test, outcome, independent second token, process, bridge, or
trajectory occurs.
\end{proposition}

\begin{proof}
The interpretations of every structural sort, operation, relation, and axiom
are unchanged in the two expansions.  Since the new occurrence predicate is
absent from those axioms, both expansions satisfy exactly the same structural
theory, while \(\exists y\,\operatorname{Occ}_R(y)\) has opposite truth
values.  Any occurrence conclusion therefore needs either an explicit token
or an occurrence-preserving constructor from an already occurring witness.
\end{proof}

By Proposition~\ref{p4:prop:structure-does-not-entail-occurrence}, G0 does
not select a preparation, performed test, outcome, second independent
occurrence, actual map, or actual physical trajectory.  GP1--GP3
constrain and, where uniqueness is available, close the structural realization
of a constituted role.  All results
below distinguish those antecedents from the finite theorem that represents
them.

\begin{remark}[Status partition for Paper~IV]
\label{p4:rem:selector-status-partition}
The logical statuses used below are separated as follows.
\begin{enumerate}[label=(\roman*)]
\item \emph{Generated/derived/co-realized finite quantum structure:}
\(\mathcal A_{\rm kin}\), its state/effect/Born calculus, relation-free
tensor composition, CP/Kraus form after process occurrence, finite reductions,
spectral calculus, generated kinematic flows, apparatus-induced operator
systems, gauge averaging, finite statistics-selected constructions, and finite
Gibbs/KMS identities.
\item \emph{G0 tokens/occurrences:} one actual preparation/measurement package
containing the apparatus or channel, prepared state, performed test, physical
label ownership, outcome/readout, and any actual map or trajectory; plus
independent composite roots, field bridges, and process histories where used.
The package is bookkeeping: its fields are not generated from one another,
and only fields consumed by the claim are required.
\item \emph{Historical generic Paper-IV interface families:}
\(\mathsf{DYN}\), \(\mathsf{ACC}\), \(\mathsf{1P}\),
\(\mathsf{STAT}\), and \(\mathsf{OCC}\).  Their no-selection theorems remain
valid scope results for an arbitrary bare finite algebra, but they are not all
active framework-wide roots through Paper~VI-B.
\item \emph{Active through-IX physical residues:} the occurrence package in
(ii), a genuine reservoir/bath and Markov or kinetic limit when irreversibility
is claimed, and the role-specific EM, vacuum, flavor, neutral, thermal, or
cosmological operator/state/history data supplied by later papers.
\item \emph{Calibrations and analytic domains:} \(\mathrm{CAL}_t\),
\(\hbar\), temperature and energy references, central energy-zero
representatives, local-QFT/Wick domains, and controlled continuum limits.
\end{enumerate}
On the active stationary-free and visible-matter branches, one-particle
semantics and visible half-spin Fermi statistics are derived.  A numerical
bosonic occupation ceiling remains only an optional branch and has no active
consumer through Paper~VI-B.  Later physical owners do not retroactively change
the historical generic no-go theorems; they only remove those generic families
from the active current frontier.
\end{remark}

\begin{remark}[Prerequisites are relative to the consumed finite readout]
\label{p4:rem:v2187-claim-relative-finite-prerequisites}
A finite Born probability consumes its actual generated test algebra,
normalized state and effect. It does not require a pure state, a full
carrier-resolving \(\mathsf{ENV}\) extension, a graded field envelope,
one-particle semantics or a generic exchange-statistics selector. The last
roles enter only a claim that actually uses the corresponding event completion,
field or many-body construction. When a native branch already generates those
roles, they are reused rather than introduced again as independent inputs.

For a stationary tagged pole on one certified exact support, the consumed
object is the complete action-induced source-cyclic quotient of
Lemma~\ref{p4:lem:v2186-exact-source-cyclic-action} and
Theorem~\ref{p4:v2169:b10:thm:support}. A unique total Hamiltonian, unique
occupation member or unique pure vacuum is not a universal prerequisite.
An independently licensed family can determine that pole through
Theorem~\ref{p4:thm:v2186-common-source-cyclic-polynomial}; actual-pairing
isometry is additionally required only when common residues are claimed.
All mixed, boundary, constraint and interaction returns that enter the cyclic
action or change its support remain retained. This is not permission to delete
an unevaluated return or infer the actual electron root from the constructor.

On this fixed source chart a scalar \(H\mapsto H+cI\) is invisible to the
pole and normalized Gibbs/full-ground state, as in
Corollary~\ref{p4:cor:v2186-scalar-center-pole-invisibility}. Central
metric/source derivatives, absolute energies and sector-dependent centers
remain on their own consumed domains. Actual particle/rest identification,
nonstationary histories and each claimed limiting comparison are preserved.
\end{remark}



\paragraph{Paper-IV ownership and downstream import rule}
\label{p4:subsec:quantum-ownership-rule}

Paper~IV is the unique owner in the series of the finite quantum-mechanical
calculus over the realized carrier and the interface boundary of
Section~\ref{p4:subsec:constitutive-role-boundary}.  Its native theorem domain
comprises the generated finite test algebra, its full positive-normalized
state space, effects and POVMs, the blockwise Born pairing, tensor composition
for certified relation-free pairs, CP/Kraus form for certified processes,
partial evaluation and reduction, the representation of supplied
finite unitary and open-system dynamics, conserved quantities, finite
statistics-selected Fock sectors, exact CAR, the finite bosonic CCR boundary
defect, finite gauge projection, and finite Gibbs/KMS identities.  These
finite representation theorems are not rederived in later papers.

Papers~V-A--VII-B may supply sector selectors, Hamiltonians, particle identifications, couplings, backgrounds, transfer kernels, and readout functionals and then specialize the Paper-IV calculus to those data.  Paper~VII-C owns the controlled continuum reconstruction from compatible finite towers.  It may formulate selector-classification and continuum-limit consequences of the finite calculus, but the underlying finite algebra, state/Born rule, conservation, gauge-projection, thermal, and CCR statements remain Paper-IV imports.  Thus a later restatement of one of these finite results is to be read only as an explicit import or corollary, never as an independent foundation.



\paragraph{Imported \texorpdfstring{$\mathfrak F_{\le3}$}{F<=3} data}
\label{p4:subsec:imported-prefix-data}

\begin{table}[H]
\centering
\scriptsize

\begingroup
\setlength{\tabcolsep}{3pt}
\renewcommand{\arraystretch}{1.16}
\begin{tabular}{@{}L{3.0cm}L{3.4cm}L{4.4cm}L{3.0cm}@{}}
\toprule
\textbf{Imported package} & \textbf{Source layer} & \textbf{Use in Paper~IV} & \textbf{Status here} \\
\midrule
Round \(S^3\), Hopf/contact, \(\Spin(4)\), finite-carrier reference data & Paper~I-A & geometric and representation reference normalization & \textsc{Reference import} \\
Realized vector A-branch, scalar A-branch classification, B-profile target, and compact momentum truncation grammar & Paper~II & vector transfer topology, scalar realizability class, underdetermined reciprocal target, and sector bookkeeping & \textsc{Kinematic import} \\
Level-one diagonal-center grading & Paper~II-D & canonical
\(11\oplus4\) split, even/odd operator typing, and invariant-odd obstruction
& \textsc{Derived kinematic import} \\
Hopf/Reeb phase line & Papers~I-A--II & ray phase, charge-sector, Berry/Reeb interface & \textsc{Geometric import} \\
Born-reciprocal completion, initial reciprocal-position role, and scalar clock pair & Paper~III over the Paper-I-A Reeb primitive and Paper-II vector/scalar constitution & Born background and inherited scalar time/energy orientation; physical local realization uses \((\mathrm{GL}_{xp})\), while \((\mathrm{POL})\) is the explicit downstream selector assigning the position role to macroscopic metric variation and the momentum role to the fixed representation; static compressed-observable locality separately uses \(({\rm QLOC})\) & \textsc{Core plus bridge and selector interfaces} \\
Structural momentum-role carrier \(\mathcal H_p\cong\mathcal H_1\), \(m_p=1\) & Paper~II-B & fixed protected compact carrier co-realized with the occurring transfer operations; only particular readouts/calibrations add G0 tokens & \textsc{Generated physical-role import} \\
\(\mathsf{Supp}_{2T}\) Hurwitz-core role and its generated
\(C_6^{\rm grp}/C_6^{\rm link}\) descendants & Paper~I-A support atlas, with
the Paper-II-A occurrence bridge when used & optional finite internal
supports only where the core role is actualized or conditionally generated &
\textsc{Conditional support import} \\
\(\Dcl\) closed spectral-readout form class & theorem-level form after a G0 closed-process root or a certified generated process role; an arbitrary process law remains conditional & finite closed-loop spectral consequences on the displayed process domain & \textsc{Typed dynamics import} \\
Generated gravitational heat role and protected octet weights & Papers~III-C and III-E on
\(\mathsf C_{\rm IR}^{\rm der}+(\mathrm{GRAV\!-\!ASYM})\) &
\(F_g(u)=e^{-u}\), \(f_{2,g}=(4\pi)^{-2}\),
\(w_{\rm vac}=w_{\rm EH}=1/8\), \(w_{\rm nl}(U)=1\), and
\(w_{\rm nl}(C)=1/8\) with vanishing terminal MRESP visibility &
\textsc{Derived process/weight import}; MRESP coupling, renormalized vacuum,
cosmological solutions, and continuation remain tagged \\
Sharp scalar spectral frontier & Paper~III-F under
\(\mathsf{SF}_{\Lambda_F}\) & exact instantaneous/cumulative
reciprocal-position resolution staircase & \textsc{Conditional operational
frontier import} \\
Fixed-shell family \(\mathrm{REG}_\sigma\) & Papers~III-C--III-F &
\(\sigma_{\rm fin}\in\mathbb R\),
\(\delta_\sigma=1/(6\sqrt2)+\sigma_{\rm fin}\); the source value
\(1/(6\sqrt2)\) is specifically \(\mathrm{REG}_0\) &
\textsc{Unresolved selector family; source member explicit} \\
Common-rate hierarchy and Reeb use of \(\delta_\sigma\) & Paper~III, for each frozen \(\mathrm{REG}_\sigma\) member & the derived functional hierarchy and scheme-labelled value \(L_H^{(\sigma)}=4\pi^2+\delta_\sigma\) & \textsc{Derived conditional on frozen scheme} \\
GP1 operational speed ceiling \(c\) & Papers~II-A and III-A & sharp bound \(c=2\pi R_x/T_t\) on the generated clocked-Hopf doctrine & \textsc{Derived core import}; \((\mathrm{CAL}_t):T_t>0\) and numerical radius remain calibrations \\
Lorentzian operational cone and sign & Paper~III-E & causal ordering and conformal metric class & \textsc{Derived core import} \\
Local second-order free law and Euclidean reconstruction & Paper~III-E and stated Paper-IV blocks & finite dispersion only on the local free domain; real-time reconstruction only under its positivity/\(({\rm WICK}_\bullet)\) domain & \textsc{Tagged interface} \\
\bottomrule
\end{tabular}
\caption[Paper-IV import ledger]{Paper-IV import ledger.}
\label{tab:p4-framework-import-ledger}
\endgroup
\end{table}



\paragraph{Carrier convention: compact momentum versus reciprocal position unfolding}
\label{p4:sec:carrier-convention-mp-mx}

Throughout this paper, \(m\) indexes the mathematical finite-carrier family
or a representation, readout, or presentation level on the theorem domain
under discussion.  Paper~II does not physically realize every member of this
family merely by defining it.  The physical compact momentum carrier is
written explicitly as \(m_p\) and is imported from Papers~II--III as
\[
        m_p=1.
\]
Large laboratory, thermodynamic, local-net, scattering, or continuum-looking presentations belong to the reciprocal-position/interface side,
\[
        m_x=1 \quad\hbox{initially},\qquad
        m_x\gg1 \quad\hbox{on generated or supplied large-resolution branches}.
\]
Thus Peter--Weyl exhaustions, inductive local-net families, and large-window limits below are auxiliary or interface constructions unless explicitly identified with the protected microscopic carrier.  They do not replace the fixed compact momentum datum \(m_p=1\).

For unequal labelled resolutions, the level labels are retained marks, so use
Paper~II-A, Proposition~\ref{p2a:prop:typed-born-branch-transport}, first to
type the physical map to the separately certified reciprocal presentation.
Then Paper~I-C, Theorem~\ref{p1b:thm:unequal-exchange}, transports the
projector and operator together and gives the least invariant comparison
carrier.  Its compressed exchange need not square to the identity on an
asymmetric window.  This comparison preserves the preceding \(m_p,m_x\)
roles and does not turn unequal levels into an internal symmetry.

\begin{corollary}[Exact homogeneous reciprocal-position resolution import]
\label{p4:cor:exact-reciprocal-position-resolution}
Import Theorem~\ref{p3:thm:exact-spectral-resolution-staircase} and define
\[
 \ell_F:=\Lambda_F^{-1}.
\]
On a nondecreasing homogeneous
\(\mathsf{SF}_{\Lambda_F}\) branch with \(R_x=a\), the exact resolution is
\[
 \boxed{
 m_x
 =\max\!\left\{1,
 \left\lfloor
 \sqrt{1+(R_x/\ell_F)^2}-1
 \right\rfloor\right\}.}
\]
For \(m_x\ge2\), the exact inverse interval is
\[
 \boxed{
 \ell_F\sqrt{m_x(m_x+2)}
 \le R_x
 <\ell_F\sqrt{(m_x+1)(m_x+3)}.}
\]
Writing \(r=R_x/\ell_F\), the pre-floor cutoff degree satisfies
\[
 q(r):=\sqrt{1+r^2}-1=r-1+O(r^{-1}),
 \qquad
 m_x=\lfloor q(r)\rfloor=r+O(1).
\]
Thus the formerly used linear cell formula is neither exact nor an
asymptotically equal integer identity.  The frontier scale \(\ell_F\) is not
the smooth action length \(\Lambda_{\rm cut}^{-1}\) unless
\((\mathrm{ID}_{F/A})\) is imposed, and it is not the Born length
\(\ell_{\rm B}\) unless a further calibration identifies them.  Outside
\(\mathsf{SF}_{\Lambda_F}\), a finite position trajectory remains
a separately supplied/generated interface.
\end{corollary}

\begin{proof}
The first formula is the nondecreasing specialization \(a_*=a\) of the
Paper-III-F theorem.  For \(m_x\ge2\), the floor condition
\[
 m_x\le\sqrt{1+r^2}-1<m_x+1
\]
is equivalent, after adding one, squaring, and restoring
\(\ell_F\), to the displayed interval.  The asymptotic expansion is
the binomial expansion of \(\sqrt{1+r^2}\); applying the floor introduces a
bounded error.
\end{proof}



\section{The generated finite test algebra and its physical role}
\label{p4a:sec:v238-a2}
\noindent\textit{Source and claim boundary.} The native $m_p=1$ algebra is a finite direct sum of matrix algebras. Its state and test roles are constitutive consequences, not an independent textbook Hilbert-space or full-matrix-algebra postulate.

\subsection{Finite Core and Structural Readout Criterion}
\label{p4:sec:part1-scope}

\paragraph{Purpose.}
Part~I closes the finite state/effect/operator calculus on a fixed member
\(\mathcal H_m\) of the mathematical truncation family, or on a finite level
separately realized by a named interface.  The
standing GP1--GP3 readout convention separates the native finite core from
representation and interface structures: GP1 bounds constituted carriers and
operations, GP2 enforces chronological geometric origin and relation
exactness, and GP3 retains the global minimal faithful realization.  Apparatus
algebras, sector restrictions, scattering windows, thermodynamic limits, and
continuum QFT are handled by their own theorem domains after the core calculus
is closed.

\begin{definition}[Prefix-structural and role-conditioned readout criterion]
\label{p4:def:structural-criterion}
A mathematical datum is \emph{prefix-structural} when, at fixed finite
truncation level \(m\), it is intrinsic to, generated by, or theorem-fixed
from \(\mathfrak F_{\le3}^{\rm core}\) under the distinct GP1, GP2, and GP3
roles just stated.  A finite quantum representation is
\emph{interface-conditioned structural} when it is theorem-fixed from that
mathematical datum relative to its displayed residual set \(\mathcal R_T\).
Those interfaces select actual preparation/readout tokens or certify
independent occurrences, maps, and realized dynamics; the generated test
algebra already supplies event/probability-role meaning, and the interfaces do
not supply the finite algebraic forms proved below.  A datum is
\emph{representation/interface} when it is a
representation choice, accessible-sector restriction, apparatus description,
regulator representative, reciprocal-position window, continuum
presentation, or downstream sector identification whose alteration leaves
the closed finite state/effect calculus intact.
\end{definition}

\begin{theorem}[Boundary-conditioned specialization of the finite readout atlas]
\label{p4:thm:boundary-conditioned-readout-specialization}
Let
\[
 \mathcal A=\bigoplus_\lambda M_{d_\lambda}(\mathbb C),
 \qquad
 \rho=\bigoplus_\lambda\rho_\lambda\ge0,
 \qquad
 \sum_\lambda\operatorname{Tr}\rho_\lambda=1.
\]
Let \(\Phi:\mathcal A\to\mathcal A\) be an occurring completely positive
trace-preserving process and let \(0\le E\le1\) be a generated effect.  The
boundary-conditioned diagram
\[
 \Rbdry\ni
 I\xrightarrow{\rho}\mathcal A
 \xrightarrow{\Phi}\mathcal A
 \xrightarrow{e_E}I,
 \qquad e_E(a):=\operatorname{Tr}_{\mathcal A}(Ea),
\]
has the unique finite scalar evaluation
\[
 \boxed{p(E\mid\rho,\Phi)
 =\operatorname{Tr}_{\mathcal A}\!\bigl(E\Phi(\rho)\bigr)\in[0,1].}
\]
Conversely, every normalized affine effect valuation at \(\Phi=\mathrm{id}\)
is represented by a unique block density \(\rho\).  The actual process,
preparation, performed test, and outcome/readout remain the displayed
\(\mathsf{PROC}\) and token-level \(\mathsf{PROB}\) antecedents; the physical
evaluation role and algebraic form are theorem-level.
\end{theorem}

\begin{proof}
Paper~II-A Theorem~\ref{p2:thm:loop-boundary-readout-exhaustion} places a
prepared, processed, and tested finite diagram in \(\Rbdry\).
Theorem~\ref{p4:thm:density-matrix} below proves that every normalized affine
effect valuation has a unique positive unit-trace block density and is
evaluated by the block trace pairing.  Theorems~\ref{p4:thm:cp-forced}
and~\ref{p4:thm:finite-choi-kraus-instruments} prove that the occurring
same-type-stable finite process is CP and has a trace-preserving Kraus form,
so \(\Phi(\rho)\ge0\) and
\(\operatorname{Tr}\Phi(\rho)=1\).  Since \(0\le E\le1\),
\[
 0\le
 \operatorname{Tr}\!\bigl(E\Phi(\rho)\bigr)
 \le\operatorname{Tr}\Phi(\rho)=1.
\]
Trace-pairing nondegeneracy gives uniqueness.  None of these steps erases the
preparation or effect boundary.
\end{proof}

\begin{corollary}[Pure amplitude and Born square]
\label{p4:cor:pure-amplitude-born-square}
For unit vectors \(\psi,\phi\in H\) and a finite contraction
\(T:H\to H\), define
\[
 A_{\phi\leftarrow\psi}(T):=\langle\phi,T\psi\rangle .
\]
With
\(\rho_\psi=\lvert\psi\rangle\langle\psi\rvert\),
\(E_\phi=\lvert\phi\rangle\langle\phi\rvert\), and
\(\Phi_T(\rho)=T\rho T^*\),
\[
\boxed{
\begin{aligned}
\operatorname{Tr}\!\left(E_\phi\Phi_T(\rho_\psi)\right)
&=\operatorname{Tr}\!\left(
\lvert\phi\rangle\langle\phi\rvert
T\lvert\psi\rangle\langle\psi\rvert T^*\right)\\
&=\langle\phi,T\psi\rangle
  \langle\psi,T^*\phi\rangle
=\lvert\langle\phi,T\psi\rangle\rvert^2 .
\end{aligned}}
\]
Thus the complex amplitude is a boundary line readout and the probability is
its positive state/effect specialization, not an alternative arbitrary rule.
\end{corollary}

\begin{proof}
Use the rank-one identity
\(\operatorname{Tr}(\lvert u\rangle\langle v\rvert)=\langle v,u\rangle\)
and associate the four rank-one factors in the displayed order.  The second
factor is the complex conjugate of the first.  If \(T\) is an isometry this is
a channel probability; for a contraction it is the corresponding
trace-nonincreasing instrument weight.
\end{proof}

\begin{theorem}[Finite-core admissibility on the premise-minimal interface boundary]
\label{p4:thm:no-hidden-interface-hypotheses}
The native finite quantum sector of Part~I is generated from
\(\mathfrak F_{\le3}^{\rm core}\), Paper~II-C's generated finite test algebra,
and Paper~II-A's standard full state space.  Finite functional calculus,
effects, block densities, Born representation, relation-free tensor
monoidality, affine extension, factor-local same-type Choi amplification, and direct-sum
Kraus/instrument form are theorem-level consequences.  The displayed
\(\mathsf{SUB}\) and \(\mathsf{PROC}\) labels contribute only the independent
occurrence antecedents defined above; \(\mathsf{PROB}\) is token shorthand,
and \(\mathsf{QEV}\) has no residual role-level content.  Native nontrivial
kinematic automorphism flows
are generated algebraically; \(\mathsf{DYN}\) is required when one is
identified with actual elapsed evolution or when an actual open evolution is
invoked.  The full carrier endomorphism algebra may be used as a graded field
algebra only on \(\mathsf{FENV}\).  It becomes the physical full-carrier test
algebra either when a generated carrier-separating family is present or on an
independently adjoined \(\mathsf{ENV}\) root.
Apparatus models, accessible algebras,
coarse-graining maps, laboratory efficiencies, scattering windows, continuum
refinements, thermal baths, control Hamiltonians, detector responses, and
downstream sector assignments enter only through their stated theorem
domains.
\end{theorem}

\begin{proof}
The generated finite algebra is fixed upstream before an apparatus,
environment, or continuum presentation is selected.  Paper~II-A's canonical
state--test theorem proves the positive functional, effect, and block-density
structure.  Theorems below prove the Born representation, tensor products,
partial evaluations, complete positivity, instruments, and, on the
\(\mathsf{DYN}\) domain, the representation of supplied closed generators.
Paper~II-C's coupling theorem determines exactly when generated bridges enlarge
the operator algebra to a full carrier envelope; Paper~II-D's central grading
and the \(\mathsf{FENV}/\mathsf{ENV}\) distinction determine whether that
envelope is a graded field algebra or a parity-breaking event algebra.  A datum
outside this closure selects a basis, sector, channel, profile, response
coefficient, Hamiltonian, reservoir, or limiting system and therefore belongs
to the corresponding representation or interface theorem domain.  Hence the
finite calculus is closed without concealing the semantics that make it a
quantum calculus.
\end{proof}



\paragraph*{Central hypothesis ledger}
\label{p4:sec:central-proof-discharge-ledger}

\begin{proposition}[Paper-IV centralized hypothesis rule]
\label{p4:prop:central-proof-discharge}
Membership in a finite carrier or in a named interface domain is a type check,
not a proof of that domain's analytic or dynamical hypotheses.  Every
Paper-IV result must separately verify each property it uses, including
positivity, normalization, complete positivity, semigroup composition,
spectral isolation, KMS analyticity, reflection or graded-reflection
positivity, coercivity, locality/leakage estimates, convergence, and
reservoir or continuum compatibility.  A result may instead cite an earlier
theorem that proves the exact property on the same domain.
\end{proposition}

\begin{proof}
Finiteness makes all everywhere-defined linear maps bounded and reduces many
representation questions to matrix algebra, but it does not imply positivity
of a chosen functional, complete positivity of a map, the semigroup law,
reflection positivity, convergence of an integral, locality of a Hamiltonian,
or convergence of a sequence of finite models.  Explicit finite witnesses
already separate these properties:
\begin{enumerate}[label=(\alph*), leftmargin=2.2em]
\item \(-\mathbf1\in M_d(\CC)\) is finite but not positive, and a positive
      functional need not be normalized;
\item the transpose map
      \(\tau:M_2(\CC)\to M_2(\CC)\) is positive, unital, and trace preserving
      but is not completely positive, because
      \((\operatorname{id}\otimes\tau)(\ket{\Omega}\bra{\Omega})\) is the flip
      operator and has eigenvalue \(-1\);
\item for a non-scalar self-adjoint \(H\), the finite automorphisms
      \(\alpha_t=\operatorname{Ad}(e^{it^2H})\) do not satisfy
      \(\alpha_{t+s}=\alpha_t\alpha_s\) in general;
\item the scalar Euclidean kernel \(C(\tau)=-e^{-|\tau|}\) fails reflection
      positivity, since its one-point reflected quadratic form at any
      \(\tau>0\) is \(C(2\tau)<0\);
\item for the trivial dynamics on \(M_2(\CC)\), the vector state
      \(\omega(a)=\langle e_1,ae_1\rangle\) is not KMS: with
      \(A=\ket{e_1}\bra{e_2}\) and \(B=A^*\), one has
      \(\omega(AB)=1\ne0=\omega(BA)\);
\item the bounded finite sequence \(X_m=(-1)^m\mathbf1\) does not converge.
\end{enumerate}
Likewise, a matrix algebra carries no support or locality relation until that
relation is supplied, and a finite-dimensional integral can diverge when its
weight is noncoercive.  Therefore finite membership licenses the relevant
matrix operations but cannot discharge a hypothesis that is an additional
inequality, composition law, compatibility condition, or limit statement.
This is precisely the separation between finite type correctness and theorem
correctness.
\end{proof}

\begin{definition}[Paper-IV exhaustion and readout ledger]
\label{p4:def:exhaustion-readout-ledger}
The centralized Paper-IV ledger has three rows.
\begin{enumerate}[label=(\roman*), leftmargin=2em]
\item \emph{Finite-core scope:} states, effects, POVMs, blockwise Born
      evaluation, relation-free composites, certified channels, reductions,
      and represented supplied dynamics comprise Paper~IV's structural finite
      calculus on any generated finite direct-sum algebra.  The legacy role
      labels record only empirical meaning and actual-domain existence.  This
      is an ownership and scope statement, not a claim that an arbitrary map,
      dynamics, locality, KMS state, or limit follows from finite membership.
\item \emph{Representation exhaustion:} rays, wavefunctions, bases, projective selective-readout representations, path sums, scattering windows, local nets, thermal models, and finite-QFT modules are admitted as representation or interface realizations of the finite core.
\item \emph{Readout chronology:} detector efficiencies, sector assignments, SM/QCD/cosmological quantities, and continuum limits use the finite quantum calculus after it is closed.
\end{enumerate}
\end{definition}

\paragraph{Standing algebra convention.}
Part~I starts with the smallest unital \(*\)-algebra generated by the
constituted finite tests.  At the physical level \(m_p=1\) this is the exact
\(43\)-dimensional algebra \(\mathcal A_{\rm kin}\) of Paper~II-C.  The full
endomorphism algebra \(\End(\mathcal H_m)\) remains a mathematical presentation
algebra.  It becomes a graded field algebra only on the
\(\mathsf{FENV}\) domain and a physical event algebra only on the
\(\mathsf{ENV}\) extension domain.  Throughout this paper, a displayed use of
\(\mathcal A_m^{\rm env}\) therefore carries the appropriate field or event
domain unless the sentence explicitly says that only mathematical
presentation availability is meant.
All finite state, effect, tensor, and process representation theorems are
stated for a general finite direct sum of matrix algebras and do not require
\(\mathsf{FENV}\) or \(\mathsf{ENV}\).

\begin{definition}[Generated test algebra, presentation envelope, and accessible algebra]
\label{p4:def:resolved-vs-accessible-algebra}
At a finite carrier level \(m\), let \(\mathcal T_m\) be the constituted
self-adjoint test family and put
\[
 \mathcal A_m^{\rm gen}
 :=\operatorname{alg}_{\CC}^{*}(\mathbf1,\mathcal T_m)
 \subseteq\End(\mathcal H_m).
\]
This is the \emph{generated test algebra}.  Abstractly it has finite
Wedderburn form \(\bigoplus_\lambda M_{d_\lambda}(\CC)\).  At \(m_p=1\),
\[
 \mathcal A_1^{\rm gen}=\mathcal A_{\rm kin}
 \cong\CC\oplus M_2\oplus M_2\oplus M_4\oplus M_3\oplus M_3.
\]
The \emph{presentation envelope}
\[
 \mathcal A_m^{\rm env}:=\End_{\CC}(\mathcal H_m)
\]
exists mathematically on the same carrier.  It is a physical test algebra only
after \(\mathsf{ENV}\) is certified; \(\mathsf{FENV}\) may instead type it as
a graded field envelope without making its odd part observable.  An
\emph{accessible algebra} is any
generated or theorem-fixed unital \(*\)-subalgebra used by a stated regime or
apparatus.  It may equal \(\mathcal A_m^{\rm gen}\), a generated enlargement,
or a restriction selected by an actual charge, pointer, or superselection
interface.  Accordingly, when such a proper algebra is not already generated
or theorem-fixed, it belongs to the residual accessibility realization family
\[
 \mathfrak A_{\rm ACC}
 :=\{\mathcal A_{\rm acc}\subseteq\mathcal A_m^{\rm env}:
 \mathcal A_{\rm acc}\text{ is a physically supplied admissible unital
 }*\text{-subalgebra}\}.
\]
Selecting a member is then an \(\mathsf{ACC}\) interface; if the algebra is
generated by charge/test/apparatus data, its status is instead theorem-level.
Direct sums are therefore native generated algebras when their
central projectors are themselves generated; they are not automatically
post hoc restrictions of a prior full envelope.
\end{definition}


\begin{lemma}[Simultaneous presentation covariance of the imported envelope]
\label{p4:lem:unitary-naturality-resolved-events}
On the mathematical presentation domain, and on the physical
\(\mathsf{ENV}\) domain when certified, the matrix envelope of
Definition~\ref{p3:def:block-algebra-envelope}, every unitary change of
finite presentation acts by simultaneous conjugation:
\[
u\mathcal A_m^{\mathrm{env}} u^{-1}=\mathcal A_m^{\mathrm{env}}
\qquad
\text{for all }u\in U(\mathcal H_m).
\]
The carrier action, Dirac/spatial operators, projectors, states, and effects
must all be conjugated together.  A generic \(u\) need not preserve the
distinguished \(\Spin(4)\), Casimir, coordinate, Toeplitz, grading, or
protected-block data and is not thereby an upstream geometric or dynamical
symmetry.
\end{lemma}

\begin{proof}
For a finite Hilbert space,
\(u\End(\mathcal H_m)u^{-1}=\End(\mathcal H_m)\) algebraically.  Simultaneous
conjugation preserves every equation among the constituted tuple because it
is an isomorphism of presentations.  Holding that tuple fixed while
conjugating only one component is instead a symmetry claim and is valid only
for the actual stabilizer of the fixed data.
\end{proof}


\begin{lemma}[Subalgebras require selectors]
\label{p4:lem:proper-subalgebras-require-selectors}
\leavevmode\par\noindent
Let \(\mathcal H\) be a finite-dimensional Hilbert carrier and let
\(\mathcal B\subseteq\End(\mathcal H)\) be a unital \(*\)-subalgebra.  If
\(\mathcal B\neq\CC\Id\) and \(\mathcal B\neq\End(\mathcal H)\), then the
inclusion of \(\mathcal B\) determines nontrivial selector data: equivalently,
after a unitary identification one has a finite Wedderburn decomposition
\[
\mathcal H\cong\bigoplus_\alpha (V_\alpha\otimes M_\alpha),
\qquad
\mathcal B\cong\bigoplus_\alpha
\bigl(\End(V_\alpha)\otimes\Id_{M_\alpha}\bigr),
\]
so \(\mathcal B\) carries central sector projectors, multiplicity and commutant
data, a preferred decomposition, or a preferred tensor-factor structure.
Therefore a proper finite subalgebra is an accessible or sector algebra inside the
maximal unreduced endomorphism envelope unless its selector data is intrinsic to,
generated by, or theorem-fixed by the constituted carrier.
\end{lemma}

\begin{proof}
The finite-dimensional structure theorem for unital \(*\)-subalgebras of a full
matrix algebra gives the displayed decomposition.  If there is more than one summand,
the central projections onto the summands are nontrivial sector selectors.  If there
is one summand but \(M_\alpha\) is nontrivial, the commutant
\(\Id_{V_\alpha}\otimes\End(M_\alpha)\) is a multiplicity/tensor-factor selector.  If
\(V_\alpha\) is a proper factor, the chosen factorization itself is extra structure.
The only cases in which no such selector remains are the scalar algebra
\(\CC\Id\) and the full algebra \(\End(\mathcal H)\).  By exact GP2, projector families,
sector labels, commutants, preferred bases, tensor factors, apparatus decompositions,
and superselection data cannot be appended as structural selectors.  Hence every
proper non-scalar algebra is accessible or interface-level unless its selector is
already supplied by the constituted carrier.
\end{proof}

\begin{remark}[Generated sector data are not inserted selectors]
Paper~II-C's central projectors in \(\mathcal A_{\rm kin}\) are polynomials in
the constituted chiral Casimirs.  Their direct-sum decomposition is therefore
generated structure, not a later accessible-sector choice.  The preceding
lemma diagnoses provenance; it does not prefer \(\End(\mathcal H)\) over a
proper algebra whose Wedderburn data have already been derived.
\end{remark}

\begin{definition}[Carrier-resolving full-envelope extension]
\label{p4:def:carrier-resolving-event-role}
Let \(\mathcal B\subseteq\End(\mathcal H_m)\) be the unital \(*\)-algebra
assigned to a proposed \(\mathsf{ENV}\) extension.  The extension is
\emph{carrier resolving} when every two distinct density operators on the
resolved carrier can be separated by a self-adjoint test in \(\mathcal B\):
\[
\rho\ne\sigma
\quad\Longrightarrow\quad
\exists\,b=b^*\in\mathcal B:
\Tr(\rho b)\ne\Tr(\sigma b).
\]
Carrier resolution alone forces the full envelope by
Lemma~\ref{p4:lem:event-separation-full-envelope}.  It is stronger than the
native requirement that states be separated as functionals on the generated
algebra.  Two full-carrier density matrices that differ only in the
Hilbert--Schmidt annihilator of \(\mathcal A_m^{\rm gen}\) represent the same
state on that algebra and need not be separated.  Thus carrier resolution is
not a prerequisite for finite quantum kinematics.  It is a theorem condition
on a generated test family when one exists, and a genuine extension only when
that separating family must be independently adjoined.
\end{definition}

\begin{lemma}[Finite event separation forces the full envelope]
\label{p4:lem:event-separation-full-envelope}
If a unital \(*\)-subalgebra
\(\mathcal B\subseteq\End(\mathcal H_m)\) is carrier resolving in the sense of
Definition~\ref{p4:def:carrier-resolving-event-role}, then
\[
\mathcal B=\End(\mathcal H_m).
\]
\end{lemma}

\begin{proof}
Put \(N=\dim\mathcal H_m\).  If the real self-adjoint part
\(\mathcal B_{\rm sa}\) were a proper subspace of
\(\End(\mathcal H_m)_{\rm sa}\), finite-dimensional Hilbert--Schmidt duality
would give a nonzero \(h=h^*\) such that
\[
\Tr(hb)=0\qquad(b=b^*\in\mathcal B).
\]
Because \(\mathbf1\in\mathcal B\), one also has \(\Tr h=0\).  Choose
\[
0<\varepsilon<\frac{1}{N\|h\|}.
\]
Then
\[
\rho_\pm:=\frac{\mathbf1}{N}\pm\varepsilon h
\]
are distinct positive trace-one matrices, while
\[
\Tr(\rho_+b)-\Tr(\rho_-b)=2\varepsilon\Tr(hb)=0
\qquad(b=b^*\in\mathcal B).
\]
Thus \(\mathcal B\) would fail to separate two resolved states, contrary to
the hypothesis.  Hence
\(\mathcal B_{\rm sa}=\End(\mathcal H_m)_{\rm sa}\), and complexification
gives \(\mathcal B=\End(\mathcal H_m)\).
\end{proof}


\begin{theorem}[Generated finite test algebra and carrier-resolving event extension]
\label{p4:thm:full-finite-observable-algebra}
At the physical level \(m_p=1\), the prefix generates
\[
 \boxed{
 \mathcal A_1^{\rm gen}=\mathcal A_{\rm kin}
 \cong\CC\oplus M_2(\CC)\oplus M_2(\CC)\oplus M_4(\CC)
       \oplus M_3(\CC)\oplus M_3(\CC),
 \qquad \dim_{\CC}\mathcal A_{\rm kin}=43.}
\]
Its self-adjoint part, projections, and effects form the native finite
physical state--test/event calculus by Paper~II-A's generated-role theorem;
no \(\mathsf{QEV}\) assignment remains.

If a generated test algebra \(\mathcal B\subseteq\End(\mathcal H_1)\) is
carrier resolving in the sense of
Definition~\ref{p4:def:carrier-resolving-event-role}, then its event algebra
is forced without an additional ENV witness to be
\[
 \boxed{\mathcal A_1^{\rm env}
 =\End(\mathcal H_1)\cong M_{15}(\CC).}
\]
If no such family is generated, an independently adjoined
\(\mathsf{ENV}\) root may impose it.  On the level-one grading domain that
root is a parity-breaking full-event extension and is not a premise of the
contained core.
\end{theorem}


\begin{proof}
Paper~II-C, Theorem~\ref{p2c:thm:exact-level-one-kinematic-algebra}, proves
the displayed block equality by recovering the two chiral generators from
\(\widehat J_i,\widehat P_i\), separating the six sectors with a polynomial
in their Casimirs, and applying irreducibility on each central block.  Its
Corollary~\ref{p2c:cor:contained-quantum-kinematic-existence} uses
noncommuting block effects to establish the contained finite
quantum-kinematic system.

Lemma~\ref{p4:lem:event-separation-full-envelope} proves that an explicitly
carrier-resolving extension cannot be proper.  Theorem~\ref{p4:thm:test-graph-env-discharge}
below gives the complementary generated route: a connected graph of
constituted physical tests first forces the full matrix algebra in the test
slot, after which carrier resolution is automatic.  Paper~III-B's
\(\End(\mathcal H_1)\) remains mathematically available throughout; neither
mathematical presentation availability nor a field-only bridge licenses it as
the full physical event algebra.
\end{proof}

\begin{corollary}[Observable symmetry and the amount of resolved state data]
\label{p4:cor:operational-symmetry-state-dimensions}
On the respective domains where the displayed algebras are available as
tests, their effective conjugation symmetries and normalized state-space
affine dimensions are
\[
\begin{array}{c|c|c}
\text{test algebra}&\text{effective symmetry}&\text{affine dimension}\\\hline
\mathcal A_{\rm kin}&SO(3)\times SO(3)&42\\
\mathcal A_X=M_{11}(\CC)\oplus M_4(\CC)&SO(4)&136\\
M_{15}(\CC)\ \text{if generated}&\Spin(4)&224
\end{array}
\]
Thus resolving a central geometric/protected population is not full state
tomography.  In particular, native block-diagonal Hamiltonians and CPTP
channels whose Kraus operators lie in \(\mathcal A_X\) conserve that
population and do not force \(p_g=11/15\).
\end{corollary}

\begin{proof}
Paper~II-C, Theorem~\ref{p2c:thm:operational-symmetry-filtration}, gives
the conjugation kernels and their effective quotients.  A finite direct sum
\(\bigoplus_iM_{d_i}(\CC)\) has normalized state-space affine dimension
\(\sum_i d_i^2-1\), giving \(43-1\), \(137-1\), and \(225-1\).
For a central projection \(P_g\) and block-diagonal Kraus operators,
\[
 \Tr\!\left(P_g\sum_\alpha K_\alpha\rho K_\alpha^*\right)
 =\Tr\!\left(\rho P_g\sum_\alpha K_\alpha^*K_\alpha\right)
 =\Tr(\rho P_g).
\]
The Hamiltonian statement follows from \([H,P_g]=0\).  Population is
therefore conserved on the domain of
Paper~II-C, Corollary~\ref{p2c:cor:native-central-population-conservation};
mere preservation of the dephased subspace is a weaker channel condition
and need not conserve its central probabilities.  The generated exchange
process considered later supplies additional dynamics when it occurs.
\end{proof}

\begin{theorem}[Test-typed coupling graph and automatic ENV discharge]
\label{p4:thm:test-graph-env-discharge}
Let
\[
 \mathcal H=\bigoplus_{\lambda}V_\lambda,
 \qquad
 \mathcal A_0=\bigoplus_{\lambda}\End(V_\lambda)
\]
be a generated finite test algebra written in its generated central blocks,
and let \(X_{\rm test}\) be a constituted family of additional self-adjoint
physical tests already generated in the test/event role.  Form the nonzero
sector-coupling graph \(G_{X_{\rm test}}\) as in Paper~II-C.  If this graph is
connected, then
\[
 \boxed{
 \operatorname{alg}_{\CC}^{*}(\mathcal A_0,X_{\rm test})
 =\End(\mathcal H),
 }
\]
and this full matrix algebra is already the generated physical test/event
algebra.  It separates every two distinct full-carrier density operators, so
no independent \(\mathsf{ENV}\) root is required.

At level one, if a generated physical-test algebra has already reached
\[
 M_{11}(\CC)\oplus M_4(\CC)
\]
and one constituted nonzero physical-test bridge connects the two central
blocks, then
\[
 \boxed{C^*(M_{11}\oplus M_4,T_{\rm test})=M_{15}(\CC)}
\]
as the generated event algebra.  A merely field-typed odd bridge satisfies the
separate \(\mathsf{FENV}\) theorem but does not, without test/event typing,
discharge \(\mathsf{ENV}\).
\end{theorem}

\begin{proof}
Paper~II-C, Theorem~\ref{p2c:thm:sector-coupling-graph-closure}, gives
\[
 \operatorname{alg}_{\CC}^{*}(\mathcal A_0,X_{\rm test})
 =\bigoplus_{C\in\pi_0(G_{X_{\rm test}})}\End(\mathcal H_C).
\]
Connectedness leaves one component and therefore the full matrix algebra.
Because \(X_{\rm test}\) entered through the constituted physical test slot,
Paper~II-A's generated-role theorem preserves that typing under the licensed
finite \(*\)-algebra closure; this is the step that would fail for a
field-only bridge.  Finally, if \(\rho\ne\sigma\) are two density operators,
then \(h:=\rho-\sigma\) is a nonzero self-adjoint element of
\(\End(\mathcal H)\) and
\[
 \Tr((\rho-\sigma)h)=\Tr(h^2)>0.
\]
Thus the generated full matrix test algebra is carrier resolving.  The
level-one statement is the two-component special case of the same graph
closure theorem.
\end{proof}

\begin{corollary}[One carrier-spanning effect discharges \(\mathsf{ENV}\)]
\label{p4:cor:one-effect-env-discharge}
In the setting of
Theorem~\ref{p4:thm:test-graph-env-discharge}, choose
\[
 v=\bigoplus_\lambda v_\lambda,
 \qquad v_\lambda\in V_\lambda\setminus\{0\},
 \qquad
 E_v:=\frac{\lvert v\rangle\langle v\rvert}{\langle v,v\rangle}.
\]
If \(E_v\) is already constituted in the physical test/effect role, then
\[
 \boxed{C^*(\mathcal A_0,E_v)=\End(\mathcal H),}
\]
so this single generated effect discharges \(\mathsf{ENV}\).  The
carrier-spanning condition \(v_\lambda\ne0\) for every generated central
block is open and dense in \(\mathcal H\); mathematical existence alone does
not constitute the effect.
\end{corollary}

\begin{proof}
For every pair of central blocks,
\[
 P_\lambda E_vP_\mu
 =\frac{\lvert v_\lambda\rangle\langle v_\mu\rvert}
 {\langle v,v\rangle}\ne0.
\]
Given \(x\in V_\lambda\) and \(y\in V_\mu\), choose
\(L_\lambda\in\End(V_\lambda)\) and
\(R_\mu\in\End(V_\mu)\) with
\(L_\lambda v_\lambda=x\) and \(R_\mu^*v_\mu=y\).  Then
\(L_\lambda(P_\lambda E_vP_\mu)R_\mu\) is the corresponding nonzero
scalar multiple of \(\lvert x\rangle\langle y\rvert\).  Hence every
rectangular matrix unit between every two blocks lies in the generated
\(*\)-algebra, which is therefore \(\End(\mathcal H)\).  The set on which
some \(v_\lambda=0\) is a finite union of proper closed subspaces, proving
the open-dense statement.
\end{proof}

\begin{definition}[Level-one graded field envelope]
\label{p4:def:graded-field-envelope}
Import \(\Gamma_c\) from Paper~II-D and put
\[
 \Theta(T):=\Gamma_cT\Gamma_c.
\]
Let \(\mathcal A_{\rm ev}\subseteq
\End(\mathcal H_1)^{\bar0}\) be the generated even algebra and let
\(Y_{\rm odd}\subseteq\End(\mathcal H_1)^{\bar1}\) be an actually generated
odd field family.  Define
\[
 \mathcal F:=C^*(\mathcal A_{\rm ev},Y_{\rm odd}).
\]
\end{definition}

\begin{theorem}[Odd-bridge criterion for the graded field envelope]
\label{p4:thm:odd-bridge-field-envelope}
Under the sector-block hypotheses of
Theorem~\ref{p2c:thm:sector-coupling-graph-closure}, the six-sector algebra
satisfies \(\mathcal F=M_{15}(\CC)\) exactly when the total nonzero coupling
graph is connected.  In particular, if the even coordinates have already
generated
\[
 \mathcal A_{\rm ev}=M_{11}(\CC)\oplus M_4(\CC),
\]
then any one nonzero central-odd bridge gives
\[
 \boxed{\mathcal F=M_{15}(\CC).}
\]
Its parity-preserving observable fixed point is
\[
 \boxed{\mathcal A_{\rm obs}=\mathcal F^\Theta
 =M_{11}(\CC)\oplus M_4(\CC).}
\]
Declaring all self-adjoint elements of \(\mathcal F\) observable is the
separate \(\mathsf{ENV}\) full-event extension.
\end{theorem}

\begin{proof}
Theorem~\ref{p2c:thm:sector-coupling-graph-closure} states that the
\(*\)-algebra generated from the six full sector blocks and any generated
operator family is the direct sum of full matrix algebras over the connected
components of its nonzero-coupling graph.  This proves the first equivalence.
If the even closure is \(M_{11}\oplus M_4\), any nonzero odd block connects
the two already internally connected parity components; the same theorem
therefore gives \(M_{15}\).

Finally, relative to
\(\mathcal H_1=\mathcal H_+\oplus\mathcal H_-\),
\(\Gamma_c=\operatorname{diag}(I_{11},-I_4)\).  Its conjugation fixes exactly
the two diagonal blocks and negates the rectangular blocks, so the fixed
point of \(M_{15}\) is \(M_{11}\oplus M_4\).  Event meaning for the odd
self-adjoint elements is not supplied by this algebraic closure.
\end{proof}


\begin{remark}[Invariant odd no-go and covariant bridge frontier]
\label{p4:rem:covariant-odd-bridge-frontier}
Paper~II-D, Lemma~\ref{p2:lem:no-invariant-central-odd-bridge}, proves that
every Spin(4)-invariant central-odd endomorphism is zero, while its generated
coordinate/functional-calculus grammar is central even.  Consequently the
current invariant/even closure cannot itself discharge the
\(\mathsf{FENV}\) occurrence.  A geometrically lawful later route would be a
generated covariant family
\[
 T(gx)=\rho(g)T(x)\rho(g)^{-1},
 \qquad T(x)\in\operatorname{Hom}(\mathcal H_-,\mathcal H_+)
 \oplus\operatorname{Hom}(\mathcal H_+,\mathcal H_-),
\]
for example from an independently constituted nonparallel spinorial
source or other actual central-odd field input.  The canonical correlated
Born connection and its holonomy alone cannot furnish a nonzero parallel
odd section (Theorem~\ref{p2b:thm:v264-no-parallel-odd}).
If the total block-coupling graph is connected,
Theorem~\ref{p4:thm:odd-bridge-field-envelope} derives the full field
envelope.  No such family or field occurrence is asserted by the present
finite core.
\end{remark}


\begin{remark}[Status of the generated core, even envelope, and field envelope]
\label{p4:rem:finite-resolution-selection-status}
The direct-sum algebra is not an accessible restriction imposed after a full
physical envelope; its sector projectors are already generated Casimir
polynomials.  States are positive normalized functionals on this algebra.
Full-carrier density matrices differing only by off-block coherences restrict
to the same state and therefore do not create a resolution defect.  At
\(m_p=1\), the annihilator of \(\mathcal A_{\rm kin}\) inside \(M_{15}\) has
complex dimension \(225-43=182\).

Current II-C coordinates preserve the \(11\)-dimensional geometric and
\(4\)-dimensional protected components, so even their nondegenerate closure is
at most \(M_{11}\oplus M_4\), not \(M_{15}\).  Once both even components are
internally full, an actual nonzero central-odd bridge certifies
\(\mathsf{FENV}\) and generates \(M_{15}\) as a graded field envelope; its
observable fixed point remains \(M_{11}\oplus M_4\).  The separate
carrier-resolving \(\mathsf{ENV}\) interface is required to declare the odd
self-adjoint elements to be physical events.
Restricted laboratory, charge-preserving, pointer, subsystem-local, and
superselection algebras remain valid when their own selectors are generated
or theorem-fixed.
\end{remark}


\begin{remark}[Finite-dimensional scope and continuum bridge maps]
\label{p4:rem:finite-scope}
Every fixed mathematical carrier \(\mathcal H_m\) is exactly finite
dimensional, and every theorem on a supplied finite level is exact on that
level.  This does not physically realize every \(m\): the compact momentum
physical momentum carrier derived upstream is \(m_p=1\), while other values occur as
mathematical symbol levels, readout/interface resolutions, or later
role-specific realizations.  When the discussion concerns laboratory
position or continuum presentation, the growing parameter belongs to the
reciprocal-position/interface side.
Continuum wave mechanics only appears later through explicit bridge maps: finite operators may be assigned classical symbol content only on a stated symbol-readout domain, retained symbols exhaust the classical continuum only in the Peter--Weyl sense, and Euclidean kernels generate real time only through finite positive Hamiltonian reconstruction.
The finite matrix algebras themselves are therefore not being identified as literal subalgebras of \(C(S^3)\).
\end{remark}


\begin{theorem}[Three-map bridge from finite carriers to continuum presentation and real time]
\label{p4:thm:three-map-bridge}
Let \(\mathcal H_m\) be the finite \(\Spin(4)\)-equivariant carrier imported from Paper~II and let
\[
\mathcal A_m^{\mathrm{env}}=\End(\mathcal H_m)
\]
be its mathematical presentation envelope.  When an arbitrary element of this
algebra is treated as a physical event, the \(\mathsf{ENV}\) domain of
Theorem~\ref{p4:thm:full-finite-observable-algebra} is understood.  The scalar
Toeplitz operators used in (i) need only belong to their explicitly generated
Toeplitz image and do not require full-envelope closure.
Then the framework licenses the following three bridge maps, and only on the stated theorem domains.

\begin{enumerate}[label=(\roman*), leftmargin=2.3em]
\item \textbf{Finite operator to retained symbol.}
Let \(\mathcal P_m\) denote the retained scalar Peter--Weyl symbol space of
Paper~II, Definition~\ref{p2:def:retained-pw-symbols}, with Haar
\(L^2\)-inner product, and let
\(\Pi_m^{\rm sc}\) project onto the scalar diagonal Peter--Weyl carrier
\(\mathcal H_m^{\rm sc}\).  For every retained scalar symbol
\(f\in\mathcal P_m\), scalar multiplication is defined on
\(L^2(S^3)\), and its finite scalar realization is the Toeplitz compression
\begin{equation}
\label{p4:eq:three-map-toeplitz-realization}
\mathsf T_m^{\rm sc}(f)
:=\Pi_m^{\rm sc}M_f\Pi_m^{\rm sc}
\in\End(\mathcal H_m^{\rm sc}).
\end{equation}
This is the scalar-diagonal domain of the Paper-II compression theorem and
its scalar Toeplitz leakage identities.  If
\(J_m^{\rm sc}:\mathcal H_m^{\rm sc}\hookrightarrow\mathcal H_m\) is the
canonical inclusion, its supported extension
\[
J_m^{\rm sc}\mathsf T_m^{\rm sc}(f)(J_m^{\rm sc})^*
\in\End(\mathcal H_m)
\]
is a mathematical full-carrier operator with scalar-diagonal support; calling
it a physical event requires that it belong to the generated test algebra or
that the relevant \(\mathsf{ENV}\) domain be certified.  A bundle-valued or
off-diagonal continuation to all of \(\mathcal H_m\) requires a separately
fixed ambient bundle representation; Paper~II,
Theorem~\ref{p2:thm:finite-toeplitz-coordinate-uniqueness}, proves uniqueness of the
finite retraction after that ambient action is fixed, not uniqueness of the
ambient continuation.
Put
\[
\mathcal T_m^{\rm sc}:=\operatorname{im}(\mathsf T_m^{\rm sc})
\subseteq\End(\mathcal H_m^{\rm sc}) .
\]
The canonical retained-symbol readout on this Toeplitz image is
\begin{equation}
\label{p4:eq:three-map-symbol-readout}
\sigma_m^{\rm sc}:\mathcal T_m^{\rm sc}\longrightarrow
\mathcal P_m\cap(\ker\mathsf T_m^{\rm sc})^\perp,
\qquad
\sigma_m^{\rm sc}(\mathsf T_m^{\rm sc} f)=f_\perp ,
\end{equation}
where \(f=f_0+f_\perp\),
\(f_0\in\ker\mathsf T_m^{\rm sc}\), and
\(f_\perp\perp\ker\mathsf T_m^{\rm sc}\).
Thus a finite operator has classical scalar symbol content only after it lies in the Toeplitz image, or after an additional covariant-symbol/interface readout has been explicitly supplied.
There is no structural identification \(\End(\mathcal H_m)\subset C(S^3)\).

\item \textbf{Retained symbols to continuum presentation.}\par
The retained symbol spaces form an increasing Peter--Weyl exhaustion,
\begin{equation}
\label{p4:eq:three-map-pw-exhaustion}
\mathcal P_0\subseteq\mathcal P_1\subseteq\cdots,
\qquad
\mathcal P_{\mathrm{poly}}:=\bigcup_{m\ge0}\mathcal P_m,
\end{equation}
and Paper~II, Theorem~\ref{p2:thm:continuum_limit}, gives
\begin{equation}
\label{p4:eq:three-map-continuum-density}
\overline{\mathcal P_{\mathrm{poly}}}^{\|\cdot\|_\infty}=C(S^3).
\end{equation}
Smooth, distributional, \(L^2\), vector, spinor, or bundle-valued continuum
presentations require their corresponding symbol spaces and topologies; they
are not obtained by treating the full matrix algebra as a commutative algebra.
The index \(m\) in this Peter--Weyl exhaustion is an auxiliary
symbol/interface index and is not a dynamical flow \(m_p=1\to\infty\) of the
protected compact momentum carrier.

\item \textbf{Positive Euclidean spectral kernel to a finite Hamiltonian pair.}
Let a finite Euclidean two-point block have positive spectral form
\begin{equation}
\label{p4:eq:three-map-positive-euclidean-kernel}
C_E(\tau)=\sum_{a=1}^N Z_a e^{-E_a|\tau|/\hbar},
\qquad
Z_a\ge0,\qquad E_a\ge0 .
\end{equation}
After quotienting possible reflection-null vectors, let
\[
\mathcal K_C:=\operatorname{span}\{e_a:Z_a>0\}
\]
be the finite reconstruction space.  Finite spectral theory defines on this
space a positive self-adjoint Hamiltonian
\begin{equation}
\label{p4:eq:three-map-hamiltonian}
H e_a=E_a e_a,\qquad H\ge0,
\end{equation}
with Euclidean contraction \(e^{-|\tau|H/\hbar}\) and unitary evolution
\begin{equation}
\label{p4:eq:three-map-real-time}
U(\ttilde)=e^{-iH\ttilde/\hbar}.
\end{equation}
The identification of the unitary parameter with the already constituted
physical scalar-cover time uses the GP1 operational cone and Lorentzian
orientation imported from Paper~III,
Theorem~\ref{p3:thm:lorentzian-signature}.  The finite kernel reconstructs
\(H\) on \(\mathcal K_C\).  It identifies that reconstruction space with an
already constituted carrier block only when a same-space semigroup
factorization or an explicit unitary intertwiner has also been proved.  It
neither selects the Lorentzian sign nor proves a general Euclidean field
reconstruction.  The latter retains the relevant Paper-III \(({\rm WICK}_\bullet)\) domain
outside the explicitly finite reflection-positive block stated here.
\end{enumerate}
\end{theorem}


\begin{proof}
For (i), scalar multiplication by \(f\) is canonical on the scalar
Peter--Weyl carrier.  Compression by \(\Pi_m^{\rm sc}\) therefore gives the
displayed scalar Toeplitz operator.  Paper~II proves separately that once an
off-diagonal ambient action has been fixed, its equivariant finite return is
the orthogonal compression; it does not infer that ambient action from the
scalar symbol.  The map
\(\mathsf T_m^{\rm sc}:\mathcal P_m\to\End(\mathcal H_m^{\rm sc})\) is a
linear map between finite-dimensional spaces.
The Haar \(L^2\)-inner product gives the orthogonal decomposition
\[
\mathcal P_m=\ker\mathsf T_m^{\rm sc}
\oplus(\ker\mathsf T_m^{\rm sc})^\perp .
\]
If \(\mathsf T_m^{\rm sc}f=\mathsf T_m^{\rm sc}g\), then
\(f-g\in\ker\mathsf T_m^{\rm sc}\), so the orthogonal component
\(f_\perp\) is independent of the chosen representative.  Hence
\(\sigma_m^{\rm sc}\) is a well-defined inverse of
\(\mathsf T_m^{\rm sc}\) on the scalar Toeplitz image after quotienting the
kernel.
This proves the finite-operator-to-retained-symbol bridge and also shows its limitation: an arbitrary element of \(\End(\mathcal H_m)\) is a finite quantum operator, not automatically a scalar function on \(S^3\).

For (ii), Paper~II, Theorem~\ref{p2:thm:continuum_limit}, identifies the union of the retained scalar coefficient spaces with the Peter--Weyl polynomial algebra on \(SU(2)\cong S^3\).
The Peter--Weyl theorem gives uniform density in \(C(S^3)\), which is exactly \eqref{p4:eq:three-map-continuum-density}.
No operator-norm convergence of the full noncommutative matrix algebras to a commutative algebra is used.

For (iii), use one orthonormal reconstruction vector for each nonzero
spectral channel.  Each vector carries the Euclidean factor
\(e^{-E_a|\tau|/\hbar}\).  Defining \(H e_a=E_a e_a\) on
\(\mathcal K_C\) gives a positive self-adjoint finite Hamiltonian
and recovers the Euclidean contraction semigroup
\(e^{-|\tau|H/\hbar}\).
The spectral theorem then gives the unitary group
\(e^{-iH\ttilde/\hbar}\).  On the Paper-III physical-time domain its
parameter is the already oriented scalar-cover time.  A scalar kernel with
two distinct exponents, for example, requires a two-dimensional
reconstruction space even if it was recorded as one scalar matrix element.
Thus the kernel alone does not put \(H\) on the original carrier block.
This is a finite spectral reconstruction, not an assumption of a
pre-existing continuum Lorentzian QFT and not an input to the Paper-III sign
theorem.
\end{proof}


\begin{remark}[Use of the three bridge maps in later papers]
\label{p4:rem:three-map-bridge-export}
Later appearances of \(C^\infty(S^3)\), \(L^2(S^3)\), continuum fields, Euclidean kernels, real-time Hamiltonians, or Lorentzian causal language are admissible only through the relevant part of Theorem~\ref{p4:thm:three-map-bridge} or through a more specific theorem that states a stronger domain.
This prevents the three common hidden assumptions: that a finite matrix algebra is literally a function algebra, that symbol density is the same as continuum QFT, or that Euclidean notation by itself supplies physical Lorentzian time.
\end{remark}


\begin{remark}[Exceptional-core finite-cell imports]
\label{p4:rem:series-selected-finite-operator-imports}
When this paper cites the selected finite-cell corollaries of Paper~I-A or the
selected-carrier graph of Paper~II, it uses them only as already-established
finite mathematical data in the compact sector fixed upstream.  A Hurwitz or
trimer representative is used only on an actualized or conditionally
generated \(\mathsf{Supp}_{2T}\) domain.  Those imports are not used here to select
the compact sector, and they are not assigned Standard-Model, cosmological,
or phenomenological roles inside Paper~IV.  Their role here is narrower:
once a finite carrier, projector family, Clifford block, shell algebra, or
\(\mathsf{Supp}_{2T}\)-relative trimer has been supplied, the constituted
quantum roles and the finite representation theorems determine the
corresponding state/effect, statistics-selected Fock, mixing, and channel
calculus.
\end{remark}




\subsection{Framework-native ingredients relevant to quantum mechanics}
\label{p4:sec:part1-native}

The upstream ingredients used in Part~I are stated in the following subsections.



\paragraph{The finite carrier}
\label{p4:sec:part1-native-carrier}

Paper~II constructs the finite carrier (Theorem~\ref{p2:thm:truncation})
\begin{equation}
\label{p4:eq:part1-Hm}
\mathcal H_m
=
\bigoplus_{j_L+j_R\le m} V_{j_L}\otimes V_{j_R},
\qquad
\dim(\mathcal H_m)<\infty,
\end{equation}
where $V_j$ denotes the spin-$j$ irreducible representation of $\SU(2)$.
This is the structural kinematic carrier of Part~I.



\paragraph{The observable algebra}
\label{p4:sec:part1-native-algebra}

The native observable/test algebra is the generated algebra
\begin{equation}
\label{p4:eq:part1-Am}
\mathcal A_m^{\mathrm{gen}}
:=\operatorname{alg}_{\CC}^{*}(\mathbf1,\mathcal T_m)
\subseteq\End(\mathcal H_m),
\end{equation}
with involution induced by the inner product on \(\mathcal H_m\).  At the
physical level \(m_p=1\), Theorem~\ref{p4:thm:full-finite-observable-algebra}
gives
\begin{equation}
\label{p4:eq:part1-Am-full-matrix}
\mathcal A_1^{\mathrm{gen}}=\mathcal A_{\rm kin}
\cong\CC\oplus M_2\oplus M_2\oplus M_4\oplus M_3\oplus M_3.
\end{equation}
The direct-sum blocks are related, not unrelated: their central projectors are
generated polynomials in the constituted chiral Casimirs.  The full
\(\mathcal A_m^{\rm env}=\End(\mathcal H_m)\) is used as a graded field
algebra only on \(\mathsf{FENV}\), and as a physical test algebra only on the
carrier-resolving \(\mathsf{ENV}\) domain.
For example, if an exact charge or pointer decomposition
\begin{equation}
\label{p4:eq:part1-sector-Hm}
\mathcal H_m=\bigoplus_{q}\mathcal H_q,
\qquad P_q:\mathcal H_m\to\mathcal H_q,
\end{equation}
is imposed, then the corresponding charge-preserving accessible algebra is
\begin{equation}
\label{p4:eq:part1-accessible-blocks}
\mathcal A_{\mathrm{acc}}
=
\{a\in\End(\mathcal H_m):[a,P_q]=0\text{ for all }q\}
\cong
\bigoplus_q \End(\mathcal H_q).
\end{equation}
This corresponding block algebra is accessible because that new charge or
pointer family is generated.  Off-diagonal spaces belong mathematically to the
presentation envelope.  Generated central-even bridges may enlarge the
observable algebra within a parity component.  Generated central-odd bridges
belong to the \(\mathsf{FENV}\) field envelope but do not become physical
tests unless the carrier-resolving \(\mathsf{ENV}\) interface supplies that
additional event meaning.



\paragraph{Presentation covariance and constituted symmetries}
\label{p4:sec:part1-native-covariance}

In the finite setting, a unitary change of presentation acts on the complete
constituted tuple by simultaneous conjugation,
\begin{equation}
\label{p4:eq:part1-inner}
\mathcal A_m^{\mathrm{gen}}\longmapsto
u\mathcal A_m^{\mathrm{gen}}u^{-1},
\qquad a\mapsto\Ad_u(a):=uau^{-1},
\qquad u\in U(\mathcal H_m).
\end{equation}
State--effect pairs transform covariantly with it.  This is why numerical
predictions are independent of the chosen presentation.  It does not promote
every unitary to a symmetry of the fixed geometric tuple.  The latter
symmetries are its actual stabilizer, consistently with the restricted
no-preferred-multiplicity-basis theorem of Paper~II,
Theorem~\ref{p2:thm:no-preferred-basis}.



\paragraph{Subsystem gluing and no preferred decomposition}
\label{p4:sec:part1-native-gluing}

For any two constituted role occurrences whose factor tests persist and whose
frozen relation slot contains no cross-relation, GP2 supplies the universal
relation-free product.  The legacy \(\mathsf{SUB}\) label records exactly this
application certificate; it adds no composition law.
Theorems~\ref{p4:thm:independent-finite-composability} and
\ref{p4:thm:tensor-product} prove the finite algebraic consequence for direct
sums as well as full matrices: independent coexistence is multiplicative,
whereas a direct sum within one algebra records generated alternatives or
central sectors.
Because the framework supplies no preferred ordering, braiding, grading, hidden interaction, or boundary datum at this level, the corresponding minimal symmetric linear realization is the ordinary tensor product.
This is the structural origin of multipartite composition and reduced states.
Complete positivity follows separately when an actual affine process is stable
on one certified relation-free same-type copy; \(\mathsf{PROC}\) is the
compatibility name for that existence/stability antecedent.



\paragraph{Hopf phase and Reeb flow}
\label{p4:sec:part1-native-hopf}

Paper~I-A identifies a Hopf \(U(1)\) fiber and the same geometry carries a
Reeb/contact flow.  Pure-state phase redundancy, however, first follows from
projective complex Hilbert geometry.  Upstream Paper~II-D retains the full
associated-character atlas \(\{L^q:q\in\mathbb Z\}\); its principal
unit-Chern result does not by itself select an associated physical charge.
Paper~IV introduces a new comparison role: the primitive positively oriented
associated line for which one principal Hopf turn is compared with one central
ray-phase turn.  That role derives the degree-one character used below.  The
generally noncentral Hopf/Reeb action on mode sectors is kept distinct.


\begin{remark}[Geometric Reeb phase versus physical sector phases]
\label{p4:rem:reeb-geometry-no-shell-selector}
The Hopf/Reeb construction supplies a geometric phase and its associated
transport, but neither a source phase at each curvature-shell landmark nor a
physical relative CP phase follows from that fact alone.  The protected
resolution clock and its finite-sector handoffs therefore import no Reeb
standing-wave phase selector.  Any physical phase must instead be derived from
its occurring source/current operator and the appropriate connection or
transport law.
\end{remark}




\section{Positive states, effects and Born evaluation}
\label{p4a:sec:v238-a3}
\noindent\textit{Source and claim boundary.} Positive normalized functionals, block-density representatives, effects and the Born pairing are algebraic consequences on the generated test algebra. Individual preparation and detector tokens retain their physical occurrence conditions.

\subsection{Quantum identifications on the generated algebra}
\label{p4:sec:part1-selected}

\begin{table}[!htbp]
\centering
\caption[Quantum identifications from native data]{Quantum identifications from native data.}
\label{p4:tab:selected-identifications}
\small
\renewcommand{\arraystretch}{1.2}
\setlength{\tabcolsep}{4pt}
\begin{tabularx}{\textwidth}{>{\raggedright\arraybackslash}p{0.23\textwidth}
                                >{\raggedright\arraybackslash}X
                                >{\raggedright\arraybackslash}p{0.16\textwidth}}
\toprule
\textbf{Quantum identification} & \textbf{Framework object} & \textbf{Status in Part~I} \\
\midrule
Generated test & element \(a\in\mathcal A\), with \(\mathcal A\) a generated finite test algebra & prefix theorem object \\
Physical observable & self-adjoint element \(a=a^*\in\mathcal A\) & generated test-role content \\

Yes/no proposition & projection \(P\in\mathcal A\) with \(P^2=P=P^*\) & generated sharp-effect content \\

State & \(\omega\in S(\mathcal A)\), a positive normalized functional & derived finite state space \\

Block density & \(\rho=\bigoplus_\lambda\rho_\lambda\ge0\), \(\sum_\lambda\Tr\rho_\lambda=1\) & derived \\

Measurement outcome type & effect \(E\in\mathcal A\), \(0\le E\le\Id\) & generated effect content; actual outcome is a G0 token \\

General measurement role & POVM \(\{E_i\}\) with \(\sum_i E_i=\Id\) & derived effect construction; performed readout uses token-level \(\mathsf{PROB}\) shorthand \\

Physical process & Schr\"odinger-picture CPTP map; dually a unital CP map & derived form after actual factor-local map certification on one same-type product, or a generated dilation \\

Composite system & tensor product \(\mathcal A_A\otimes\mathcal A_B\) & derived after two-occurrence/no-cross-relation certification \\

Ray / projective class & central \(U(1)_{\rm ray}\) quotient of a nonzero vector; degree-one Hopf comparison & derived + comparison \\

Closed dynamics & blockwise inner automorphism flow on \(\mathcal A\), with
fixed-flow generator class \([H]\in\mathcal A_{\rm sa}/Z(\mathcal A)_{\rm sa}\) &
native kinematic examples; \(\mathsf{DYN}\) selects an actual flow, scale,
and trajectory while the inner representation and central quotient are
derived \\
\bottomrule
\end{tabularx}
\TableNotes{This ledger separates algebraic theorem content from residual
empirical or existence interfaces.  No full carrier envelope is required.}
\end{table}



\paragraph{Observables and yes/no questions}
\label{p4:sec:part1-observables}


\begin{definition}[Observable and event]
\label{p4:def:observable-event}
Let \(\mathcal A\) be the generated finite test algebra on the theorem domain
under discussion.  A physical \emph{observable} role is a self-adjoint
element \(a=a^*\), a \emph{yes/no event} is a projection
\(P^2=P=P^*\), and an \emph{effect} is an operator
\(0\le E\le\Id\), all in \(\mathcal A\).  These are generated physical
test/event roles.  The definition does not perform a test or register an
outcome token.
\end{definition}


The effect language is the finite-resolution replacement for ``measurable event with possibly coarse outcome'' and will later produce POVMs.



\begin{proposition}[Exact finite operator bridge on a generated carrier]
\label{p4:prop:exact-finite-operator-bridge}
Let an upstream licensed constructor supply a finite retained spectral carrier.
On that carrier, Paper~I-C's exact spectral evaluation and exact all-level
sampling/transport theorems give a lossless finite representation of the
retained geometric operator algebra, including every operator relation whose
hypotheses lie in the selected finite band.  This is an exact mathematical
operator transport; it does not select a preferred cutoff, a state, an
effect, or a preparation.

Paper~IV therefore takes the generated finite \(C^*\)-algebra as its
structural input and derives the state/effect/Born calculus below, rather than
postulating a second finite quantum carrier.
\end{proposition}
\begin{proof}
Apply the exact finite evaluation map and all-level transport of
Theorems~\ref{p1b:thm:full-sampling} and
\ref{p1b:thm:all-level-sampling} to the retained band.
At the native carrier, Corollary~\ref{p1b:cor:v233-eight-fibre-export}
also gives an eight-fibre realization. Since the operator is already
constituted, this presentation requires no new physical point-support
occurrence; it does not generate additional support incidence or a
performed instrument.  Exact intertwining
preserves the selected operator products, adjoints, and relations.  The
sampling theorems themselves contain no rule selecting the band or a state,
so those roles remain typed separately.
\end{proof}



\paragraph{Finite states and equivalent effect valuations}
\label{p4:sec:part1-states}

Paper~II-A, in Definition~\ref{p2:def:standard-finite-state-space} and
Theorem~\ref{p2:thm:canonical-finite-state-test-evaluation}, fixes
\(S(\mathcal A)\) as the full positive-normalized state space of the generated
finite algebra.  No specifically quantum trace or
probability postulate is added here.  Restriction of a state to effects is
automatically normalized and finitely additive; conversely those finite
effect laws reconstruct the state.  The theorem below records the equivalent
block-density presentation.  Interpreting the resulting number as an observed
probability-role evaluation is already generated.  Selecting which state is
prepared, performing the test, registering an outcome, or comparing finite
frequencies with that weight are token-level uses of the compatibility label
\(\mathsf{PROB}\).


\begin{definition}[Finite state space and effect valuation]
\label{p4:def:state}
Let
\[
 \mathcal A=\bigoplus_{\lambda=1}^{r}M_{d_\lambda}(\CC),
 \qquad
 S(\mathcal A)
 :=\{\omega:\mathcal A\to\CC:
       \omega(a^*a)\ge0,\ \omega(\Id)=1\},
\]
where linearity of \(\omega\) is understood, and put
\[
 \operatorname{Eff}(\mathcal A)
 :=\{E\in\mathcal A_{\rm sa}:0\le E\le\Id\}.
\]
For \(\omega\in S(\mathcal A)\), its effect valuation is
\(p_\omega(E):=\omega(E)\).  More generally, a \emph{finite effect valuation}
means a map \(p:\operatorname{Eff}(\mathcal A)\to[0,1]\) satisfying
\[
 p(\Id)=1,
 \qquad
 p\!\left(\sum_iE_i\right)=\sum_i p(E_i)
 \quad\text{whenever }E_i\ge0\text{ and }\sum_iE_i\le\Id.
\]
Presentation changes act covariantly on the algebra, state, and effect
together.  No fixed state is required to be invariant under every unitary.
\end{definition}

For a canonical closed-loop readout, the dimension weights are
already fixed by Corollary~\ref{p2a:cor:v233-loop-coefficients}.
This is compatible with the preparation simplex explicitly calculated
in Corollary~\ref{p2b:cor:v233-native-invariant-states}.
Neither presentation invariance nor those operation symmetries
select a general preparation.

\begin{remark}[Covariance versus fixed-state invariance]
The covariance in Definition~\ref{p4:def:state} is covariance of the
algebra--state--effect triple under finite relabeling.  It is not invariance of
a fixed state under all unitary conjugations; that stronger demand would force
a blockwise tracial restriction and erase most of the finite state space.
\end{remark}

\begin{proposition}[Symmetry-selected preparation and its direct-sum limit]
\label{p4:prop:symmetry-selected-preparation}
Let \(\rho\) be a density matrix on a simple carrier \(M_d(\CC)\).  If the
constituted preparation role requires
\[
 u\rho u^*=\rho\qquad\text{for every }u\in U(d),
\]
then \(\rho=I_d/d\); hence the maximally symmetric state is unique.  More
generally, for
\(\mathcal A=\bigoplus_\lambda M_{d_\lambda}(\CC)\), invariance under every
blockwise inner automorphism gives exactly
\[
 \rho_\lambda=\frac{p_\lambda}{d_\lambda}I_{d_\lambda},
 \qquad p_\lambda\ge0,
 \qquad \sum_\lambda p_\lambda=1.
\]
Thus block symmetry alone does not select the central weights
\((p_\lambda)_\lambda\).  A physically constituted full-carrier unitary
symmetry, or independent geometry fixing those weights, is required to reduce
that residual preparation family; occurrence of a preparation remains a
token-level statement.
\end{proposition}

\begin{proof}
Unitary invariance places \(\rho\) in the commutant of the full matrix
algebra, so \(\rho=cI_d\); normalization gives \(c=1/d\).  Blockwise
application gives \(\rho_\lambda=c_\lambda I_{d_\lambda}\).  Positivity gives
\(c_\lambda\ge0\), and with
\(p_\lambda:=\operatorname{Tr}(\rho_\lambda)=d_\lambda c_\lambda\), trace-one
normalization gives the displayed simplex of residual central weights.
\end{proof}

\begin{proposition}[Canonical evaluation and selector exclusion]
\label{p4:prop:no-insertion-finite-readout-uniqueness}
For every generated finite algebra, \(p_\omega(E)=\omega(E)\) is the canonical
state--effect evaluation and obeys the effect laws of
Definition~\ref{p4:def:state}.  A second rule depending on a chosen
decomposition of the same effect, or assigning extra weights to otherwise
unlabelled outcomes, requires an additional basis, context, apparatus, or
sector selector.  Such a rule is interface data rather than a second native
state evaluation.
\end{proposition}


\begin{proof}
Positivity and normalization of \(\omega\) give \(0\le\omega(E)\le1\), and
linearity gives finite additivity.  Naturality is
\((\eta\circ\phi)(a)=\eta(\phi(a))\) for a unital \(*\)-homomorphism \(\phi\).
A decomposition-, basis-, channel-, or apparatus-dependent modification
contains the additional selector data just named and belongs to that
interface domain.
\end{proof}


\begin{theorem}[Finite direct-sum state and effect-valuation representation]
\label{p4:thm:density-matrix}
Let \(\mathcal A=\bigoplus_\lambda M_{d_\lambda}(\CC)\).  Restriction to
effects is a bijection between \(S(\mathcal A)\) and the finite effect
valuations of Definition~\ref{p4:def:state}.  Every such state has a unique
block density
\[
 \rho=\bigoplus_\lambda\rho_\lambda,
 \qquad \rho_\lambda\ge0,
 \qquad \sum_\lambda\Tr(\rho_\lambda)=1,
\]
such that
\begin{equation}
\label{p4:eq:omega-trace}
 \omega(a)=\Tr_{\mathcal A}(\rho a)
 :=\sum_\lambda\Tr(\rho_\lambda a_\lambda),
 \qquad a\in\mathcal A,
\end{equation}
and
\[
 p(E)=\omega(E)
 =\sum_\lambda\Tr(\rho_\lambda E_\lambda)
\qquad
(0\le E\le\Id).
\]
Conversely every positive block density of total trace one defines such a
state and valuation.
\end{theorem}


\begin{proof}
\ProofPlan{We prove the finite direct-sum representation in four steps:
affine extension on effects, extension to a positive linear functional,
trace-dual density-matrix representation, and separation of
selector-dependent alternatives.}
\begin{proofsteps}

\item \textbf{Additivity, homogeneity, and monotonicity on effects.}
Finite additivity on effects first gives rational homogeneity on the interval
\(0\le E\le\Id\): subdividing an effect into \(n\) identical pieces gives
\(p(E/n)=p(E)/n\), and finite sums give \(p(qE)=q p(E)\) whenever
\(0\le qE\le\Id\) and \(q\in\mathbb Q_{\ge0}\).  Monotonicity follows from
additivity because \(0\le E\le F\le\Id\) implies
\(p(F)=p(E)+p(F-E)\ge p(E)\).  For fixed \(E\), the function
\(f_E(t):=p(tE)\) is additive on its allowed interval and monotone.  A monotone
solution of the Cauchy equation is real-linear, so
\(p(tE)=t\,p(E)\) for every real \(t\ge0\) for which \(tE\le\Id\).


\item \textbf{Extension to a positive linear functional.}
For a positive operator \(X\ge0\), choose \(r>0\) with \(0\le X/r\le\Id\) and set
\[
\omega_+(X):=r\,p(X/r).
\]
The preceding homogeneity shows that this definition is independent of the chosen
\(r\).  Additivity of \(p\) gives
\(\omega_+(X+Y)=\omega_+(X)+\omega_+(Y)\) for positive \(X,Y\).  Every self-adjoint
finite matrix has a Jordan decomposition \(h=h_+-h_-\) with \(h_\pm\ge0\) and
\(h_+h_-=0\), so define
\[
\omega(h):=\omega_+(h_+)-\omega_+(h_-).
\]
If \(h=X-Y=X'-Y'\) are two differences of positive operators, then
\(X+Y'=X'+Y\), and additivity of \(\omega_+\) gives the same value; hence the
extension is well defined and real-linear on self-adjoint elements.  Complex-linear
extension by \(a=h_1+ih_2\) defines a linear functional on
\(\mathcal A\).  Positivity is immediate from construction, and
normalization gives \(\omega(\Id)=p(\Id)=1\).


\item \textbf{Trace-dual density-matrix representation.}
Conversely, any positive normalized linear functional restricts to a normalized
additive probability assignment on the effect interval.  Finite-dimensional
trace pairing on each matrix summand gives a unique block matrix
\(\rho=\bigoplus_\lambda\rho_\lambda\) with
\(\omega(a)=\sum_\lambda\Tr(\rho_\lambda a_\lambda)\).  Testing on rank-one
projections inside each block gives \(\rho_\lambda\ge0\), and normalization is
\(\sum_\lambda\Tr\rho_\lambda=1\).  This proves
\eqref{p4:eq:omega-trace} and the effect formula.  Uniqueness follows from
nondegeneracy of the block trace pairing.


\item \textbf{Selector-dependent alternatives.}
Nonlinear, basis-weighted, outcome-weighted, or context-dependent alternatives
either fail the defining finite coarse-graining law or require selector data
of the type excluded by exact GP2 from the unlabelled domain.  They may occur
as apparatus/interface readouts after the relevant selector has been supplied.
\end{proofsteps}

\end{proof}

\begin{corollary}[State-root and test-root equivalence at finite resolution]
\label{p4:cor:state-test-root-equivalence}
For a finite generated algebra \(\mathcal A\), either of the following
descriptions determines the other without a new role-level occurrence:
\[
 \mathcal A_{\rm sa}
 \xrightarrow{\ a\mapsto\widehat a\ }
 \operatorname{Aff}(S(\mathcal A),\RR),
 \qquad
 S(\mathcal A)
 \xrightarrow{\ \omega\mapsto(a\mapsto\omega(a))\ }
 \mathcal A_{\rm sa}^*.
\]
The first arrow is an order-unit isomorphism by Paper~II-A, and the state
space is the normalized positive base of the dual cone.  Hence one may take
an occurring separating state/preparation family or an occurring finite test
algebra as the root; the dual role, projection/effect lattice, spectral
squaring, and Jordan observable product are generated.  The associative
orientation, complete positivity, and tensor factorization are not.  A named preparation or
performed test remains a token.
\end{corollary}

\begin{proof}
Paper~II-A Theorem~\ref{p2:thm:canonical-finite-state-test-evaluation}
proves the first assertion, including surjectivity onto affine functions.
Positivity and normalization define exactly the displayed base of the dual
cone.  GP2 retains these canonical dual constructions, while G0 still governs
selection of an individual token.
\end{proof}



\paragraph{Born evaluation law at finite resolution}
\label{p4:sec:part1-born}


\begin{theorem}[Finite blockwise Born evaluation theorem]
\label{p4:thm:born}
Let \(\mathcal A=\bigoplus_\lambda M_{d_\lambda}(\CC)\), let
\(\omega\in S(\mathcal A)\) have block density \(\rho\), and let
\(P\in\mathcal A\) be a projection.  Its state evaluation is
\begin{equation}
\label{p4:eq:born-projection}
\omega(P)=\sum_\lambda\Tr(\rho_\lambda P_\lambda).
\end{equation}
More generally, for every effect \(E\in\operatorname{Eff}(\mathcal A)\),
\begin{equation}
\label{p4:eq:born-effect}
\omega(E)=\sum_\lambda\Tr(\rho_\lambda E_\lambda)\in[0,1].
\end{equation}
These normalized evaluations are the finite probability-role values and may
be denoted \(\operatorname{Prob}_\omega(P)\) and
\(\operatorname{Prob}_\omega(E)\).  A \(\mathsf{PROB}\)-tagged token specifies
which preparation, test, outcome record, or calibration is actual.  No extra
normalization, additivity, trace-form, or semantic hypothesis is required.
\end{theorem}


\begin{proof}
Theorem~\ref{p4:thm:density-matrix} supplies the unique block density of
\(\omega\).  For a projection \(P=P^*=P^2\), positivity gives
\[
0\le \omega(P),
\qquad
0\le \omega(\Id-P)=1-\omega(P),
\]
so \(0\le\omega(P)\le1\).  By the block trace representation,
\[
\omega(P)=\sum_\lambda\Tr(\rho_\lambda P_\lambda).
\]
If \(E\) is any effect, then \(0\le E\le\Id\), so the same positivity and
normalization give \(0\le\omega(E)\le1\), and the same representation gives
the displayed block sum.  Linearity supplies finite coarse-graining
additivity.  A token-level \(\mathsf{PROB}\) occurrence selects an actual
boundary instance of this generated rule; it does not change the rule or its
role typing.
\end{proof}


\begin{remark}[Status of finite Born evaluation]
\label{p4:rem:born-evaluation-status}
Theorem~\ref{p4:thm:born} is not a Gleason reconstruction from projection
probabilities alone and it is not an independent trace-form postulate.  The
full positive-normalized state space is already part of Paper~II-A's finite
algebra architecture; finite trace duality gives its block density and effect
evaluation.  GP2 excludes unsupported contextual weights from the unlabelled domain.
Alternative probability rules are accessible/apparatus readouts only
when their selector data has been explicitly supplied.
\end{remark}

\begin{remark}[Projection-only Gleason routes are optional and block-sensitive]
\label{p4:rem:gleason-available}
On a full matrix block of dimension at least three, a projection-only Gleason
argument can reconstruct the state from orthogonal-resolution additivity.
That route is not uniform on \(\mathcal A_{\rm kin}\), which contains
one- and two-dimensional blocks.  The effect-valuation theorem above works on
every finite direct sum, including those blocks, and therefore requires no
full-envelope detour.  Whenever a later \(\mathsf{ENV}\) domain is separately
certified, the ordinary full-matrix Gleason theorem remains available there.
\end{remark}





\paragraph{The structural finite quantum system}
\label{p4:sec:part1-full-carrier}

The physical level-one structural system is therefore the triple
\begin{equation}
\label{p4:eq:part1-carrier-pair}
(\mathcal H_1,\mathcal A_{\rm kin},S(\mathcal A_{\rm kin})),
\qquad
\mathcal A_{\rm kin}
\cong\CC\oplus M_2\oplus M_2\oplus M_4\oplus M_3\oplus M_3,
\end{equation}
and the analogous system on any finite theorem domain is
\((\mathcal H_{\mathcal A},\mathcal A,S(\mathcal A))\), where
\(\mathcal H_{\mathcal A}=\bigoplus_\lambda\CC^{d_\lambda}\) is the minimal
faithful carrier of \(\mathcal A=\bigoplus_\lambda M_{d_\lambda}\).  No full
carrier envelope or separate fundamental wavefunction postulate is required.
Wavefunctions, bases, and state vectors enter later as representations or
presentation choices.  \(\mathsf{FENV}\) is an optional graded field-algebra
growth domain, whereas \(\mathsf{ENV}\) is the distinct full-event extension.



\section{Subsystem composition and completely positive operations}
\label{p4a:sec:v238-a4}
\noindent\textit{Source and claim boundary.} Commuting local images and tensor gluing are obtained on the certified relation-free GP2 composition domain; local processes and dilations yield complete positivity and Kraus representatives without a second autonomous channel axiom.

\subsection{Structural operators and their action}
\label{p4:sec:part1-operators}

\begin{table}[!htbp]
\centering
\caption[Structural quantum operators]{Structural quantum operators.}
\label{p4:tab:structural-operators}
\renewcommand{\arraystretch}{1.2}
\begin{tabularx}{\textwidth}{>{\raggedright\arraybackslash}p{0.20\textwidth}
                                >{\raggedright\arraybackslash}p{0.22\textwidth}
                                >{\raggedright\arraybackslash}X}
\toprule
\textbf{Operator / map} & \textbf{Carrier} & \textbf{Action in Part~I} \\
\midrule
Block density \(\rho=\bigoplus_\lambda\rho_\lambda\) & \(\mathcal H_{\mathcal A}\) & state representative; evaluates tests by the block trace pairing \\

Projection $P$ & $\mathcal H_m$ & sharp yes/no event \\

Effect $E$ & $\mathcal H_m$ & coarse event / POVM element \\

Unitary $U$ & $\mathcal H_m$ & presentation change or native kinematic
conjugation; implements actual closed time evolution only for a selected
\(\mathsf{DYN}\) automorphism flow \\

Channel \(\Phi\) & \(\mathcal A\) or \(S(\mathcal A)\) & CPTP form derived after actual factor-local map certification on one same-type product, or from a generated dilation \\

Instrument $\{\Phi_i\}$ & state space & outcome-resolved measurement update \\

Partial trace $\TrB$ & $\mathcal A_A\otimes\mathcal A_B$ & reduced state / local prediction map \\

Closed generator $H$ & $\mathcal H_m$ & a self-adjoint algebra element
generates a native kinematic flow; its identification and scale as an actual
Hamiltonian require a named closed-dynamics domain \\

Lindbladian $\mathcal L$ & state space & generator of a theorem-domain reduced Markovian semigroup; not native closed-sector data \\
\bottomrule
\end{tabularx}
\TableNotes{This operator ledger separates derived role/algebraic forms from
their token and root interfaces.  \(\mathsf{PROB}\) abbreviates particular
preparation/readout tokens; \(\mathsf{SUB}\), \(\mathsf{PROC}\), and \(\mathsf{DYN}\) record the
existence certificates for composites, maps, and actual evolution.}
\end{table}



\paragraph{Effects and POVMs}
\label{p4:sec:part1-povm}

For a generated finite test algebra \(\mathcal A\), a positive element
\(E\in\mathcal A\) with \(0\le E\le\mathbf1_{\mathcal A}\) is an effect.
A finite family $\{E_i\}_{i=1}^n$ with
\begin{equation}
\label{p4:eq:povm}
E_i\ge0,
\qquad
\sum_{i=1}^n E_i=\mathbf1_{\mathcal A}
\end{equation}
is a POVM \cite{Naimark1940,Holevo1982,Kraus1983}.  If
\(\mathcal A=\bigoplus_\lambda\End(\mathcal H_\lambda)\), a state is
represented by block densities \(\rho=(\rho_\lambda)_\lambda\), and its
canonical normalized outcome weights are
\begin{equation}
\label{p4:eq:povm-prob}
p_i=\omega_\rho(E_i)
=\sum_\lambda\Tr(\rho_\lambda E_{i,\lambda}).
\end{equation}
The effects and normalized weights are already physical event/probability-role
content.  A particular preparation, performed POVM, registered outcome, or
frequency/calibration claim uses token-level G0 (compatibly denoted
\(\mathsf{PROB}\) where needed).  Projective measurements are the special case
in which the $E_i$ are pairwise orthogonal projections.



\paragraph{Unconditional finite channels and complete positivity}
\label{p4:sec:part1-cp}

Ordinary positivity is not enough for an unconditional local finite process.
The residual \(\mathsf{PROC}\) interface supplies an actual affine,
normalization-preserving map and certifies that it acts factor-locally when
one occurrence is transformed while an independently usable same-type
occurrence is untouched.
The tensor-product law for the two occurrences is mathematical once the
\(\mathsf{SUB}\) occurrence certificate has supplied their joint application
domain.  The factor-local lift is, by the product-operation constructor,
\(\operatorname{id}_{\mathcal A}\otimes\Phi\); preservation of the admitted
joint states is process well-typedness, not a second stability axiom.
Positivity on that one same-type copy then forces complete
positivity, block by block; arbitrary mathematical ancilla amplifications
follow afterward.  Thus neither a full matrix envelope nor the physical
existence of every possible ancilla is assumed.  The
Choi--Kraus--Stinespring results
\cite{Stinespring1955,Choi1975,Kraus1983} are used only as
finite-dimensional representation theorems.


\begin{theorem}[Factor-local same-type action forces complete positivity]
\label{p4:thm:cp-forced}
Let
\[
\mathcal A=\bigoplus_{\lambda\in\Lambda}
\End(\mathcal H_\lambda)
\]
be a finite generated test algebra, let
\(\tau(X):=\sum_\lambda\Tr(X_\lambda)\), and let
\(F:S(\mathcal A)\to S(\mathcal A)\) be an affine map supplied on an actual
\(\mathsf{PROC}\) domain.  Its unique complex-linear extension
\(\Phi:\mathcal A\to\mathcal A\) is obtained on the positive cone from
\[
\Phi(X)=
\begin{cases}
\tau(X)\,F\!\left(X/\tau(X)\right),&X\ne0,\\
0,&X=0,
\end{cases}
\qquad X\ge0,
\]
and then by real and complex linearity.  Suppose \(\mathsf{SUB}\) certifies
an independently usable same-type occurrence, so that its joint algebra is
\(\mathcal A\otimes\mathcal A\), and suppose \(\mathsf{PROC}\) certifies
that \(F\) acts only on the second factor.  Its product-typed joint action is
then \(\operatorname{id}_{\mathcal A}\otimes\Phi\), which is positive
because an actual process preserves the admitted joint state cone.  Then
\(\Phi\) is completely positive.  The displayed extension is
\(\tau\)-preserving and hence is a CPTP map of the finite direct-sum algebra.
\end{theorem}


\begin{proof}
Affinity of \(F\) makes the displayed positive-cone extension additive:
for \(X,Y\ge0\), normalization of \(X+Y\) is the convex combination of the
normalizations of \(X\) and \(Y\) with weights
\(\tau(X)/\tau(X+Y)\) and \(\tau(Y)/\tau(X+Y)\).  The same formula gives positive
homogeneity.  Every self-adjoint element is a difference of positive
elements.  If \(X-Y=X'-Y'\) with all four elements positive, then
\(X+Y'=X'+Y\); positive-cone additivity gives
\(\Phi(X)-\Phi(Y)=\Phi(X')-\Phi(Y')\).  Hence the difference extension is
well defined and real linear on the self-adjoint part of \(\mathcal A\), and
then complex linear.  Positivity is immediate, while
\(\tau(\Phi(X))=\tau(X)\) first for \(X\ge0\) and hence by linearity for all
\(X\).

Let \(\iota_\lambda:M_{d_\lambda}\to\mathcal A\) and
\(q_\mu:\mathcal A\to M_{d_\mu}\) denote the canonical block inclusion and
projection, and set
\(\Phi_{\mu\lambda}:=q_\mu\Phi\iota_\lambda\).  Restricting
\(\operatorname{id}_{\mathcal A}\otimes\Phi\) to the
\((\lambda,\lambda)\) input block and projecting the output to the
\((\lambda,\mu)\) block shows that
\[
\operatorname{id}_{M_{d_\lambda}}\otimes\Phi_{\mu\lambda}
\]
is positive.  Applied to the maximally entangled rank-one operator in
\(M_{d_\lambda}\otimes M_{d_\lambda}\), this makes the Choi matrix of every
\(\Phi_{\mu\lambda}\) positive.  The finite Choi theorem,
Theorem~\ref{p4:thm:finite-choi-kraus-instruments}, makes each block component
completely positive.  Since
\(q_\mu\Phi(X)=\sum_\lambda\Phi_{\mu\lambda}(X_\lambda)\), their finite sum
is completely positive on \(\mathcal A\).  Hence every finite mathematical
amplification is positive, without assuming each amplification as an
upstream physical occurrence.
\end{proof}

\begin{remark}[Relation of \(\mathsf{PROC}\) to \(\mathsf{SUB}\) and \(\mathsf{DYN}\)]
\label{p4:rem:proc-relative-to-sub-dyn}
The representation content formerly placed in \(\mathsf{PROC}\) is now a
theorem: once an actual affine normalized map occurs as a well-typed
factor-local action on one same-type product,
Theorem~\ref{p4:thm:cp-forced} gives complete positivity and
Theorem~\ref{p4:thm:finite-choi-kraus-instruments} gives block Kraus form.
What remains irreducible in \(\mathsf{PROC}\) is application-level existence:
the algebra does not select a map or assert that it occurs.  A supplied
\(\mathsf{DYN}\) dilation, environment state, and subsystem occurrence can
instantiate such a map by unitary evolution and partial evaluation, but
\(\mathsf{SUB}\) and \(\mathsf{DYN}\) alone do not manufacture those data for
an arbitrary channel.  Thus the residual interface is an occurrence
certificate, not an additional mathematical process axiom.
\end{remark}




\paragraph{Kraus form and instruments}
\label{p4:sec:part1-kraus}

\begin{theorem}[Finite Choi--Kraus representation and instruments]
\label{p4:thm:finite-choi-kraus-instruments}
Let
\[
\mathcal A=\bigoplus_{\lambda\in\Lambda}\End(\mathcal H_\lambda),
\qquad
\mathcal B=\bigoplus_{\mu\in M}\End(\mathcal K_\mu)
\]
be finite-dimensional \(C^*\)-algebras and let
\(\Phi:\mathcal A\to\mathcal B\) be linear.  Write
\(\Phi_{\mu\lambda}=q_\mu\Phi\iota_\lambda\) for its matrix-block
components.  Then \(\Phi\) is completely positive if and only if every
ordinary Choi matrix \(C_{\mu\lambda}\) of
\(\Phi_{\mu\lambda}:\End(\mathcal H_\lambda)\to
\End(\mathcal K_\mu)\) is positive.  Equivalently, there are finitely many
operators
\(K_{\mu\lambda,a}:\mathcal H_\lambda\to\mathcal K_\mu\) such that
\begin{equation}
\label{p4:eq:kraus}
\bigl(\Phi(\rho)\bigr)_\mu
=\sum_{\lambda,a}
K_{\mu\lambda,a}\rho_\lambda K_{\mu\lambda,a}^*.
\end{equation}
Moreover \(\Phi\) preserves total block trace if and only if, for every
input block \(\lambda\),
\begin{equation}
\label{p4:eq:kraus-tp}
\sum_{\mu,a}K_{\mu\lambda,a}^*K_{\mu\lambda,a}
=\Id_{\mathcal H_\lambda} .
\end{equation}
A finite family \(\{\Phi_i\}\) is an instrument exactly when every
\(\Phi_i\) is CP and \(\sum_i\Phi_i\) is trace preserving.  For a state
\(\rho=(\rho_\lambda)_\lambda\) it has normalized outcome weights and
posterior states
\begin{equation}
\label{p4:eq:instrument}
p_i=\sum_\mu\Tr\bigl((\Phi_i(\rho))_\mu\bigr),
\qquad
\rho_i'=\frac{\Phi_i(\rho)}{p_i}
\quad (p_i>0).
\end{equation}
Thus a supplied instrument represents its normalized outcome weights and
conditional posterior states inside the finite CP calculus.  The algebra
does not select which instrument is physically realized or which outcome
occurs in an individual trial.
\end{theorem}

\begin{proof}
If \(\Phi\) is completely positive, composing it with a block inclusion and
a block projection shows that every \(\Phi_{\mu\lambda}\) is completely
positive.  The ordinary finite Choi argument therefore gives
\(C_{\mu\lambda}\ge0\).

Conversely, diagonalize each positive Choi matrix as
\(C_{\mu\lambda}=\sum_a
\ket{\xi_{\mu\lambda,a}}\bra{\xi_{\mu\lambda,a}}\).
Blockwise vectorization uniquely supplies maps
\(K_{\mu\lambda,a}:\mathcal H_\lambda\to\mathcal K_\mu\) by
\[
\ket{\xi_{\mu\lambda,a}}
=\sum_j e_{\lambda,j}\otimes K_{\mu\lambda,a}e_{\lambda,j}.
\]
Comparison of matrix-unit blocks gives
\(\Phi_{\mu\lambda}(X)=\sum_aK_{\mu\lambda,a}X
K_{\mu\lambda,a}^*\).  Summing over \(\lambda\) yields
\eqref{p4:eq:kraus}.  Every matrix amplification is then a finite sum of
positive conjugation maps, so \(\Phi\) is completely positive.

For a positive input supported in block \(\lambda\), cyclicity of trace gives
\[
\sum_\mu\Tr\bigl((\Phi(\rho))_\mu\bigr)
=\Tr\!\left(
\rho_\lambda\sum_{\mu,a}K_{\mu\lambda,a}^*K_{\mu\lambda,a}
\right).
\]
Equality with \(\Tr(\rho_\lambda)\) for all such inputs is equivalent to
\eqref{p4:eq:kraus-tp}.  Applying this criterion to
\(\sum_i\Phi_i\) gives the instrument statement, and
\eqref{p4:eq:instrument} is normalization of the subnormalized CP output.
Empirical outcome and update meaning still require the corresponding
\(\mathsf{PROB}\) and \(\mathsf{PROC}\) interfaces.
\end{proof}



% BEGIN v2.296: OUTCOME-WISE QUOTIENT OF FINITE GENERATED INSTRUMENTS
\begin{theorem}[Outcome-wise descent of finite instruments and adaptive protocols]
\label{p4:thm:v296-outcome-wise-instrument-descent}
Let $j:\mathcal A_0\simeq\mathcal A\subseteq\End(H)$ be a
specified finite-dimensional unital $C^*$-test algebra with the
canonical state restriction $Q$ of
Theorem~\ref{p2b:thm:v293-operational-multiplicity-quotient}.
Use the finite instrument grammar of
Theorem~\ref{p4:thm:finite-choi-kraus-instruments}. Let a registered
family of finite instruments on $H$ have completely
positive, trace-nonincreasing outcome maps $\Psi_{a,i}$, with
$\sum_i\Psi_{a,i}$ trace preserving for each setting $a$. Consider
finite adaptive protocols built from these instruments, classical
outcome-dependent choices of later settings, and final readouts in
$\mathcal A$. The following are equivalent:
\begin{enumerate}[label=(\roman*),leftmargin=2em]
\item For every registered setting and each \emph{individual} outcome,
$\Psi_{a,i}^*(\mathcal A)\subseteq\mathcal A$.
\item There are unique completely positive reduced outcome maps
$\widetilde\Psi_{a,i}$ on the quotient state algebra such that
$Q\Psi_{a,i}=\widetilde\Psi_{a,i}Q$ on all positive, possibly
subnormalized, states; they constitute normalized instruments when
summed, and every finite adaptive protocol has the same full joint
outcome probabilities, reduced posterior states when nonzero, and
final $\mathcal A$-readouts before and after restriction.
\end{enumerate}
Equivalently, (i) can be tested by the complete real Hermitian basis
of $\mathcal A$, \emph{separately for each outcome map}. In the exact
kinematic and full-coordinate multiplicity-blind algebras of
Proposition~\ref{p2c:prop:v294-finite-process-descent}, those bases
have respectively $43$ and $137$ elements.
The weaker condition
$(\sum_i\Psi_{a,i})^*(\mathcal A)\subseteq\mathcal A$ is not
sufficient: put $H=\mathbb C^2\otimes\mathbb C^2$,
$\mathcal A=M_2\otimes I_2$, $P_i=|i\rangle\langle i|$ and
$\Psi_i(\rho)=(I_2\otimes P_i)\rho(I_2\otimes P_i)$.
The sum is an $\mathcal A$-preserving dephasing channel, but
$\Psi_i^*(I_4)=I_2\otimes P_i\notin\mathcal A$, so the actual outcome
weights distinguish states with identical $Q$-images.
These conclusions concern an \emph{already specified} instrument
family, not the physical occurrence of a particular measurement or an
outcome in a trial.
\end{theorem}
\begin{proof}
If (i) holds, $j^{-1}\Psi_{a,i}^*j$ is completely positive, since
$j$ is a $*$-isomorphism and the unitality of each outcome map is not
needed. Duality on all block states defines the unique positive
subnormalized quotient map and gives the stated intertwining. Since
$\sum_i\Psi_{a,i}^*(I)=I$, its reduced sum preserves normalization.
By induction on the length of an outcome word,
$Q\Psi_{a_n,i_n}\cdots\Psi_{a_1,i_1}
=\widetilde\Psi_{a_n,i_n}\cdots\widetilde\Psi_{a_1,i_1}Q$,
also when later $a_k$ depends on previous outcomes. The trace of a
subnormalized state is its outcome-word probability because $j$ is
unital and $Q$ preserves trace. Dividing by a nonzero word probability
preserves the quotient posterior and every subsequent test expectation.
Conversely, the asserted one-step outcome probabilities and posterior
$\mathcal A$-expectations for all initial states imply that
$\rho\mapsto\operatorname{Tr}(\Psi_{a,i}(\rho)A)$ depends only on
$Q(\rho)$ for every $A\in\mathcal A$. The exact effect converse in
Theorem~\ref{p2b:thm:v293-operational-multiplicity-quotient} yields
$\Psi_{a,i}^*(A)\in\mathcal A$. Real-basis testing and the displayed
projective counterexample prove the remaining assertions; for the
counterexample, $|\psi\rangle\langle\psi|\otimes P_0$ and
$|\psi\rangle\langle\psi|\otimes P_1$ have the same $Q$-state but
opposite outcome probabilities for outcome $0$.
\end{proof}
% END v2.296: OUTCOME-WISE QUOTIENT OF FINITE GENERATED INSTRUMENTS



\paragraph{Composite systems, reduced states, and partial trace}
\label{p4:sec:part1-partial-trace}

Let
\(\mathcal A=\bigoplus_\lambda\End(\mathcal H_{A,\lambda})\) and
\(\mathcal B=\bigoplus_\mu\End(\mathcal H_{B,\mu})\).  By
Theorem~\ref{p4:thm:tensor-product}, their relation-free composite algebra is
\[
\mathcal A\otimes\mathcal B
\cong
\bigoplus_{\lambda,\mu}
\End(\mathcal H_{A,\lambda}\otimes\mathcal H_{B,\mu}).
\]
Every joint state is represented by positive block densities
\(\rho_{\lambda\mu}\) with total trace one.  Its marginals are the block
densities
\begin{equation}
\label{p4:eq:partial-trace}
 (\rho_A)_\lambda
=\sum_\mu\Tr_{\mathcal H_{B,\mu}}(\rho_{\lambda\mu}),
\qquad
 (\rho_B)_\mu
=\sum_\lambda\Tr_{\mathcal H_{A,\lambda}}(\rho_{\lambda\mu}).
\end{equation}
For every local effect \(E_A=(E_{A,\lambda})_\lambda\), partial evaluation
recovers the local weight:
\begin{equation}
\label{p4:eq:local-prob}
\sum_{\lambda,\mu}\Tr\!\left[
\rho_{\lambda\mu}
(E_{A,\lambda}\otimes\Id_{\mathcal H_{B,\mu}})
\right]
=\sum_\lambda\Tr\bigl((\rho_A)_\lambda E_{A,\lambda}\bigr).
\end{equation}
This is the canonical finite direct-sum local-evaluation rule.  Reading it as
a physical reduced-state prediction requires the supplied joint occurrence
and probability interfaces.



\subsection{Tensor-product composition and entanglement}
\label{p4:sec:part1-entanglement}


\begin{proposition}[Common-witness descendants co-occur without a second token]
\label{p4:prop:common-witness-cooccurrence}
Let \(x\) be an already occurring G0 witness in a generated role \(X\), and
let \(F:X\to Y\) and \(G:X\to Z\) be typed host-natural constructors which
preserve the complete source witness and its declared distinctions.  Then the
unique product pairing
\[
 \langle F,G\rangle:X\longrightarrow Y\times Z,
 \qquad
 x\longmapsto(F(x),G(x)),
\]
is an occurrence-preserving constructor.  Hence \(F(x)\) and \(G(x)\)
co-occur without an independent \(\mathsf{SUB}\) token.

This conclusion does not assert relation-free independence.  Every relation
forced by their common ancestry remains in the image of
\(\langle F,G\rangle\); an independent tensor-product branch must still pass
the chronological no-cross-relation test.
\end{proposition}

\begin{proof}
The categorical product property makes \(\langle F,G\rangle\) the unique map
whose two projections are \(F\) and \(G\), so no member, parameter, or second
source witness is chosen.  Host naturality and complete witness preservation
pass componentwise to the pairing.  Applying this generated constructor to
the occurring \(x\) therefore produces the joint descendant token.  Its image
is generally the graph of the common-source relation, which proves the final
restriction.
\end{proof}


\begin{theorem}[Exact GP2 product-operation interchange]
\label{p4:thm:gp2-product-operation-interchange}
Let \(\mathcal A_A,\mathcal A_B\) be generated finite test algebras and let
\(\mathcal C\) be a joint finite test realization.  Suppose the complete
unitary operation roles of the GP2 relation-free product are represented by
one exact operation-preserving map
\[
 \varrho:
 \mathcal U(\mathcal A_A)\times\mathcal U(\mathcal A_B)
 \longrightarrow \mathcal U(\mathcal C),
\]
whose restrictions are the unitary parts of faithful unital
\(*\)-embeddings \(i_A:\mathcal A_A\hookrightarrow\mathcal C\) and
\(i_B:\mathcal A_B\hookrightarrow\mathcal C\):
\[
 \varrho(u,\mathbf1)=i_A(u),
 \qquad
 \varrho(\mathbf1,v)=i_B(v).
\]
Then the two local algebra images commute elementwise:
\[
 \boxed{[i_A(a),i_B(b)]=0
 \qquad(a\in\mathcal A_A,\ b\in\mathcal A_B).}
\]
This is the typed product-operation realization of the complete GP2 product,
not an independently inserted commutation relation.
\end{theorem}

\begin{proof}
Componentwise multiplication in the relation-free product gives
\[
 (u,\mathbf1)(\mathbf1,v)=(u,v)
 =(\mathbf1,v)(u,\mathbf1).
\]
Exact preservation of the generated operation law therefore gives
\[
 i_A(u)i_B(v)=\varrho(u,v)=i_B(v)i_A(u)
\]
for all local unitaries.  If \(x=x^*\) and \(\|x\|\le1\), finite functional
calculus gives the unitary
\[
 u_x=x+i(\mathbf1-x^2)^{1/2},
 \qquad x=\frac12(u_x+u_x^*).
\]
After scaling, every self-adjoint element is a linear combination of two
unitaries, and every algebra element is a linear combination of at most four.
Bilinearity of the commutator proves the displayed conclusion.
\end{proof}

\begin{remark}[Why exactness is required]
\label{p4:rem:product-operation-exactness}
A projective realization may satisfy
\(\varrho(g)\varrho(h)=\omega(g,h)\varrho(gh)\) and turn commuting factor
operations into phase-commuting Weyl pairs.  The multiplier \(\omega\) is a
cross-role \(2\)-cocycle.  It belongs only on a separately generated
braided/interaction branch; inserting it on the unbraided relation-free
branch would violate GP2 no-insertion.  Conversely, Cartesian support freedom
without an operation-preserving realization is insufficient to imply
commutation.
\end{remark}


\begin{theorem}[Relation-free independent-occurrence criterion]
\label{p4:thm:independent-finite-composability}
Let
\[
\mathcal A_A=\bigoplus_{\lambda\in\Lambda}\End(\mathcal H_{A,\lambda}),
\qquad
\mathcal A_B=\bigoplus_{\mu\in M}\End(\mathcal H_{B,\mu})
\]
be generated finite test algebras.  The residual \(\mathsf{SUB}\) interface
certifies the application fact that an occurrence of each type exists jointly.
If GP2 supplies no generated cross-relation between them and its complete
product-operation role is represented exactly as in
Theorem~\ref{p4:thm:gp2-product-operation-interchange}, its relation-exact
joint realization is a finite algebra \(\mathcal C\) with:
\begin{enumerate}[label=(\roman*), leftmargin=2em]
\item there are faithful unital \(*\)-embeddings
      \(i_A:\mathcal A_A\hookrightarrow\mathcal C\) and
      \(i_B:\mathcal A_B\hookrightarrow\mathcal C\);
\item the embedded local algebras commute elementwise,
      \([i_A(a),i_B(b)]=0\) for all local observables;
\item \(\mathcal C\) is generated as a unital \(*\)-algebra by
      \(i_A(\mathcal A_A)\) and \(i_B(\mathcal A_B)\);
\item for every pair of minimal central projections
      \(z_\lambda\in\mathcal A_A\) and \(w_\mu\in\mathcal A_B\),
      \(i_A(z_\lambda)i_B(w_\mu)\ne0\).
\end{enumerate}
Condition (iv) is the exact finite-dimensional meaning of ``no deleted
cross-sector pair.''  A diagonal sector match, graph constraint, or selected
quotient would set at least one such product to zero and therefore insert a
cross-relation.  Noncommuting or braided gluing is a different supplied
interaction/ordering role.  None of these algebraic facts asserts that the two
occurrences exist; that is precisely the remaining \(\mathsf{SUB}\) content.
\end{theorem}


\begin{proof}
The \(\mathsf{SUB}\) certificate supplies only joint occurrence.
Theorem~\ref{p4:thm:gp2-product-operation-interchange} derives
commutation from the exact componentwise operation law before any tensor
algebra is used.  GP2 no-omission retains both faithful local test images, and
constructor-relative closure retains the algebra they generate.  In a finite
direct sum, every ideal is a sum of simple blocks.
Consequently any nonzero kernel of the multiplication map from
\(\mathcal A_A\otimes\mathcal A_B\) deletes at least one joint central block
\(z_\lambda\otimes w_\mu\), exactly contradicting (iv).  Conversely a
specified deletion, projective multiplier, or noncommuting relation is
additional cross-role data.  Thus the displayed conditions are the finite
algebraic form of the exact unbraided GP2 product realization, while the
occurrence itself is not derived.
\end{proof}


\begin{theorem}[Finite direct-sum tensor monoidality]
\label{p4:thm:tensor-product}
Let \(A\) and \(B\) satisfy
Theorem~\ref{p4:thm:independent-finite-composability}.  Multiplication
induces an isomorphism
\[
\mu:\mathcal A_A\otimes\mathcal A_B
\overset{\sim}{\longrightarrow}\mathcal C,
\qquad
\mu(a\otimes b)=i_A(a)i_B(b).
\]
Every finite faithful representation of this joint algebra has, for integers
\(r_{\lambda\mu}\ge1\), the form
\[
\mathcal K\cong
\bigoplus_{\lambda,\mu}
\bigl(\mathcal H_{A,\lambda}\otimes
\mathcal H_{B,\mu}\otimes\mathbb C^{r_{\lambda\mu}}\bigr).
\]
Its minimal faithful carrier is therefore, uniquely up to unitary
equivalence,
\begin{equation}
\label{p4:eq:tensor-product-carrier}
\bigoplus_{\lambda,\mu}
\mathcal H_{A,\lambda}\otimes\mathcal H_{B,\mu},
\end{equation}
and their joint observable algebra is
\begin{equation}
\label{p4:eq:tensor-product}
\mathcal C=\mathcal A_A\otimes\mathcal A_B
\cong\bigoplus_{\lambda,\mu}
\End(\mathcal H_{A,\lambda}\otimes\mathcal H_{B,\mu}).
\end{equation}
Product states are the functionals \(\omega_A\otimes\omega_B\).  A state on
the tensor product is \emph{entangled} precisely when it is not a convex
combination of product states.
\end{theorem}

\begin{proof}
Because the two faithful images commute, multiplication defines a unital
\(*\)-homomorphism
\[
\mu:\mathcal A_A\otimes\mathcal A_B\longrightarrow\mathcal C,
\qquad
\mu(a\otimes b)=i_A(a)i_B(b).
\]
Generation makes \(\mu\) surjective.  The domain is the direct sum
\[
\mathcal A_A\otimes\mathcal A_B
\cong\bigoplus_{\lambda,\mu}
\End(\mathcal H_{A,\lambda}\otimes\mathcal H_{B,\mu}).
\]
Its kernel is an ideal and hence a sum of some simple
\((\lambda,\mu)\) blocks.  If one such block lay in the kernel, the
corresponding central product
\(i_A(z_\lambda)i_B(w_\mu)\) would vanish, contrary to
Theorem~\ref{p4:thm:independent-finite-composability}(iv).  Thus the kernel
is zero and \(\mu\) is an isomorphism.  The irreducible representations are
the defining representations of the displayed simple blocks, so every
faithful representation contains each with multiplicity
\(r_{\lambda\mu}\ge1\), and minimal faithful representation sets every
multiplicity to one.  This conclusion uses no carrier-resolving event role or
full-envelope assumption.  The state statements are then the standard
definitions on the resulting tensor product.
\end{proof}


\begin{remark}[Even-observable locality in fermionic sectors]
The tensor-composition theorem concerns independently addressable physical observable
algebras.  In fermionic presentations, odd field generators may graded-commute, while
the physical even local observable algebras commute ordinarily.  A supplied grading or
braiding is additional constituted structure and therefore lies outside the unbraided
hypothesis of Theorem~\ref{p4:thm:independent-finite-composability}.
\end{remark}

This tensor law is the algebraic expression of relation-exact independent
typing.  Its proof uses the already generated finite algebras, a
\(\mathsf{SUB}\) occurrence certificate, and GP2 relation exactness; minimal
faithful representation then follows from the direct-sum algebra itself.  It
is independent of the complete-positivity theorem, which needs only one
independently usable same-type copy.



\subsection{Non-signaling}
\label{p4:sec:part1-nonsignaling}


\begin{theorem}[Non-signaling]
\label{p4:thm:nonsignaling}
Let \(\mathcal A\) and \(\mathcal B\) be finite direct-sum algebras, let
\(\rho_{AB}\) be the block density of a state on
\(\mathcal A\otimes\mathcal B\), and let
\(\Phi_B\) be any total-block-trace-preserving CP map acting only on
\(\mathcal B\).  Then the marginal state of \(A\) is unchanged:
\begin{equation}
\label{p4:eq:nonsignaling}
\operatorname{marg}_A\!\left[
(\operatorname{id}_A\otimes\Phi_B)(\rho_{AB})\right]
=\operatorname{marg}_A(\rho_{AB}).
\end{equation}
Hence no choice of nonselective local operation on $B$ can alter the outcome statistics of measurements performed only on $A$.  Selective remote instruments can change conditional local states after conditioning on the remote outcome; without outcome communication, the averaged map is CPTP and the reduced state is unchanged.
\end{theorem}


\begin{proof}
Let \(\Phi_B^*\) be the adjoint with respect to total block trace.  Trace
preservation is equivalent to
\(\Phi_B^*(\mathbf1_{\mathcal B})=\mathbf1_{\mathcal B}\).  For every
\(X_A\in\mathcal A\), the marginal after the local operation evaluates as
\begin{align*}
\omega_{AB}'(X_A\otimes\mathbf1_{\mathcal B})
&=
\omega_{AB}\!\left(
X_A\otimes\Phi_B^*(\mathbf1_{\mathcal B})\right)\\
&=
\omega_{AB}(X_A\otimes\mathbf1_{\mathcal B}).
\end{align*}
Finite states are separated by their evaluations, so the two marginals are
equal.  Hence local normalized weights on \(A\) are independent of the chosen
nonselective CPTP map on \(B\).  Selective conditioning is a different,
communicated outcome interface.

\end{proof}



\section{Finite dynamics, automorphisms, symmetries and generators}
\label{p4a:sec:v238-a5}
\noindent\textit{Source and claim boundary.} Constituted finite one-parameter flows have exact self-adjoint generator representations; finite Noether and open GKLS statements are qualified by their actual generator, process and symmetry domains. No arbitrary dynamical law is selected by the algebra alone.

\subsection{Closed and open generators}
\label{p4:sec:part1-generators}

\begin{corollary}[Native nontrivial kinematic automorphism flows]
\label{p4:cor:native-kinematic-flows}
For every self-adjoint \(X\in\mathcal A_{\rm kin}\),
\[
 \alpha_s^X(a):=e^{isX}ae^{-isX}
\]
is a continuous automorphism flow of \(\mathcal A_{\rm kin}\).  Noncentral
self-adjoint choices in any noncommutative block give nonconstant flows;
self-adjoint representatives of the constituted kinematic generators provide
examples.  Hence nontrivial kinematic automorphism flows exist without a new
dynamics postulate.  Identifying one flow with elapsed physical succession,
fixing its physical scale, and asserting that an actual system follows it
remain the \(\mathsf{DYN}\) occurrence/selection interface.  Once the
parametrized flow is fixed, its inner generator is determined only up to the
central gauge proved below.
\end{corollary}

\begin{proof}
Finite functional calculus gives \(e^{isX}\in\mathcal A_{\rm kin}\), so
conjugation is a norm-continuous one-parameter group of
\(*\)-automorphisms.  Its derivative at \(s=0\) is \(i[X,a]\).  If the flow
were constant for every \(a\), then \(X\) would lie in the center; therefore
every noncentral self-adjoint \(X\) gives a nonconstant flow.
\end{proof}

A norm-continuous one-parameter automorphism group of a finite direct-sum
algebra fixes each minimal central projection: its induced permutation of the
finite block set is continuous and therefore constant.  On every full matrix
block the restricted flow is inner.  Once a closed-dynamics theorem
domain has supplied such a flow on
\(\mathcal A=\bigoplus_\lambda\End(\mathcal H_\lambda)\), the finite
one-parameter generator closure of Paper~III-A, together with blockwise
innerness, gives self-adjoint \(H_\lambda\) and
\(H=\bigoplus_\lambda H_\lambda\) such that
\begin{equation}
\label{p4:eq:heisenberg-flow}
\alpha_t(a)=e^{itH}ae^{-itH}.
\end{equation}
No separate infinite-dimensional Stone/Hille--Yosida hypothesis is involved
in this finite representation statement.

\begin{corollary}[Central gauge of a fixed finite closed generator]
\label{p4:cor:central-gauge-closed-generator}
Let \(\alpha_t\) be one fixed parametrized norm-continuous automorphism flow
of the finite direct-sum algebra \(\mathcal A\).  If self-adjoint
\(H,H'\in\mathcal A\) both implement that same flow, then
\[
 \boxed{H-H'\in Z(\mathcal A)_{\rm sa}.}
\]
Conversely, adding any \(Z\in Z(\mathcal A)_{\rm sa}\) leaves the flow
unchanged.  Thus the observable-flow generator is intrinsically the inner
derivation
\[
 \delta_H(a)=i[H,a]
\]
or, equivalently, the quotient class
\[
 \boxed{[H]\in \mathcal A_{\rm sa}/Z(\mathcal A)_{\rm sa}.}
\]
This quotient is for a fixed flow parameterization.  It does not quotient a
physical rescaling of time, and it does not erase central energy data when a
later theorem uses \(H\) itself as an energy observable, Gibbs weight, phase
reference, or coupling datum.
\end{corollary}

\begin{proof}
If \(H\) and \(H'\) implement the same parametrized flow, differentiating at
\(t=0\) gives \(i[H,a]=i[H',a]\) for every \(a\in\mathcal A\).  Hence
\(H-H'\) commutes with all of \(\mathcal A\) and lies in its self-adjoint
center.  Conversely, for central self-adjoint \(Z\),
\(e^{it(H+Z)}=e^{itH}e^{itZ}\), and the central factor cancels in conjugation.
The statements about scale and later uses of \(H\) follow because neither is
part of equality of the fixed observable automorphism flow.
\end{proof}

\begin{proposition}[Bare-algebra naturality selects no nontrivial dynamics]
\label{p4:prop:bare-algebra-dynamics-no-selector}
Let
\(\mathcal A=\bigoplus_\lambda M_{d_\lambda}(\CC)\), and suppose a proposed
canonical rule assigns a generator class
\[
 [H(\mathcal A)]\in
 \mathcal A_{\rm sa}/Z(\mathcal A)_{\rm sa}
\]
naturally under every \(*\)-automorphism of \(\mathcal A\).  Then
\[
 \boxed{[H(\mathcal A)]=0.}
\]
Consequently the bare finite algebra, its state space, and its effect
interval cannot select a nontrivial \(\mathsf{DYN}\) member.
\end{proposition}

\begin{proof}
Naturality under every blockwise inner automorphism gives
\[
 uH_\lambda u^*-H_\lambda=c_\lambda(u)I_{d_\lambda}
\]
on each simple block.  Taking the trace gives \(c_\lambda(u)=0\).
Therefore \(H_\lambda\) commutes with every unitary and is scalar.  Hence
\(H\in Z(\mathcal A)_{\rm sa}\), so its quotient class and inner derivation
vanish.
\end{proof}

\begin{remark}[Exact geometric criterion for reducing \(\mathsf{DYN}\)]
\label{p4:rem:one-ray-dynamics-criterion}
If additional constituted geometry acts on
\(\operatorname{Der}(\mathcal A)\), it selects a unique nonzero dynamical ray
exactly when the invariant subspace is one-dimensional.  A constituted time
orientation then fixes the sign; a dimensionful rate remains calibration.
If the invariant subspace is zero no nontrivial flow is selected, and if its
dimension exceeds one a selector remains.  Reeb, connection, Dirac, modular,
or timed-operation data may discharge this criterion only when their
constructor and invariant-ray proof are already generated.  An actual
trajectory remains an occurrence statement.
This is precisely the finite linear one-ray branch of
Theorem~\ref{p3a:thm:finite-generator-one-ray-rigidity}: the sector-specific
geometry supplies the parent signature, while the universal theorem supplies
the selection logic.
\end{remark}

Exact GP2 explicitly lists a Hamiltonian generator among data that cannot be
inserted as constitutive structure.  Therefore \(H\) is not selected by a
closed Euclidean loop.  A Hamiltonian is imported from an upstream generated
timed-operation law, obtained by Paper~III-A's finite generator closure from
an already occurring \(\mathsf{DYN}\) one-parameter group together with its
Hamiltonian interpretation/scale data, or supplied on an explicit theorem
domain.  The generator theorem represents a flow; it does not select or
actualize one.  A closed Euclidean loop does not by itself select real-time
evolution; the Euclidean-to-Hamiltonian bridge retains its reconstruction
hypotheses.
Here $H$ is the generator in units with $\hbar=1$; if $H_{\rm phys}$ denotes a physical Hamiltonian, then $e^{itH}$ is replaced by $e^{itH_{\rm phys}/\hbar}$.
On block-density states this becomes
\begin{equation}
\label{p4:eq:schrodinger-flow}
\rho_\lambda(t)=e^{-itH_\lambda}\rho_\lambda(0)e^{itH_\lambda}.
\end{equation}
For a fixed product environment assignment, unitary system--environment
evolution followed by partial trace gives an exact CPTP reduced map.  Initial
correlations require a separately supplied CP assignment map or a restricted
compatibility domain.  Only when the reduced family is additionally a
norm-continuous Markov semigroup does it have the GKLS form stated below.

\begin{corollary}[Schr\"odinger and Heisenberg dynamics share one occurrence]
\label{p4:cor:schrodinger-heisenberg-one-corealized-dynamics}
Once a finite closed flow \(\alpha_t\) is generated or supplied on a
\(\mathsf{DYN}\) domain, its state and test forms
\[
 F_t(\rho)=e^{-itH}\rho e^{itH},
 \qquad
 \alpha_t(a)=e^{itH}ae^{-itH}
\]
satisfy the separating evaluation identity
\[
 \boxed{\Tr(F_t(\rho)a)=\Tr(\rho\,\alpha_t(a)).}
\]
They are therefore one co-realized dynamics in evaluation-dual
presentations; a second \(\mathsf{DYN}\) or \(\mathsf{PROC}\) occurrence for
the other picture is redundant.

Likewise, once an open CPTP semigroup \(\Phi_t\) is supplied, its unique
Heisenberg dual
\[
 \Tr(\Phi_t(\rho)a)=\Tr(\rho\,\Phi_t^\sharp(a))
\]
is a co-realized unital CP semigroup on the same domain \(t\ge0\).  This
duality does not supply a negative-time map, inverse channel, Markovian
semigroup when only isolated channels are given, system--environment
assignment, generator selection, initial state, or actual history.
\end{corollary}

\begin{proof}
The closed identity is trace cyclicity.  In the open case, finite
state--test separation gives the unique adjoint map; the Kraus representation
shows that trace preservation and complete positivity of \(\Phi_t\) become
unitality and complete positivity of \(\Phi_t^\sharp\).  Family composition
is contravariant, and the additive one-parameter law yields the dual
semigroup.  Co-realization and the scope boundary are
Theorem~\ref{p2a:thm:evaluation-dual-corealization} and
Corollary~\ref{p2a:cor:finite-state-test-dual-dynamics}.
\end{proof}



\subsection{Closed structural quantum outputs}
\label{p4:sec:part1-outputs}

\begin{table}[!htbp]
\centering
\caption[Closed structural quantum outputs]{Closed structural quantum outputs.}
\label{p4:tab:structural-outputs}
\renewcommand{\arraystretch}{1.2}
\begin{tabularx}{\textwidth}{>{\raggedright\arraybackslash}p{0.35\textwidth}
                                >{\raggedright\arraybackslash}X}
\toprule
\textbf{Structural quantum output} & \textbf{Derived result in Part~I} \\
\midrule
Generated finite test algebra & at \(m_p=1\), \(\mathcal A_{\rm kin}\cong
\mathbb C\oplus M_2\oplus M_2\oplus M_4\oplus M_3\oplus M_3\), with
the full \(\End(\mathcal H_1)\) typed as a graded field envelope under
\(\mathsf{FENV}\), and as a full event algebra when generated tests are
carrier resolving or under an independently adjoined \(\mathsf{ENV}\) root \\

State space & full \(S(\mathcal A)\) derived as positive normalized
functionals \(\Leftrightarrow\) block-density tuples \\

Canonical evaluation & \(\omega_\rho(E)=\sum_\lambda\Tr(\rho_\lambda E_\lambda)\);
generated event/probability-role meaning; \(\mathsf{PROB}\) only selects actual tokens \\

Subsystem composition & a \(\mathsf{SUB}\) occurrence certificate plus GP2
relation exactness \(\Rightarrow\) finite direct-sum tensor product \\

Entanglement & non-factorization of global states on tensor products \\

Reduced local predictions & partial trace and local event formula \\

Non-signaling & local statistics independent of remote choice of CPTP map \\

Generalized measurements & algebraic effects and POVMs derived; empirical
meaning and actual instruments remain residual interfaces \\

Closed dynamics & nontrivial kinematic automorphism flows are native; a
supplied \(\mathsf{DYN}\) identification of an actual flow has an inner
unitary representation \\

Open dynamics & a supplied \(\mathsf{DYN}\) norm-continuous CPTP semigroup has GKLS/Lindblad form \\

No-go structure & uncertainty, complementarity, and no-cloning follow from noncommutativity and linearity \\
\bottomrule
\end{tabularx}
\TableNotes{This Part-I output ledger separates exact finite representation
theorems from residual application certificates.  In particular, the state,
evaluation, tensor, CP/Kraus, and supplied-flow forms hold for finite direct
sums; \(\mathsf{FENV}\) types a full carrier field envelope, while
\(\mathsf{ENV}\) is needed when later arguments use it as a full event
algebra.}
\end{table}



\paragraph{Supplied closed automorphism flows}
\label{p4:sec:part1-unitary-flow}


\begin{theorem}[Supplied closed automorphism flows are unitary]
\label{p4:thm:inner-flow}
\leavevmode\par\noindent
Let \(\{\alpha_t\}_{t\in\RR}\) be a norm-continuous one-parameter group of
\(*\)-automorphisms of a finite generated algebra
\[
\mathcal A=\bigoplus_{\lambda\in\Lambda}\End(\mathcal H_\lambda).
\]
Then every minimal central block is invariant, and there are self-adjoint
operators \(H_\lambda\in\End(\mathcal H_\lambda)\) such that, with
\(H=\bigoplus_\lambda H_\lambda\),
\begin{equation}
\label{p4:eq:inner-flow}
\alpha_t(a)=e^{itH}ae^{-itH}
\end{equation}
for all \(a\in\mathcal A\) and all \(t\in\RR\).  Thus the implementation
form is derived for a supplied \(\mathsf{DYN}\) flow.  The scalar parameter
orientation is imported upstream, and nontrivial kinematic flows already
exist; what is not derived from the algebra is the selection of this
particular flow as actual elapsed evolution, its physical scale, or its
trajectory.
\end{theorem}


\begin{proof}
\ProofPlan{We first fix the finite center, then apply innerness of derivations
on each full matrix block.}
\begin{proofsteps}

\item \textbf{The center is fixed.}
Every automorphism permutes the minimal central projections.  The induced map
from the connected group \(\RR\) to this finite permutation group is
continuous, hence constant.  Since \(\alpha_0=\operatorname{id}\), every
minimal central projection is fixed and each full matrix block is invariant.

\item \textbf{Block derivations are inner.}
On each block, norm continuity of the finite-dimensional one-parameter group
gives a derivation generator
\[
\delta_\lambda(a):=
\left.\frac{d}{dt}\right|_{t=0}\alpha_t(a),
\qquad a\in\End(\mathcal H_\lambda).
\]
Every derivation of a full matrix algebra is inner, and the \(*\)-condition
therefore gives a self-adjoint \(H_\lambda\) with
\[
\delta_\lambda(a)=i[H_\lambda,a].
\]

\item \textbf{Integration.}
The finite-dimensional operator ODE
\[
\frac{d}{dt}\beta_{\lambda,t}(a)
=i[H_\lambda,\beta_{\lambda,t}(a)],
\qquad
\beta_{\lambda,0}(a)=a,
\]
has unique solution
\[
\beta_{\lambda,t}(a)=e^{itH_\lambda}ae^{-itH_\lambda}.
\]
It coincides with \(\alpha_t\) on that block.  Taking the finite direct sum
over \(\lambda\) proves \eqref{p4:eq:inner-flow}.
\end{proofsteps}

\end{proof}



\paragraph{Open dynamics and GKLS form}
\label{p4:sec:part1-gkls}

\begin{theorem}[Supplied finite Markovian semigroups have GKLS form]
\label{p4:thm:gkls}
On a reduced-open-dynamics theorem domain, let \(\mathcal H\) be finite dimensional and let \(\{\Phi_t\}_{t\ge0}\) be a norm-continuous CPTP Markovian semigroup on states in \(\End(\mathcal H)\).  Then its generator has the form
\begin{equation}
\label{p4:eq:gkls}
\frac{d\rho}{dt}
=
-i[H,\rho]
+
\sum_\alpha\left(L_\alpha\rho L_\alpha^*-
\frac12\{L_\alpha^*L_\alpha,\rho\}\right)
\end{equation}
for some self-adjoint \(H\) and noise operators \(L_\alpha\).
\end{theorem}

This is the finite representation theorem for a supplied reduced Markovian semigroup.  The system--environment split and Markovian regime name the theorem domain; they are not premises of the closed envelope algebra.

\begin{proof}
Let \(\mathcal V=\End(\mathcal H)\).  Since \(\mathcal V\) is finite dimensional and \(t\mapsto\Phi_t\) is norm continuous, the semigroup has a bounded generator
\[
\mathcal L(\rho)=\lim_{t\downarrow0}\frac{\Phi_t(\rho)-\rho}{t}.
\]
For \(t>0\), Theorem~\ref{p4:thm:finite-choi-kraus-instruments} gives a Kraus representation of \(\Phi_t\).  Because \(\Phi_t\to\operatorname{id}\) in norm, the Choi matrix of \(\Phi_t\) converges to the rank-one Choi matrix of the identity.  Diagonalizing the positive Choi matrix in a gauge adapted to this rank-one block gives
\[
K_0(t)=\Id+tG+o(t),
\qquad
K_\alpha(t)=\sqrt t\,L_\alpha+o(\sqrt t).
\]
Expanding the Kraus sum gives
\[
\mathcal L(\rho)=G\rho+\rho G^*+\sum_\alpha L_\alpha\rho L_\alpha^*.
\]
Trace preservation gives
\[
G+G^*+\sum_\alpha L_\alpha^*L_\alpha=0.
\]
Writing the skew-adjoint part of \(G\) as \(-iH\), with \(H=H^*\), gives
\[
G=-iH-\frac12\sum_\alpha L_\alpha^*L_\alpha .
\]
Substitution yields \eqref{p4:eq:gkls}.  No unbounded-operator or domain subtlety is present in finite dimension.
\end{proof}



\paragraph{Uncertainty, complementarity, and no-cloning}
\label{p4:sec:part1-no-go}

Noncommutativity of any generated finite test algebra \(\mathcal A\)
immediately yields the Robertson inequality.  For self-adjoint
\(A,B\in\mathcal A\) and state \(\omega\),
\begin{equation}
\label{p4:eq:uncertainty}
\mathrm{Var}_\omega(A)\,\mathrm{Var}_\omega(B)
\ge
\frac14\,|\omega([A,B])|^2.
\end{equation}
This is the structural statement of complementarity: when
\([A,B]\ne0\), the two observables admit no joint spectral measure and no
complete common eigenbasis.  Particular common eigenvectors may still exist,
so noncommutation alone does not exclude every state that is sharp for both.

Likewise, linearity and tensorial composition imply the no-cloning obstruction.
There is no CPTP map that clones all input states,
\begin{equation}
\label{p4:eq:nocloning}
\Phi(\rho\otimes\sigma)=\rho\otimes\rho
\qquad
\text{for all }\rho.
\end{equation}
Indeed, if such a fixed blank state $\sigma$ and linear channel $\Phi$ cloned every density matrix, then for two distinct states $\rho_1,\rho_2$ and $0<p<1$ linearity would give
\[
\Phi\bigl((p\rho_1+(1-p)\rho_2)\otimes\sigma\bigr)
=
p\rho_1\otimes\rho_1+(1-p)\rho_2\otimes\rho_2,
\]
whereas cloning the mixture itself would require
\[
(p\rho_1+(1-p)\rho_2)\otimes(p\rho_1+(1-p)\rho_2),
\]
which contains the cross terms $p(1-p)(\rho_1\otimes\rho_2+\rho_2\otimes\rho_1)$.  These expressions disagree except in trivial restricted cases.



\paragraph{Symmetries, Hamiltonian Pictures, and Transport}
\label{p4:part:symmetry-dynamical-pictures}

\paragraph{Orientation.}
This part collects the symmetry, conserved-quantity, dynamical-picture, perturbative, and adiabatic transport results that organize finite quantum evolution once a Hamiltonian or supplied automorphism flow has been selected.



\subsection{Symmetries, generators, and conserved quantities}
\label{p4:sec:part2-symmetry}



\paragraph{Transition-probability preserving symmetries}
\label{p4:sec:part2-wigner}

Choose any finite Hilbert-space realization of the structural sector, for example a GNS presentation or the defining carrier \(\mathcal H_m\).
A physical symmetry of pure states is a bijection of the projective space \(\PP(\mathcal H_m)\) that preserves Born transition probabilities.
Equivalently, for rays \([\psi],[\phi]\in\PP(\mathcal H_m)\), one requires preservation of
\begin{equation}
\label{p4:eq:transition-prob}
P([\psi],[\phi])
=
\frac{|\langle\psi,\phi\rangle|^2}{\|\psi\|^2\,\|\phi\|^2}.
\end{equation}
Thus the physically meaningful content of symmetry is formulated directly in the projective language already native to the Hopf-phase structure.


\begin{theorem}[Finite Wigner theorem]
\label{p4:thm:wigner-finite}
Let \(\mathcal H_m\ne0\), put \(d=\dim\mathcal H_m\), and let
\(\mathsf S:\PP(\mathcal H_m)\to\PP(\mathcal H_m)\) be a bijection preserving
the transition probabilities \eqref{p4:eq:transition-prob}.  If \(d\ge2\),
then there exists a unitary or antiunitary operator \(W\) on
\(\mathcal H_m\), unique up to an overall phase, such that
\begin{equation}
\label{p4:eq:wigner-implementer}
\mathsf S([\psi])=[W\psi]
\end{equation}
for every nonzero \(\psi\in\mathcal H_m\).  Its semilinear type---unitary or
antiunitary---is also uniquely determined by \(\mathsf S\).

If \(d=1\), the projective space has one point.  Both unitary and antiunitary
implementers exist, so their type is not determined; within either fixed type
the implementer is still unique up to phase.
\end{theorem}


\begin{proof}
\ProofPlan{We normalize the images of an orthonormal basis, determine the
induced real isometry on each coordinate projective line, prove that its
identity/conjugation type is global, extend from a dense affine chart, and
then prove phase and type uniqueness.}
\begin{proofsteps}

\item \textbf{Basis normalization and coordinate projective lines.}
Assume first that \(d\ge2\).  Choose an orthonormal basis
\((e_1,\dots,e_d)\), and choose unit representatives \(f_j\) of
\(\mathsf S([e_j])\).  Since
\[
P([e_i],[e_j])=\delta_{ij},
\]
the \(f_j\) form an orthonormal basis.  Compose \(\mathsf S\) with the
projectivization of the unitary \(f_j\mapsto e_j\).  It is therefore enough
to prove the theorem when every coordinate ray is fixed.

For \(j\ge2\), the projective line
\[
L_j:=\PP(\operatorname{span}\{e_1,e_j\})
\]
is exactly the set of rays orthogonal to every \([e_k]\) with
\(k\notin\{1,j\}\).  Hence \(\mathsf S(L_j)=L_j\).  There is consequently a
unique bijection \(\sigma_j:\CC\to\CC\) such that
\[
\mathsf S([e_1+ze_j])=[e_1+\sigma_j(z)e_j].
\]
The phase of the original representative \(f_j\) may be chosen so that
\(\sigma_j(1)=1\).

\item \textbf{Each coordinate map is identity or conjugation.}
Comparison with the coordinate rays gives
\(|\sigma_j(z)|=|z|\).  Comparing the transition probabilities of
\([e_1+ze_j]\) and \([e_1+we_j]\), and expanding the squared numerator, gives
\[
\Re\!\bigl(\overline{\sigma_j(z)}\sigma_j(w)\bigr)
=\Re(\bar zw)
\qquad(z,w\in\CC).
\]
Thus \(\sigma_j\) is a surjective isometry of the real inner-product space
\(\CC\cong\RR^2\) fixing the origin.  The displayed identity and surjectivity
show directly that it is real linear: for every \(t\in\CC\), the vectors
\(\sigma_j(z+w)-\sigma_j(z)-\sigma_j(w)\) and
\(\sigma_j(rz)-r\sigma_j(z)\), \(r\in\RR\), have zero real inner product with
every \(\sigma_j(t)\), hence vanish.  A real-linear orthogonal map fixing
\(1\) sends \(i\) to \(i\) or \(-i\).  Therefore
\[
\sigma_j(z)=z
\qquad\text{or}\qquad
\sigma_j(z)=\bar z.
\]

\item \textbf{Coordinate reconstruction and global type.}
Consider a ray in the affine chart with nonzero first coordinate:
\[
\ell_z=[e_1+\textstyle\sum_{j=2}^d z_je_j].
\]
Its image has a unique representative
\([e_1+\sum_{j=2}^d w_je_j]\).  Transition probabilities with the coordinate
rays give \(|w_j|=|z_j|\).  Comparing with every
\([e_1+te_j]\) gives
\[
|1+\bar t z_j|
=
|1+\overline{\sigma_j(t)}w_j|
\qquad(t\in\CC),
\]
and hence \(w_j=\sigma_j(z_j)\).

The identity/conjugation choice is the same for all \(j\).  Indeed, if
\(\sigma_j=\operatorname{id}\) and
\(\sigma_k=\operatorname{conj}\), transition probability between the rays of
\[
e_1+e_j+e_k
\qquad\text{and}\qquad
e_1+ie_j+ie_k
\]
would change: the squared inner-product numerator is \(5\) before the map and
\(1\) afterward.  This is impossible.  For \(d=2\) there is only one
coordinate map, so no compatibility step is needed.

It follows that on this affine chart \(\mathsf S\) is induced either by the
identity or by coordinatewise conjugation.  Transition probability determines
the Fubini--Study distance
\[
d_{\rm FS}([\psi],[\phi])
=\arccos\sqrt{P([\psi],[\phi])},
\]
so \(\mathsf S\) is continuous.  The affine chart is dense, and the same
formula therefore holds on all rays.  Undoing the initial basis-normalizing
unitary yields a unitary or antiunitary implementer \(W\).


\item \textbf{Phase uniqueness, type uniqueness, and the one-dimensional case.}
Suppose first that \(W_1\) and \(W_2\) have the same semilinear type and
implement the same projective map.  Then
\(T:=W_2^{-1}W_1\) is linear and fixes every ray.  Writing
\(Te_j=\lambda_je_j\) and applying \(T\) to \(e_j+e_k\) gives
\(\lambda_j=\lambda_k\).  Hence
\[
T=\lambda\mathbf1,
\qquad
|\lambda|=1.
\]
If \(W_1\) and \(W_2\) had opposite types, then \(T\) would be antiunitary and
would still fix every ray.  For two orthonormal vectors \(e_1,e_2\), the rays
of \(e_1,e_2,e_1+e_2\) force
\(Te_1=\lambda e_1\) and \(Te_2=\lambda e_2\).  Antilinearity then gives
\[
T(e_1+ie_2)=\lambda e_1-i\lambda e_2,
\]
which is not collinear with \(e_1+ie_2\).  Thus opposite types cannot
implement the same ray map when \(d\ge2\).

Finally, if \(d=1\), \(\PP(\mathcal H_m)\) has one element.  Every unitary and
every antiunitary fixes it, so the type is not determined.  The quotient of
two implementers of the same type is a unitary scalar, proving phase
uniqueness within each type.
\end{proofsteps}

\end{proof}

This theorem classifies the Hilbert-space representation of a supplied
transition-probability-preserving ray symmetry.  It does not generate the
physical symmetry map \(\mathsf S\).  For \(d\ge2\), the ray redundancy
\(\psi\sim e^{i\theta}\psi\) is exactly the remaining phase ambiguity of its
implementer; in \(d=1\), semilinear type remains additionally undetermined.

\begin{remark}[Status of the finite Wigner theorem]
\label{p4:rem:wigner-status}
Theorem~\ref{p4:thm:wigner-finite} is a classical finite-dimensional result \cite{Wigner1959,Uhlhorn1963}.
The proof above records the coordinate-line compatibility and the
one-dimensional exception explicitly because later symmetry arguments
repeatedly use the theorem in sector-specific form.
\end{remark}



\paragraph{Projective group actions and one-parameter subgroups}
\label{p4:sec:part2-projective-groups}

If a group \(G\) acts by transition-probability preserving maps on rays and
\(\dim\mathcal H_m\ge2\), choose for each \(g\in G\) an implementing unitary
or antiunitary \(W(g)\) from Theorem~\ref{p4:thm:wigner-finite}.  Type
uniqueness defines
\[
\varepsilon:G\longrightarrow\mathbb Z_2,
\qquad
\varepsilon(g)=
\begin{cases}
0,&W(g)\text{ unitary},\\
1,&W(g)\text{ antiunitary}.
\end{cases}
\]
Semilinear types multiply, so \(\varepsilon\) is a homomorphism and
\(G_+:=\ker\varepsilon\) is the unitary subgroup.  For all \(g,h\in G\),
\begin{equation}
\label{p4:eq:projective-rep}
W(g)W(h)=\omega(g,h)\,W(gh),
\qquad
|\omega(g,h)|=1.
\end{equation}
On \(G_+\) this is an ordinary projective representation.  Globally,
associativity gives the twisted cocycle law
\[
\omega(g,h)\omega(gh,k)
=
\begin{cases}
\omega(h,k)\omega(g,hk),&\varepsilon(g)=0,\\
\overline{\omega(h,k)}\,\omega(g,hk),&\varepsilon(g)=1.
\end{cases}
\]
The multiplier \(\omega\) records only the unobservable central ray-phase
ambiguity of Hilbert-space implementers.  The degree-one character
\(\chi_1\) of Theorem~\ref{p4:thm:hopf-ray-phase} supplies a geometric
Hopf-to-ray comparison when that associated line is the theorem domain; it
does not identify an arbitrary projective multiplier with the noncentral
Hopf action on the whole carrier.
In particular, a continuous one-parameter subgroup through the identity lies
in \(G_+\) and is represented by unitary families; antiunitary realizations
occur only outside that identity component, for example in discrete
time-reversal-type symmetries.  In dimension one the ray action does not
determine \(\varepsilon\); a semilinear-type choice must be supplied before
this group-level statement is used.



\paragraph{One-parameter symmetries and self-adjoint generators}
\label{p4:sec:part2-symmetry-generators}

Textbook conserved charges arise when a continuous symmetry subgroup acts on
the finite carrier.  Let \(\{U(s)\}_{s\in\RR}\) be an already constituted
strongly continuous one-parameter unitary group on \(\mathcal H_m\).  In
finite dimension strong continuity is norm continuity, so
Corollary~\ref{p3a:cor:v149-finite-unitary-generator} yields, without a
separate Stone-type physical premise,
\begin{equation}
\label{p4:eq:Us-generator}
U(s)=e^{-isG}
\end{equation}
for a unique self-adjoint generator \(G\).  On observables,
\begin{equation}
\label{p4:eq:observable-symmetry-flow}
\beta_s(A)=U(s)AU(s)^*=e^{-isG}Ae^{isG},
\qquad
\frac{d}{ds}\beta_s(A)\Big|_{s=0}=-i[G,A].
\end{equation}
Thus every constituted continuous unitary group determines a Hermitian
infinitesimal generator.  Conversely, every self-adjoint \(G\) defines the
mathematical unitary group \(e^{-isG}\); it is a \emph{physical symmetry}
only on a theorem domain where that group preserves the complete
physical/action-owning presentation or, equivalently for the relevant
closed-dynamics claim, the generated or supplied dynamics.  The Part~I
Hamiltonian evolution is the special case in which the one-parameter group is
the selected closed-dynamics flow and the generator is the Hamiltonian.

\begin{remark}[Finite generator closure is internal; symmetry status is typed]
\label{p4:rem:v458-finite-generator-status}
The exponential/self-adjoint representation above is finite-dimensional
mathematics applied to an already constituted one-parameter unitary group; it
does not actualize that group, select its physical parameterization, or make
it a symmetry of an independently supplied Hamiltonian.  Those are distinct
\(\mathsf{DYN}\), calibration, and invariance claims.  On native same-owner
action branches, Paper~II-A's
Corollary~\ref{p2a:cor:v368-native-noether-antecedent} and its concrete
Paper~III-F realization may already derive the action-invariance antecedent.
On external/supplied \(H,U\) domains no such provenance inference is made.
\end{remark}



\paragraph{Constants of motion and the commutant criterion}
\label{p4:sec:part2-constants-of-motion}

The textbook notion of a conserved observable is immediate in the finite Born framework.
For a closed Hamiltonian flow
\begin{equation}
\label{p4:eq:state-flow-H}
\rho(t)=e^{-itH}\rho(0)e^{itH},
\end{equation}
the expectation of an observable \(Q\) is conserved when it is constant along the flow for every initial state.


\begin{theorem}[Conserved quantities and commutants]
\label{p4:thm:constants-of-motion}
On a closed-dynamics theorem domain, let \(H\) be the supplied self-adjoint finite generator and let \(Q\in\mathcal A_m^{\mathrm{env}}\) be self-adjoint.
The following are equivalent:
\begin{enumerate}[label=(\roman*), leftmargin=2em]
\item for every density matrix \(\rho\), the expectation \(\Tr(\rho(t)Q)\) is independent of \(t\),
\item the Heisenberg-evolved observable satisfies
\begin{equation}
\label{p4:eq:heisenberg-invariant}
Q(t):=e^{itH}Qe^{-itH}=Q
\end{equation}
for all \(t\),
\item \([H,Q]=0\).
\end{enumerate}
Equivalently, the conserved observables are precisely the self-adjoint elements of the commutant \(\{H\}'\subset\mathcal A_m^{\mathrm{env}}\).
\end{theorem}


\begin{proof}
Differentiating \(Q(t)=e^{itH}Qe^{-itH}\) gives
\begin{equation}
\label{p4:eq:heisenberg-derivative}
\frac{dQ(t)}{dt}=i[H,Q(t)].
\end{equation}
Hence \(Q(t)=Q\) for all \(t\) iff \([H,Q]=0\).
Using cyclicity of the trace,
\(
\Tr(\rho(t)Q)=\Tr(\rho(0)Q(t))
\),
so invariance of every expectation value is equivalent to invariance of \(Q(t)\) itself.
\end{proof}

\begin{corollary}[Finite Noether-I criterion for continuous symmetries]
\label{p4:cor:noether-finite}
Let \(U(s)=e^{-isG}\) be an already constituted one-parameter unitary group
and let \(H\) be the self-adjoint generator of a closed-dynamics theorem
domain on the same carrier.  Suppose the dynamics is invariant under the
group,
\begin{equation}
\label{p4:eq:H-symmetry-invariance}
U(s)HU(s)^*=H
\qquad
\text{for all }s\in\RR.
\end{equation}
Then
\[
 \boxed{[H,G]=0,}
\]
and therefore \(G\) is a conserved observable.  Conversely, if
\([H,G]=0\), then the unitary group generated by \(G\) commutes with the
closed time evolution.  Calling that group a physical symmetry still requires
its stated physical theorem-domain provenance.
\end{corollary}

\begin{proof}
Differentiating \(U(s)HU(s)^*=H\) at \(s=0\) gives
\(-iGH+iHG=0\), hence \([H,G]=0\).  Conversely, if \([H,G]=0\), then
\(G\) commutes with every polynomial in \(H\) and therefore with
\(e^{-itH}\) by finite-dimensional functional calculus.
Theorem~\ref{p4:thm:constants-of-motion} then makes \(G\) conserved.  The
final typing clause is not part of the commutator algebra: a mathematical
unitary group becomes a physical symmetry only on the declared
symmetry/action domain.
\end{proof}

\begin{remark}[Native versus supplied Noether antecedent]
\label{p4:rem:v458-native-supplied-noether}
Corollary~\ref{p4:cor:noether-finite} itself is an algebraic conditional and
does not postulate conservation.  On a native same-owner action branch,
Paper~II-A's same-owner unmarked-automorphism theorem and native Noether
antecedent, concretely realized for the local \(A+\Phi+g\) action in
Paper~III-F, may already supply the action-invariance relation whose finite
Hamiltonian representative is \eqref{p4:eq:H-symmetry-invariance}, provided
the required Hamiltonian/Legendre or closed-dynamics bridge has also been
certified.  For an externally supplied \(H\) and \(U\), or when the group
changes a generated source, coupling, boundary datum, orientation, or other
retained mark, invariance remains an explicit hypothesis and need not hold.
Thus not every generated operation is a Noether symmetry.
\end{remark}

\begin{remark}[Noether-I/Noether-II typing boundary]
\label{p4:rem:v176-noether-one-two-boundary}
Corollary~\ref{p4:cor:noether-finite} is the global one-parameter/Noether-I
logic: a continuous physical symmetry of the closed dynamics yields a
conserved generator.  The local gauge and diffeomorphism automorphisms of the
Paper-III-F \(A+\Phi+g\) action instead enter the same-owner variational
closure through Noether II and primarily yield Euler--Lagrange identities,
moment maps, and constraints.  Gauge redundancy therefore must not be
promoted to an independent physical charge merely by applying the finite
Noether-I criterion.
\end{remark}

This is the finite-resolution form of the Noether-I commutant logic.  It
requires no underlying classical phase space, but it does require the typed
symmetry-of-dynamics antecedent; the existence of a unitary group alone is not
that antecedent.  Rotational invariance is the prime example.  When
\([H,\hat J_i]=0\) for the diagonal generators introduced in
Section~\ref{p4:sec:part2-angular}, the corresponding angular-momentum
components are conserved; when only
\([H,\hat J^2]=[H,\hat J_z]=0\), the carrier decomposes into
\(\ket{J,M}\) multiplets.



\paragraph{Good quantum numbers and symmetry sectors}
\label{p4:sec:part2-good-quantum-numbers}

If a commuting family of observables \(H,Q_1,\dots,Q_r\) is given, finite-dimensional spectral theory yields a simultaneous orthogonal decomposition of the carrier into joint eigenspaces.
The corresponding eigenlabels are the framework version of textbook good quantum numbers.
Symmetry therefore organizes the carrier into blocks of fixed conserved data, while degeneracies within a block are explained by residual multiplet structure.
In particular, the bounded angular-momentum sectors derived later are not additional decorations: they are the rotation-symmetry decomposition of the same finite carrier.




\subsection{Stationary states and dynamical pictures}
\label{p4:sec:part2-dynamical-pictures}



\paragraph{Finite spectral resolution of the Hamiltonian}
\label{p4:sec:part2-spectral-resolution}

On a closed-dynamics theorem domain, let $H\in\mathcal A_m^{\mathrm{env}}$ be the supplied self-adjoint finite generator.
Because the carrier is finite-dimensional, the spectral theorem is exact and purely algebraic: there is a finite set of real eigenvalues $\Spec(H)=\{E_\alpha\}$ and pairwise orthogonal projectors $P_\alpha$ such that
\begin{equation}
\label{p4:eq:H-spectral-resolution}
H=\sum_{\alpha} E_\alpha P_\alpha,
\qquad
P_\alpha P_\beta=\delta_{\alpha\beta}P_\alpha,
\qquad
\sum_\alpha P_\alpha=\mathbf 1.
\end{equation}
Consequently the exact unitary flow is
\begin{equation}
\label{p4:eq:U-spectral-resolution}
U(t):=e^{-itH}=\sum_\alpha e^{-itE_\alpha}P_\alpha.
\end{equation}
All textbook energy decompositions therefore appear internally once the closed generator is supplied by the theorem domain: there is no need for additional continuum assumptions to obtain spectral projectors, energy measurements, or phase evolution, and the spectral projectors do not select the upstream compact sector.


\begin{theorem}[Finite spectral evolution and Bohr frequencies]
\label{p4:thm:finite-spectral-evolution}
For every observable $A\in\mathcal A_m^{\mathrm{env}}$, its Heisenberg evolution under $H$ admits the exact expansion
\begin{equation}
\label{p4:eq:A-heis-bohr}
A_H(t):=e^{itH}Ae^{-itH}=\sum_{\alpha,\beta} e^{it(E_\alpha-E_\beta)}P_\alpha A P_\beta.
\end{equation}
Hence all transition phases are built from the finite set of Bohr frequencies
\begin{equation}
\label{p4:eq:bohr-frequencies}
\omega_{\alpha\beta}=E_\alpha-E_\beta.
\end{equation}
In particular, the diagonal blocks $P_\alpha A P_\alpha$ are time-independent, while off-diagonal coherences oscillate with frequencies fixed by energy differences.
\end{theorem}


\begin{proof}
Insert the spectral decomposition \eqref{p4:eq:U-spectral-resolution} on both sides of $A$ and use orthogonality of the projectors.
\end{proof}



\paragraph{Stationary states and energy sectors}
\label{p4:sec:part2-stationary-states}

A density matrix $\rho$ is stationary when it is invariant under the Hamiltonian flow,
\begin{equation}
\label{p4:eq:rho-stationary-def}
\rho(t)=e^{-itH}\rho e^{itH}=\rho
\qquad\text{for all }t.
\end{equation}
Pure energy eigenvectors are the simplest examples, but the finite framework also allows mixed stationary states supported on degenerate eigenspaces.


\begin{theorem}[Characterization of stationary states]
\label{p4:thm:stationary-states}
On the same closed-dynamics theorem domain, let $H$ be self-adjoint on the finite carrier.
For a density matrix $\rho$, the following are equivalent:
\begin{enumerate}[label=(\roman*), leftmargin=2em]
\item $\rho$ is stationary in the sense of \eqref{p4:eq:rho-stationary-def},
\item $[H,\rho]=0$,
\item $\rho$ is block-diagonal in the spectral decomposition of $H$,
\begin{equation}
\label{p4:eq:rho-energy-blocks}
\rho=\sum_\alpha P_\alpha\rho P_\alpha.
\end{equation}
\end{enumerate}
A pure ray $[\psi]$ is stationary precisely when $\psi$ lies in a single energy eigenspace of $H$.
Within a degenerate eigenspace, every density matrix is stationary.
\end{theorem}


\begin{proof}
The equivalence of (i) and (ii) is the state-side version of the adjoint invariance criterion.
Using \eqref{p4:eq:H-spectral-resolution}, one has
\[
[H,\rho]=\sum_{\alpha,\beta}(E_\alpha-E_\beta)P_\alpha\rho P_\beta,
\]
which vanishes iff all off-diagonal blocks with $E_\alpha\neq E_\beta$ vanish.
That is exactly the block-diagonal condition \eqref{p4:eq:rho-energy-blocks}.
For a pure state $\rho_\psi=\lvert\psi\rangle\!\langle\psi\rvert$ this means that $\psi$ is supported in a single eigenspace.
\end{proof}

Thus textbook stationary spectroscopy survives unchanged in finite form.
The only structural difference from the infinite-dimensional presentation is that the energy resolution is always discrete and finite at fixed truncation.



\paragraph{Schr\"odinger and Heisenberg pictures}
\label{p4:sec:part2-sch-heis}

The two standard dynamical pictures are exactly equivalent formulations of the same adjoint Hamiltonian flow.
In the Schr\"odinger picture the observables are fixed and the states evolve,
\begin{equation}
\label{p4:eq:sch-picture}
\rho_S(t)=U(t)\rho_S(0)U(t)^*,
\qquad
A_S(t)=A.
\end{equation}
In the Heisenberg picture the state is fixed and observables evolve,
\begin{equation}
\label{p4:eq:heis-picture}
\rho_H=\rho_S(0),
\qquad
A_H(t)=U(t)^*AU(t)=e^{itH}Ae^{-itH}.
\end{equation}
The equality of predictions is immediate from trace cyclicity:
\begin{equation}
\label{p4:eq:picture-equivalence}
\Tr(\rho_S(t)A)=\Tr(\rho_H A_H(t)).
\end{equation}
Differentiating the Heisenberg observable gives the usual equation of motion
\begin{equation}
\label{p4:eq:heisenberg-equation}
\frac{dA_H}{dt}=i[H,A_H],
\end{equation}
with $\hbar=1$ in the present convention; restoring dimensions yields $\dot A_H=(i/\hbar)[H,A_H]$.
Thus the equivalence between evolving states and evolving observables follows from the finite algebraic formalism.



\paragraph{Interaction picture and perturbing a reference Hamiltonian}
\label{p4:sec:part2-interaction-picture}

Suppose first that the Hamiltonian is split as
\begin{equation}
\label{p4:eq:H-split}
H=H_0+V,
\end{equation}
where $H_0$ is a solvable reference part and $V$ is a time-independent interaction.  The time-dependent case is obtained by replacing the full exponential below by the exact propagator of $H_0+V(t)$.
The interaction picture is defined by removing the exact $H_0$ flow:
\begin{equation}
\label{p4:eq:interaction-picture-def}
\rho_I(t):=e^{itH_0}\rho_S(t)e^{-itH_0},
\qquad
A_I(t):=e^{itH_0}Ae^{-itH_0}.
\end{equation}
Its propagator in the time-independent case is
\begin{equation}
\label{p4:eq:UI-def}
U_I(t):=e^{itH_0}e^{-itH}.
\end{equation}
More generally, if $U(t)$ solves $i\dot U(t)=(H_0+V(t))U(t)$, then the interaction-picture propagator is
\[
U_I(t):=e^{itH_0}U(t),
\]
so that $\rho_I(t)=U_I(t)\rho_I(0)U_I(t)^*$.
The interaction Hamiltonian is the Heisenberg transform of the interaction with respect to $H_0$,
\begin{equation}
\label{p4:eq:VI-def}
V_I(t):=e^{itH_0}V(t)e^{-itH_0},
\end{equation}
with $V(t)=V$ in the time-independent case,
and one obtains the finite interaction-picture evolution equation
\begin{equation}
\label{p4:eq:UI-evolution}
i\frac{dU_I}{dt}=V_I(t)U_I(t),
\qquad
U_I(0)=\mathbf 1.
\end{equation}
Because the carrier is finite-dimensional, this equation always has a unique global solution and the corresponding Dyson series converges absolutely whenever $V_I(t)$ is norm-continuous, or at least norm-integrable, on the finite time interval under consideration.
Once the split \(H=H_0+V\) is retained, the interaction picture is another
generated presentation of the same co-realized dynamics, not a third process
occurrence.  The split itself, and any claim that \(H_0\) is the physically
preferred reference, is not produced by picture duality.

\begin{corollary}[Commuting split]
\label{p4:cor:commuting-split}
If $[H_0,V]=0$, then $V_I(t)=V$ is constant and
\begin{equation}
\label{p4:eq:commuting-factorization}
e^{-itH}=e^{-itH_0}e^{-itV}=e^{-itV}e^{-itH_0}.
\end{equation}
Hence the interaction picture becomes exact factorization rather than an approximation scheme.
\end{corollary}

\begin{proof}
If \([H_0,V]=0\), then \(e^{itH_0}Ve^{-itH_0}=V\), so the interaction-picture generator is constant.  Since commuting finite matrices have commuting exponentials, \(e^{-it(H_0+V)}=e^{-itH_0}e^{-itV}=e^{-itV}e^{-itH_0}\).
\end{proof}


This is the precise finite-resolution setting in which time-dependent perturbation theory is later formulated.
The framework therefore contains the full textbook triad of dynamical pictures --- Schr\"odinger, Heisenberg, and interaction --- before any appeal to continuum wave mechanics.




\subsection{Generated process rigidity on simple finite event algebras}
\label{p4:sec:generated-process-rigidity-update}

Let $\tau(A)=d^{-1}\Tr A$ on $M_d(\CC)$ and define
$P_1(A)=\tau(A)I$, $P_0=I-P_1$.

\begin{theorem}[Automorphism-neutral primitive bistochastic semigroup]
\label{p4:thm:neutral-primitive-qms-rigidity}
Let $(T_t)_{t\ge0}$ be norm-continuous, unital, completely positive and trace
preserving, and suppose
\[
 T_t(UAU^*)=UT_t(A)U^*
 \qquad(U\in U(d)).
\]
If one nonzero traceless observable decays, then for some $\kappa>0$,
\[
 \boxed{T_t=P_1+e^{-\kappa t}P_0.}
\]
The dual semigroup has unique stationary state $I/d$.
\end{theorem}

\begin{proof}
Under conjugation, $M_d(\CC)=\CC I\oplus M_d(\CC)_0$ with irreducible
traceless adjoint module.  Schur's lemma gives $T_t=P_1+b(t)P_0$.
The semigroup law gives $b(s+t)=b(s)b(t)$; continuity and complete-positive
contractivity give $b(t)=e^{-\kappa t}$, and the decay witness makes
$\kappa>0$.
\end{proof}

\begin{corollary}[Neutral quantum semigroup as a universal-continuation instance]
\label{p4:cor:neutral-qms-universal-continuation}
On the theorem domain of
Theorem~\ref{p4:thm:neutral-primitive-qms-rigidity}, the continuation role is
the completely positive bistochastic process on
\(\CC I\oplus M_d(\CC)_0\), the unmarked symmetry is conjugation by
\(U(d)\), and the admissible generator cone on the nontrivial irreducible
module is the single ray
\[
 \RR_{\ge0}(-P_0).
\]
The decay witness excludes its zero endpoint.  Hence
Theorem~\ref{p3a:thm:finite-generator-one-ray-rigidity} recovers the process
law up to the positive clock calibration \(\kappa\).  The corollary identifies
the law once this typed process role is generated and occurring; it does not
create an actual process token from the matrix algebra alone.
\end{corollary}

\begin{theorem}[Generated reversible mixing graph and common stationary state]
\label{p4:thm:generated-reversible-mixing-graph}
Let \(\Gamma=(V,E)\) be a finite generated undirected graph, let
\(\mu_i>0\) be a generated vertex weight, and let \(G\) be its generated
unmarked automorphism group preserving \(\mu\).  On a declared reversible
Markov-process domain, every strictly positive \(G\)-invariant conductance
family
\[
 c_{ij}=c_{ji}>0\qquad(\{i,j\}\in E)
\]
defines
\[
 (L_cf)(i)=\frac1{\mu_i}\sum_{j:\{i,j\}\in E}
 c_{ij}\bigl(f(j)-f(i)\bigr),
 \qquad T_t^{(c)}=e^{tL_c}.
\]
Then \(T_t^{(c)}\) is a positive unital contraction semigroup, is reversible
in \(\ell^2(V,\mu)\), and has the normalized stationary measure
\[
 \boxed{\pi_i=\frac{\mu_i}{\sum_{k\in V}\mu_k}.}
\]
If \(\Gamma\) is connected, this stationary probability is unique for every
such conductance family.  If \(E/G\) has \(q\) edge orbits, the admissible
strictly positive generators form a \(q\)-parameter cone; after quotienting
the common time calibration, \(q-1\) conductance ratios remain when \(q>1\).
For \(q=1\), the generator is a single positive ray and the process law is
unique up to clock calibration.  If \(\Gamma\) has \(r>1\) connected
components, the stationary probabilities form an \((r-1)\)-simplex of
component weights, with the relative weights inside each component fixed by
\(\mu\).
\end{theorem}

\begin{proof}
For the weighted inner product
\(\langle f,g\rangle_\mu=\sum_i\mu_i\overline{f(i)}g(i)\), symmetry of the
conductances gives
\[
 \langle f,-L_cf\rangle_\mu
 =\sum_{\{i,j\}\in E}c_{ij}|f(i)-f(j)|^2
 =\frac12\sum_i\sum_{j:\{i,j\}\in E}
       c_{ij}|f(i)-f(j)|^2\ge0.
\]
Thus \(L_c\) is self-adjoint and dissipative, while its off-diagonal rates
are nonnegative and \(L_c1=0\); finite-dimensional Markov-semigroup theory
gives the claimed positive unital contractions.  Detailed balance reads
\(\pi_i(L_c)_{ij}=c_{ij}/\sum_k\mu_k=\pi_j(L_c)_{ji}\).
The kernel of the displayed Dirichlet form consists exactly of functions
constant on connected components, proving both stationary statements.
Finally, a \(G\)-invariant conductance is constant on each edge orbit, so the
cone and calibration counts follow.
\end{proof}

\begin{corollary}[State information common to an unresolved process family]
\label{p4:cor:reversible-family-common-stationary-state}
Under the hypotheses of
Theorem~\ref{p4:thm:generated-reversible-mixing-graph}, assume the exact
continuation family is generated and nonempty.  If the graph is connected,
the stationary-state readout is constant on the entire conductance moduli
space.  Therefore
Corollary~\ref{p3a:cor:common-process-invariant} derives \(\pi\) at role level
even when \(q>1\) and the dynamical law remains unresolved.  Convergence of a
particular physical history still requires an occurring process member and
its actual history.  In particular, the current \(11\oplus4\) central grading
does not satisfy this theorem without a generated inter-block transition
graph and positive edge support.
\end{corollary}

\begin{corollary}[Transition-typed \(11\oplus4\) one-edge relaxation]
\label{p4:cor:eleven-four-one-edge-relaxation}
Assume the generated central blocks have vertex measures
\(\mu_+=11\), \(\mu_-=4\), and that a constituted physical
\emph{transition/process} bridge supplies one positive undirected edge between
them on the reversible Markov domain of
Theorem~\ref{p4:thm:generated-reversible-mixing-graph}.  Then, for one
conductance \(c>0\), the generator on central observables is
\[
 L_c=
 \begin{pmatrix}
  -c/11 & c/11\\
  c/4 & -c/4
 \end{pmatrix}.
\]
It has eigenvalues \(0\) and \(-15c/44\).  Hence, writing
\(\kappa=15c/44\), the law is the unique nonstationary reversible one-edge
law up to clock calibration, its stationary probability is
\[
 \boxed{\pi=\left(\frac{11}{15},\frac4{15}\right),}
\]
and every occurring probability history satisfies
\[
 p_+(t)=\frac{11}{15}
  +e^{-\kappa t}\left(p_+(0)-\frac{11}{15}\right),
 \qquad
 p_-(t)=1-p_+(t).
\]
The bare grading and a merely field-typed central-odd bridge do not satisfy
the transition/process hypothesis.
\end{corollary}

\begin{proof}
The graph has two vertices and one edge orbit, so
Theorem~\ref{p4:thm:generated-reversible-mixing-graph} gives one positive
conductance ray and stationary weights proportional to \(11,4\).  Substitution
in its generator formula gives the displayed matrix.  Its rows sum to zero,
its determinant vanishes, and its trace is \(-c/11-c/4=-15c/44\), proving the
spectrum.  The dual probability line orthogonal to \(\pi\) therefore decays
by \(e^{-\kappa t}\), which gives the formulas.  The final boundary follows
from the theorem's declared process type and from the field/event distinction
in Theorem~\ref{p4:thm:test-graph-env-discharge}.
\end{proof}

% BEGIN FRAMEWORK v2.264: OBSERVABLE-MINIMAL PAIRED EXCHANGE AND CONVERSES
\subsection{Paired-exchange realization and exact observable converses}
\label{p4:sec:v264-paired-population}
\noindent
The following mathematical realization strengthens the \emph{conditional}
transition theorem above without upgrading the existing native test algebra
to $M_{15}$ or postulating an individually physical central-odd effect.
It imports the canonical \emph{two-copy operator candidate} of
Theorem~\ref{p2b:thm:v264-canonical-paired-exchange}; actual exchange
remains a separate action/process occurrence obligation.

\begin{theorem}[Conditional exact $11/15$ paired-exchange population law]
\label{p4:thm:v264-paired-population-law}
Let an actually licensed, identified, relation-free two-carrier occurrence
consume $H_{\times}=\hbar\kappa K$, with $K$ the mathematical
exchange of Theorem~\ref{p2b:thm:v264-canonical-paired-exchange}.
Assume a corresponding dynamical token licenses
$U_\theta=e^{-i\theta K}$ with $\theta=\kappa t$, and assume
that the second carrier is \emph{actually prepared} in
$\sigma_E=I_{15}/15$.  Then for every full-carrier density
$\rho$ (including off-block coherence), the resulting reduced CPTP
channel $\Phi_\theta$ satisfies
\[
\boxed{\operatorname{Tr}(P_g\Phi_\theta(\rho))
  =\cos^2\theta\operatorname{Tr}(P_g\rho)
   +\sin^2\theta\frac{11}{15}},
\qquad
\Phi_\theta(\rho)=\operatorname{Tr}_E\!\left[
 U_\theta(\rho\otimes\sigma_E)U_\theta^*\right].
\]
In particular, $\theta=\pi/2$ yields $p_g'=11/15$ for every input
population, \emph{not} necessarily the full output density $I_{15}/15$.
For $n$ independent, equally prepared collisions the population is
\[
 p_g^{(n)}=\frac{11}{15}+
  \bigl(p_g^{(0)}-\tfrac{11}{15}\bigr)\cos^{2n}\theta.
\]
If an independent occurring Poisson collision law has rate $r>0$,
then $p_g(t)=11/15+(p_g(0)-11/15)e^{-r\sin^2\theta\,t}$.
The mathematical existence of $K$ \emph{does not} establish any of
these action, preparation, or process antecedents.
\end{theorem}
\begin{proof}
From $K^3=K$ and $K^2=Q_{\times}$ follows
$U_\theta=I+(\cos\theta-1)Q_{\times}-i\sin\theta K$.
In the bath trace state, the bath block probabilities are
$q_g=11/15$ and $q_s=4/15$. The crossed sector
$\mathcal H_g\otimes\mathcal H_s$ loses
$p_gq_s\sin^2\theta$ of system-geometric population, whereas
$\mathcal H_s\otimes\mathcal H_g$ gains
$(1-p_g)q_g\sin^2\theta$; subtracting and adding gives the box.
Terms linear in off-block coherence vanish because the bath and the
measured projector are central-even and the channel intertwines the
central conjugation. Partial trace of unitary conjugation with a
density bath is CPTP; since the bath state is scalar and $U_\theta$
commutes with the diagonal $\Spin(4)$ action, it is covariant.
Iterating $p\mapsto q_g+\cos^2\theta(p-q_g)$ proves the $n$-collision
formula. Poisson averaging uses
$\mathbb E(a^{N_t})=\exp(rt(a-1))$ at $a=\cos^2\theta$.
\end{proof}

\begin{proposition}[All-state population conservation: Hamiltonian and Kraus converses]
\label{p4:prop:v264-population-converse}
Let $P_g$ be any fixed orthogonal block projector. For a
self-adjoint Hamiltonian $H$ on a finite composite carrier,
$\operatorname{Tr}((P_g\otimes I)e^{-itH}\rho e^{itH})$
is constant in $t$ for \emph{every} joint density $\rho$ if and only if
\[
 \boxed{[H,P_g\otimes I]=0.}
\]
For a CPTP map $\Phi(\rho)=\sum_jL_j\rho L_j^*$ on a \emph{single}
carrier, $\operatorname{Tr}(P_g\Phi(\rho))=\operatorname{Tr}(P_g\rho)$
for every $\rho$ if and only if all its Kraus operators are block
diagonal: $P_gL_jP_s=P_sL_jP_g=0$ for every $j$.
Thus even test algebras may coexist with population-changing dynamics,
but Hamiltonians in $\mathcal A_{\rm ev}\otimes\mathcal A_{\rm ev}$
conserve the first-system central population identically.
\end{proposition}
\begin{proof}
The Hamiltonian derivative at $t=0$ is
$i\operatorname{Tr}([H,P_g\otimes I]\rho)$; vanishing for every
density forces the commutator to vanish, and the converse follows
by exponentiation. For the channel, insert a pure input in
$\operatorname{Ran}(P_s)$. Its initial geometric population is zero,
while its final population is the nonnegative sum
$\sum_j\|P_gL_j\psi_s\|^2$; each term must vanish. With a pure
$\operatorname{Ran}(P_g)$ input, trace preservation and conservation
of geometric population force
$\sum_j\|P_sL_j\psi_g\|^2=0$.
These are precisely the vanishing off-block Kraus entries; the
converse is immediate.
\end{proof}

\begin{proposition}[Explicit covariant GKLS alternative; not action selection]
\label{p4:prop:v264-covariant-gkls-alternative}
For every \emph{chosen} $\gamma>0$, the following mathematical generator
on $M_{15}(\CC)$ is of GKLS form:
\[
 \boxed{\mathcal L_\gamma(\rho)=\gamma\left[
 \operatorname{Tr}(P_g\rho)P_s+
 \operatorname{Tr}(P_s\rho)P_g-
 \tfrac12\{4P_g+11P_s,\rho\}\right].}
\]
It is trace preserving, unital and covariant under
$U(11)\times U(4)$, hence under the represented $\Spin(4)$.
Its unique stationary full density is $I_{15}/15$, and
\[
 \dot p_g=\gamma(11-15p_g),\qquad
 p_g(t)=\frac{11}{15}+
 (p_g(0)-\tfrac{11}{15})e^{-15\gamma t}.
\]
This construction demonstrates mathematical compatibility of a
population-changing \emph{process} with a central-even native
\emph{test algebra}; it does not assert that the common action
selects $\gamma$, a bath, jump operators, or this process.
\end{proposition}
\begin{proof}
Choose orthonormal bases $g_i$ ($i=1,\ldots,11$) and $s_j$
($j=1,\ldots,4$). The $88$ jumps
$\sqrt\gamma|s_j\rangle\langle g_i|$ and
$\sqrt\gamma|g_i\rangle\langle s_j|$ give the displayed GKLS
operator after summation. Its formula is basis independent and
$\mathcal L_\gamma(I)=0$; trace preservation follows because the
jump-rate operator is $4\gamma P_g+11\gamma P_s$.
The geometric population obeys the displayed first-order law.
Within the geometric traceless, protected traceless, off-diagonal
and central-population-difference subspaces the strictly positive
decay rates are respectively $4\gamma$, $11\gamma$,
$15\gamma/2$ and $15\gamma$. Only the normalized identity is fixed,
proving unique stationarity.
\end{proof}

\begin{theorem}[Complete native protocol blindness under covariant instruments]
\label{p4:thm:v264-all-history-native-blindness}
Let $\Theta(\rho)=\Gamma_c\rho\Gamma_c$ and
$\mathcal E=(\operatorname{id}+\Theta)/2$. For a finite sequence
of CP instrument-outcome maps $\mathcal I_{o_k}$ satisfying
$\mathcal I_{o_k}\Theta=\Theta\mathcal I_{o_k}$ and a final
even effect $F=\Theta(F)$, every finite history probability is
unchanged upon replacing the initial density $\rho$ by
$\mathcal E(\rho)$:
\[
\boxed{\operatorname{Tr}\!\left[
 F\mathcal I_{o_n}\cdots\mathcal I_{o_1}(\rho)\right]
=\operatorname{Tr}\!\left[
 F\mathcal I_{o_n}\cdots\mathcal I_{o_1}(\mathcal E(\rho))
\right].}
\]
Thus even arbitrarily long \emph{finite} native covariant protocols
cannot distinguish the $88$ real off-block coherence coordinates.
This does not prohibit a new, independently licensed parity-breaking
instrument or reference, or a different test algebra.
\end{theorem}
\begin{proof}
Every instrument map commutes with $\mathcal E$; induction moves
$\mathcal E$ through any finite composition. Since $F$ is even,
$\operatorname{Tr}(FA)=\operatorname{Tr}(F\mathcal E(A))$ for all
$A$, proving the identity. On Hermitian matrices the kernel of
$\mathcal E$ consists exactly of the two conjugate $11\times4$
rectangular blocks, a real vector space of dimension $88$.
\end{proof}

\begin{remark}[The exact pending physical contraction]
\label{p4:rem:v264-paired-action-firewall}
Under the printed common action of Paper~III-F, the locally uninserted
fermion term is bilinear $\bar\Psi D(B;Y)\Psi$; on its zero-field
branch the first spinorial variation is zero, while its second
variation is the paired-index fermion operator. Paper~V-A's nonzero
rooted \emph{chiral} coefficient is central-even under the
momentum $11/4$ grading and has a different operator owner
(Proposition~\ref{p2b:prop:v264-chiral-central-distinction}).
Neither supplies an already evaluated $\mathcal H_1\otimes\mathcal H_1$
action Hessian. If an action owner later supplies a normalized
self-adjoint two-carrier Hessian $H_{\rm act}$, the observable-targeted
test is
\[
 \boxed{\mathcal X_{\rm pair}(S_*)=
 \left\|(P_s\otimes P_g)H_{\rm act}
 (P_g\otimes P_s)\right\|_{\rm HS}^2.}
\]
A positive value certifies a nonzero crossed block, \emph{not}
the special $K$ law, environment preparation, temporal rate, or a
completed $\mathsf{PROC}$/$\mathsf{DYN}$ incidence. Even $\mathcal
X_{\rm pair}>0$ need not by itself prove population relaxation for
a particular prepared history. This expression remains an
\emph{unexecuted same-action target}, not an evaluated certificate
in the current common action. The full field-envelope
$\mathsf{FENV}$ still independently requires a typed central-odd
one-system bridge; it is unnecessary for the paired population
observable under the stated alternative process antecedents.
\end{remark}
% END FRAMEWORK v2.264: OBSERVABLE-MINIMAL PAIRED EXCHANGE AND CONVERSES

\begin{proposition}[Modular blindness to central weights]
\label{p4:prop:modular-central-weight-blindness}
For
\[
 \mathcal A=\bigoplus_iM_{d_i}(\CC),
 \qquad
 \rho=\bigoplus_i p_i\sigma_i,
\]
with \(p_i>0\) and faithful block densities \(\sigma_i>0\), the modular
automorphism group is
\[
 \rho^{it}A\rho^{-it}
 =\bigoplus_i\sigma_i^{it}A_i\sigma_i^{-it},
\]
and is independent of every central weight $p_i$.  Modular flow therefore
cannot by itself select the missing block probabilities.

For a nonfaithful state, let \(s_i=\mathbf1_{(0,\infty)}(\sigma_i)\) and
use the support algebra
\(\bigoplus_i s_iM_{d_i}(\mathbb C)s_i\), omitting zero-weight blocks.
The restricted density is faithful there and the same cancellation holds.
This support restriction is a statement about the state representation;
it does not delete a direction used by a different claim-complete owner
or define an invertible modular automorphism on the original full algebra.
\end{proposition}

\begin{proof}
On each faithful block,
\((p_i\sigma_i)^{it}=p_i^{it}\sigma_i^{it}\) by positive spectral
calculus, and the two scalar factors cancel in conjugation.  On its support,
a nonzero finite positive matrix has strictly positive eigenvalues, so the
same calculation applies after restriction.  If a block has a kernel,
its negative imaginary powers are not invertible on the full block;
positivity of \(p_i\) alone cannot remove that kernel.
\end{proof}

\begin{theorem}[Finite simple-algebra KMS uniqueness]
\label{p4:thm:finite-simple-kms-uniqueness}
For the inner dynamics $\alpha_t(A)=e^{itH}Ae^{-itH}$ on $M_d(\CC)$, a
normalized finite $\beta$-KMS density is uniquely
\[
 \rho_\beta=\frac{e^{-\beta H}}{\Tr e^{-\beta H}}.
\]
On a direct sum the same calculation determines the Gibbs density only within
each block and leaves the positive central mixture free.
\end{theorem}

\begin{proposition}[Effects do not determine instruments]
\label{p4:prop:effect-instrument-nonselection}
For a projection $P_i$, both
\[
 \mathcal I_i(\rho)=P_i\rho P_i
 \quad\text{and}\quad
 \mathcal J_i(\rho)=\Tr(P_i\rho)\sigma_i
\]
are completely positive operations with the same outcome probability
$\Tr(P_i\rho)$, but generically have different post-measurement states.  Thus
a POVM/effect family does not select an apparatus instrument; a generated
pointer/Stinespring process is required.
\end{proposition}



% BEGIN v2.36 STATE/SCATTERING ATLAS


\section{Rays, projective phases and finite sector representations}
\label{p4a:sec:v238-a6}
\noindent\textit{Source and claim boundary.} Finite GNS, Hopf phase, statistics, central blocks and conditional event extensions are representational consequences of supplied state and exchange data, not a new continuum field ontology.

\paragraph{Representation Languages and Finite Many-Body Sectors}
\label{p4:part:derived-sector}

\paragraph{Orientation.}
This part starts the representation/interface sequence.  It records how the closed finite algebra appears as rays, phases, wavefunction presentations, tensor powers, statistics sectors, and superselection blocks before finite-field or scattering-specific bridges are introduced.



\subsection{Representation and finite-QFT bridge convention}
\label{p4:sec:part2-purpose}



\paragraph{Purpose and bridge role}
\label{p4:sec:part2-purpose-role}

This part begins the representation and interface sequence built from the closed quantum sector.  It covers rays, phases, wavefunction presentations, tensor powers, statistics, and superselection; the subsequent internal parts carry the finite-QFT, symmetry, transport, mixing, reciprocal-position, measurement, scattering, and internal-sector interfaces.



\paragraph{Status labels and carrier convention}
\label{p4:sec:part2-status-labels}

We use the following compact labels in this part:
\begin{enumerate}[label=(\roman*), leftmargin=2em]
\item \emph{finite theorem}: an exact finite-algebra statement;
\item \emph{interface theorem}: an exact finite statement plus a stated compatibility, limit, or presentation map;
\item \emph{presentation}: a basis, kernel, wavefunction, ray, or path-sum representation of the same finite data;
\item \emph{illustrative}: a standard protocol or textbook normal form expressed in finite algebra;
\item \emph{test channel}: a framework-specific finite readout channel not needed by the core calculus.
\end{enumerate}
Continuum, local-net, thermal, and scattering language always uses the carrier convention of Section~\ref{p4:sec:carrier-convention-mp-mx}: \(m_p=1\) remains the fixed compact momentum datum, while large-position, thermodynamic, and continuum-looking limits belong to the reciprocal-position/interface unfolding \(m_x\), or to compatible auxiliary presentation ladders for that unfolding.



\subsection{Derived Representation and Interface Ingredients}
\label{p4:sec:part2-inventory}



\paragraph{Representation and interface ledger}
\label{p4:sec:part2-inventory-ledger}

\begingroup
\small
\setlength{\tabcolsep}{4pt}
\renewcommand{\arraystretch}{1.12}

\begin{xltabular}{\linewidth}{@{}
  >{\raggedright\arraybackslash}p{0.23\linewidth}
  >{\raggedright\arraybackslash}p{0.14\linewidth}
  >{\raggedright\arraybackslash}p{0.12\linewidth}
  >{\raggedright\arraybackslash}X
@{}}
\caption[Representation and interface ingredients]{Representation and interface ingredients.}
\label{p4:tab:part2-inventory}\\

\toprule
\textbf{Ingredient}
& \textbf{Status}
& \textbf{Coverage}
& \textbf{Role / remaining limitation} \\
\midrule
\endfirsthead

\caption[]{Representation and interface ingredients (continued).}\\
\toprule
\textbf{Ingredient}
& \textbf{Status}
& \textbf{Coverage}
& \textbf{Role / remaining limitation} \\
\midrule
\endhead

\midrule
\multicolumn{4}{r}{\emph{Continued on next page}}\\
\bottomrule
\endfoot

\bottomrule
\endlastfoot

GNS realization of a state
& finite theorem
& covered
& Recovers the textbook Hilbert-space presentation from an algebraic state; no
extra axiom is introduced. \\

Projective/ray language and Hopf phase
& presentation
& covered
& Identifies pure states with projective classes and ties phase redundancy to
the native Hopf fiber. \\

Wavefunctions from maximal abelian subalgebras
& interface theorem
& covered
& Recovers basis/position-style presentations; still basis-dependent rather than
structural. \\

Many-body tensor powers and identical-particle sectors
& finite theorem + supplied sector
& conditional
& Theorem~\ref{p4:thm:permutation-sector-decomposition} derives the full
Young-sector decomposition.  Bose or Fermi Fock calculus follows after the
corresponding statistics selector.  On the bosonic branch GP1 already forces
finite occupation support, while the standard initial-sublevel construction
additionally requires a generated occupation boundary
\(\mathsf P_{\rm gbd}(\mathrm{occ},N_{\max})\); neither the statistics
selector nor its numerical boundary follows from permutation covariance. \\

Superselection sectors
& presentation + interface theorem
& covered
& Gives an exact finite charge-sector theorem for accessible algebras and an
effective decoherence-stabilized refinement; exhaustive charge catalogues are
left to later sector-specific papers. \\

Spin and finite internal multiplets
& presentation + interface theorem
& covered
& Arises through fuzzy \(\CP^1\simeq S^2\) and finite blocks; bounded addition,
recoupling, and Zeeman/Rabi model sectors are explicit. \\

Symmetry generators and conserved quantities
& presentation
& covered
& Unitary/projective symmetry actions, one-parameter generators, and
conserved-quantity criteria are explicit; systematic classification of all
framework symmetry groups is left to later sector-specific papers. \\

Stationary states, spectral decomposition, and dynamical pictures
& presentation
& covered
& Finite spectral resolution plus Schr\"odinger, Heisenberg, interaction, and
perturbative picture changes are explicit; finite-window scattering and
resonance language are added later in the representation/interface sequence. \\

Three-level spectral geometry and Weyl invariants
& presentation + interface theorem
& covered
& Organizes rank-three hierarchy and mixing data; relative diagonalization,
finite unitary mixing, and rephasing-invariant data are explicit. \\

Local fuzzy \(\CP^1\simeq S^2\) patch and canonical observables
& interface theorem
& covered
& Gives local noncommuting coordinates and the source of spin-like sectors. \\

Heisenberg contraction limit
& interface theorem
& covered
& The ambient commutator is canonical after a low-sector sandwich, while the
commutator of compressed observables retains its own top-edge defect;
continuum CCR require a controlled core/limit. \\

Interferometric, Stern--Gerlach, Bell, and teleportation protocols
& illustrative
& covered
& Shows standard laboratory protocols inside the finite algebraic formalism. \\

Measurement theory beyond sharp projection
& illustrative + interface theorem
& covered
& Includes POVMs, instruments, weak monitoring, stochastic filtering, weak
values on a specified weak-probe model, and feedback-compatible trajectory
language; apparatus-specific noise and disturbance require the full indirect
measurement model. \\

Wave-mechanical kernels, discrete path sums, and path-integral limits
& finite theorem + open interface
& finite sums covered
& The finite Trotter path sums are exact.  Schr\"odinger/kernel convergence
requires a supplied operator limit, and no continuum path measure or
real-time oscillatory-integral theorem is claimed. \\

Continuum/QFT, curved-background, and thermodynamic bridges
& interface theorem
& covered
& Realizes the upstream Paper-III local free dispersion as a positive finite
Hamiltonian on stable isolated blocks, and develops inductive local finite-to-field
limits, compatible algebraic local nets, quasi-free GNS reconstruction, weak
smeared-field distribution limits, exact finite coarse-graining and threshold
bookkeeping, finite Bogoliubov transport on time-dependent mode spaces, and
quasi-local many-body limits.  All large, local, or continuum parameters here
are reciprocal-position/interface data, not flows of \(m_p=1\); fully
interacting renormalized relativistic continuum QFT remains beyond scope. \\

Bound states, solvable models, and spectral examples
& illustrative
& covered
& Exact supplied two-level/Zeeman calculations, an explicit finite oscillator,
and conditional well, barrier, Coulomb, and spin--orbit representatives are
presented with their coefficient-extraction or comparison domains. \\

Scattering, resonances, and decay channels
& interface theorem
& covered
& Includes finite-window channel scattering, threshold opening,
Feshbach-type resonances, unstable-block reduction, decay-channel
coarse graining, and a derived open-channel large-\(m_x\)
M{\o}ller/cross-section bridge; full exact infinite-volume transport theory is
not developed. \\

Berry phase from Hopf holonomy
& presentation
& covered
& Gives a geometric account of phase transport in cyclic sectors. \\

Quantum control, estimation, and error-correction protocols
& illustrative
& moderate
& States informationally complete tomography, finite controllability criteria,
quantum Fisher information, and Knill--Laflamme protected-code conditions; full
computation and fault-tolerant architectures remain beyond scope. \\

\end{xltabular}
\TableNotes{This Part-II roadmap records which ingredients are exact finite theorems, interface theorems, presentations, or illustrative protocols, and states the remaining limitation of each language.}
\endgroup



\subsection{Rays, phases, and representation choices}
\label{p4:sec:part2-rays}



\paragraph{Finite GNS realization}
\label{p4:sec:part2-gns}

Given a state $\omega$ on $\mathcal A_m^{\mathrm{env}}$, the GNS construction produces a Hilbert-space realization \cite{GelfandNaimark1943,Segal1947,BratteliRobinson1987}
\begin{equation}
\label{p4:eq:gns}
(\pi_\omega,\mathcal H_\omega,\Omega_\omega)
\end{equation}
such that
\begin{equation}
\label{p4:eq:gns-state}
\omega(a)=\langle\Omega_\omega,\pi_\omega(a)\Omega_\omega\rangle.
\end{equation}


\begin{proposition}[Finite GNS realization equals density-matrix representation]
\label{p4:prop:finite-gns-density}
Let \(\mathcal A_m^{\mathrm{env}}\cong \End(\mathcal H_m)\) and let \(\omega(a)=\Tr(\rho a)\) be a state.  The GNS construction applied to \((\mathcal A_m^{\mathrm{env}},\omega)\) gives a cyclic Hilbert-space representation \((\pi_\omega,\mathcal H_\omega,\Omega_\omega)\) satisfying
\[
\omega(a)=\langle\Omega_\omega,\pi_\omega(a)\Omega_\omega\rangle .
\]
If \(\rho\) has support projection \(S\), then the GNS space is naturally isomorphic to the completion of \(\mathcal A_m^{\mathrm{env}}S\) in the inner product
\[
\langle x,y\rangle_\omega=\Tr(\rho x^*y).
\]
In particular, for a pure state \(\rho=\ket{\psi}\bra{\psi}\), the GNS representation is unitarily equivalent to the defining representation on \(\mathcal H_m\) with cyclic vector \(\psi\).  Thus the usual Hilbert-space state-vector picture is recovered as a finite representation of the structural algebraic state, not as an additional primitive axiom.
\end{proposition}


\begin{proof}
Let \(S\) be the support projection of \(\rho\).  The GNS inner product on the algebra is
\[
  \langle x,y\rangle_\omega=\omega(x^*y)=\Tr(\rho x^*y).
\]
Its null space
\[
N_\omega=\{x:\omega(x^*x)=0\}
\]
is a left ideal: if \(x\in N_\omega\) and \(a\in\mathcal A_m^{\mathrm{env}}\),
then \(a^*a\le\|a\|^2\mathbf1\) gives
\(\omega((ax)^*(ax))\le\|a\|^2\omega(x^*x)=0\).
Moreover \(N_\omega=\{x:xS=0\}\), so the quotient map
\([x]\mapsto xS\) identifies the GNS quotient with the finite-dimensional
space \(\mathcal A_m^{\mathrm{env}}S\) equipped with the displayed inner
product.  Left multiplication gives the representation \(\pi_\omega\), and
the class of \(\mathbf1\) (equivalently its image \(S\)) is cyclic.  If
\(\rho=\ket{\psi}\bra{\psi}\) is pure, then \(S\) is the rank-one projection
onto \(\CC\psi\); the map \(aS\mapsto a\psi\) is unitary onto
\(\mathcal H_m\) and intertwines left multiplication with the defining
representation.  This proves the stated finite GNS realization.
\end{proof}


In finite dimension this is equivalent to the density-matrix presentation already obtained in Part~I in the sense that every algebraic state is represented by a density matrix. The GNS Hilbert space need not be the original carrier $\mathcal H_m$: for mixed states it is generally a representation space built from the algebra and the state, while pure states recover the usual cyclic vector representation on a minimal carrier. Thus ordinary Hilbert-space quantum mechanics is recovered as a representation of the structural algebraic data rather than as an independent starting axiom.



\paragraph{Rays and Hopf phase}
\label{p4:sec:part2-hopf-rays}

For a finite direct sum, pure states are supported on one simple block and
have a rank-one density there.  On the full-envelope representation domain
this specializes to
\(\rho_\psi=\ket{\psi}\bra{\psi}\) for a unit vector \(\psi\).
Their physical content depends only on the projective class \([\psi]\), so
\[
\psi\sim e^{i\theta}\psi
\]
is an exact theorem of the finite state calculus.  The geometric Hopf circle
enters through the degree-one associated character supplied by Paper~II-D; it
must not be confused with the generally noncentral geometric Hopf/Reeb action
on the complete mode carrier.


\begin{theorem}[Projective ray phase and the degree-one Hopf comparison]
\label{p4:thm:hopf-ray-phase}
Let \(\mathcal H_m\) be a nonzero finite complex Hilbert carrier and, for every
unit vector \(\psi\in\mathcal H_m\), put
\[
\rho_\psi=\ket{\psi}\bra{\psi}.
\]
Then
\[
\rho_\psi=\rho_\phi
\quad\Longleftrightarrow\quad
\phi=e^{i\theta}\psi
\]
for unit vectors \(\psi,\phi\).  Consequently the central action
\[
U(1)_{\rm ray}\curvearrowright\mathcal H_m,
\qquad
z\cdot\psi=z\psi,
\]
leaves every finite effect probability and expectation value invariant.

Retain first the upstream associated-character family
\[
\{L^q:q\in\mathbb Z\}.
\]
Now constitute the Paper-IV ray-comparison role as the primitive associated
line compatible with the frozen positive Hopf orientation and with one
principal Hopf turn mapped to one central ray-phase turn.  Its signature
requires
\[
q\in\mathbb Z,\qquad |q|=1,\qquad q>0,
\]
so uniquely \(q=+1\).  Hence the comparison character is
\[
\chi_1:U(1)_{\rm Hopf}\longrightarrow U(1)_{\rm ray},
\qquad
\chi_1(e^{i\alpha})=e^{i\alpha}.
\]
Thus \(\chi_1\) is \emph{DERIVED-LATER} relative to the upstream
associated-character atlas: it identifies one turn of the oriented principal
Hopf fiber with one central ray-phase turn for this comparison role only.
This does not assert that the geometric Hopf action on \(\mathcal H_m\) is
scalar; on mode sectors that action may carry distinct integer or half-integer
weights, and no physical Reeb, matter, hypercharge, or gauge charge is selected
by \(\chi_1\).
\end{theorem}

\begin{remark}[Ray comparison is not a physical charge selector]
\label{p4:rem:ray-comparison-not-charge-selector}
The chronology is
\[
\begin{aligned}
&\text{Paper II-D: }\{L^q\}_{q\in\mathbb Z}\text{ retained}\\
&\qquad\longrightarrow\text{ Paper IV primitive positive ray-comparison role}\\
&\qquad\longrightarrow q=+1.
\end{aligned}
\]
The last value belongs only to the comparison character \(\chi_1\).  It does
not select a physical Hopf/Reeb charge sector; those sectors remain a full
atlas until a later accessibility/sector theorem independently distinguishes a
member.
\end{remark}


\begin{proof}
For any effect \(0\le E\le\mathbf1\),
\[
\Tr(\rho_{e^{i\theta}\psi}E)=\Tr(\rho_\psi E),
\]
because \(\ket{e^{i\theta}\psi}\bra{e^{i\theta}\psi}=\ket{\psi}\bra{\psi}\).
Thus the phase action is operationally invisible to all probabilities and expectation values of the finite theory.
Conversely, equality of rank-one projectors implies equality of their
one-dimensional ranges, so \(\phi=z\psi\) for some \(z\ne0\).  Unit
normalization gives \(|z|=1\), hence \(z=e^{i\theta}\).

Paper~II-D, Theorem~\ref{p2:thm:forced-hopf-chern-unit}, gives the oriented
primitive Hopf bundle \(U(1)\to S^3\to S^2\) with \(c_1=1\).
Its associated-character construction,
Proposition~\ref{p2:prop:hopfion-charge}, assigns to
\(\rho_q(e^{i\alpha})=e^{iq\alpha}\) the line \(L^q\); the primitive
degree-one line is therefore \(L^1\) with character \(\chi_1(z)=z\).
Composition with the central scalar action on a pure-state line gives the
displayed comparison representation.  Since a noncentral geometric mode
representation is not used in this construction, no equality between that
representation and scalar multiplication on the entire carrier follows.
\end{proof}

No independent quantum phase postulate is introduced: phase redundancy is
projective Hilbert geometry, and the primitive Hopf line supplies its oriented
degree-one geometric comparison.  Hopf/Reeb charges on mode sectors remain
distinct representation data and may produce the superselection structure
described below.



\paragraph{Wavefunctions from maximal abelian subalgebras}
\label{p4:sec:part2-wavefunctions}

Choose a maximal abelian subalgebra (MASA) of $\mathcal A_m^{\mathrm{env}}$.
Diagonalizing that MASA gives a preferred presentation basis and identifies a state vector with its coordinate components in that basis.
These components are the finite-resolution analogue of a wavefunction.
They are therefore \emph{presentation-dependent} rather than structural: different MASAs give different but unitarily equivalent wavefunction pictures.



\subsection[Many-body systems, statistics, and superselection]
{Many-body systems, statistics,\\ and superselection}
\label{p4:sec:part2-manybody}



\paragraph{Tensor powers and identical particles}
\label{p4:sec:part2-identical}

Repeated independently instantiated \(\mathsf{SUB}\) roles produce tensor powers
\(\mathcal H^{\otimes n}\) and \(\mathcal A^{\otimes n}\).  When the occurrences
have the same constituted type, the symmetric group acts canonically by
\[
U_\pi(\psi_1\otimes\cdots\otimes\psi_n)
=
\psi_{\pi^{-1}(1)}\otimes\cdots\otimes\psi_{\pi^{-1}(n)} .
\]

\begin{theorem}[Permutation-sector decomposition and statistics boundary]
\label{p4:thm:permutation-sector-decomposition}
Let \(\mathcal H\) be finite dimensional and let \(U\) be the preceding
permutation representation of \(S_n\) on \(\mathcal H^{\otimes n}\).  For each
partition \(\lambda\vdash n\), let \(\chi_\lambda\) be the corresponding
irreducible character and \(d_\lambda=\chi_\lambda(e)\).  Then
\[
P_\lambda
:=
\frac{d_\lambda}{n!}
\sum_{\pi\in S_n}\chi_\lambda(\pi^{-1})U_\pi
\]
is the orthogonal projector onto the \(\lambda\)-isotypic sector, possibly the
zero sector, and
\[
\mathcal H^{\otimes n}
=
\bigoplus_{\lambda\vdash n}P_\lambda\mathcal H^{\otimes n}.
\]
The bosonic and fermionic spaces are respectively
\[
\operatorname{Sym}^n\mathcal H=P_{(n)}\mathcal H^{\otimes n},
\qquad
\Lambda^n\mathcal H=P_{(1^n)}\mathcal H^{\otimes n},
\]
with the latter zero when \(n>\dim\mathcal H\).  Permutation covariance derives
this complete Young-sector decomposition but does not select either the
trivial or sign sector as the physically realized statistics.  Thus the
mathematical decomposition is co-realized, while the generic physical
statistics realization family
\[
 \mathfrak A_{\rm STAT}(n)
 :=\{P_\lambda\mathcal H^{\otimes n}:\lambda\vdash n,\ P_\lambda\ne0\}
\]
remains unresolved until a supplied or later-derived statistics discriminator
is available.
\end{theorem}

\begin{proof}
The operators \(U_\pi\) form a unitary representation of the finite group
\(S_n\).  The character-orthogonality relations state that the elements
\[
e_\lambda
=
\frac{d_\lambda}{n!}
\sum_{\pi\in S_n}\chi_\lambda(\pi^{-1})\pi
\in\CC[S_n]
\]
are mutually orthogonal central self-adjoint idempotents and sum to the
identity of \(\CC[S_n]\).  Applying the representation \(U\) gives the
projectors \(P_\lambda\), their orthogonality, and the displayed direct sum.
For \(\lambda=(n)\), the character is trivial, so \(P_{(n)}\) symmetrizes all
tensor factors.  For \(\lambda=(1^n)\), the character is the sign character,
so \(P_{(1^n)}\) antisymmetrizes them.  The group-algebra decomposition
contains every nonzero Young sector and has no criterion that chooses one
central idempotent as the physical carrier.  Such a choice therefore does not
follow from permutation covariance.
\end{proof}

Accordingly, a finite Bose or Fermi construction below begins only after a
statistics-sector selector has been supplied.  Alternatively, such a selector
could be derived from a separate spin--statistics theorem on a named local
relativistic QFT domain; no such selection theorem belongs to the finite
Paper-I-A--III prefix.



\paragraph{Superselection sectors}
\label{p4:sec:part2-superselection}

If a conserved algebraic or topological charge commutes with all observables accessible in a given regime, the full state space decomposes into superselection blocks.
Interference between distinct blocks is then operationally inaccessible.
This is the algebraic form of superselection in the Born framework.


\begin{theorem}[Finite charge superselection]
\label{p4:thm:charge-superselection}
Let $Q_1,\dots,Q_r$ be commuting self-adjoint charges on $\mathcal H_m$, let $\{P_q\}$ be the joint spectral projectors, and let $\mathcal A_{\mathrm{acc}}\subseteq \End(\mathcal H_m)$ be the algebra of observables accessible in a given regime.
If
\[
[A,Q_i]=0
\qquad
\text{for all }A\in \mathcal A_{\mathrm{acc}},\ i=1,\dots,r.
\]
then:
\begin{enumerate}[label=(\roman*)]
\item the carrier decomposes as
\[
\mathcal H_m=\bigoplus_q \mathcal H_q,
\qquad
\mathcal H_q:=P_q\mathcal H_m;
\]
\item every accessible observable is block diagonal, equivalently the compression map embeds the accessible algebra as a block-diagonal subalgebra,
\[
\mathcal A_{\mathrm{acc}}
\hookrightarrow
\bigoplus_q P_q\mathcal A_{\mathrm{acc}}P_q,
\qquad
A\longmapsto (P_qAP_q)_q;
\]
\item for every state $\rho$, all accessible expectation values are reproduced by the dephased state
\[
\rho_{\mathrm{ss}}:=\sum_q P_q\rho P_q.
\]
\end{enumerate}
Hence off-diagonal coherences between distinct charge sectors are operationally invisible.
If, in addition, every admissible Hamiltonian commutes with the charges and
every admissible reduced channel preserves block-diagonal states---for
example, it satisfies
\(\Phi\Delta_Q=\Delta_Q\Phi\Delta_Q\)---then the sectors are dynamically
stable.  Here
\[
\Delta_Q(\rho)=\sum_q P_q\rho P_q,
\]
is the dephasing conditional expectation.  The stronger equality
\(\Phi\Delta_Q=\Delta_Q\Phi\) is sufficient but is not claimed equivalent to
block preservation.  With the stated accessible-algebra restriction and
dynamical stability, the sectors are exact superselection sectors.  If the
off-diagonal blocks are only suppressed by environmental decoherence or
restricted access, the same sector split is effective rather than exact.
\end{theorem}


\begin{proof}
Because the charges commute, the finite spectral theorem gives a joint decomposition
\[
\mathcal H_m=\bigoplus_q P_q\mathcal H_m.
\]
If $A$ commutes with every $Q_i$, it commutes with every joint spectral projector $P_q$, hence
\[
A=\sum_q P_qAP_q.
\]
Thus the compression map above is injective and realizes $\mathcal A_{\mathrm{acc}}$ as a block-diagonal subalgebra of the displayed direct sum, without requiring the compressed blocks to vary independently.
Therefore
\[
\Tr(\rho A)=\Tr\!\Bigl(\sum_q P_q\rho P_q\,A\Bigr)=\Tr(\rho_{\mathrm{ss}}A)
\]
for every $A\in\mathcal A_{\mathrm{acc}}$.
The last statement only distinguishes whether the vanishing of cross-sector observables is exact at the algebra/dynamics level or merely generated effectively by decoherence.
\end{proof}

\begin{corollary}[Level-one central-parity superselection]
\label{p4:cor:central-parity-superselection}
Let \(\Gamma_c\) and \(P_\pm\) be the Paper-II-D central grading.  For every
full-carrier density representative \(\rho\) and every native accessible
effect \(E\),
\[
 \boxed{
 \Tr(\rho E)
 =\Tr\!\left[(P_+\rho P_++P_-\rho P_-)E\right].}
\]
Thus cross-parity coherences are inaccessible to the native observable
algebra.  If every admissible Hamiltonian is central-even and every admissible
channel preserves the dephased subspace, the two sectors are dynamically
stable exact superselection sectors.  Without that dynamics clause, the
displayed identity is an exact kinematic/accessibility statement.
\end{corollary}

\begin{proof}
Theorem~\ref{p2:thm:diagonal-center-grading} shows that every native
accessible effect commutes with \(\Gamma_c\), and hence with \(P_\pm\).  The
two cross terms \(P_+\rho P_-E\) and \(P_-\rho P_+E\) have zero trace, which
gives the displayed identity.  A central-even Hamiltonian preserves the two
eigenspaces, and the stated channel condition preserves their dephased state
space, so Theorem~\ref{p4:thm:charge-superselection} supplies the dynamical
claim.
\end{proof}

\begin{remark}[Central parity, statistics, and odd fields]
\label{p4:rem:central-parity-field-boundary}
This is diagonal-center tensor/spinorial parity, not yet a boson/fermion or
spin--statistics theorem.  An odd field may connect the two sectors in the
graded field envelope \(\mathcal F\) without becoming an observable in
\(\mathcal A_{\rm obs}\).
\end{remark}

So superselection is not a new collapse postulate.
It is a statement about which block structure is visible to the observables and channels actually admitted in the regime under study.


\begin{remark}[Hopf half-turn does not discharge the central-odd bridge]
\label{p4:rem:hopf-half-turn-not-central-odd}
Paper~II-E and Paper~III-A identify the provenance-fixed Hopf half-turn on
the protected block with \(\pm\Gamma_{\rm sp}\) and with the geometric
antipodal operation.  This finite involution preserves the Paper-II-D
\(11\oplus4\) diagonal-center sectors and therefore does not supply the
central-odd operator required by any later odd-field/FENV bridge.  Finite
Hopf parity and diagonal-center oddness remain different typed gradings.
\end{remark}

\begin{corollary}[Hopf/Reeb charge sectors as finite superselection data]
\label{p4:cor:hopf-reeb-superselection}
Let \(Q_\xi\) be the finite self-adjoint generator representing the Paper~I-A
Reeb/Hopf circle action on a resolved carrier block, with spectral projectors
\(P_q\).  Put
\[
U_\xi(\theta):=e^{-i\theta Q_\xi},
\qquad
\Delta_\xi(X):=
\int_{U(1)}U_\xi(\theta)XU_\xi(\theta)^*\,d\theta
=\sum_qP_qXP_q .
\]
Assume
\(\mathcal A_{\rm acc}\subseteq\{Q_\xi\}'\) and every admissible reduced
channel is \(U_\xi\)-covariant:
\[
\Phi\!\left(U_\xi(\theta)\rho U_\xi(\theta)^*\right)
=
U_\xi(\theta)\Phi(\rho)U_\xi(\theta)^* .
\]
Then
\[
\Phi\Delta_\xi=\Delta_\xi\Phi,
\]
and
\[
\mathcal H_m=\bigoplus_q P_q\mathcal H_m
\]
is an exact finite superselection decomposition.  The weaker, correctly typed
condition
\(\Phi\Delta_\xi=\Delta_\xi\Phi\Delta_\xi\) is sufficient merely for
preservation of block-diagonal states.  If block preservation holds only
after environmental dephasing or restricted access, the decomposition is
effective.
\end{corollary}


\begin{proof}
The accessible-algebra claim follows from
Theorem~\ref{p4:thm:charge-superselection}.  Integrating the covariance
identity over the compact circle gives
\(\Phi\Delta_\xi=\Delta_\xi\Phi\).  Hence a block-diagonal state remains
block diagonal and all accessible expectations are insensitive to its
off-diagonal coherences.  The weaker intertwining identity gives the same
block-preservation conclusion on the range of \(\Delta_\xi\).
\end{proof}

Thus Hopf/Reeb weights may serve as quantum charge-sector data on the
explicit accessible-algebra/\(U_\xi\)-covariant-channel domain of the
corollary.
Geometry alone does not make every such weight an exact superselection rule.

\begin{remark}[Dephasing compatibility does not imply charge covariance]
\label{p4:rem:dephasing-not-charge-covariance}
The implication used above is one-way.  For example, on \(\CC^2\) let
\(Q=\operatorname{diag}(0,1)\), let \(S\) exchange the two charge
eigenvectors, and put \(\Phi(\rho)=S\rho S^*\).  This channel commutes
with charge dephasing, but sends \(\lvert0\rangle\langle1\rvert\) to
\(\lvert1\rangle\langle0\rvert\), reversing its nonzero \(U(1)\)
weight.  It therefore fails charge covariance.  A selection rule that uses
individual charge differences requires the actual covariance identity;
dephasing compatibility certifies only the stated block preservation.
\end{remark}

\begin{remark}[CR/Szeg\H{o} scalar polarization and the finite quantum carrier]
\label{p4:rem:szego-polarization-interface}
Paper~I-A also supplies a canonical scalar boundary polarization, the Hardy/Szeg\H{o}
projector \(\Pi_{\rm Sz}:L^2(S^3)\to H^2(S^3)\).  In scalar boundary sectors
this gives the upstream Hardy/Szeg\H{o} Toeplitz interface
\[
T_f^{\rm CR}=\Pi_{\rm Sz}M_f\Pi_{\rm Sz}.
\]
This is a continuum scalar-boundary polarization, not the finite Born carrier and
not the finite carrier projector \(\Pi_m\).  It therefore tightens the geometric
meaning of scalar wavefunction/Toeplitz language, while the quantum state calculus
of Part~I and the finite QFT bridge below remain based on the constituted finite
carrier, its algebra of observables, and its sector projectors.
\end{remark}

\begin{remark}[Basic and contact analytic scope]
\label{p4:rem:basic-contact-analytic-scope}
The Paper-I-A Hopf basic sector has one scalar harmonic ground slot and no basic
\((0,1)\) harmonic zero mode, so Paper~IV's finite quantum calculus does not
silently inherit an extra transverse zero-mode multiplet.  The Paper-I-A contact
homogeneous dimension \(Q=4\) controls Heisenberg/contact analytic scaling, but it
is not the Lorentzian scalar-cover positive-frequency condition and is not an
additional continuum spacetime postulate.  The flat hyper-K\"ahler cone similarly
contributes scale-rigidity language, not an extra curvature source.
\end{remark}




\section{Reciprocal coordinates and canonical quantum presentations}
\label{p4a:sec:v238-a7}
\noindent\textit{Source and claim boundary.} All position-role and Heisenberg-contraction statements retain fixed bounded momentum and their explicitly stated spatial-resolution domains; finite textbook Hamiltonians are comparison presentations, not primitive premises.

\paragraph{Reciprocal-Position and Textbook Hamiltonian Presentations}
\label{p4:part:reciprocal-position-presentations}

\paragraph{Orientation.}
This part separates reciprocal-position and textbook Hamiltonian presentations from the microscopic compact-momentum datum.  Fuzzy angular momentum, Born-pair sectors, oscillators, wells, barriers, and central-potential representatives are presentation tools for finite quantum data, not new upstream axioms.



\subsection{Emergent canonical sector on the reciprocal position side}
\label{p4:sec:part2-canonical}



\paragraph{Microscopic compact momentum and the large reciprocal-position unfolding}
\label{p4:sec:part2-momentum-position-split}

The canonical presentation requires an explicit distinction between the two
reciprocal sides of the framework.  Papers~II--III identify a compact momentum
geometry whose physically selected microscopic compact sector is the minimal
truncation level
\begin{equation}
\label{p4:eq:m-momentum-physical}
m_p=1.
\end{equation}
Paper~II supplies the finite \(\Spin(4)\)-equivariant carrier family.
Paper~II-B's generated-carrier theorem fixes the realized
compact momentum sector \(m_p=1\) without invoking \(\Dcl\).  On the
separate \(\Dcl\)-tagged dynamics branch,
Theorem~\ref{p3:thm:no-momentum-m-flow} then proves that evolution inside one
fixed spectral-action sector cannot change that already selected label.  The
diagonal \(\SU(2)\) content of
that imported sector therefore gives the exact microscopic angular-momentum
content of the compact momentum sector.

On the Paper-III generated reciprocal-localization branch, the
reciprocal-position cotangent generator germ is a derived localization
certification; on the occurring minimal local-action parent,
Theorem~\ref{p3f:thm:minimal-configuration-tangent-vdom} derives the
admissible position-metric variational domain without a separate
\((\mathrm{POL})\) selector.  Paper~III-F's separate
\(({\rm QLOC})\) predicate supplies static compressed-observable locality
when that analytic domain is used.  Macroscopic locality is therefore carried
by the reciprocal-position presentation rather than by the fixed compact
momentum sector.
Accordingly, textbook particle mechanics with extended position space belongs
to the reciprocal position-side presentation, whose realized refinement level
can become large only through a generated or supplied resolution trajectory.
In the primitive unresolved stage
this unfolding begins at the minimal value
\begin{equation}
\label{p4:eq:m-position-initial}
m_x=1,
\end{equation}
and its cumulative value is the derived Paper~II-C frontier history
\((\mathrm{HIST}_{x}^{K,t})\), not an independent tower, trajectory, or
monotonicity postulate.  On a sharp physical frontier the scale is
\(\Lambda_F\); it is identified with an action cutoff only under the separate
\((\mathrm{ID}_{F/A})\) match.
On the homogeneous \(\mathsf{SF}_{\Lambda_F}\) branch,
Corollary~\ref{p4:cor:exact-reciprocal-position-resolution} gives that
trajectory exactly; large \(m_x\) occurs only if the actual scale factor
crosses the corresponding spectral thresholds.  Other coarse-grained or
laboratory regimes require a separately supplied compatible refinement.  The ordinary Schr\"odinger
picture of motion in \(\mathbb R^3\), when a convergence theorem to that
presentation has been supplied, is therefore a large-\(m_x\) reciprocal-position
interface, not a large-\(m\) limit of the compact momentum sector.

The general finite-\(m\) representation theory derived below remains useful in two distinct ways.
First, it gives exact mathematical statements for every member of the
finite-carrier family, without claiming that every \(m\) is physically
realized.
Second, at large realized reciprocal-position refinement, the same bounded
formulas organize candidate finite approximants.  Convergence to a canonical
Schr\"odinger operator requires compatible embeddings and strong- or
norm-resolvent convergence; Papers~II--III do not supply that further theorem.
This separation keeps the framework architecture explicit: compact momentum
stays microscopic and minimally truncated, while reciprocal-position
resolution follows the exact spectral staircase on
\(\mathsf{SF}_{\Lambda_F}\), or a separately supplied compatible trajectory on
another branch.

\begin{proposition}[Momentum-role realization dictionary]
\label{p4:prop:momentum-role-realization-dictionary}
The upstream and downstream momentum statements fit together without an
operator limit on a common Hilbert space:
\begin{enumerate}[label=(\roman*),leftmargin=2.2em]
\item Paper~II-C gives the exact first-order bundle isomorphism
\[
 \mathfrak B_a:\mathscr K_a\overset\sim\longrightarrow T_b^*S_B^3,
 \qquad
 \mathfrak B_a(\widehat P_i(a))=e_b^i,
\]
for the moving compact transvection-operator fibre.
\item On a reciprocal-position sphere of radius $a$, Paper~II-D gives
\[
 \Pi_i^{(a)}=\frac{\hbar}{a}Q_i,
 \qquad
 \sigma_{\rm phys}(\Pi_i^{(a)})(x,p)=p(e_i(x)),
 \qquad
 \sum_i(\Pi_i^{(a)})^2=\hbar^2\Delta_a.
\]
Thus $\Pi_i^{(a)}$ quantizes the local cotangent component functions but is
not the same finite operator as $\widehat P_i$ on $\mathcal H_1$.
\item Along the large-$m_x$, large-radius sharp-frontier family, Paper~VII-C
proves that the $\Pi_i^{(a)}$ become commuting only in trace moments, with
limiting variables $p_i$ on the momentum ball.  The microscopic label remains
$m_p=1$ throughout.
\end{enumerate}
Accordingly the physical dictionary is
\[
 \boxed{
 \widehat P_i\ \xleftrightarrow{\ \text{exact local germ}\ }
 T^*U
 \ \xleftrightarrow{\ \text{principal symbol}\ }
 \Pi_i^{(a)}
 \ \xRightarrow{\ \text{trace moments}\ }
 p_i,}
\]
where the middle double arrow denotes role/symbol correspondence rather than
finite-operator equality.  No step identifies the global compact momentum
carrier with $T^*U$ or sends $m_p$ to infinity.
\end{proposition}

\begin{theorem}[Fixed-parent spectral restrictions and resolvent convergence]
\label{p4:thm:fixed-parent-resolvent-convergence}
Let \(H=H^*\) on a Hilbert space \(\mathcal K\), and let \(P_N\) be increasing
spectral projections of \(H\) with \(P_N\to\mathbf1\) strongly.  Set
\(\mathcal K_N=P_N\mathcal K\), let
\(j_N:\mathcal K_N\hookrightarrow\mathcal K\) be the canonical inclusion,
and let \(H_N=H|_{\mathcal K_N}\).  Then for every
\(z\in\mathbb C\setminus\mathbb R\),
\[
 \boxed{
 j_N(H_N-z)^{-1}P_N
 \longrightarrow(H-z)^{-1}
 \quad\text{strongly}.}
\]
If the omitted spectral energies escape,
\[
 E_N^{\rm omit}:=
 \inf\{|\lambda|:\lambda\in\operatorname{spec}(H|_{(1-P_N)\mathcal K})\}
 \longrightarrow\infty,
\]
then the convergence is in norm; for spectral restrictions,
\[
 \bigl\|(H-z)^{-1}
 -j_N(H_N-z)^{-1}P_N\bigr\|
 =
 \sup_{\lambda\in\operatorname{spec}(H|_{(1-P_N)\mathcal K})}
 |\lambda-z|^{-1}.
\]
Thus exact restrictions on one fixed parent have canonical comparison maps.
Norm-resolvent language for full-space wall Hamiltonians additionally
requires a self-adjoint extension on the complement, for example a positive
wall \(W_N\) whose complementary energies tend to \(+\infty\).
\end{theorem}

\begin{proof}
Functional calculus identifies the embedded finite resolvent with
\(P_N(H-z)^{-1}\).  Strong convergence follows from \(P_N\to1\).  Because
\(P_N\) is spectral, the difference is
\((1-P_N)(H-z)^{-1}\), whose norm is the displayed spectral supremum; the
escape condition makes it vanish.
\end{proof}

\begin{theorem}[Expanding-sphere target maps and Mosco interface]
\label{p4:thm:expanding-sphere-mosco-interface}
Let \(a_N\to\infty\), let
\(\mathcal K_N=L^2(S^3_{a_N},d\mu_{a_N})\), and let
\(\mathcal K_\infty=L^2(\mathbb R^3,dx)\).  Choose base points
\(p_N\in S^3_{a_N}\), radii
\[
 r_N\longrightarrow\infty,
 \qquad
 \frac{r_N}{a_N}\longrightarrow0,
\]
and cutoffs \(\chi_N\in C_c^\infty(B_{r_N}(0))\) with
\(0\le\chi_N\le1\) and \(\chi_N=1\) on \(B_{r_N/2}(0)\).  For all large
\(N\), the exponential map at \(p_N\) is injective on \(B_{r_N}(0)\).
Write its volume density as
\[
 d\mu_{a_N}(\exp_{p_N}x)=\vartheta_N(x)\,dx
\]
and define a contraction
\(J_N:\mathcal K_\infty\to\mathcal K_N\) by
\begin{equation}
\label{p4:eq:expanding-sphere-identification-map}
 (J_Nf)(\exp_{p_N}x)
 :=
 \chi_N(x)\vartheta_N(x)^{-1/2}f(x)
 \quad (|x|<r_N),
\end{equation}
with value zero outside that chart.  Then
\begin{equation}
\label{p4:eq:expanding-sphere-asymptotic-isometry}
 \langle J_Nf,J_Ng\rangle_{\mathcal K_N}
 =
 \int_{\mathbb R^3}\chi_N^2\overline f g\,dx
 \longrightarrow
 \langle f,g\rangle_{\mathcal K_\infty}.
\end{equation}

Declare \(u_N\to u\) strongly when
\(\lVert u_N-J_Nu\rVert_{\mathcal K_N}\to0\), and declare
\(u_N\rightharpoonup u\) weakly when
\(\sup_N\lVert u_N\rVert<\infty\) and
\[
 \langle u_N,J_Nf\rangle_{\mathcal K_N}
 \longrightarrow
 \langle u,f\rangle_{\mathcal K_\infty}
 \quad\text{for every }f\in\mathcal K_\infty.
\]
Let \(q_N\) and \(q_\infty\) be densely defined, closed quadratic forms on
\(\mathcal K_N\) and \(\mathcal K_\infty\), respectively, with one common
lower bound
\[
 q_N[u]\ge-\omega\lVert u\rVert^2,
 \qquad
 q_\infty[u]\ge-\omega\lVert u\rVert^2.
\]
Extend each form by \(+\infty\) off its form domain, and suppose
\begin{enumerate}[label=(M\arabic*), leftmargin=2.7em]
\item whenever \(u_N\rightharpoonup u\),
\[
 q_\infty[u]\le\liminf_{N\to\infty}q_N[u_N];
\]
\item for every \(u\in\mathcal K_\infty\), there is a sequence
\(u_N\to u\) strongly such that
\[
 \limsup_{N\to\infty}q_N[u_N]\le q_\infty[u].
\]
\end{enumerate}
Let \(H_N\) and \(H_\infty\) be the self-adjoint operators associated with
these forms.  Then, for every \(\alpha>\omega\) and
\(f\in\mathcal K_\infty\),
\begin{equation}
\label{p4:eq:expanding-sphere-strong-resolvent}
 \bigl\lVert
 (H_N+\alpha)^{-1}J_Nf
 -J_N(H_\infty+\alpha)^{-1}f
 \bigr\rVert_{\mathcal K_N}
 \longrightarrow0.
\end{equation}
Equivalently, the analogous statement holds for every
\(z\in\mathbb C\setminus\mathbb R\).  Conditions (M1)--(M2) are the complete
generalized Mosco interface used here.  They imply no operator-norm estimate:
norm-resolvent convergence requires a separate proof of
\[
 \left\lVert
 (H_N-z)^{-1}J_N-J_N(H_\infty-z)^{-1}
 \right\rVert_{\mathcal B(\mathcal K_\infty,\mathcal K_N)}
 \longrightarrow0
\]
and, if a full-space comparison is intended, a corresponding estimate on
the complementary spherical sector.  Neither \(a_N\to\infty\) nor an
integer band history alone supplies (M1), (M2), or those norm estimates.
\end{theorem}

\begin{proof}
The change-of-variables formula in the exponential chart gives
\eqref{p4:eq:expanding-sphere-asymptotic-isometry}; dominated convergence
applies because the cutoffs are bounded by one and equal one on balls whose
radii tend to infinity.  Fix \(\alpha>\omega\) and \(f\in\mathcal K_\infty\).
The vectors
\[
 u_N=(H_N+\alpha)^{-1}J_Nf
\]
are the unique minimizers of the strictly convex, coercive functionals
\[
 \mathcal F_N(v)
 :=q_N[v]+\alpha\lVert v\rVert^2
 -2\operatorname{Re}\langle J_Nf,v\rangle.
\]
By (M1), (M2), and
\eqref{p4:eq:expanding-sphere-asymptotic-isometry}, these functionals Mosco
converge to
\[
 \mathcal F_\infty(v)
 :=q_\infty[v]+\alpha\lVert v\rVert^2
 -2\operatorname{Re}\langle f,v\rangle.
\]
The common lower bound makes their sublevel sets uniformly bounded.  The
liminf inequality gives weak convergence of every minimizing subsequence to
the unique minimizer
\(u=(H_\infty+\alpha)^{-1}f\); applying a recovery sequence to \(u\) gives
convergence of the minimum values.  Strict convexity of the shifted forms
then upgrades the minimizing convergence to
\(\lVert u_N-J_Nu\rVert\to0\), which is
\eqref{p4:eq:expanding-sphere-strong-resolvent}.  The resolvent identity and
the uniform self-adjoint resolvent bound extend the conclusion from
\(-\alpha\) to every nonreal \(z\).  Operator-norm convergence is stronger
than this pointwise conclusion and therefore requires the separate estimate
stated in the theorem.
\end{proof}

\begin{remark}[GP1--GP3 status of the momentum carrier and position-resolution labels]
\label{p4:rem:mp-mx-selection-status}
The labels \(m_p\) and \(m_x\) have different logical roles.  The compact
momentum label \(m_p\) is a finite representation label of the ultraviolet
momentum carrier, so changing it would change the physical carrier fixed by
Paper~II-B; Paper~III's no-momentum-\(m\)-flow theorem additionally shows that
the selected label cannot change within the stated \(\Dcl\) dynamics.  The reciprocal
position label \(m_x\) is an infrared harmonic-resolution label associated
with the realized spatial size; changing it changes the number of resolved
position modes, not the compact momentum carrier.  Thus GP1 permits large but
finite \(m_x\) windows at realized stages but does not generate them.
Paper~III-F supplies the exact \(\mathsf{SF}_{\Lambda_F}\) staircase when that
operational frontier occurs.  GP2 forbids treating those windows as a hidden
flow \(m_p\to m_p'\) or as a new selector for the microscopic compact sector.
\end{remark}

\begin{remark}[Compatibility with Paper~III momentum no-flow]
\label{p4:rem:mp-mx-compatible-with-p3}
On the \(\Dcl\) domain on which Paper~III,
Theorem~\ref{p3:thm:no-momentum-m-flow}, is proved, this large-position
presentation is compatible with its momentum-label preservation result.  The
continuum-looking wavefunction, canonical particle, and laboratory scattering
languages below are large-\(m_x\) reciprocal-position interfaces.  They do
not require a dynamical flow of the compact momentum truncation.  Throughout
this paper the microscopic compact momentum carrier remains the fixed minimal
compact sector
\[
        m_p=1,
\]
while the realized reciprocal position refinement may grow to \(m_x\gg1\).
\end{remark}

\begin{remark}[Regime dictionary: what applies where]
\label{p4:rem:regime-dictionary}
To avoid confusion, the framework distinguishes three levels of statement.

\emph{(i) Exact finite-quantum statements on every admitted carrier.}
These include the Hilbert-space realization of states and observables, density
matrices, POVMs and instruments, tensor products, reduced states, unitary and
GKLS dynamics, spectral decomposition, perturbation theory, adiabatic transport,
sector projectors, and finite many-body/Fock constructions.  They also include
the three standard dynamical pictures: Schr\"odinger, Heisenberg, and
interaction.  These statements apply on every finite sector admitted by the
framework, including the microscopic compact momentum sector with \(m_p=1\).

\emph{(ii) Exact statements specific to the microscopic compact momentum sector.}
These include the physically selected minimal truncation \(m_p=1\), the corresponding bounded diagonal-angular-momentum content, and the compact momentum-side interpretation of the underlying \(\Spin(4)\) kinematics.
Such statements are exact microscopic structure, not large-scale approximations.

\emph{(iii) Emergent reciprocal-position statements.}
These include canonical coordinate--momentum language, contraction to the Heisenberg commutation relations, textbook wavefunction models on effective spatial domains, partial-wave orbital towers, and large-flight-time scattering language.
They belong to the reciprocal-position unfolding once its realized refinement
reaches large \(m_x\), either through the exact
\(\mathsf{SF}_{\Lambda_F}\) staircase or through a separately supplied compatible
trajectory.

In particular, the \emph{Heisenberg picture} is exact on every finite carrier, whereas the \emph{Heisenberg canonical algebra} \([X_i,P_j]\approx i\hbar\,\delta_{ij}\) is only an emergent reciprocal-position contraction.
Thus each theorem is asserted only in the sector and regime specified by the framework architecture.
\end{remark}



\paragraph{Controlled contraction to the Heisenberg algebra}
\label{p4:sec:part2-heisenberg}

Let \(\mathcal H_d=\operatorname{span}\{|0\rangle,\ldots,|d-1\rangle\}\), and define
truncated ladder operators by
\[
a_d|n\rangle=\sqrt n\,|n-1\rangle,
\qquad
a_d^\dagger|n\rangle=\sqrt{n+1}\,|n+1\rangle
\]
whenever the displayed vector lies in \(\mathcal H_d\), with zero action at the two
truncation boundaries.

\begin{theorem}[Exact finite oscillator boundary defect and low-sector CCR]
\label{p4:thm:finite-heisenberg-contraction}
The truncated ladder operators satisfy
\[
[a_d,a_d^\dagger]=\Id_d-d|d-1\rangle\langle d-1|.
\]
For any \(\kappa>0\), define
\[
X_d=\sqrt{\frac{\hbar}{2\kappa}}(a_d+a_d^\dagger),
\qquad
P_d=-i\sqrt{\frac{\hbar\kappa}{2}}(a_d-a_d^\dagger).
\]
Then
\begin{equation}
\label{p4:eq:heisenberg}
[X_d,P_d]=i\hbar\bigl(\Id_d-d|d-1\rangle\langle d-1|\bigr).
\end{equation}
If \(Q_K=\sum_{n=0}^{K}|n\rangle\langle n|\) with \(K\le d-2\), then the compressed
ambient commutator is exact on the low sector:
\[
Q_K[X_d,P_d]Q_K=i\hbar Q_K.
\]
This must be distinguished from the commutator of the compressed observables,
which has its own top-edge defect:
\begin{equation}
\label{p4:eq:compressed-heisenberg-boundary}
[Q_KX_dQ_K,Q_KP_dQ_K]
=
i\hbar\bigl(Q_K-(K+1)|K\rangle\langle K|\bigr).
\end{equation}
Consequently, if \(Q_R=\sum_{n=0}^{R}|n\rangle\langle n|\) and \(R<K\), then
\[
Q_R[Q_KX_dQ_K,Q_KP_dQ_K]Q_R=i\hbar Q_R.
\]
For every normalized \(\psi\in\mathcal H_d\),
\[
\left|\langle\psi,([X_d,P_d]-i\hbar\Id_d)\psi\rangle\right|
=\hbar d\,|\langle d-1|\psi\rangle|^2.
\]
Hence the ambient commutator agrees exactly with the canonical relation after
low-sector compression, whereas the algebra generated by the compressed
observables retains the boundary term
\eqref{p4:eq:compressed-heisenberg-boundary}.  Canonical matrix elements are
recovered on fixed finite-support vectors that remain below the moving
boundary, and for sequences satisfying
\(d|\langle d-1|\psi_d\rangle|^2\to0\).  It does not converge in full operator
norm, consistently with the finite-dimensional trace obstruction.
\end{theorem}

\begin{proof}
Direct action on the number basis gives
\(a_da_d^\dagger|n\rangle=(n+1)|n\rangle\) for \(n<d-1\) and zero at the top level,
whereas \(a_d^\dagger a_d|n\rangle=n|n\rangle\).  Their difference is therefore
\(\Id_d-d|d-1\rangle\langle d-1|\).  Substitution into the definitions of
\(X_d,P_d\) gives \eqref{p4:eq:heisenberg}.  Compression by \(Q_K\) removes the
original boundary projector exactly.  The operators \(Q_Ka_dQ_K\) and
\(Q_Ka_d^\dagger Q_K\) are instead the ladder operators truncated at \(K\);
the same basis calculation with dimension \(K+1\) gives
\eqref{p4:eq:compressed-heisenberg-boundary}.  A further sandwich by \(Q_R\),
\(R<K\), removes that new boundary projector.  The matrix-element estimate is
immediate.  The full
operator-norm defect has norm \(\hbar d\), so no full-norm CCR limit is asserted.
\end{proof}

Canonical commutators are therefore controlled reciprocal-position contractions of
the finite oscillator calculus, not exact primitive finite-dimensional relations.



\paragraph{Born reciprocity and dual sectors}
\label{p4:sec:part2-reciprocity}

The framework retains both momentum-like and reciprocal-position sectors.
The compact momentum side is physically fixed at the minimal compact sector
\(m_p=1\), whereas the reciprocal position side carries the framework's
resolution unfolding.  It starts at the derived seed \(m_x=1\).  On
\(\mathsf{SF}_{\Lambda_F}\), it grows by the exact Paper-III-F staircase when the
actual homogeneous radius crosses successive spectral thresholds; otherwise
a large-\(m_x\) trajectory must be supplied separately.
The canonical pair language is therefore best regarded as the large-\(m_x\) dual-sector presentation of an underlying compact finite geometry, not as the deepest level of the theory.


\paragraph{Typed reciprocity at unequal finite levels.}
The generated Born map is inherited by licensed unmarked descendants through
Paper~II-A, Theorem~\ref{p2a:thm:generated-born-inheritance}.  In the
\emph{auxiliary mathematical exchange family}, unequal total-weight labels are
retained marks, so exchange is not an internal automorphism of one marked
presentation.  Paper~II-A,
Proposition~\ref{p2a:prop:typed-born-branch-transport}, together with
Paper~I-C, Theorem~\ref{p1b:thm:unequal-exchange}, gives the comparison
\[
 \mathcal H_{m_x,x}\otimes\mathcal H_{m_p,p}
 \longleftrightarrow
 \mathcal H_{m_p,x}\otimes\mathcal H_{m_x,p}
\]
and transports every licensed finite operator block with the exchanged
projection.  This displayed \(\mathcal H_{m,x}\) family is a comparison
construction; it is not the physical reciprocal-position resolution object.
On the physical branch, the position-side scalar datum is the Born-derived
spectral window \(\mathcal P_x(m_x)\), while the microscopic compact momentum
carrier remains \(\mathcal H_{1,p}\) with \(m_p=1\).  Thus the shorthand
level pair \((m_x,1)\) has a reciprocal comparison branch \((1,m_x)\) only in
the typed Born-transport sense; Born reciprocity does not create a second
microscopic position carrier, force \(m_x=1\), create a second momentum
occurrence, or identify the two differently typed cutoffs.  Internal Born
parity statements require an actually Born-stable subspace or comparison
object.




\paragraph{Scalar Born-pair compact-boson quantization}
\label{p4:sec:scalar-born-pair-compact-boson}

The scalar branch classified in Paper~II and realized in Paper~III, together with its Paper-III reciprocal completion and physical orientation, supplies the compact scalar Born pair and its integer character lattice.  The zero-mode lattice is framework-derived.  The chiral oscillator tower, Euclidean world-sheet torus, and Dedekind factor belong to the standard compact-boson CFT completion and are stated below as an explicit interface hypothesis rather than derived from two compact circles alone.

\begin{lemma}[Integer character lattice of the scalar Born pair]
\label{p4:lem:scalar-born-character-lattice}
Let the scalar Born pair be the compact conjugate pair
\(S^1_t\times S^1_E\).  Its physical periods obey the Paper-III relation
\[
T_{\rm per}E_{\rm per}=2\pi\hbar.
\]
After dividing by this phase unit, write the dimensionless reciprocal radii
as \(\widehat R_t\widehat R_E=1\).  The finite Born quantization of the pair
has intrinsic character labels
\[
  (n,w)\in\mathbb Z\oplus\mathbb Z,
\]
where \(n\) labels the character of one scalar circle and \(w\) labels the character of the conjugate scalar circle.  The exchange of the Born pair acts by
\[
  (n,w;R)\longmapsto (w,n;R^{-1}).
\]
\end{lemma}

\begin{proof}
Paper~1B, Theorem~\ref{p1b:thm:group}, gives the primitive
$\mathbb Z^2$ character lattice and factor exchange. The physical period
calibration remains $T_{\rm per}E_{\rm per}=2\pi\hbar$ from Paper~III.
In the resulting reciprocal-radius convention
$\widehat R_t\widehat R_E=1$, factor exchange sends $R$ to $R^{-1}$,
giving the displayed action on the locally declared labels.
\end{proof}


\begin{theorem}[Native Born-eigencharge lattice and scalar deck parity]
\label{p4:thm:v190-born-eigencharge-lattice}
For the native scalar character $(n,w)\in\mathbb Z^2$, define
\[
 p_+:=n+w,\qquad p_-:=n-w.
\]
Then
\[
 \boxed{p_+\equiv p_-\pmod2,}
\]
and the image of the native character lattice in the Born eigencharge plane is
exactly
\[
 \boxed{\{(p_+,p_-)\in\mathbb Z^2:p_+\equiv p_-\pmod2\}.}
\]
The common parity is the scalar deck parity of the two-fold Born-eigenphase
cover.  The plus quotient $n=w$ lies on $p_-=0$, while the relative quotient
$n=-w$ lies on $p_+=0$; both quotient axes are deck-even because their
nonzero eigencharge is even.  Therefore the native relative scalar comparison
band does not see the deck-odd parent coset.
\end{theorem}

\begin{proof}
The integer map and same-parity image are Paper~1B,
Theorem~\ref{p1b:thm:eigenphase}; its deck and relative-quotient
interpretations are Corollary~\ref{p1b:cor:relative-characters}.
For the quotient axes used here, $n=w$ gives $p_-=0$ and $p_+=2n$,
whereas $n=-w$ gives $p_+=0$ and $p_-=2n$. Thus the comparison band
on the relative quotient has even deck parity.
\end{proof}

\begin{theorem}[Conditional scalar Born-pair compact-boson interface]
\label{p4:thm:scalar-born-compact-boson-partition}
Assume, in addition to the scalar Born character lattice of
Lemma~\ref{p4:lem:scalar-born-character-lattice}, the standard free-scalar chiral
completion on a Euclidean world-sheet torus with modulus \(\tau\), and choose the
action normalization in which
\begin{equation}
\label{p4:eq:scalar-born-left-right-charges}
p_L=\frac nR+wR,
\qquad
p_R=\frac nR-wR,
\qquad
(n,w)\in\mathbb Z^2.
\end{equation}
Then the compact-boson interface partition function is
\begin{equation}
\label{p4:eq:scalar-born-compact-boson-partition}
Z_{\rm Born}(R;\tau)
=
\frac{1}{|\eta_{\rm D}(\tau)|^2}
\sum_{n,w\in\mathbb Z}
q^{p_L^2/4}\,\bar q^{p_R^2/4},
\qquad q=e^{2\pi i\tau},
\end{equation}
where \(\eta_{\rm D}(\tau)=q^{1/24}\prod_{k\ge1}(1-q^k)\).  Finite character and
oscillator windows approximate this expression on the declared CFT interface; the
infinite oscillator completion is not primitive finite-carrier data.
\end{theorem}

\begin{proof}
The native scalar lattice is supplied by
Lemma~\ref{p4:lem:scalar-born-character-lattice}. The specified chiral
completion supplies the oscillator trace $|\eta_{\rm D}(\tau)|^{-2}$,
the modulus and the quadratic normalization in
\eqref{p4:eq:scalar-born-left-right-charges}. Its zero-mode trace is
therefore exactly the theta-series completion in Paper~1B,
Theorem~\ref{p1b:thm:theta}, giving
\eqref{p4:eq:scalar-born-compact-boson-partition} on this declared
interface.
\end{proof}



\begin{corollary}[Self-dual left/right zero modes as Born eigencharges]
\label{p4:cor:v190-selfdual-born-eigencharges}
On the declared compact-boson interface at the self-dual presentation $R=1$,
\[
 \boxed{p_L=p_+=n+w,\qquad p_R=p_-=n-w.}
\]
Hence the native same-parity condition becomes $p_L\equiv p_R\pmod2$,
scalar Born exchange preserves $p_L$ and reverses $p_R$, the plus quotient is
the $p_R=0$ axis, and the relative quotient is the $p_L=0$ axis.  The
identification of these native eigencharges with chiral left/right CFT momenta
remains conditional on the compact-boson interface; the underlying lattice
and parity statement are framework-native.
This is the $R=1$ specialization of Paper~1B,
Theorems~\ref{p1b:thm:eigenphase} and~\ref{p1b:thm:theta}, with the
left/right momentum interpretation supplied by the present completion.
\end{corollary}

\begin{corollary}[Modular invariance of the scalar Born-pair interface partition function]
\label{p4:cor:scalar-born-modular-invariance}
For every scalar radius representative \(R>0\), the compact-symbol interface partition function
\(Z_{\rm Born}(R;\tau)\) of \eqref{p4:eq:scalar-born-compact-boson-partition} is invariant under the modular generators
\[
  T:\tau\longmapsto \tau+1,
  \qquad
  S:\tau\longmapsto -\frac1\tau .
\]
Equivalently,
\[
  Z_{\rm Born}(R;\tau+1)=Z_{\rm Born}(R;\tau),
  \qquad
  Z_{\rm Born}(R;-1/\tau)=Z_{\rm Born}(R;\tau).
\]
\end{corollary}

\begin{proof}
The partition function has exactly the oscillator factor, modulus
and radius normalization of Paper~1B, Theorem~\ref{p1b:thm:theta}.
Its Poisson-summation and Dedekind-eta transformation identities establish
invariance under both generators for every $R>0$.
\end{proof}

\begin{corollary}[Scalar Born exchange as circle T-duality]
\label{p4:cor:scalar-born-t-duality}
Under scalar Born exchange,
\[
  R\mapsto R^{-1},
  \qquad
  n\leftrightarrow w,
\]
the partition function \eqref{p4:eq:scalar-born-compact-boson-partition} is invariant.  Thus the scalar Born-pair exchange is ordinary compact-boson circle T-duality in the scalar sector.
\end{corollary}

\begin{proof}
Apply the radius-inversion identity in Paper~1B,
Theorem~\ref{p1b:thm:theta}, to the declared completion. The native
Born exchange swaps $n,w$ and sends $R$ to $R^{-1}$, precisely the
transformation fixing $p_L$ and reversing $p_R$ in that theorem.
\end{proof}

\begin{remark}[Scope]
\label{p4:rem:scalar-born-duality-scope}
The integer character lattice and Born exchange are native scalar Born-pair results.  Modular invariance, the Dedekind oscillator factor, and the full circle T-duality statement apply after the standard free-scalar chiral CFT completion and its normalization have been supplied.  No bare radius inversion for the full \(S^3\) compact momentum carrier, WZW background, string compactification, or change of the protected selection \(m_p=1\) is asserted.
\end{remark}




\paragraph{Conclusion and Export Ledger}
\label{p4:part:conclusion-export-ledger}

\paragraph{Orientation.}
The final part records what Paper~IV exports to the later sector and continuum papers: the finite quantum core, representation/interface tools, resolution-change mechanisms, and the boundaries of experimental or phenomenological specialization.


% BEGIN SC v1.80 CLOCK-FREE-READOUT


\subsection{Finite clock readout, free generator, and comparison-band trace}
\label{p4:sec:v180-clock-free-readout}

\begin{proposition}[Exact Hopf-charge state quotient]
\label{p4:prop:exact-hopf-charge-state-quotient}
Let $P_q$, $q\in\{-2,-1,0,1,2\}$, be the generated spectral projectors of
the fixed $m_p=1$ momentum Hopf charge.  For a positive parent density
$\rho$ with $\operatorname{Tr}\rho=1$, put
\[
 R(\rho)_q=\operatorname{Tr}(\rho P_q),\qquad
 \sigma(p)=\sum_q\frac{p_q}{d_q}P_q,
 \quad(d_q)=(1,3,7,3,1).
\]
Then $R\sigma=\mathrm{id}$ on the probability simplex and every charge effect
$E=\sum_qe_qP_q$ obeys
$\operatorname{Tr}(\sigma(R\rho)E)=\operatorname{Tr}(\rho E)$.
This is the complete state quotient for the charge-readout algebra only; it is
not a preparation process or a realization of the full phase-sensitive scalar
role.
\end{proposition}

\begin{corollary}[Clock-family free Hamiltonian and same-action Noether charges]
\label{p4:cor:clock-family-free-hamiltonian}
On the existing admitted stationary scalar free block,
\[
 \boxed{H_{\rm free}
 =\sqrt{\sum_iG_i^2+\mu^2c_{\rm op}^4I}
 =\sqrt{\hbar^2c_{\rm op}^2\Delta_a+\mu^2c_{\rm op}^4I}.}
\]
Hence $[H_{\rm free},Q_i]=0$ and the represented chiral spatial charges are
conserved by native Noether on this same action domain.  On degree $k$,
\[
 \mathcal E_k^2=E_\star^2k(k+2)+\mu^2c_{\rm op}^4.
\]
For a Hopf weight $q$ in that degree,
\[
 \mathcal E_k^2=(qE_\star)^2
 +E_\star^2[k(k+2)-q^2]+\mu^2c_{\rm op}^4,
\]
where the second term is spatial kinetic content rather than a new rest mass.
\end{corollary}

\begin{corollary}[Matched spatial and free-energy window]
\label{p4:cor:matched-spatial-free-energy-window}
On the same free spectral block, monotone functional calculus gives
\[
 \boxed{
 \mathbf1_{[0,\Lambda_F^2]}(\Delta_a)
 =\mathbf1_{[0,\mathcal E_F]}(H_{\rm free}),\qquad
 \mathcal E_F=\sqrt{\hbar^2c_{\rm op}^2\Lambda_F^2+
 \mu^2c_{\rm op}^4}.}
\]
Thus a specified spatial sharp frontier fixes its free-Hamiltonian energy
window; it does not truncate independent Fourier tests on the compact energy
circle.  The retained primitive degree-one frontier requires the already
stated threshold correction before $a\Lambda_F\ge\sqrt3$.
\end{corollary}

\begin{remark}[Inherited traces at the level-two algebra comparison]
\label{p4:rem:chiral-algebra-inherited-trace-comparison}
At the position comparison level $m_{\rm cmp}=2$, the abstract clock algebra
is the same $\mathbb C\oplus M_2\oplus M_3$ as on the fixed momentum
carrier, but the represented traces differ:
\[
 \tau_p(a_0,a_1,a_2)=
 \frac{6a_0+3\operatorname{tr}_2a_1+\operatorname{tr}_3a_2}{15},
\]
\[
 \tau_{x,2}(a_0,a_1,a_2)=
 \frac{a_0+2\operatorname{tr}_2a_1+3\operatorname{tr}_3a_2}{14}.
\]
This is why the exact algebra isomorphism at $m_{\rm cmp}=2$ cannot be used
to identify the $14$-dimensional position band with the $15$-dimensional
$m_p=1$ momentum Hilbert space or to replace either parent action trace.
The algebra and its two inherited traces are Paper~1B,
Theorem~\ref{p1b:prop:degree-two}, applied to the two carriers
identified in Paper~II-C.
The exact conversion is now explicit in Paper~I-C,
Corollary~\ref{p1b:cor:degree-two-trace-density}. Under the same
generator-preserving identification,
\[
 h=\left(\frac{28}{5},\frac75I_2,\frac{14}{45}I_3\right),
 \qquad \tau_p(a)=\tau_{x,2}(ha).
\]
A state with density \(\rho_p\) relative to \(\tau_p\) has
density \(\rho_{x,2}=h\rho_p\) relative to \(\tau_{x,2}\).
Thus the same functional has exactly specified density representatives
in the two inherited normalizations. The element \(h\) is central
and supplies no basis rotation inside an irreducible action block.
\end{remark}


\begin{corollary}[Abstract clock-state equivalence at the degree-two comparison]
\label{p4:cor:v185-clock-state-space-equivalence}
Let
\(
 \Phi_2:\mathcal A_{{\rm clk},x}(2)\overset\sim\longrightarrow
 \mathcal A_{{\rm clk},p}
\)
be the generator-preserving $*$-isomorphism of Paper~II-C.  Pullback gives an
affine bijection of full algebraic state spaces,
\[
 \boxed{
 \Phi_2^*:S(\mathcal A_{{\rm clk},p})
 \overset\sim\longrightarrow S(\mathcal A_{{\rm clk},x}(2)),
 \qquad \omega\longmapsto\omega\circ\Phi_2.}
\]
The inverse is pullback by $\Phi_2^{-1}$.  This abstract state-space
correspondence does not identify the inherited normalized carrier traces
$\tau_p$ and $\tau_{x,2}$ displayed above and therefore does not identify the
$15$- and $14$-dimensional parent Hilbert carriers.
\end{corollary}
\begin{proof}
Paper~1B, Theorem~\ref{p1b:prop:degree-two}, gives the state-space
bijection for the isomorphism constructed in Paper~II-C. Explicitly,
composition with a unital $*$-isomorphism and its inverse preserves
positive normalized functionals in both directions. The preceding
inherited trace formulas remain distinct, so this affine equivalence
preserves the separate fifteen- and fourteen-dimensional carrier roles.
\end{proof}

In density coordinates this state-space bijection is
\(\rho_{x,2}=h\rho_p\), after the displayed algebra
identification, with inverse \(\rho_p=h^{-1}\rho_{x,2}\).
This uses the fixed trace conversion above and preserves the
separate carrier interpretations.

\begin{corollary}[Frontier-shell entry at one free-energy ceiling]
\label{p4:cor:frontier-shell-entry-free-energy}
For every later degree admitted by the existing nondecreasing sharp frontier,
\[
 a_k=\frac{\sqrt{k(k+2)}}{\Lambda_F}
 \quad\Longrightarrow\quad
 \mathcal E_k(a_k)=\mathcal E_F.
\]
At a later radius in the same comparison family,
\[
 \mathcal E_k(a)^2-\mu^2c^4
 =(\mathcal E_F^2-\mu^2c^4)(a_k/a)^2.
\]
Spectral admission is not particle creation; occupation, transition,
backreaction, and the actual background history remain separately typed.
\end{corollary}
% END SC v1.80 CLOCK-FREE-READOUT



\subsection{Native coordinate control and complete finite readout}
\label{p4:sec:v226-control-readout}
\begin{corollary}[Exact native control algebras and their conditional reach]
\label{p4:cor:v226-native-controls}
Hamiltonian controls proportional only to the four intrinsic coordinates
have Lie algebra $\mathfrak{so}(11)\oplus\mathfrak{usp}(4)$ of dimension 65.
They do not have the full 137-dimensional native unitary Lie algebra merely
because their associative envelope is $M_{11}\oplus M_4$.
If the actual reversible control family also contains the centered square
$Q_0$ and the oriented quartic word $\Omega_4$, its Lie algebra is
$\mathfrak{s}(\mathfrak u(11)\oplus\mathfrak u(4))$, of dimension 136.
Under the stated concatenation and reversal conditions, the closure of
reachable operations is the corresponding compact connected group.
The coordinates with $\Omega_4$ alone already generate all connected inner
automorphisms of the native observable algebra
(Corollary~\ref{p1:cor:v226-native-inner-controls}); the extra $Q_0$ direction
is a relative central phase. A global scalar control extends the Hilbert-space
Lie algebra to $\mathfrak u(11)\oplus\mathfrak u(4)$.
\end{corollary}
\begin{proof}
Theorem~\ref{p1:thm:v226-control-lie} gives the exact algebras and their
block representations. These particular connected groups are compact and
closed; apply the supplied reversible finite-control theorem. A control
coefficient, time cost or pulse protocol is still part of the actual process.
\end{proof}

\begin{corollary}[Optimal-radius localization has zero protected chirality information]
\label{p4:cor:v226-chirality-fisher}
For the normalized protected family $\rho_s=(I+s\chi)/4$, a measurement
with exact intrinsic first moments and outcome radius $r$ has classical
Fisher information at $s=0$ at most $1-3R_{\rm mom}^2/r^2$.
At $r=\sqrt3R_{\rm mom}$ every such measurement gives identical
probabilities for $\rho_+$ and $\rho_-$.
\end{corollary}
\begin{proof}
This is Theorem~\ref{p2c:thm:v226-chirality-radius}, including arbitrary
noncovariant outcomes and additional labels. It concerns this classical
readout, not the quantum Fisher information of every available instrument.
\end{proof}
The complete native algebra instead admits the 137-outcome reconstruction
of Theorem~\ref{p1:thm:v226-tomography}, with explicit dual operators in
Theorem~\ref{p1b:thm:v226-dual-coefficients}. That is full native-state
tomography, while the sharp-radius position task has an unattained
informational-completeness infimum as stated in
Theorem~\ref{p2c:thm:v226-ic-radius-infimum}.
An effect construction does not determine its post-measurement instrument,
prepare a particular state, or actualize all independent Hamiltonian controls.

% END FRAMEWORK v2.226: KINEMATIC CLOSURES

% BEGIN FRAMEWORK v2.227: DOWNWARD CLOSURES


\section{Measurement, tomography and structural export}
\label{p4a:sec:v238-a8}
\noindent\textit{Source and claim boundary.} Operational state-update, measurement and tomography certificates are distinguished from the occurrence of instruments, frequencies, histories and external detector response. Paper~IV-B owns physical quantum processes; Paper~IV-C preserves full auxiliary derivations.

\paragraph{Measurement, Protocols, Scattering, and Internal-Sector Interfaces}
\label{p4:part:measurement-scattering-internal}

\paragraph{Orientation.}
This part groups finite apparatus and protocol interfaces, scattering windows, unstable channels, internal sectors, and path-sum presentations.  These results explain how ordinary laboratory and sector language is represented inside the finite calculus after a channel, window, or distinguished finite sector has been specified.



\subsection{Measurement update and effective classicality}
\label{p4:sec:part2-classicality}



\paragraph{Accessible subalgebras and coarse observables}
\label{p4:sec:part2-coarse-subalgebras}

A laboratory setup does not generally access the full algebra $\mathcal A_m^{\mathrm{env}}$.
Instead it first induces a system-side operator system; a physical
\(*\)-algebra follows only when the apparatus test grammar also generates the
required adjoint and sequential-product closure.
Classical appearance arises when the accessible algebra is approximately commuting and the environment suppresses off-diagonal terms relative to that subalgebra.


\begin{theorem}[Accessible operator systems and algebras are induced by the apparatus]
\label{p4:prop:accessible-algebras-selector-domain}
Let a laboratory channel be represented by a finite apparatus subsystem
\(\mathcal H_A\), initial apparatus state \(\sigma_A\), apparatus pointer
operator system \(S_A\subseteq\End(\mathcal H_A)\) containing
\(\{\Pi_\alpha^A\}\), and system--apparatus unitary \(U_{SA}\) on
\(\mathcal H_S\otimes\mathcal H_A\).  Define the Heisenberg readout map
\[
 R_{\mathcal E}(X)
 :=
 \operatorname{Tr}_A\!\left[
 (\mathbf1_S\otimes\sigma_A)U_{SA}^*XU_{SA}
 \right].
\]
Then \(R_{\mathcal E}\) is unital and completely positive.  The induced
system effects are
\begin{equation}
\label{p4:eq:induced-apparatus-effects}
  E_\alpha
  =R_{\mathcal E}(\mathbf1_S\otimes\Pi_\alpha^A).
\end{equation}
They satisfy \(E_\alpha\ge0\) and
\(\sum_\alpha E_\alpha=\mathbf1_S\).  If
\(Q_i^S=(Q_i^S)^*\) are the system-side charges or exact system selectors
resolved by the same channel, its exact
induced system readout is the operator system
\begin{equation}
\label{p4:eq:derived-accessible-operator-system}
 S_{\rm acc}(\mathcal E)
 :=
 \operatorname{span}_{\CC}
 \left\{
 \mathbf1_S,\ Q_i^S,\
 R_{\mathcal E}(\mathbf1_S\otimes b):b\in S_A
 \right\}
 \subseteq\End(\mathcal H_S).
\end{equation}
On a theorem domain whose constituted physical test grammar closes these
system-side induced tests under adjoints and sequential products, the
accessible algebra is the finite unital \(*\)-algebra
\begin{equation}
\label{p4:eq:derived-accessible-algebra}
  \mathcal A_{\rm acc}(\mathcal E)
  =
  C^*\!\left(S_{\rm acc}(\mathcal E)\right)
  \subseteq\End(\mathcal H_S).
\end{equation}
Without that product-closure doctrine,
\(S_{\rm acc}(\mathcal E)\) is the exact physical readout object and
\(C^*(S_{\rm acc})\) is only its mathematical presentation envelope.  Thus
every generator lies on the common system carrier: a literal apparatus
pointer \(b\) enters only through
\(R_{\mathcal E}(\mathbf1_S\otimes b)\).  An accessible system algebra is
therefore computed from the finite readout geometry; it is not an arbitrary
subalgebra and not an alternative maximal structural algebra.
\end{theorem}


\begin{proof}
Write
\(\sigma_A=\sum_rq_r\lvert r\rangle\langle r\rvert\) and define
\[
 V_r:\mathcal H_S\longrightarrow\mathcal H_S\otimes\mathcal H_A,
 \qquad
 V_r\psi=\sqrt{q_r}\,\psi\otimes\lvert r\rangle .
\]
Then
\[
 R_{\mathcal E}(X)
 =\sum_rV_r^*U_{SA}^*XU_{SA}V_r,
 \qquad
 \sum_rV_r^*V_r=\mathbf1_S.
\]
This is a Kraus representation of a unital completely positive map.
For any system state \(\rho_S\), the probability of pointer outcome \(\alpha\) is
\[
\operatorname{Tr}_{SA}\!\left[
U_{SA}(\rho_S\otimes\sigma_A)U_{SA}^*
(\mathbf1_S\otimes\Pi_\alpha^A)
\right]
=
\operatorname{Tr}_S(\rho_S E_\alpha).
\]
This proves the induced-effect formula.  Positivity and normalization also
follow directly from the UCP property and
\(\sum_\alpha\Pi_\alpha^A=\mathbf1_A\).  Since \(S_A\) is unital and
self-adjoint, its transported image together with the system-side \(Q_i^S\)
has the same properties, proving that \(S_{\rm acc}(\mathcal E)\) is an
operator system on \(\mathcal H_S\).

If the physical test grammar supplies the stated product closure, GP2
no-omission retains the generated algebra
\(C^*(S_{\rm acc}(\mathcal E))\).  Without that constructor, mathematical
multiplication cannot promote the operator system to a physical algebra.
When the generated algebra is proper and non-scalar,
Lemma~\ref{p4:lem:proper-subalgebras-require-selectors} identifies the
apparatus/sector distinctions that produced it; they are fixed before the
observed result and are not inserted into the structural envelope.
Consequently changing \(S_{\rm acc}(\mathcal E)\) or
\(\mathcal A_{\rm acc}(\mathcal E)\) changes the experimental channel, not
the upstream generated algebra or the Born readout of Part~I.
\end{proof}


\begin{corollary}[Channel-specific readout of finite observables]
\label{p4:cor:channel-specific-readout}
A finite experiment \(\mathcal E\) measures the readout
\[
  O_{\mathcal E}=\ell_{\mathcal E}(O_{\rm geom}),
\]
where \(O_{\rm geom}\) is the framework-native finite observable or invariant and \(\ell_{\mathcal E}\) is the finite functional induced by the apparatus effects, cuts, reconstruction map, transport kernel, and accessible operator system or generated algebra of Theorem~\ref{p4:prop:accessible-algebras-selector-domain}.  Different experiments may therefore read different finite projections of the same geometric structure.  A comparison with data is GP2-admissible only after the channel functional \(\ell_{\mathcal E}\) has been specified independently of the observed outcome.
\end{corollary}

\begin{proof}
Theorem~\ref{p4:prop:accessible-algebras-selector-domain} constructs the accessible effects and operator system from the finite apparatus channel, and its algebra on the stated product-closure domain.  Any reported number is obtained by applying the corresponding finite probability or estimator functional to the underlying observable.  If two apparatuses induce different effects, coarse grainings, or transport kernels, the resulting functionals may differ even when the native observable is the same.  Because \(\ell_{\mathcal E}\) is fixed by the apparatus geometry before comparison, this is not a post-hoc fitted channel choice.
\end{proof}




\paragraph{Projective supports and selective-readout basins}
\label{p4:sec:part2-projective-actualization}

The preceding theorem fixes the accessible effects and hence the Born weights of a
finite apparatus channel.  This subsection records the corresponding selective
readout geometry.  It does not add a subquantum local value assignment.  It shows
that, once a finite theorem-domain readout support is supplied, the Hopf/projective
structure already determines a canonical unresolved hit space, a canonical simplex
measure, and a deterministic basin rule whose unresolved frequencies are the Born
weights.  The result is a canonical GP2-compatible representation of selective
readout, not a uniqueness theorem for all measurable partitions and not an ontic
collapse mechanism.

\begin{definition}[Finite readout support]
\label{p4:def:finite-readout-support}
A \emph{finite readout support} for a theorem-domain channel is a finite-dimensional
complex Hilbert space \(\mathcal R\) generated by the already constituted carrier,
by a declared apparatus/interface, or by a later sector support.  A resolved
readout context on \(\mathcal R\) is an orthonormal channel family
\[
        \mathcal C=\{e_1,\ldots,e_d\},
        \qquad d=\dim_\CC\mathcal R,
\]
whose one-dimensional projectors generate the corresponding finite pointer
algebra.  The index \(d\) is a readout-support rank.  It is not a new compact
momentum truncation; when the microscopic compact sector is meant, the fixed
input remains \(m_p=1\).
\end{definition}

\begin{definition}[Projective hit space and pointer simplex]
\label{p4:def:projective-hit-simplex}
For a finite readout support \(\mathcal R\), the unresolved hit space is the
projective ray space
\[
        \PP(\mathcal R)
        :=(\mathcal R\setminus\{0\})/\CC^\times .
\]
For a resolved context \(\mathcal C=\{e_i\}_{i=1}^d\), define the pointer simplex
\[
        \Delta_{d-1}
        :=\Bigl\{(x_1,\ldots,x_d)\in\RR_{\ge0}^d:
        \sum_i x_i=1\Bigr\}
\]
and the squared-amplitude moment map
\begin{equation}
\label{p4:eq:projective-hit-moment-map}
        \mu_{\mathcal C}([u])
        :=
        \left(
        \frac{|\langle e_1,u\rangle|^2}{\langle u,u\rangle},
        \ldots,
        \frac{|\langle e_d,u\rangle|^2}{\langle u,u\rangle}
        \right)
        \in \Delta_{d-1} .
\end{equation}
\end{definition}

\begin{lemma}[State-independent unitary-natural hit measure]
\label{p4:lem:no-insertion-hit-measure}
Suppose the unresolved hit measure on \(\PP(\mathcal R)\) is fixed before the
prepared state, Born point, outcome labels, and apparatus response are inserted, and
is natural under every unitary relabeling of \(\mathcal R\).  Then it is the normalized
Fubini--Study measure \(\mu_{\rm FS}\).
\end{lemma}

\begin{proof}
The projective space \(\PP(\mathcal R)\) is the compact homogeneous space
\(U(d)/(U(1)\times U(d-1))\), which has a unique normalized invariant probability
measure.  That measure is Fubini--Study.  State- or apparatus-dependent covariant
densities can be constructed only after those data are supplied; the stated role
separation excludes them from the unresolved hit measure and places all such
dependence in the Born point and the basin map.
\end{proof}

\begin{lemma}[Complex Hopf rays give uniform simplex volume]
\label{p4:lem:fs-simplex-pushforward}
For every resolved context \(\mathcal C\), the pushforward of normalized
Fubini--Study measure by \(\mu_{\mathcal C}\) is the normalized affine volume
measure \(\nu_{\Delta}\) on \(\Delta_{d-1}\):
\[
        (\mu_{\mathcal C})_*\mu_{\rm FS}=\nu_{\Delta}.
\]
\end{lemma}

\begin{proof}
Choose coordinates \(u_i=\langle e_i,u\rangle\) and write, on the unit sphere in
\(\mathcal R\cong\CC^d\),
\[
        u_i=\sqrt{x_i}\,e^{i\phi_i},
        \qquad x_i\ge0,
        \qquad \sum_i x_i=1 .
\]
The complex Euclidean measure satisfies
\[
        \prod_i d^2u_i
        =
        \prod_i r_i\,dr_i\,d\phi_i
        =
        2^{-d}\prod_i dx_i\,d\phi_i .
\]
After imposing \(\sum_i x_i=1\) and quotienting the common Hopf phase, the remaining
measure in the variables \(x_i=|u_i|^2\) is proportional to
\[
        \delta\!\left(1-\sum_i x_i\right)\prod_i dx_i,
\]
which is the affine volume form on the simplex.  Normalization gives the displayed
pushforward identity.
\end{proof}

\begin{definition}[Barycentric selective-readout basins]
\label{p4:def:barycentric-actualization-basins}
Let \(v_i\) be the vertices of \(\Delta_{d-1}\).  For a Born point
\[
        b=(b_1,\ldots,b_d)\in\Delta_{d-1},
\]
define the \(i\)-th barycentric basin by
\[
        R_i(b):=\operatorname{conv}\bigl(b,
        v_1,\ldots,\widehat{v_i},\ldots,v_d\bigr).
\]
The hat means that \(v_i\) is omitted.
\end{definition}

\begin{lemma}[Barycentric basin volume]
\label{p4:lem:barycentric-basin-volume}
The regions \(R_i(b)\) triangulate \(\Delta_{d-1}\) up to common faces, and
\[
        \frac{\operatorname{Vol}(R_i(b))}
             {\operatorname{Vol}(\Delta_{d-1})}
        = b_i .
\]
\end{lemma}

\begin{proof}
The point \(b\) has barycentric coordinates \((b_i)_i\) relative to the vertices
\((v_i)_i\).  Replacing the vertex \(v_i\) of the full simplex by \(b\) gives a
subsimplex whose determinant is linear in the replaced vertex and whose normalized
volume is the corresponding barycentric coordinate.  Equivalently, translate to a
basis of edge vectors of \(\Delta_{d-1}\); the determinant for \(R_i(b)\) is the
full determinant multiplied by the coefficient \(b_i\).  The subsimplices meet only
on lower-dimensional faces and fill the simplex, proving both the triangulation and
the volume formula.
\end{proof}

\begin{theorem}[Canonical projective selective-readout probability representation]
\label{p4:thm:projective-readout-actualization}
Let \(\mathcal R\) be a constituted or explicitly supplied finite readout
support with resolved context
\(\mathcal C=\{e_i\}_{i=1}^d\).  Let a finite apparatus channel induce effects
\(E_i\ge0\), \(\sum_iE_i=\Id\), and let \(\rho\) be the finite state on the
corresponding system algebra.  Put
\[
        b_i:=\Tr(\rho E_i),
        \qquad b=(b_i)_i\in\Delta_{d-1}.
\]
For an unresolved hit \([u]\in\PP(\mathcal R)\), define
\[
        x:=\mu_{\mathcal C}([u])\in\Delta_{d-1}
\]
and select the realized channel by
\begin{equation}
\label{p4:eq:projective-actualization-selector}
        \Sigma_{\mathcal E,\mathcal R}(\rho,[u])=i
        \quad\Longleftrightarrow\quad
        x\in R_i(b),
\end{equation}
with boundary ties assigned by any fixed theorem-domain convention; the boundary has
\(\nu_\Delta\)-measure zero.  Then the unresolved hit average gives
\[
        \mu_{\rm FS}\bigl\{[u]:
        \Sigma_{\mathcal E,\mathcal R}(\rho,[u])=i\bigr\}
        =
        \Tr(\rho E_i).
\]
Thus the selective branch can be represented canonically as a deterministic basin
hit on the projective readout support, while the unresolved frequencies are exactly
the finite Born weights.  This asserts existence and canonical status under the stated state-independent
unitary-natural measure premise; it does not assert that no other measurable partition
of the simplex can have the same volumes.
\end{theorem}

\begin{proof}
By Lemma~\ref{p4:lem:no-insertion-hit-measure}, the unresolved hit measure on
\(\PP(\mathcal R)\) is \(\mu_{\rm FS}\).  By
Lemma~\ref{p4:lem:fs-simplex-pushforward}, its image under \(\mu_{\mathcal C}\) is
uniform affine simplex volume.  Therefore
\[
\begin{aligned}
\mu_{\rm FS}\{[u]:\Sigma_{\mathcal E,\mathcal R}(\rho,[u])=i\}
&=
\nu_\Delta(R_i(b))\\
&=
\frac{\operatorname{Vol}(R_i(b))}{\operatorname{Vol}(\Delta_{d-1})}\\
&= b_i\\
&=\Tr(\rho E_i),
\end{aligned}
\]
where the third equality is Lemma~\ref{p4:lem:barycentric-basin-volume}.  The rule
therefore selects one branch for a fixed unresolved hit and reproduces the finite
Born readout after unresolved averaging.
\end{proof}

\begin{corollary}[Projective obstruction and the actual frequency route]
\label{p4:cor:unique-ergodic-born-frequencies}
For \(\dim_{\mathbb C}\mathcal R=d>1\), there is no continuous map
\(T:\PP(\mathcal R)\to\PP(\mathcal R)\) whose unique invariant probability
measure is Fubini--Study volume.  Hence the former continuous uniquely
ergodic projective implementation is retired as an empty nontrivial domain.
The one-shot probability representation of
Theorem~\ref{p4:thm:projective-readout-actualization} is unchanged.

For an actual stationary successive finite instrument, the event-sector
certificate of Theorem~\ref{p4:thm:v2210-lit-instrument-mixing} instead
gives quantitative finite frequency bounds, and
Corollary~\ref{p4:cor:v2211-event-frequency-strong-law} gives almost-sure
Born frequencies whenever the computed event-cyclic transfer radius is
strictly less than one.  Neither independence of successive outcomes nor
mixing of unused sectors is required.
\end{corollary}

\begin{proof}
Put \(n=d-1\).  The integral cohomology ring is
\(H^*(\mathbb{CP}^n;\mathbb Z)=\mathbb Z[h]/(h^{n+1})\),
\(\deg h=2\).  A continuous self-map acts by \(T^*h=kh\) for an integer
\(k\), and by \(k^j\) on degree \(2j\).  Thus its square has Lefschetz number
\[
 L(T^2)=\sum_{j=0}^{n}k^{2j}\ge1.
\]
The Lefschetz fixed-point theorem \cite[Theorem~2C.3]{hatcher2002algebraic}
therefore gives a fixed point of \(T^2\).  Its \(T\)-orbit has period one
or two and supports an invariant atomic probability.  This differs from
non-atomic Fubini--Study volume for \(n\ge1\), excluding the asserted
unique invariant law.  At \(d=1\) projective space is a point and the
statement is trivial.  The finite instrument claims follow from the cited
transfer-covariance calculation and the strong-law proof in Paper~IV-B; they use
actual instrument cylinder probabilities rather than a projective map.
\end{proof}

\begin{remark}[Scope of the selective-readout representation]
\label{p4:rem:actualization-status}
Theorem~\ref{p4:thm:projective-readout-actualization} is a one-shot
representation after the state, apparatus effects, readout support and
channel are fixed.  It does not define collapse dynamics, an actual hit law
or noncontextual pre-existing values.  The selector is contextual to
\((\mathcal E,\mathcal R,\mathcal C)\) and need not be the unique partition
with those basin volumes.  The continuous Fubini--Study uniquely ergodic
route is excluded by Corollary~\ref{p4:cor:unique-ergodic-born-frequencies}.
The live repeated-frequency route consumes an actual stationary instrument
and its computed event-sector contraction or covariance bound; its finite
error and almost-sure consequences are derived in Paper~IV-B. The state/effect
Born evaluation itself consumes no such mixing premise.
\end{remark}



Theorem~\ref{p4:thm:v255-sharp-event-frequency} sharpens that sufficient
mixing test: for every occurring finite stationary instrument, the exact
deterministic frequency criterion is $d_1=\ell(Qv)=0$, where $Q$ is the
fixed-mode projector of its event-cyclic transfer. Periodic event modes
are allowed. The theorem proves both directions, the exact variance
floor $d_1$, and a finite $O(N^{-1})$ sampling bound when it vanishes.

Theorem~\ref{p4:thm:v258-nonstationary-frequency-law} removes the stationary
preparation restriction for every repeated time-homogeneous finite instrument.
For an arbitrary initial state it derives an almost-sure limiting frequency
and every moment from $\mathcal K=\mathcal C\mathcal I_1\mathcal C$,
where $\mathcal C$ is the finite channel's Ces\`aro projector.
Corollary~\ref{p4:cor:v258-asymptotic-frequency-povm} gives its complete
finite atomic law as a harmonic POVM on the initial state. A deterministic
rate and its equality to the initial one-shot probability are then explicit
tests, rather than universal assumptions.

\paragraph{State update as noncommutative conditioning}
\label{p4:sec:part2-conditioning}

For the supplied ideal L\"uders instrument associated with a sharp event
\(P\), the outcome operation is
\(\mathcal I_P(\rho)=P\rho P\).  When \(\Tr(\rho P)>0\), its normalized
selective update is
\begin{equation}
\label{p4:eq:luders}
\rho\mapsto \rho'=\frac{P\rho P}{\Tr(\rho P)}
\end{equation}
and is the noncommutative analogue of Bayesian conditioning.  If
\(\Tr(\rho P)=0\), that conditioned branch is not realized.  The event \(P\)
does not by itself select this instrument.  If \(U\) is any unitary preserving
\(\operatorname{Ran}P\), then
\[
\mathcal I_{P,U}(\rho)=UP\rho PU^*
\]
has the same effect because \((UP)^*(UP)=P\), and its output is again supported
in \(\operatorname{Ran}P\).  Hence repeatability alone does not make the
L\"uders operation unique; its ideal minimal-disturbance or nondemolition
status is additional instrument-domain data.  General outcome-dependent
updates are described by \eqref{p4:eq:instrument}.  What is structural is the
CP-instrument calculus of conditional updates, not a unique ``collapse'' map
attached to \(P\).

The resolution-change, accessibility, and decoherence conditions for effective classicality are derived in Part~\ref{p4:part:part3}.



\subsection{Tomography, metrology, controllability, and protected codes}
\label{p4:sec:part3-control}



\paragraph{Informationally complete tomography}
\label{p4:sec:part3-tomography}

A finite POVM $\{E_a\}_{a=1}^N$ is informationally complete if the probabilities
\begin{equation}
\label{p4:eq:ic-probs}
p_a=\Tr(\rho E_a)
\end{equation}
determine $\rho$ uniquely.


\begin{theorem}[Finite IC reconstruction]
\label{p4:thm:ic-reconstruction}
\leavevmode\par\noindent
Let \(\mathcal H\) be a \(d\)-dimensional sector and let
\(\{E_a\}_{a=1}^N\) be a POVM on \(\mathcal H\).  The following are equivalent:
\begin{enumerate}[label=(\roman*),leftmargin=2.2em]
\item \(\{E_a\}\) is informationally complete;
\item the linear map
\[
T:\mathrm{Mat}(d,\CC)\to \CC^N,
\qquad
T(X):=(\Tr(XE_a))_{a=1}^N,
\]
is injective;
\item the operators \(E_a\) span \(\mathrm{Mat}(d,\CC)\).
\end{enumerate}
When these conditions hold, one may choose dual-frame operators \(\{F_a\}_{a=1}^N\) such that
\begin{equation}
\label{p4:eq:tomography}
X=\sum_{a=1}^N \Tr(XE_a)\,F_a
\qquad
\text{for all }X\in\mathrm{Mat}(d,\CC).
\end{equation}
In particular,
\[
\rho=\sum_{a=1}^N p_aF_a.
\]
The reconstruction operators \(F_a\) are dual-frame operators; they need not be positive and need not themselves form a POVM.
\end{theorem}


\begin{proof}
By definition, the POVM is informationally complete exactly when \(T(\rho_1)=T(\rho_2)\) implies \(\rho_1=\rho_2\) for all states \(\rho_1,\rho_2\).
Since states span the Hermitian operator space and hence the full complex matrix algebra, this is equivalent to injectivity of \(T\) on \(\mathrm{Mat}(d,\CC)\).
Thus (i) and (ii) are equivalent.

Now \(T(X)=0\) if and only if \(\Tr(XE_a)=0\) for every \(a\), which means that \(X\) annihilates \(\operatorname{span}\{E_a\}\) under the nondegenerate bilinear trace pairing \((X,Y)\mapsto\Tr(XY)\).
Therefore \(T\) is injective exactly when this trace annihilator is \(\{0\}\), i.e.\ exactly when \(\{E_a\}\) spans \(\mathrm{Mat}(d,\CC)\).  Equivalently, for physical density matrices, informational completeness means that the real linear span of the Hermitian effects contains the whole real vector space of Hermitian operators.
This proves the equivalence of (ii) and (iii).

Now suppose the equivalent conditions hold.
Because \(T\) is an injective linear map from a \(d^2\)-dimensional space into \(\CC^N\), its image is a \(d^2\)-dimensional subspace.
Choose any linear left inverse \(S:\operatorname{im}(T)\to \mathrm{Mat}(d,\CC)\) with \(S\circ T=\operatorname{id}\).
Writing \(e_a\) for the standard basis vectors of \(\CC^N\), define \(F_a:=S(e_a)\) after extending \(S\) arbitrarily to all of \(\CC^N\).
Then
\[
X=(S\circ T)(X)=\sum_a \Tr(XE_a)\,F_a,
\]
which is exactly \eqref{p4:eq:tomography}.
\end{proof}

Quantum state tomography is repeated Born evaluation on an informationally complete finite observable family.



\paragraph{Finite metrology and Fisher information}
\label{p4:sec:part3-metrology}

If a parameter $\theta$ is encoded by
\begin{equation}
\label{p4:eq:param-encoding}
\rho_{\theta}=e^{-i\theta G}\rho\,e^{i\theta G},
\end{equation}
then the quantum Fisher information $F_Q(\rho,G)$ gives the intrinsic precision scale and, for locally unbiased estimators satisfying the theorem-domain regularity conditions, the Cram\'er--Rao bound reads \cite{Helstrom1976,Holevo1982,BraunsteinCaves1994}
\begin{equation}
\label{p4:eq:qcrb}
\mathrm{Var}(\widehat\theta)\ge \frac{1}{nF_Q(\rho,G)}
\end{equation}
for $n$ independent repetitions.


\begin{theorem}[Pure-state quantum Fisher information on the finite carrier]
\label{p4:thm:pure-qfi}
Let \(\rho=\ket{\psi}\!\bra{\psi}\) be a pure state on a finite sector and let \(G=G^*\) generate the unitary encoding \eqref{p4:eq:param-encoding}.
Then
\begin{equation}
\label{p4:eq:qfi-pure}
F_Q(\ket{\psi},G)=4\,\mathrm{Var}_{\psi}(G).
\end{equation}
Consequently the precision scale for pure-state metrology is controlled by exactly the same variance appearing in the uncertainty and speed-limit bounds.
\end{theorem}


\begin{proof}
\ProofPlan{We read pure-state quantum Fisher information by expanding the unitary-encoded state, extracting the fidelity-loss coefficient, and applying the pure-state QFI formula.}
\begin{proofsteps}

\item \textbf{Pure-state unitary expansion.}
Write \(\ket{\psi_\theta}=e^{-i\theta G}\ket{\psi}\).
Then
\[
\ket{\psi_{\theta+d\theta}}
=
\Bigl(
\mathbf 1-iG\,d\theta-\frac12G^2\,d\theta^2
\Bigr)\ket{\psi}
+
O(d\theta^3).
\]

\item \textbf{Fidelity-loss coefficient.}
Hence
\[
\langle \psi_\theta,\psi_{\theta+d\theta}\rangle
=
1-i\langle G\rangle_\psi d\theta-\frac12\langle G^2\rangle_\psi d\theta^2+O(d\theta^3),
\]
and therefore
\[
|\langle \psi_\theta,\psi_{\theta+d\theta}\rangle|^2
=
1-\bigl(\langle G^2\rangle_\psi-\langle G\rangle_\psi^2\bigr)d\theta^2+O(d\theta^3)
=
1-\mathrm{Var}_\psi(G)\,d\theta^2+O(d\theta^3).
\]

\item \textbf{Quantum Fisher information readout.}
For pure states, the quantum Fisher information is four times the coefficient of \(d\theta^2\) in the fidelity loss:
\[
F_Q(\ket{\psi},G)
=
4\lim_{d\theta\to 0}
\frac{1-|\langle \psi_\theta,\psi_{\theta+d\theta}\rangle|^2}{d\theta^2}.
\]
Substituting the expansion above gives \eqref{p4:eq:qfi-pure}.
\end{proofsteps}

\end{proof}

So metrological sensitivity is again controlled by the same variance structure that underlies the speed-limit and uncertainty bounds.



\subsection{Canonical states and the reduced readout certificate}
\label{p4:sec:v227-state-response}
\begin{corollary}[Canonical symmetry is inherited from the actual Hamiltonian]
\label{p4:cor:v227-state-covariance}
For the normalized full-ground or Gibbs state of an actual finite $H$,
a unitary symmetry of $H$ is a symmetry of the state. With
$\epsilon=\|[H,U]\|$, the respective covariance defects are at most
$2\min(1,\epsilon/\Delta)$ and $2\min(1,\beta\epsilon)$ in trace norm.
Fixed occupied-block canonical mixtures use the weighted within-block
bounds, rather than a gap between different central sectors.
\end{corollary}
\begin{proof}
Apply Theorem~\ref{p1b:thm:v227-state-symmetry} and
Corollary~\ref{p1b:cor:v227-block-covariance}. Proper pure substates of a
degenerate ground eigenspace are outside the full-projector hypothesis.
\end{proof}
At fixed state/support, the independent native central offsets leave even
Liouvillian readouts unchanged; blockwise normalized canonical states with
fixed weights are unchanged too. The globally normalized Gibbs ensemble
can change its weights, and odd sources retain relative offsets
(Proposition~\ref{p1b:prop:v227-central-boundary}). Thus no preferred
central weight follows from this additional invariance.

For pole claims the existing source-cyclic action, moment-pencil and
cofactor certificates already avoid selecting an unused total occupation
member. The finite complete-response certificate of
Theorem~\ref{p1b:thm:v227-finite-response} provides another exact route,
including an unknown reducing dark complement. It requires actual complete
sample data and source incidence; it supplies no missing charged source
word, physical clock or numerical pole by itself.

% END FRAMEWORK v2.227: DOWNWARD CLOSURES



\subsection{Conclusion of Part I}
\label{p4:sec:part1-conclusion}

Part~I establishes the exact generated finite algebra, its full state/effect
order and canonical Born evaluation, finite direct-sum tensor composition,
reduced-state calculus, CP/Kraus representation of supplied processes, and
native kinematic flows together with the blockwise inner form of any supplied
closed flow and the central-gauge classification of its fixed-flow generator.  The full carrier envelope is an optional
\(\mathsf{FENV}\) graded field extension or \(\mathsf{ENV}\) full-event
extension, not a premise of these results.
The remaining parts use the appropriate generated algebra or explicitly
declare an envelope, representation, or application interface.




\subsection{Conclusion}
\label{p4:sec:conclusion}
\label{p4:sec:summary}



\paragraph{Structural finite quantum closure}
\label{p4:sec:conclusion-structural-compatibility}

Paper~IV closes the finite state/effect/operator calculus on the generated
finite test algebra rather than presupposing a full carrier envelope.  At the
protected compact momentum level \(m_p=1\), Paper~II-C gives the exact result
\[
\mathcal A_{\rm kin}
\cong
\CC\oplus M_2\oplus M_2\oplus M_4\oplus M_3\oplus M_3,
\qquad
\dim_\CC\mathcal A_{\rm kin}=43.
\]
Its noncommuting blocks already contain a finite quantum-kinematic system.
For every finite direct sum \(\mathcal A\), the full state space
\(S(\mathcal A)\), unique block-density representation, effects, and
canonical Born evaluation are mathematical theorems.  The
generated-role theorem additionally gives physical event and probability-role
meaning.  \(\mathsf{QEV}\) is eliminated; \(\mathsf{PROB}\) retains only
preparation, performed-test, outcome/readout, frequency, and calibration
tokens.

Once \(\mathsf{SUB}\) certifies two occurrences with no cross-relation and
the complete product-operation role is represented exactly, factor
interchange derives commuting local images and GP2 relation exactness gives
the direct-sum tensor product, its marginals, and non-signaling identity.
Once \(\mathsf{PROC}\) supplies an actual affine
normalized factor-local map, its positive same-type product action forces
complete positivity and block Kraus form.  Nontrivial kinematic automorphism
flows exist natively.  An
actually supplied/actualized \(\mathsf{DYN}\) member has the blockwise inner representation
proved here, and a supplied full-matrix Markov semigroup has the stated GKLS
form.  The level-one central grading makes the current even observable
envelope at most \(M_{11}\oplus M_4\).  An actual odd bridge plus total graph
connectivity licenses \(M_{15}\) as the \(\mathsf{FENV}\) graded field
envelope, whose parity-preserving observable fixed point remains
\(M_{11}\oplus M_4\).  The separate carrier-resolving \(\mathsf{ENV}\) certificate is required
to use \(M_{15}\) as a full event algebra only when neither a connected
generated physical-test coupling graph nor another carrier-separating
generated test family is available.
For a specified apparatus, the system-side accessible operator system and,
on its generated product-closure domain, accessible algebra are derived by
the UCP readout map; the residual \(\mathsf{ACC}\) content is the actual
apparatus/channel and closure domain, not an arbitrary subalgebra choice.
Projective selective-readout sampling is a representation theorem built after
the state and apparatus effects are fixed; it is neither a stochastic
collapse mechanism nor an additional finite-core axiom.

The remaining preparation/test/outcome tokens, independent composite roots,
processes, trajectories, field bridges, and ungenerated event extensions are
specialized G0 witnesses.  Generated physical event roles are descendants,
not new roots.  The witnesses certify stage-safe geometric actualization of
distinguishability without
altering the generated algebra or selecting an outcome-dependent apparatus
map.



\paragraph{Interface tools}
\label{p4:sec:conclusion-representation-bridge}

Wavefunctions, projective rays, the primitive positive degree-one
Hopf/ray comparison (\textsc{derived-later} from the retained associated-line
atlas and not a physical charge selector), finite readout supports,
angular-momentum bases, worked Hamiltonians, path sums, finite Fock sectors
with the historical generic \(\mathsf{1P}\), \(\mathsf{STAT}\), and optional
bosonic \(\mathsf{OCC}\) interfaces, and finite Euclidean kernels,
finite-window
scattering, local or graded-local pre-nets, quasi-free fields, finite thermal/KMS identities,
and connection-equipped time-dependent blocks are finite representation or
interface tools.  On the active visible branch, \(\mathsf{1P}\) and Fermi
\(\mathsf{STAT}\) are later discharged; no independent numerical bosonic
occupation ceiling is used through Paper~VI-B.  The GP1-calibrated speed \(c\),
operational cone, and
Lorentzian sign are Paper-III imports.  On
\(\mathsf{SF}_{\Lambda_F}\), reciprocal-position resolution obeys the exact
square-root spectral staircase; otherwise any large-\(m_x\) trajectory is a
separately supplied interface.  Finite relativistic dispersion further
requires the stated local second-order free sector.  Large-window M\o ller
limits, complex resonance widths, irreversible GKLS limits, thermal kinetic
equations, Wick reconstruction, and continuum QFT retain their
Cauchy/reservoir/plasma/\(({\rm WICK}_\bullet)\)/locality hypotheses.  A
single fixed finite bath cannot supply a nonunitary exact all-time semigroup
by Theorem~\ref{p4:thm:fixed-finite-bath-semigroup-rigidity}; reset, a
genuine reservoir, a controlled scaling limit, or an explicitly finite-time
approximation is required.  In particular,
the graded-fermionic \(({\rm WICK}_{\rm OS,fin})\) branch uses the complete reflection
matrix, scalar population kinetics uses the invariant occupation algebra, and
history-dependent return uses the two-time
\((\mathrm{DLOC}_{\rm pkt})\) instrument rather than a final-time projector.  Their
large/local/continuum behavior belongs to the reciprocal-position or
presentation side \(m_x\), not to a change of the fixed microscopic compact
momentum datum \(m_p=1\).

\begin{remark}[Historical generic interfaces and active through-IX count]
\label{p4:rem:paper4-selector-count}
The native finite state/effect/tensor/CP closure has no structural selector.
For a completely generic bare finite algebra the five interface classes
\[
 \boxed{\mathsf{DYN},\qquad\mathsf{ACC},\qquad
 \mathsf{1P},\qquad\mathsf{STAT},\qquad\mathsf{OCC}.}
\]
remain useful no-selection diagnostics.  They must not, however, be counted as
five active structural roots of the current Papers-I-A--VI-B chain.

In that chain, \(\mathsf{ACC}\) and physical localization are one actual
preparation/measurement occurrence package; \(\mathsf{1P}\) and visible
\(\mathsf{STAT}\) are closed on their active branches; no theorem consumes an
independent numerical bosonic \(\mathsf{OCC}\) ceiling; and generic
\(\mathsf{DYN}\) is replaced by the role-specific physical owner/state
packages used by later papers.  Therefore
\[
 \boxed{N_{\rm Paper\ IV}^{\rm active\ structural,\le IX}=0.}
\]
Actual source/state and cyclic-response evaluation remains wherever a
role-specific readout has not been evaluated. An actual apparatus/preparation
is consumed when a measurement is claimed; a genuine reservoir/Markov/kinetic
domain is consumed when irreversible dynamics is claimed. These domains do
not become universal premises of every finite probability or pole constructor.
\end{remark}


\begin{remark}[Joint compatibility of the generic many-body interfaces]
\label{p4:rem:joint-manybody-compatibility}
The five-family display above is an owner/status ledger, not a claim that every
Cartesian tuple is a physical candidate.  A proposed one-particle block
\(\mathcal K=\operatorname{Ran}P\) and a physical one-particle generator
\(H_1\) are jointly admissible only on the stability locus
\[
 \boxed{[P,H_1]=0},
\]
with the additional reducing/isolation conditions of the corresponding free
field theorem.  Likewise an exact accessible charge-sector algebra with
minimal central projectors \(P_q\) can coexist with a continuous closed
dynamics only when
\[
 \boxed{[H,P_q]=0\quad\text{for every retained sector }q},
\]
and an open channel must preserve the same block-diagonal state algebra on the
declared theorem domain.  These are joint GP2 compatibility conditions, not
new selectors.

The occupation family is branch-dependent rather than an unconditional fifth
coordinate: on a Fermi branch finite occupation follows from
\(\Lambda^n\mathcal K=0\) for \(n>\dim\mathcal K\), so a bosonic
\(\mathsf{OCC}\) datum is type-inapplicable; on a Bose branch a generated
finite occupation boundary may remain live.  Consequently, on Paper~V-A's
visible-matter branch where \(\mathsf{1P}\) and Fermi \(\mathsf{STAT}\) are
derived later, the generic five-family owner ledger specializes to the joint
\(\mathsf{DYN}+\mathsf{ACC}\) problem.  This later branchwise reduction does
not rewrite the generic Paper-IV chronology.
\end{remark}



\paragraph{Experimental and phenomenological boundary}
\label{p4:sec:conclusion-experimental-boundary}

Concrete detector efficiencies, prepared-state selections, apparatus windows,
weak-coupling hypotheses, transport pipelines, sector assignments,
flavor/mass closures, Standard-Model identifications, Reeb transfer channels,
and full renormalized continuum QFT are downstream uses of the finite quantum
layer.  They may test, specialize, or validate Paper~IV; they do not alter the
canonical state evaluation or reopen the compact carrier.  Any claim that
they enlarge the generated algebra must separately satisfy the graph or
\(\mathsf{FENV}/\mathsf{ENV}\) criterion with the claimed field/event typing.



\paragraph{Finite-QM compatibility claim}
\label{p4:sec:conclusion-final-closure}

Thus Paper~IV establishes the finite-domain mathematical closure of ordinary
finite quantum state, effect, tensor, and CP representation theory inside the
Born--Hopf framework without a hidden full-envelope premise.  It does not
derive which preparation, performed test, outcome/frequency record, second
independent occurrence, or actual process occurs; nor does its generic
physical representation layer select a \(\mathsf{DYN}\), an actual
\(\mathsf{ACC}\) apparatus/closure domain, \(\mathsf{1P}\),
\(\mathsf{STAT}\), or \(\mathsf{OCC}\)
member.  It
does derive the physical state/test/event roles and their normalized
probability evaluation.  The remaining ordinary token/occurrence interfaces are
\(\mathsf{PROB}\), \(\mathsf{SUB}\),
\(\mathsf{PROC}\), and \(\mathsf{DYN}\).  The distinct
\(\mathsf{FENV}/\mathsf{ENV}\) witness families type an odd-field bridge and,
when not already generated, event completion.  Thus the combined residual
witness ledger contains all six families when their respective branches are
live; the four-item occurrence list is not an exhaustive bridge/completion
ledger.  Under geometric primacy G0 these are specialized
occurrence, bridge, preparation, process, trajectory, and readout
certificates, not independent foundations.  G0 requires faithful intrinsic
actualization but does not choose which preparation, process, outcome, or
trajectory occurs.  Local-field,
irreversible, continuum, and phenomenological languages remain controlled
interface presentations on their stated domains.

The upstream compact momentum result is chronologically firewalled from all
of these later interfaces:
\[
\begin{aligned}
\mathcal H_p\cong\mathcal H_1,\quad m_p=1
&\quad\text{are independent of}\quad
\mathsf{ENV},\mathsf{FENV},\mathsf{DYN},\mathsf{ACC},\mathsf{1P},\\
&\hspace{6em}\mathsf{STAT},\mathsf{OCC},\mathrm{REG}_\sigma,\mathrm{POL}.
\end{aligned}
\]
Likewise mathematical availability of \(M_{15}=\End(\mathcal H_1)\) is not
physical event generation: it becomes the generated event algebra only when a
test-typed connected closure (or an explicitly occurring carrier-resolving
event extension) supplies the required marks.

Within the narrower generated-process domain, a fully automorphism-neutral
primitive bistochastic semigroup on a simple matrix algebra is rigid:
\(T_t=P_1+e^{-\kappa t}P_0\), with the maximally mixed stationary state and
only the positive rate left for time placement.  A specified finite simple
Hamiltonian and inverse temperature likewise have a unique normalized
Gibbs/KMS state; this is not a unique pure-vacuum prerequisite for finite
readouts. These theorems do not select arbitrary central block weights on
direct sums, and an effect family still does not select its measurement
instrument.  Generic apparatus, reservoir, and
actual-process ownership therefore remain exactly where the main ledger puts
them.

The universal continuation update sharpens this boundary.  A generated
connected reversible mixing graph with generated vertex measure fixes its
stationary probability even when several invariant conductance ratios leave
the process law unresolved; a one-edge-orbit graph fixes the law itself up to
clock calibration.  Conversely, a disconnected graph retains central
component weights, and the bare \(11\oplus4\) grading supplies no transition
edges.  Thus common invariant state information can descend from a family of
sector-specific continuations without identifying those continuations with
one another or assuming one literal dynamical parent.



% BEGIN FRAMEWORK v2.266: FINITE CONTROL ORBITS AND LUEDERS HISTORY OBSTRUCTIONS
\subsection{Exact native control-orbit span and repeated-readout blindness}
\label{p4a:sec:v266-readout-orbit}
\begin{theorem}[One central direction is missing under all inner controls]
\label{p4a:thm:v266-controlled-span}
Assume that the full connected block-preserving inner conjugations
$U(11)\times U(4)$ are available as mathematical reversible controls.
For any exact-position five-effect intrinsic POVM, the real linear span
of all its conjugated effects is exactly
\[
\boxed{\mathbb RI\oplus\operatorname{Herm}_0(11)
\oplus\operatorname{Herm}_0(4),\qquad \dim_{\mathbb R}=136.}
\]
The single absent native Hermitian direction is the central
geometric/protected population contrast. An independently accessible
central projector $P_S$ extends the span to all $137$ real dimensions
of $M_{11}\oplus M_4$. Full C$^*$-algebra generation by the original
five effects does not already give this linear test span.
\end{theorem}
\begin{proof}
Theorem~\ref{p2c:thm:v266-minimal-protected-readout} writes every
five-effect exact-position POVM as
$E_i=w_iI+\sum_ag_{ia}X_a^D/r$. Each $X_a^D$ is traceless
\emph{separately} in both blocks. Consequently every conjugated effect
has normalized block traces $w_i$ in both sectors and is orthogonal to
$Z=P_G/11-P_S/4$. This excludes at least one central direction.
Conversely the span contains nonzero coordinate operators with nonzero
traceless parts in both blocks. Averaging the conjugation orbit over
one unitary factor isolates the traceless part in the other, and the
real adjoint representation of $SU(n)$ on $\operatorname{Herm}_0(n)$
is irreducible for $n=11,4$. Independent block conjugations therefore
span both traceless factors. The normalization $I$ is already present,
so the resulting dimension is $1+120+15=136$. The line spanned by
$P_S$ is independent of that space and completes it.
\end{proof}

\begin{theorem}[L\"uders histories of five exact-position effects remain chirality blind]
\label{p4a:thm:v266-lueders-blindness}
For a fixed five-effect exact-intrinsic-position POVM, repeat its
\emph{canonical L\"uders instrument} any finite number of times,
without any intervening controls. Every protected compressed history
effect belongs to the real five-dimensional Clifford spin factor
\[
\mathcal J=\operatorname{span}_{\mathbb R}
\{I_4,\gamma_1,\gamma_2,\gamma_3,\gamma_4\}.
\]
Thus the two protected chiral states
$\rho_\pm=(I_4\pm\chi)/4$ have identical full outcome-string
probabilities at every finite history length. In contrast, the
same-label two-shot event can extract the quadratic fingerprint $Q$.
\end{theorem}
\begin{proof}
The five-effect inverse-moment theorem gives protected effects
$E_i^S=w_iI_4+u_i\cdot\gamma$. Their positive square roots have
form $K_i^S=a_iI_4+b_i\gamma(n_i)$ with real coefficients (including
the scalar special case). For $A=cI_4+\gamma(v)$,
Clifford anticommutation gives
\[
\begin{split}
&(aI+b\gamma(n))(cI+\gamma(v))(aI+b\gamma(n))\\
&=\bigl[c(a^2+b^2)+2ab\,n\cdot v\bigr]I\\
&\quad+\gamma\bigl((a^2-b^2)v+2b^2(n\cdot v)n+2abc\,n\bigr).
\end{split}
\]
Therefore the Heisenberg update $A\mapsto K_i^SAK_i^S$
preserves $\mathcal J$. Starting from $I_4$, arbitrary compositions
stay in $\mathcal J$. The chirality $\chi$ is trace-orthogonal to
$\mathcal J$, so all protected outcome-string probabilities agree
between $\rho_+$ and $\rho_-$. Proposition~
\ref{p2c:prop:v266-repeat-fingerprint} supplies the separate
quadratic two-shot observable. This no-go does not apply to other
post-measurement instruments or to histories with additional controls.
\end{proof}

\begin{corollary}[Sharp chirality information with reduced support]
\label{p4a:cor:v266-small-chirality-readout}
The upper total-variation and classical Fisher-information bounds of
Corollary~\ref{p4:cor:v226-chirality-fisher} are attained by seven
exact-position outcomes for $r\ge\sqrt3\Rmom$ and by six for
$r\ge\sqrt{34}\Rmom/3$. Five cannot read protected chirality
at any radius. The explicit six-effect pure chirality construction
at $r\ge2\Rmom$ avoids geometric-block contamination on arbitrary
native preparations; the smaller-radius six-trine contrast need only
be interpreted as a chiral test for a known protected preparation.
\end{corollary}
\begin{proof}
Theorem~\ref{p2c:thm:v266-sharp-chirality-support} gives the sharp
saturators; Theorem~\ref{p2c:thm:v266-minimal-protected-readout}
gives the five-effect no-go and the six-effect global chirality contrast.
\end{proof}
No access to a continuous control group, calibrated instrument, or
physical measurement occurrence is inferred from finite Lie closure
or unitary orbit algebra alone.
% END FRAMEWORK v2.266: FINITE CONTROL ORBITS AND LUEDERS HISTORY OBSTRUCTIONS



% BEGIN FRAMEWORK v2.297: FINITE OBSERVABLE-SPECIFIC OUTCOME-HISTORY SATURATION
\begin{theorem}[Finite Heisenberg saturation of registered observable histories]
\label{p4a:thm:v297-finite-outcome-word-saturation}
Let $H$ be a specified finite Hilbert carrier of complex dimension $N$.
Fix a finite set of \emph{registered} terminal Hermitian effects
$E_1,\ldots,E_s$ and a finite set of CP trace-nonincreasing outcome maps
$\Psi_1,\ldots,\Psi_b$ on $\End(H)$, with the usual instrument
normalization when the maps are grouped into settings. Put
\[
 \mathcal O_0=\operatorname{span}_{\mathbb R}
 \{I_H,E_1,\ldots,E_s\}\subseteq\operatorname{Herm}(H),\qquad
 \mathcal O_{n+1}=\mathcal O_n+
 \sum_{i=1}^b\Psi_i^*(\mathcal O_n),
 \quad d=\dim_{\mathbb R}\mathcal O_0.
\]
Then the increasing sequence stabilizes by $L=N^2-d$ steps:
$\mathcal O_L=\mathcal O_{L+1}=\mathcal O_\infty$.
It is exactly the real span of all pulled-back terminal effects
$\Psi_{i_1}^*\cdots\Psi_{i_m}^*E_a$ and outcome-word effects
$\Psi_{i_1}^*\cdots\Psi_{i_m}^*I_H$, for every finite word
(including the empty word). Adaptive choice of registered settings
introduces no additional effect beyond these words.

Let $\mathcal A\subseteq\End(H)$ be a finite-dimensional represented
unital $C^*$-subalgebra and let $Q$ be its state restriction as in
Theorem~\ref{p2b:thm:v293-operational-multiplicity-quotient}.
The following are equivalent:
\begin{enumerate}[label=(\roman*),leftmargin=2em]
\item Every registered finite outcome-word probability and terminal
$E_a$-readout is determined by $Q(\rho)$ for \emph{all} initial density
matrices $\rho$.
\item $\mathcal O_L\subseteq\mathcal A_{\rm sa}$.
\end{enumerate}
Thus the condition is decidable after at most $N^2-d$ closure rounds,
using real finite-dimensional linear algebra. If it fails, there is a
word of length at most $L$ with Hermitian pulled-back effect
$T\notin\mathcal A$. Writing the ordinary-trace conditional expectation
of Theorem~\ref{p2c:thm:v295-complete-port-twirl} as
$\mathsf E_{\mathcal A}$, the two density matrices
\[
 \rho_\pm=\frac{I_H}{N}\pm\varepsilon
 (T-\mathsf E_{\mathcal A}T),\qquad
 0<\varepsilon\le
 \frac{1}{N\|T-\mathsf E_{\mathcal A}T\|_{\rm op}},
\]
have identical $Q$-images, but differ on that finite history by
$2\varepsilon\|T-\mathsf E_{\mathcal A}T\|_{\rm HS}^2>0$.

This is an \emph{observable-specific} criterion. It need not imply
$\Psi_i^*(\mathcal A)\subseteq\mathcal A$ for all $i$, which is the
stronger complete-process condition of
Theorem~\ref{p4:thm:v296-outcome-wise-instrument-descent}.
It is a mathematical test for an already specified grammar, not a
license for occurrence of an effect, instrument, or GP3 comparison.
\end{theorem}
\begin{proof}
Each Heisenberg outcome map is real linear on $\operatorname{Herm}(H)$, so induction
identifies $\mathcal O_n$ with the real span of words of length at most
$n$ applied to the terminal effects and $I_H$. Since
$\dim_{\mathbb R}\operatorname{Herm}(H)=N^2$, each strict inclusion raises the
dimension by at least one. Once an inclusion is equality the
recursion preserves equality thereafter. Hence it stabilizes by
$N^2-d$ rounds, also for any adaptive subgrammar.
If (ii) holds, every pulled-back effect belongs to $\mathcal A$,
so its expectation is determined by state restriction $Q$; this proves
(i), including each finite outcome word. Conversely, suppose a word
effect $T$ lies outside $\mathcal A$. The trace-preserving conditional
expectation gives the nonzero Hermitian
$X=T-\mathsf E_{\mathcal A}T$, orthogonal to $\mathcal A$ in
Hilbert--Schmidt pairing. The displayed bound makes $\rho_\pm$
positive and trace-one. They have identical $\mathcal A$-expectations,
but the specified pulled-back history separates them by
$2\varepsilon\operatorname{Tr}(X^2)>0$. Thus (i) fails.
If $\mathcal O_L$ is not a subset of $\mathcal A$, some actual word
spanning it fails membership; the stabilization bound supplies the
finite length. The finite word test is narrower than demanding
invariance of all of $\mathcal A$ under each instrument outcome.
\end{proof}

\begin{proposition}[Observable-specific quotient can survive a non-descending process]
\label{p4a:prop:v297-specific-versus-full-process}
On $H=\mathbb C^2\otimes\mathbb C^2$ let
$\mathcal A=M_2(\mathbb C)\otimes I_2$ and let
$U=|0\rangle\langle0|\otimes I_2+
|1\rangle\langle1|\otimes Z$ be controlled-$Z$.
Register the channel $\Psi(\rho)=U\rho U^*$ and the two terminal effects
$E_\pm=(I_4\pm Z\otimes I_2)/2$. Then all their finite-history
readouts are multiplicity blind,
$\mathcal O_\infty=\operatorname{span}_{\mathbb R}
\{I_4,Z\otimes I_2\}\subset\mathcal A$, although
\[
 \Psi^*(X\otimes I_2)=X\otimes Z\notin\mathcal A.
\]
If instead the declared readout is
$F_+=(I_4+X\otimes I_2)/2$, then
$\Psi^*(F_+)=(I_4+X\otimes Z)/2\notin\mathcal A$, and the
quotient-equivalent states
$|+\rangle\langle+|\otimes|0\rangle\langle0|$ and
$|+\rangle\langle+|\otimes|1\rangle\langle1|$ yield respectively
probabilities $1$ and $0$ for the post-channel $F_+$ event.
Consequently one cannot infer the admissibility of a whole process
quotient from the preservation of a particular registered observable.
\end{proposition}
\begin{proof}
Controlled-$Z$ commutes with $Z\otimes I_2$ and conjugates
$X\otimes I_2$ to $X\otimes Z$. The first registered grammar is
therefore fixed by all word pullbacks. For the second readout the
conjugated effect contains $X\otimes Z$ and distinguishes the two
states, while their partial traces over the second qubit coincide.
The last inference is exactly the distinction between the restricted
criterion of Theorem~\ref{p4a:thm:v297-finite-outcome-word-saturation}
and the full-process quotient of
Theorem~\ref{p4:thm:v296-outcome-wise-instrument-descent}.
\end{proof}
% END FRAMEWORK v2.297: FINITE OBSERVABLE-SPECIFIC OUTCOME-HISTORY SATURATION


% BEGIN FRAMEWORK v2.298: ONE-STEP INFORMATIONALLY COMPLETE OUTCOME CERTIFICATE
\begin{theorem}[One-step outcome-descent certificate from an informationally complete readout]
\label{p4a:thm:v298-ic-one-step-outcome-descent}
Let $H$ be a finite-dimensional complex Hilbert space and let
$\mathcal A\subseteq\End(H)$ be a represented unital finite-dimensional
$C^*$-algebra. Suppose $E_1,\ldots,E_s\in\mathcal A$ are positive,
$\sum_jE_j=I_H$, and span $\mathcal A_{\rm sa}$ over $\mathbb R$.
Fix a finite registered family of CP trace-nonincreasing outcome maps
$\Psi_1,\ldots,\Psi_b$ on $\End(H)$; registration is a premise, not a
consequence of the existence of a Kraus representation.
Write $\mathsf E_{\mathcal A}$ for the ordinary-trace orthogonal conditional
expectation onto $\mathcal A$ from
Theorem~\ref{p2c:thm:v295-complete-port-twirl} and $Q$ for state restriction
to $\mathcal A$ from
Theorem~\ref{p2b:thm:v293-operational-multiplicity-quotient}.
The following are equivalent:
\begin{enumerate}[label=(\roman*),leftmargin=2em]
\item Every finite registered outcome-word probability followed by a
terminal $E_j$ readout depends on the initial state only through $Q$.
\item $\Psi_i^*(E_j)\in\mathcal A$ for every $1\le i\le b$ and
$1\le j\le s$.
\item $\Psi_i^*(\mathcal A)\subseteq\mathcal A$ for every $i$.
\end{enumerate}
Equivalently, the finite nonnegative diagnostic
\[
\Delta_{\rm IC}:=\sum_{i=1}^b\sum_{j=1}^s
\bigl\|\Psi_i^*(E_j)-\mathsf E_{\mathcal A}(\Psi_i^*(E_j))\bigr\|_{\rm HS}^2
\]
vanishes exactly when all finite histories descend. If it is nonzero,
some \emph{single} outcome $i$ followed by one readout $E_j$ distinguishes
two density matrices with identical restrictions $Q(\rho)$.
This strengthens the finite-length bound of
Theorem~\ref{p4a:thm:v297-finite-outcome-word-saturation} only when the
registered terminal readout is informationally complete for $\mathcal A$.
It does not register an instrument or establish a GP3 whole-candidate
comparison.
\end{theorem}
\begin{proof}
Since the $E_j$ span the Hermitian part of $\mathcal A$, real linearity of
each Heisenberg outcome map gives (ii)$\Leftrightarrow$(iii), including the
complex-linear extension to $\mathcal A$. Under (iii), every iterated
pullback of $I_H$ or an $E_j$ belongs to $\mathcal A$, proving (i) by
induction for every outcome word, including adaptive selections.
If (ii) fails for indices $(i,j)$, put
$T=\Psi_i^*(E_j)$ and $X=T-\mathsf E_{\mathcal A}(T)\ne0$.
The trace-preserving conditional expectation satisfies
$\Tr(X)=0$ and $\Tr(XA)=0$ for all $A\in\mathcal A$.
For sufficiently small $\varepsilon>0$ the matrices
$\rho_{\pm}=I_H/\dim H\pm\varepsilon X$ are density matrices with
$Q(\rho_+)=Q(\rho_-)$, yet
\[
\Tr(E_j\Psi_i(\rho_+))-\Tr(E_j\Psi_i(\rho_-))
 =2\varepsilon\Tr(X^2)>0.
\]
Hence (i) implies (ii), with a one-outcome witness. Finally each term of
$\Delta_{\rm IC}$ is nonnegative, and the sum vanishes exactly when (ii)
holds. No numerical value of this diagnostic is claimed to be independent
of the chosen informationally complete effect frame; its zero locus is.
\end{proof}

\begin{corollary}[Minimal mathematical informationally complete kinematic readout]
\label{p4a:cor:v298-kinematic-43-effect-frame}
The native six-type kinematic test algebra
$\mathcal A_{\rm kin}\cong\mathbb C\oplus M_2\oplus M_2\oplus M_4\oplus
M_3\oplus M_3$ has real Hermitian dimension $43$. On every faithful
represented implementation it admits a mathematical POVM with exactly
$43$ strictly positive effects spanning $\mathcal A_{{\rm kin},\rm sa}$,
and no POVM with fewer effects can span this entire space.
If such a readout is registered, the one-step test of
Theorem~\ref{p4a:thm:v298-ic-one-step-outcome-descent} requires only
$43b$ operator-membership checks for $b$ registered outcome maps.
Neither an apparatus occurrence nor intrinsic position-moment matching
is asserted by the abstract frame construction.
\end{corollary}
\begin{proof}
Put $d=43$ and choose a real basis $F_1,\ldots,F_{d-1}$ of the
ordinary-traceless Hermitian subspace of the represented algebra.
For sufficiently small $\varepsilon>0$, the $d$ effects
\[
E_j=\frac1dI_H+\varepsilon F_j\quad(1\le j<d),\qquad
E_d=\frac1dI_H-\varepsilon\sum_{j<d}F_j
\]
are strictly positive and sum to $I_H$. Their real span contains $I_H$
and all $F_j$, hence equals $\mathcal A_{{\rm kin},\rm sa}$.
The span of fewer than $d$ effects has dimension below $d$, proving the
lower bound. The dimension count is $1+4+4+16+9+9=43$, independent of
multiplicity-blind representations. The last assertion follows from the
preceding theorem.
\end{proof}
% END FRAMEWORK v2.298: ONE-STEP INFORMATIONALLY COMPLETE OUTCOME CERTIFICATE


% BEGIN FRAMEWORK v2.299: CONTINUOUS IC READOUT FINITE TEST
\begin{theorem}[Finite one-step certificate from a continuous informationally complete POVM]
\label{p4a:thm:v299-continuous-ic-event-test}
Let $\mathcal A\subset\operatorname{End}(H)$ be a finite-dimensional
represented unital $C^*$-algebra with
$d=\dim_{\mathbb R}\mathcal A_{\rm sa}$, and let $E$ be a specified
normalized POVM on a measurable outcome space $\Omega$ satisfying
\[
 E(B)\in\mathcal A,\qquad
 \operatorname{span}_{\mathbb R}\{E(B):B\ \mathrm{measurable}\}
 =\mathcal A_{\rm sa}.
\]
Choose measurable $B_1,\ldots,B_d$ with
$F_j=E(B_j)$ a real basis. For any specified finite family of
completely positive trace-nonincreasing outcome maps
$\{\Psi_i\}_{i=1}^b$, the following are equivalent:
\begin{enumerate}[label=(\roman*),leftmargin=2em]
\item Every registered finite outcome history, followed by any
terminal event $B\subset\Omega$, has a probability depending on
the initial density matrix only through restriction to $\mathcal A$.
\item $\Psi_i^*(F_j)\in\mathcal A$ for every $i,j$.
\item $\Psi_i^*(\mathcal A)\subseteq\mathcal A$ for every $i$.
\end{enumerate}
Thus the $bd$ exact operator-membership tests suffice even if the
terminal outcome space is uncountable. Equivalently the sum
\[
 \Delta_{E,\Psi}=\sum_{i,j}\bigl\|\Psi_i^*(F_j)-
 \mathsf E_{\mathcal A}(\Psi_i^*(F_j))\bigr\|_{\rm HS}^{2}
\]
vanishes exactly when all those histories descend, where
$\mathsf E_{\mathcal A}$ is the ordinary-trace conditional
expectation. A nonzero defect yields a witness consisting of one
specified outcome $i$ and one measurable terminal event $B_j$.
The $B_j$ need not be disjoint or sum to $\Omega$; this finite test
does not replace the continuous POVM by a finite, covariant POVM.
It also does not assert physical registration or GP3 whole-candidate
fidelity.
\end{theorem}
\begin{proof}
Finite dimension and the span hypothesis supply a basis of $d$
event effects. Complex linearity of each Heisenberg map shows
(ii)$\Leftrightarrow$(iii). Under (iii), any finite composition of
these pullbacks maps all event effects, as well as $I_H$, into
$\mathcal A$, so every word probability is determined by the
restricted input state. Conversely, if (ii) fails, choose
$T=\Psi_i^*(F_j)$ outside $\mathcal A$ and put
$X=T-\mathsf E_{\mathcal A}(T)$. The conditional expectation is
trace-preserving and Hilbert--Schmidt orthogonal, so the density
matrices $I_H/\dim H\pm\epsilon X$ for sufficiently small positive
$\epsilon$ agree on $\mathcal A$ but differ on the one-outcome
probability $\operatorname{Tr}(F_j\Psi_i(\rho))$ by
$2\epsilon\operatorname{Tr}(X^2)>0$. Each term of the displayed
finite defect is nonnegative, giving its zero criterion. The
proof is the event-span extension of
Theorem~\ref{p4a:thm:v298-ic-one-step-outcome-descent}; it does
not require discretizing or selecting a new POVM.
\end{proof}
% END FRAMEWORK v2.299: CONTINUOUS IC READOUT FINITE TEST

% BEGIN FRAMEWORK v2.300: COVARIANT SQUARE-ROOT INSTRUMENT AND NORMALIZER
\begin{theorem}[Canonical covariant square-root instrument and exact normalizer criterion]
\label{p4a:thm:v300-covariant-root-instrument}
Let $\mathcal A\subseteq\End(H)$ be a represented unital
finite-dimensional $C^*$-algebra and $(\Omega,\mu)$ a finite
measure space. Suppose $D:\Omega\to\mathcal A_+$ is measurable,
$\int_\Omega D(\omega)\,d\mu(\omega)=I_H$, and set
$K(\omega)=D(\omega)^{1/2}\in\mathcal A$. Then
\[
 \mathcal I_B(\rho)=\int_B K(\omega)\rho K(\omega)\,d\mu(\omega)
\]
defines a normalised completely positive, countably additive
instrument with POVM effects $\mathcal I_B^*(I_H)=\int_B D\,d\mu$.
For every measurable event $B$, the Heisenberg action obeys
$\mathcal I_B^*(\mathcal A)\subseteq\mathcal A$; thus the complete
outcome-history grammar descends to the restriction of states to
$\mathcal A$ as in
Theorem~\ref{p4a:thm:v299-continuous-ic-event-test}.
If a compact group $G$ preserves $\mu$, acts measurably on $\Omega$,
acts unitarily on $H$ and $D(g\omega)=U_gD(\omega)U_g^*$ almost
everywhere, this instrument is exactly $G$-covariant.

Suppose additionally that $D(\omega)$ is strictly positive for
almost every $\omega$. Let $V:\Omega\to U(H)$ be any measurable
post-update unitary field and define
\[
 \mathcal I_B^V(\rho)=\int_B
 V(\omega)K(\omega)\rho K(\omega)V(\omega)^*\,d\mu(\omega).
\]
It has precisely the same POVM effects. Its \emph{every individual}
event map preserves $\mathcal A$ in the Heisenberg picture if and only
if $V(\omega)$ normalizes $\mathcal A$ almost everywhere:
\[
 \boxed{\ \forall B:\ (\mathcal I_B^V)^*(\mathcal A)\subseteq\mathcal A
 \iff V(\omega)^*\mathcal A V(\omega)=\mathcal A\quad\mu\text{-a.e.}\ }
\]
If normalizer membership fails on a set of positive measure, a single
measurable event and a Hermitian terminal test in $\mathcal A$ give
unequal probabilities for two states that have identical restrictions
to $\mathcal A$. When the terminal POVM is informationally complete
for $\mathcal A$, one of its finitely many basis event effects can
be used as that terminal test.
The strict-positivity hypothesis in the reverse implication is
essential: a rank-one diagonal PVM can conceal nonnormalizing
post-unitaries after compression.
This proves mathematical instrument availability and a conditional
process-descent criterion; it does not register a physical instrument.
\end{theorem}
\begin{proof}
In finite dimension the continuous square-root functional calculus
makes $K$ measurable and $K\in\mathcal A$; the integral of completely
positive Kraus maps is completely positive, with normalisation
$\int K^*K=I_H$. Countable additivity holds in trace norm. The
Heisenberg pullback is $\int_B KAK\,d\mu\in\mathcal A$ for
$A\in\mathcal A$. The full adaptive-history conclusion is the
outcome-wise descent principle of
Theorem~\ref{p4:thm:v296-outcome-wise-instrument-descent}.
Functional calculus also gives
$K(g\omega)=U_gK(\omega)U_g^*$, proving covariance by a change of
variables in the integral.
For the post-unitary family the Heisenberg integrand at $A\in\mathcal A$
is $K(\omega)V(\omega)^*AV(\omega)K(\omega)$.
If $V(\omega)$ normalizes $\mathcal A$, it is in $\mathcal A$.
Conversely suppose every event pullback is in $\mathcal A$. Apply
the Hilbert--Schmidt orthogonal projection onto $\mathcal A^\perp$.
For each element of a fixed finite Hermitian basis of $\mathcal A$,
the integral of the projected integrand is zero over every event;
hence that projected integrand vanishes almost everywhere. Strict
positivity makes $K(\omega)$ invertible \emph{in} $\mathcal A$, so
$V(\omega)^*\mathcal A V(\omega)\subseteq\mathcal A$ a.e.; equality
follows in finite dimension because unitary conjugation preserves
algebra dimension. If the condition fails, the contrapositive gives
an event pullback outside $\mathcal A$. The two-state witness follows
by subtracting the trace-preserving conditional expectation as in
Theorem~\ref{p4a:thm:v299-continuous-ic-event-test}; its event-basis
argument supplies a terminal POVM effect when that POVM spans
$\mathcal A_{\rm sa}$. For sharpness, let $\mathcal A$ be the
diagonal matrices on $\mathbb C^2$, $D_i=P_i$ the rank-one diagonal
projections and $V_i$ arbitrary unitaries. Every compressed adjoint
$P_iV_i^*AV_iP_i$ lies in $\mathcal A$, although $V_i$ need not
normalize it. The full-rank assumption is therefore necessary for
the stated equivalence.
\end{proof}

\begin{proposition}[Identical POVM effects admit opposite quotient outcomes]
\label{p4a:prop:v300-swap-instrument-witness}
Let $H=\mathbb C^2\otimes\mathbb C^2$ and
$\mathcal A=M_2(\mathbb C)\otimes I_2$. Put
$P_i=|i\rangle\langle i|$ for $i=0,1$ and let $S$ swap the two
qubits. The two normalized instruments
\[
 \Phi_i(\rho)=(P_i\otimes I)\rho(P_i\otimes I),\qquad
 \widetilde\Phi_i(\rho)=S\Phi_i(\rho)S^*
\]
have the same POVM effects $P_i\otimes I$, but only $\Phi_i$ preserves
$\mathcal A$ outcome by outcome. In fact, for the Pauli matrix $Z$,
\[
 \widetilde\Phi_i^*(Z\otimes I)=P_i\otimes Z\notin\mathcal A.
\]
Thus a POVM's effects, including an informationally complete family,
do not select its physical postmeasurement instrument.
\end{proposition}
\begin{proof}
For both instruments the adjoint on the identity is $P_i\otimes I$,
and the two effects sum to $I_H$. Directly,
$S^*(Z\otimes I)S=I\otimes Z$, so compression by
$P_i\otimes I$ gives $P_i\otimes Z$.
The first instrument is an instance of the square-root construction
of Theorem~\ref{p4a:thm:v300-covariant-root-instrument} for a finite
PVM (its density need not be invertible for the forward implication).
For $0<\epsilon<1/4$ the states
$\rho_\pm=I_4/4\pm\epsilon(P_i\otimes Z)$ agree on every test in
$\mathcal A$ but differ on the displayed one-step readout.
\end{proof}
% END FRAMEWORK v2.300: COVARIANT SQUARE-ROOT INSTRUMENT AND NORMALIZER

% BEGIN FRAMEWORK v2.301: GAUGE-NATURAL OUTCOMES AND COVARIANCE COUNTERMODEL
\begin{theorem}[Gauge-natural instruments descend without an assumed Lueders update]
\label{p4a:thm:v301-gauge-natural-instrument-descent}
Let $H$ and $K=U(k_g)\times U(k_s)$ be the multiplicity-amplified
intrinsic carrier and its implementation gauge of
Theorem~\ref{p2c:thm:v301-intrinsic-gauge-fixedpoints}. Let
$\mathcal A_D=\End(H)^{\operatorname{Ad}K}$. Fix an actual or purely
mathematical instrument with CP outcome maps $\Psi_B$, where each
registered outcome label/event $B$ is \emph{gauge-trivial} and
\[
 \Psi_B(W\rho W^*)=W\Psi_B(\rho)W^*
 \quad\text{for all }W\in K,\ \rho,\ B.
\]
Then every outcome map preserves $\mathcal A_D$ in the Heisenberg
picture:
\[
 \boxed{\ \Psi_B^*(\mathcal A_D)\subseteq\mathcal A_D
 \quad\text{for every }B.\ }
\]
Hence every finite registered adaptive outcome history followed by
any $\mathcal A_D$-effect depends only on the restricted initial
state. The conclusion holds for arbitrary gauge-equivariant CP
updates; no square-root/L\"uders, unitary-normalizer or Kraus-form
selection is required. It applies separately to every registered
outcome; covariance of the nonselective sum is not enough.

The class of gauge-equivariant CP maps is closed under composition,
finite nonnegative sums and finite convex mixtures. Thus a
\emph{specified} generator grammar built entirely from gauge-equivariant
primitive outcome maps using those operations inherits this
property. This is a conditional naturality certificate: merely
calling a multiplicity space ``unmarked'' does not prove that every
occurring physical source or outcome map obeys the displayed
intertwining identity.
\end{theorem}
\begin{proof}
For $A\in\mathcal A_D$, $W^*AW=A$ for all $W\in K$ by
Theorem~\ref{p2c:thm:v301-intrinsic-gauge-fixedpoints}. Gauge
equivariance of the CP map implies the same relation for its
Heisenberg adjoint: $W^*\Psi_B^*(A)W=\Psi_B^*(W^*AW)=\Psi_B^*(A)$.
Thus the adjoint output is again fixed by $K$, hence in $\mathcal
A_D$. The outcome-wise finite-history conclusion follows by
composition, as in
Theorem~\ref{p4:thm:v296-outcome-wise-instrument-descent}.
The displayed intertwining equations are preserved by composition,
positive linear combinations and convex mixtures. Induction over
a licensed grammar using those operations proves the last assertion;
no claim is made for an uninspected constructor or a new
physical registration.
\end{proof}

\begin{proposition}[Spin(4) covariance and identical full-rank POVM effects do not force quotient descent]
\label{p4a:prop:v301-spin4-same-povm-leak}
Consider the same strict-positive, exactly $\Spin(4)$-covariant
algebra-valued directed-tangent density and its square-root
instrument of
Theorem~\ref{p4a:thm:v300-covariant-root-instrument}, on
$H=(\mathbb C^{11}\otimes\mathbb C^{k_g})\oplus
(\mathbb C^4\otimes\mathbb C^2)$ with
$\mathcal A_D=(M_{11}\otimes I_{k_g})\oplus(M_4\otimes I_2)$,
as in Theorem~\ref{p2c:thm:v301-intrinsic-gauge-fixedpoints}.
There exists a trace-preserving $\Spin(4)$-covariant CP channel
$\mathcal C$ for which
\[
 \widetilde{\mathcal I}_B=\mathcal C\circ\mathcal I_B^{\rm root}
\]
is a normalised, exactly $\Spin(4)$-covariant instrument with
\emph{identical POVM effects for every event} $B$, but whose
nonselective Heisenberg channel already fails to preserve
$\mathcal A_D$. Hence exact geometric covariance, full-rank
effects, informational completeness and the entire effect profile
do not determine quotient-faithful physical updates.
\end{proposition}
\begin{proof}
Let $P_0,P_1$ be orthogonal rank-one projections in the spinorial
multiplicity space $\mathbb C^2$. On the spinorial block set
$L_0=I_4\otimes P_0$, and for $a,b=1,\ldots,4$ set
$L_{ab}=\frac12|a\rangle\langle b|\otimes P_1$. Include the
geometric identity in the single Kraus operator
$\widetilde L_0=P_g\oplus L_0$ and extend $L_{ab}$ by zero on the
geometric block. The sum of $L^*L$ is $I_H$, so this defines a
trace-preserving CP channel $\mathcal C$. On a spinorial test
$X\in M_4$, its adjoint acts by
\[
 \mathcal C^*(0\oplus(X\otimes I_2))
  =0\oplus\left(X\otimes P_0+
                 \frac{\operatorname{Tr}X}{4}I_4\otimes P_1\right).
\]
The channel is $\Spin(4)$-covariant: $\widetilde L_0$
commutes with its represented action, and the $P_1$ component
is the completely depolarising channel on $\mathbb C^4$,
which commutes with conjugation by every unitary on that block.
Consequently $\mathcal C\circ\mathcal I_B^{\rm root}$ is a
covariant instrument. Since $\mathcal C$ is trace preserving,
$\widetilde{\mathcal I}_B^*(I)=\mathcal I_B^{\rm root*}(I)$
for every event $B$.

Choose a nonzero traceless Hermitian $X\in M_4$ and use the whole
event space $\Omega$. The displayed adjoint identity and
blockwise positive square root give
\[
 \widetilde{\mathcal I}_\Omega^*(0\oplus(X\otimes I_2))
 =0\oplus(F(X)\otimes P_0),\qquad
 F(X)=\int_\Omega D_s(\omega)^{1/2}X
                       D_s(\omega)^{1/2}\,d\mu(\omega).
\]
Strict positivity of the density implies
\[
 \operatorname{Tr}(XF(X))
 =\int_\Omega \operatorname{Tr}\!
 \bigl[(D_s(\omega)^{1/4}X D_s(\omega)^{1/4})^2\bigr]
 \,d\mu(\omega)>0,
\]
so $F(X)\ne0$. Since $P_0$ is not scalar, the output is not in
$\mathcal A_D$. The finite-dimensional conditional-expectation
witness of
Theorem~\ref{p4a:thm:v299-continuous-ic-event-test}
then supplies two quotient-equivalent input states distinguished
by a finite terminal-event history. The chosen $P_0$ is an
\emph{additional multiplicity mark}; this is a mathematical
countermodel to inference from Spin(4) covariance alone, not a
claim that GP2 generates this marked channel.
\end{proof}
% END FRAMEWORK v2.301: GAUGE-NATURAL OUTCOMES AND COVARIANCE COUNTERMODEL

% BEGIN FRAMEWORK v2.320: PAPER IV-A PETERSEN-HOPF QUANTUM INTERFACE
\section{Petersen--Hopf quantum tests and accessibility}
\label{p4a:sec:v320-petersen-quantum}
\noindent\textit{Source and role.} This section imports only the generated,
marked-$2I$ finite data of Paper~I-A,
Theorems~\ref{p1:thm:petersen-signed-and-flux} and
\ref{p1:thm:petersen-fifteen-resolution}, and the auxiliary real
readout of Paper~II-C,
Theorem~\ref{p2c:thm:v318-real-petersen-povm}.
It does not identify a Petersen module with the $m_p=1$ momentum
carrier, adjoin a new generated test to $\mathcal A_{\rm kin}$,
or certify a physical preparation, instrument, or detector.
All statements below are exact \emph{finite comparison theorems} on the
specified marked support.

\begin{theorem}[Signed Hopf-sector projectors and their invariant test algebra]
\label{p4a:thm:v320-signed-projective-test}
Let $\mathcal L=\mathbb C^{10}$ be the complexification of the real
unoriented-axis vertex representation, with the signed Petersen matrix
$S=S^*$ defined in Paper~I-A. Then
\[
 S^3=5S,\qquad
 \operatorname{Spec}S=\{(\sqrt5)^{(3)},0^{(4)},(-\sqrt5)^{(3)}\}.
\]
Its three intrinsic spectral projections are
\[
 P_+=\frac12\left(\frac{S^2}{5}+\frac{S}{\sqrt5}\right),\quad
 P_-=\frac12\left(\frac{S^2}{5}-\frac{S}{\sqrt5}\right),\quad
 P_0=I_{10}-\frac{S^2}{5},
\]
with ranks $3,3,4$. Thus the $*$-algebra generated solely by
$I_{10},S$ is $\mathbb C^3$, and the three projections form a
mathematical three-outcome PVM.

For the signed $A_5$-action $R$ on $\mathcal L$ inherited from the
oriented Hopf axes, the representation is multiplicity free,
\[
 \mathcal L\simeq\mathbf3\oplus\mathbf4\oplus\mathbf3',
 \quad \operatorname{End}_{A_5}(\mathcal L)
 =\operatorname{span}_{\mathbb C}\{P_+,P_0,P_-\}.
\]
The group-generated complex $*$-algebra on this \emph{auxiliary}
representation has block type
$M_3(\mathbb C)\oplus M_4(\mathbb C)\oplus M_3(\mathbb C)$,
of complex dimension $34$.
There is no automatic identification of these blocks with the
numerically matching $M_3,M_4,M_3$ components of the physical
$\mathcal A_{\rm kin}$ (complex dimension $43$).
\end{theorem}
\begin{proof}
The displayed polynomial identity, multiplicities, and intertwining
Hopf-axis projectors are Theorem~\ref{p1:thm:petersen-signed-and-flux}.
Lagrange interpolation at the three distinct eigenvalues gives the
projector formulas and $C^*(I,S)\cong\mathbb C^3$.
The two three-dimensional real irreducibles of $A_5$ are inequivalent,
and the four-dimensional component occurs once. Schur's lemma gives
the displayed commutant; the finite-dimensional bicommutant theorem
(or Burnside on each irreducible summand) gives the full group-image
algebra. A shared numerical block size does not constitute an
owner-preserving map between two independently defined representations.
\end{proof}

\begin{proposition}[Odd five-label permutations exchange signed mirror sectors]
\label{p4a:prop:v320-parity-graded-mirror-exchange}
For each permutation $\pi\in S_5$ of the five tetrahedral-core
labels, the induced permutation of the ten axis vertices admits a
signed-permutation lift $M_\pi$ such that
\[
 M_\pi^* S M_\pi=\operatorname{sgn}(\pi)S.
\]
Consequently even permutations preserve each of
$P_+,P_0,P_-$, while odd permutations exchange $P_+$ and $P_-$
and preserve $P_0$. The full unsigned Petersen $S_5$ graph
automorphism group is therefore not an ordinary symmetry group of
the \emph{fixed signed} operator: its odd elements act as a
sign-reversing duality. If a selected Hamiltonian $h_+P_++h_0P_0+h_-P_-$
were also invariant under a physically realized odd lift, then
necessarily $h_+=h_-$. No such odd lift is postulated as physical.
\end{proposition}
\begin{proof}
Choose a spanning tree in the Petersen graph. The signs of a lift
on its ten vertices are fixed recursively from
$S_{\pi v,\pi w}=\operatorname{sgn}(\pi)\epsilon_v\epsilon_wS_{vw}$
on its nine edges. The remaining six edges give six finite cycle
compatibility checks. The explicit five-label/ten-axis certificate
in Paper~IV-C verifies these equalities for all $60$ even and $60$
odd permutations. Substituting $S\mapsto\pm S$ in the polynomial
formulas of the three spectral projectors gives the exchange claim.
The last statement follows by conjugating the sector-scalar
Hamiltonian with an odd lift.
\end{proof}

\begin{corollary}[Complete symmetry twirl and the mirror-population limit]
\label{p4a:cor:v320-signed-twirl}
For a density operator $\rho$ on the auxiliary $\mathcal L$, its
finite rotational twirl is
\[
 \mathscr T_{A_5}(\rho)
 :=\frac1{60}\sum_{g\in A_5}R(g)\rho R(g)^*
 =\frac{\operatorname{Tr}(P_+\rho)}3P_+
 +\frac{\operatorname{Tr}(P_0\rho)}4P_0
 +\frac{\operatorname{Tr}(P_-\rho)}3P_-.
\]
Consequently every $A_5$-invariant effect on $\mathcal L$ tests only
three sector weights (two independent probabilities). It cannot
distinguish two states with the same three weights, nor can a
sector-scalar invariant Hamiltonian transfer population between these
sectors. A difference of the $\mathbf3$ and $\mathbf3'$ scalar phases
is permitted but is not a derived frequency or physical mass splitting.
\end{corollary}
\begin{proof}
The group average is a conditional expectation onto the commutant.
The representation is multiplicity free with the three listed ranks,
so Schur averaging of each diagonal block produces its normalized
trace; inter-irreducible components average to zero. All effects
commuting with $R(A_5)$ lie in the same three-dimensional commutant.
A Hamiltonian of the form $h_+P_++h_0P_0+h_-P_-$ commutes with each
sector projection, so their populations are conserved.
\end{proof}

\begin{proposition}[Real Hopf POVM: exact antiunitary blindness of real histories]
\label{p4a:prop:v320-real-axis-blindness}
Let $F_a=\frac3{10}u_a u_a^{\mathsf T}$ be the ten positive
axis effects of Paper~II-C, now represented as real matrices on an
auxiliary complex three-state Hilbert space. Their real span is
$\operatorname{Sym}(3,\mathbb R)$, of dimension six, not the
nine-dimensional real vector space $\operatorname{Herm}(3,\mathbb C)$.
In particular the distinct density matrices
\[
 \rho_\pm=\frac13I_3\pm\frac{i}{6}
 \begin{pmatrix}0&1&0\\-1&0&0\\0&0&0\end{pmatrix}
\]
are positive with eigenvalues $1/6,1/3,1/2$ and obey
\[
 \operatorname{Tr}(\rho_+F_a)
 =\operatorname{Tr}(\rho_-F_a)\qquad(a=1,\ldots,10).
\]
More generally, every finite outcome history constructed exclusively
from \emph{real} Kraus matrices and real terminal effects has the
same probability for $\rho$ and $\overline\rho$. This includes
mathematical L\"uders instruments of the ten real-axis effects.
Access to complex-valued interventions can defeat this blindness;
real-axis geometry does not forbid such interventions when independently
constituted.
\end{proposition}
\begin{proof}
Paper~II-C proves exact reconstruction of every real symmetric
trace-one state. The six real symmetric directions are therefore
spanned by the ten effects. The remaining Hermitian directions have
the form $iA$ with $A$ real antisymmetric; $\operatorname{Tr}(iA E)=0$
for every real symmetric $E$. Since $iA$ in the display has eigenvalues
$-1,0,1$ before division by six, $\rho_\pm\succeq0$.
For a word $K_n\cdots K_1$ of real Kraus matrices and real symmetric
terminal effect $E$, the Heisenberg word effect
$K_1^{\mathsf T}\cdots K_n^{\mathsf T}E K_n\cdots K_1$ is again real
symmetric. The same orthogonality proves the history statement.
\end{proof}

\begin{theorem}[Exact complementary-transfer CP no-go]
\label{p4a:thm:v320-complement-kraus-no-go}
Select two tetrahedral cores and write $T_{ij}=J_i+J_j:\mathbf3\to
\mathbf3'$. The opposite three-core selection is $-T_{ij}$,
and with $P_a=u_{ij}u_{ij}^{\mathsf T}$,
$P_t=I_3-P_a$, Paper~I-A gives
\[
 T_{ij}^*T_{ij}=4P_a+\frac14P_t.
\]
The contraction $K_{ij}=T_{ij}/2$ has the \emph{mathematical} success
effect and complement
\[
 E_{\rm succ}=K_{ij}^*K_{ij}=P_a+\frac1{16}P_t,
 \qquad E_{\rm fail}=\frac{15}{16}P_t.
\]
The associated success CP operation
$\rho\mapsto K_{ij}\rho K_{ij}^*$ is identical for $T_{ij}$ and
$-T_{ij}$. The failure effect completes the two-outcome POVM;
choosing a failure-state preparation is additional instrument data.
Therefore the complement sign is not observable through an isolated
success outcome. Distinguishing it requires a coherent reference path,
a separate phase-bearing coupling, or an equally explicit operation
not invariant under $K\mapsto-K$.
\end{theorem}
\begin{proof}
The polar identity and complement rule are
Theorem~\ref{p1:thm:petersen-subset-polar}. Dividing the positive
quadratic metric by four gives the stated effects and
$E_{\rm succ}+E_{\rm fail}=I_3$.
The sign cancels identically in $K\rho K^*$ and $K^*K$, including
postselection and every subsequent measurement of the resulting
output. No occurrence or source coupling is inferred by this algebra.
\end{proof}

\begin{corollary}[Reciprocal real-product frame leaves a nine-dimensional
correlation sector invisible]
\label{p4a:cor:v320-reciprocal-product-blindness}
On the auxiliary two-factor $\mathbb C^3_x\otimes\mathbb C^3_p$
comparison space, the 100 local product effects $F_a\otimes F_b$
span only
\[
 \operatorname{Sym}(3,\mathbb R)_x\otimes
 \operatorname{Sym}(3,\mathbb R)_p,
 \qquad \dim_{\mathbb R}=36.
\]
Within all \emph{real symmetric} $9\times9$ operators (dimension
$45$), the orthogonal complement has dimension nine and is precisely
$\Lambda^2\mathbb R^3_x\otimes\Lambda^2\mathbb R^3_p$.
Within $\operatorname{Herm}(9,\mathbb C)$ (dimension $81$), the
product effects leave $45$ directions undetected.
For the real antisymmetric $A$ in
Proposition~\ref{p4a:prop:v320-real-axis-blindness}, the distinct
positive trace-one states
\[
 R_\pm=\frac19I_9\pm\frac1{18}(A\otimes A)
\]
have identical probabilities for all 100 product effects but satisfy
\[
 \operatorname{Tr}[(R_+-R_-)(A\otimes A)]=\frac49.
\]
This is a mathematically explicit reciprocal-correlation blind pair,
not an identification of the auxiliary product with the original
fifteen-dimensional $m_p=1$ event carrier or with the 225-dimensional
Petersen edge tensor packet.
\end{corollary}
\begin{proof}
The preceding proposition gives the six-dimensional local real
symmetric span, so tensoring yields $6\cdot6=36$. The standard
symmetric decomposition is
$\operatorname{Sym}^2(V\otimes W)
=\operatorname{Sym}^2V\otimes\operatorname{Sym}^2W
\oplus\Lambda^2V\otimes\Lambda^2W$.
The matrix $A\otimes A$ is symmetric, traceless and has operator norm
one, so $R_\pm$ are positive. Because $\operatorname{Tr}(AF_a)=0$,
all 100 probabilities agree; since $\operatorname{Tr}(A^2)=-2$,
\[\operatorname{Tr}[(R_+-R_-)(A\otimes A)]
 =\frac19\operatorname{Tr}(A^2)^2=\frac49.\]
The complex Hermitian dimension is $9^2=81$, giving $81-36=45$.
\end{proof}

\begin{remark}[Paper-IV event and occurrence firewall]
\label{p4a:rem:v320-petersen-events-firewall}
The preceding spectral and POVM constructions are available on marked
mathematical comparison carriers. They are \emph{not} evidence that
$\mathcal A_{\rm kin}$ enlarges, that a Petersen source acts on its
specific $M_3$ or $M_4$ blocks, or that a corresponding detector exists.
They must pass Paper~IV-A's generated/algebra/accessibility distinction,
finite Choi--Kraus theorem, and outcome-history descent tests before a
physical observable claim. The real-axis POVM is not a continuous
$SU(2)$-covariant measurement, nor is it a full complex-state IC
measurement.
\end{remark}
% END FRAMEWORK v2.320: PAPER IV-A PETERSEN-HOPF QUANTUM INTERFACE



% BEGIN FRAMEWORK v2.340: EVEN REYE OPERATOR QUANTUM ACCESS LIMIT
\section{Parity information absent from the two-step Reye operator}
\label{p4a:sec:v340-reye-parity-readout}
\begin{proposition}[Exact spectral blindness of $T=S^2-3I$]
\label{p4a:prop:v340-reye-parity-no-go}
For a consistently oriented signed Petersen matrix $S$,
$S^3=5S$ and $T=S^2-3I$ has eigenspaces
\[
 \mathscr H_{T=2}=\mathscr H_{S=\sqrt5}\oplus
                        \mathscr H_{S=-\sqrt5},
 \quad \dim\mathscr H_{T=2}=6,
 \qquad \mathscr H_{T=-3}=\ker S,\quad\dim=4.
\]
Every spectral effect or finite instrumental response that depends
\emph{only} on $T$ acts with the same scalar on the two three-dimensional
mirror channels and is invariant under replacing $S$ by $-S$.
It cannot, by this operator alone, distinguish the signs of these
channels. A chirality-sensitive readout requires an independently
licensed odd-in-$S$ effect, preparation asymmetry, or holonomy and
its own physical occurrence/action proof.
\end{proposition}
\begin{proof}
The eigenvalues $\sqrt5$ and $-\sqrt5$ of $S$ both map to
$5-3=2$, while the zero eigenspace maps to $-3$.
The functional calculus of the two-eigenvalue symmetric operator
$T$ is scalar on $\mathscr H_{T=2}$. A function of $T$ thus
cannot separate its two rank-three summands; replacing $S$ by
$-S$ leaves $T$ unchanged. This is a mathematical comparison on
the finite carrier, not the assertion of a physically implemented
Hamiltonian or measurement.
\end{proof}
% END FRAMEWORK v2.340

% BEGIN FRAMEWORK v2.341: TWO PETERSEN-COVER RESPONSE FIREWALL
\begin{proposition}[Distinguishing cover data from a physical parity test]
\label{p4a:prop:v341-doublecover-parity-response}
The two finite lifted adjacency operators of
Theorem~\ref{p1:thm:v341-desargues-rp2-parity} are
$L(S_H)\simeq A\oplus S_H$ and
$L(-A)\simeq A\oplus(-A)$.  Their different spectra
certify that the signed dodecahedral Hopf connection and the
bipartite Desargues cross-incidence do not define the same
operator.  An odd five-label permutation exchanges the two
Hodge signings and the two six-axis face systems; it does not
supply a state-independent measurement operator on a selected
physical Hilbert space.  In particular, an effect depending only
on the parity-even $T=S_H^2-3I$ cannot distinguish the two
three-dimensional Hodge signs, even if the Desargues graph is
also present mathematically.  A test requires a GP2-authorized
odd-in-$S_H$ effect, an actual preparation, and same-owner
normalization.
\end{proposition}
\begin{proof}
The sheet-even/odd spectral block reduction is the finite graph
calculation in the cited theorem.  The $T$-blindness is
Proposition~\ref{p4a:prop:v340-reye-parity-no-go};
the observable requirements are constitutional, not spectral
consequences.
\end{proof}
% END FRAMEWORK v2.341



% BEGIN FRAMEWORK v2.347: EXACT OPERATIONAL RESOLUTION THEOREMS
\section{Exact operational distinction between the three finite operator envelopes}
\label{p4a:sec:v347-operational-resolution}
This section strengthens the already established generated-algebra,
coordinate-envelope, and \(\mathsf{FENV}/\mathsf{ENV}\) results.  It
introduces no new full-carrier test postulate.  The source is the six
multiplicity-free momentum-sector projections \(P_i\) and the central
integer/half-integer-spin projections \(P_+,P_-\) of the preceding chapters.

\begin{theorem}[Nested conditional expectations and exact test distinguishability]
\label{p4a:thm:v347-tier-distance}
On the constituted level-one carrier \(\mathcal H_1\cong\mathbb C^{15}\),
let \(\mathcal A_{43}=\mathcal A_{\rm kin}\),
\(\mathcal A_{137}=\mathcal A_{\rm obs}\), and
\(\mathcal A_{225}=\operatorname{End}(\mathcal H_1)\), as mathematical
matrix algebras, with the previously proved inclusions.  Define
\[
 \mathsf E_{43}(X)=\sum_{i=1}^6P_iXP_i,\qquad
 \mathsf E_{137}(X)=P_+XP_++P_-XP_-,\qquad
 \mathsf E_{225}(X)=X.
\]
These are unital, trace-preserving completely positive conditional
expectations, orthogonal in the Hilbert--Schmidt inner product, and
\[
 \mathsf E_{43}\mathsf E_{137}
 =\mathsf E_{137}\mathsf E_{43}=\mathsf E_{43}.
\]
The real space of self-adjoint carrier operators splits orthogonally as
\[
 \underbrace{\mathcal A_{43}^{\rm sa}}_{43}
 \oplus\underbrace{(\mathcal A_{137}^{\rm sa}
      \ominus\mathcal A_{43}^{\rm sa})}_{94}
 \oplus\underbrace{(\mathcal A_{225}^{\rm sa}
      \ominus\mathcal A_{137}^{\rm sa})}_{88},
 \qquad 43+94+88=225.
\]
For any carrier density matrices \(\rho,\sigma\) and each
\(a\in\{43,137,225\}\), the largest one-shot probability difference
obtainable from the \emph{mathematically available effects} of the
corresponding algebra equals
\begin{equation}\label{p4a:eq:v347-tier-distance}
 \boxed{
 D_a(\rho,\sigma)
 :=\sup_{\substack{E\in\mathcal A_a\\0\le E\le I}}
 \big|\operatorname{Tr}((\rho-\sigma)E)\big|
 =\frac12\|\mathsf E_a(\rho-\sigma)\|_1.}
\end{equation}
In particular two states are indistinguishable by all \(\mathcal A_a\)
effects if and only if their conditional expectations \(\mathsf E_a\rho\)
and \(\mathsf E_a\sigma\) coincide.
\end{theorem}
\begin{proof}
The \(P_i\) are mutually orthogonal, sum to the identity and refine
\(P_++P_-=I\).  Each pinching is a direct sum of matrix compressions,
which proves unitality, complete positivity, trace preservation and the
conditional-expectation property.  Orthogonality follows from trace
cyclicity.  The real self-adjoint dimension of a complex matrix algebra
\(\bigoplus_jM_{d_j}(\mathbb C)\) is \(\sum_jd_j^2\), giving the stated
\(43,137,225\) and their differences.
Put \(\Delta=\rho-\sigma\), so \(\operatorname{Tr}\Delta=0\).  For any
\(E\in\mathcal A_a\), self-adjointness of the Hilbert--Schmidt
conditional expectation gives
\(\operatorname{Tr}(\Delta E)=\operatorname{Tr}((\mathsf E_a\Delta)E)\).
The Hermitian element \(\mathsf E_a\Delta\) is traceless.  Its positive
spectral projector lies in \(\mathcal A_a\) and maximizes the trace
pairing over \(0\le E\le I\).  The maximum is half the trace norm, as
claimed.  The vanishing criterion follows by nondegeneracy of the
restricted trace pairing.
\end{proof}

\begin{corollary}[Exact 43/137 and 137/225 coherence witnesses]
\label{p4a:cor:v347-tier-witness}
Let \(u,v\) be orthogonal unit vectors in two distinct irreducible
sectors, and set
\(\psi_\pm=(u\pm v)/\sqrt2\),
\(\rho_\pm=|\psi_\pm\rangle\langle\psi_\pm|\).
If the sectors have the same central parity, then
\[
 (D_{43},D_{137},D_{225})(\rho_+,\rho_-)=(0,1,1).
\]
If they have opposite central parity, then
\[
 (D_{43},D_{137},D_{225})(\rho_+,\rho_-)=(0,0,1).
\]
\end{corollary}
\begin{proof}
The state difference is
\(\rho_+-\rho_-=|u\rangle\langle v|+|v\rangle\langle u|\),
with nonzero eigenvalues \(+1,-1\).  Distinct \(P_i\) blocks eliminate
this term under \(\mathsf E_{43}\).  The central parity pinching retains
it exactly in the same-parity case and removes it in the opposite-parity
case.  The full envelope retains it in both cases.  Apply
Theorem~\ref{p4a:thm:v347-tier-distance}.
\end{proof}
\begin{remark}[Event, field, and tomography licensing]
\label{p4a:rem:v347-tier-license}
The equalities above quantify distinguishability \emph{if all effects of
the specified algebra can be employed}.  The previously proved core and
coordinate-generated even test domains remain their respective occurrence
owners.  On a field-typed \(\mathsf{FENV}\) extension the odd matrix space
is not automatically a physical event space; a carrier-resolving
\(\mathsf{ENV}\) or actual test-coupling certificate is required before
\(D_{225}\) becomes operational.  The \(94\) and \(88\) are operator-space
dimensions, not extra species or automatically realized detector channels.
The homogeneous-orbit tomography channel counts of Paper~II-C are
\emph{continuous spans}; they must not be read as a protocol with only six
to eight sample outcomes or as exact statistical reconstruction from
finitely many noisy repetitions.
\end{remark}
% END FRAMEWORK v2.347



% BEGIN FRAMEWORK v2.352: COMPLETE TWO-RATE CENTRAL-POPULATION SEMIGROUP
\clearpage
\section{Two-rate covariant transfer with an even observable domain}
\label{p4a:sec:v352-two-rate-transfer}
The prior symmetric-rate GKLS example
\ref{p4:prop:v264-covariant-gkls-alternative} is recovered below
as a special case.  The general process is \emph{mathematically
available}, not selected by G0/GP1--GP3 and not a newly occurring
reservoir or action.

\begin{theorem}[Complete block-isotropic two-rate transfer construction]
\label{p4a:thm:v352-two-rate-transfer}
Let $G=\mathbb C^{11}$, $S=\mathbb C^4$ and $r,s\geq0$.
On block-diagonal density matrices define
\[
 \mathcal L_{r,s}(\rho_G\oplus\rho_S)=
 \left(-r\rho_G+\frac{s}{11}\operatorname{Tr}(\rho_S)I_G\right)
 \oplus
 \left(-s\rho_S+\frac{r}{4}\operatorname{Tr}(\rho_G)I_S\right).
\]
Then $e^{t\mathcal L_{r,s}}$ is a CPTP semigroup preserving the even
algebra $M_{11}\oplus M_4$, and is covariant under
$U(11)\times U(4)$, hence under the represented $\Spin(4)$.
Its central populations obey
\[
 \frac{d}{dt}\binom{p_G}{p_S}
 =\begin{pmatrix}-r&s\\r&-s\end{pmatrix}
   \binom{p_G}{p_S}.
\]
If $r,s>0$, its unique stationary state and population trajectory
are
\[
 \boxed{\rho_*=
 \frac{s}{r+s}\frac{I_G}{11}
 \oplus\frac{r}{r+s}\frac{I_S}{4}},\qquad
 p_G(t)=\frac{s}{r+s}
  +\Bigl(p_G(0)-\frac{s}{r+s}\Bigr)e^{-(r+s)t}.
\]
On the full $M_{15}$ \emph{mathematical dilation}, the semigroup has
jump operators
\[
 V_{ab}=\sqrt{r/4}\,|s_a\rangle\langle g_b|,
 \qquad W_{ba}=\sqrt{s/11}\,|g_b\rangle\langle s_a|,
 \quad 1\leq a\leq4,\ 1\leq b\leq11.
\]
These Kraus/Lindblad witnesses do not establish observable odd
operators or a physical dilation.
\end{theorem}
\begin{proof}
In GKLS form the displayed jumps give gain terms
$(r/4)\operatorname{Tr}(\rho_G)I_S$ and
$(s/11)\operatorname{Tr}(\rho_S)I_G$, and loss operator
$\sum V^*V+\sum W^*W=rP_G+sP_S$.
The full GKLS semigroup is CPTP and restricts to the stated
even-algebra generator.  Each gain term depends only on the
block trace and identity, proving covariance.  The population
matrix has eigenvalues $0$ and $-(r+s)$, and the traceless
parts of $\rho_G,\rho_S$ decay at rates $r,s$ respectively.
For positive rates the fixed state above is unique.
\end{proof}

\begin{theorem}[Exact 137-mode and 225-mode relaxation spectra]
\label{p4a:thm:v352-transfer-spectrum}
For $r,s>0$, the complex linear relaxation spectrum on the even
$137$-dimensional operator algebra has invariant mode dimensions
\[
 \begin{array}{c|rrrr}
\text{rate}&0&r+s&r&s\\\hline
\text{multiplicity}&1&1&120&15
\end{array}
\qquad(1+1+120+15=137).
\]
On the specified full GKLS dilation, the additional central-odd
$\operatorname{Hom}(G,S)\oplus\operatorname{Hom}(S,G)$ directions
have rate $(r+s)/2$ and complex multiplicity $88$, completing
$137+88=225$.  Equal numerical rates combine the coincident
eigenvalue multiplicities, without altering invariant subspaces.
\end{theorem}
\begin{proof}
Write $\rho_G=p_G I_G/11+X_G$ and
$\rho_S=p_S I_S/4+X_S$ with both $X$ traceless.
The generator sends $X_G\mapsto-rX_G$, $X_S\mapsto-sX_S$;
these traceless spaces have dimensions $11^2-1=120$ and
$4^2-1=15$.  The two population directions have eigenvalues
$0,-(r+s)$.  The GKLS gain terms vanish on off-block matrices;
the anticommutator multiplies each by $-(r+s)/2$.
There are $2\cdot11\cdot4=88$ such complex matrix units.
\end{proof}

\begin{corollary}[Exact bistochasticity criterion and relation to the old model]
\label{p4a:cor:v352-transfer-bistochastic}
The above semigroup is Schr\"odinger-unital, hence bistochastic,
if and only if
\[
 \boxed{11r=4s}.
\]
For positive rates this is equivalent to the stationary populations
$(p_G^*,p_S^*)=(11/15,4/15)$.  Setting
$r=4\gamma$ and $s=11\gamma$ recovers the earlier symmetric-jump
example \ref{p4:prop:v264-covariant-gkls-alternative}, including
its rates $4\gamma,11\gamma,15\gamma,15\gamma/2$.
\end{corollary}
\begin{proof}
On the full matrix identity,
$\mathcal L_{r,s}(I_{15})=
 (4s/11-r)P_G+(11r/4-s)P_S$.
Its vanishing is equivalent to $11r=4s$.
For positive rates the unique stationary populations are
$s/(r+s),r/(r+s)$; solve the geometric ratio.
\end{proof}

\begin{remark}[No unlicensed dynamical or field promotion]
\label{p4a:rem:v352-transfer-physical-frontier}
The dimensions $11,4$ and positive rates do not derive
$r/s$, bistochasticity, $r+s$, a bath, an action-owned jump or a
measured decay.  The construction proves that an \emph{even}
observable role can support cross-sector \emph{process} transfer
without licensing a central-odd \emph{effect}.  Establishing a
physical nonzero transition and the corresponding detector remains
an independent GP2/G0 process/readout requirement.  The numeral
$120=11^2-1$ is not evidence of an independently occurring
$2I$ or 600-cell role.
\end{remark}
% END FRAMEWORK v2.352


\subsection{Ancilla-stable operational quotient of native tests}
\label{p4a:sec:v355-native-history-quotient}
\begin{proposition}[Native finite histories cannot certify hidden carrier multiplicities]
\label{p4a:prop:v355-ancilla-stable-process-quotient}
For any finite constituted observable algebra
$\mathcal A=\bigoplus_\lambda M_{d_\lambda}(\mathbb C)$,
any finite faithful multiplicity-amplified representation, and every
finite entangled reference system, the full collection of adaptive
outcome-history probabilities of completely positive instruments
\emph{defined on $\mathcal A\otimes B(\mathcal K)$} is reproduced in
the canonical minimum faithful module. No such native experiment
distinguishes arbitrary hidden multiplicities $m_\lambda\ge1$.
\end{proposition}
\begin{proof}
Apply the sectorwise dephasing and multiplicity-factor partial trace
of Theorem~\ref{p2b:thm:v354-ancilla-stable-descent} to the initial
state with reference. Because every branch map is an operation on the
abstract joint algebra, its Heisenberg composition and every adaptive
branch-event are evaluated by the same reduced state. The all-history
result specializes Theorem~\ref{p2b:thm:v353-multitime-amplification},
which treats a common amplification factor.
\end{proof}
\begin{remark}[Physical process boundary]
\label{p4a:rem:v355-native-quotient-scope}
The statement is a capacity theorem for the operator algebra and the
specified channel domain. It is not a result about arbitrary CP maps on
the entire ambient matrix algebra of a multiplicity-amplified carrier.
A physical test of multiplicities requires a \emph{separately licensed}
commutant-sensitive operation or GP3 whole-candidate comparison. No
physical detection event, source, or central-odd $11/4$ bridge is added.
\end{remark}



% BEGIN SOURCE-INTEGRATED CROSS-PAPER MATHEMATICS v3_017, 15 November 2026.
%  Derived geometric comparison modules; physical source occurrence not promoted.
%  E3/E4 independent audit of these NEW results remains pending; inherited source statements retained.
\section{Additional exact mathematical consequences from the verified source geometry}

% -------- Research module Paper4A-IcosahedralEvaluationPOVM-v1_1 (latest supplied version) --------
% Proposed research insertion for Paper IV-A (auxiliary finite measurement).
\subsection{An exact icosahedral evaluation POVM and its phase blind spot}
\label{p4a:sec:cc13-angular-povm}

\begin{proposition}[Positive $62$-effect normalized band evaluation]
\label{p4a:prop:cc13-angular-povm}
On the complex $36$-dimensional comparison Hilbert space
$\mathscr Y_5\subset L^2(S^2,d\Omega/(4\pi))$, retain the
$62$ positive weights and degree-five reproducing kernel of
Theorem~\ref{p1b:thm:cc13-band-frame}.  Let $k_i\in\mathscr Y_5$
represent evaluation $\langle k_i,f\rangle=f(u_i)$, and define
\[
 \boxed{E_i=w_i|k_i\rangle\langle k_i|\succeq0,
 \qquad\sum_{i=1}^{62}E_i=I_{36},\qquad
 \operatorname{Tr}E_i=36w_i.}
\]
This is an $A_5$-covariant auxiliary rank-one POVM.  For a normalized
pure state $f$, the probabilities are
$p_i=w_i|f(u_i)|^2$.  Although complex \emph{amplitude} samples
reconstruct $f\in\mathscr Y_5$ exactly, these intensity probabilities
do not determine an arbitrary mixed state, nor even every pure-state ray.
\end{proposition}
\begin{proof}
The sampling-isometry identity $S^*S=I$ from Paper I-B is
$\sum_iw_i|k_i\rangle\langle k_i|=I$ in Riesz notation.
The addition theorem gives
$\|k_i\|^2=K_5(u_i,u_i)=\sum_{\ell\le5}(2\ell+1)=36$.
The three node orbits and their weights are preserved by $A_5$,
proving covariance.  There are only $62$ effects, so their linear span
cannot be the real $36^2=1296$-dimensional operator space of all
Hermitian matrices; mixed-state informational completeness is impossible.
More strongly, complex conjugation preserves $\mathscr Y_5$ and
$|\overline{f(u)}|^2=|f(u)|^2$ for every $u$.  Choose two real,
orthonormal spherical harmonics $f_1,f_2$ and
$f=(f_1+i f_2)/\sqrt2$; then $f$ and $\bar f$ are different rays
(they are orthogonal), yet have identical probabilities at all
$62$ effects and indeed at every point of $S^2$.  Thus merely
increasing the number of unoriented position-intensity probes does
not repair the phase blind spot.
\end{proof}

\begin{remark}[Readout boundary]
The $62$-outcome frame POVM is a mathematically licensed positive
measurement on the named auxiliary band.  Physical apparatus ownership,
phase-sensitive readout, or any implication for the actual cosmological
sky requires the separately constituted source/detector interface.
\end{remark}
% END SOURCE-INTEGRATED CROSS-PAPER MATHEMATICS v3_017.

% BEGIN FRAMEWORK v3.067: CONDITIONAL SCALAR COORDINATE CONTROL
\section{Mathematical scalar-coordinate control and typed readout}
\label{p4a:sec:v3067-coordinate-control}
% RESEARCH INSERTION CANDIDATE ONLY — NOT YET MERGED INTO PAPER IV-A.
% Based on Paper1A-Geometry-v5_180, Paper2C-ReciprocalLocalization-v2_173,
% Paper4A-FiniteQuantumMechanics-v4_216, and supplementary exact proof
% Paper00-SharpFrontierOccurrence-v1_4.md.
% DO NOT PROMOTE TO OCCURRING DYN, TEST, PROC, ACC OR FRONTIER WITHOUT SOURCE PROOF.

\begin{theorem}[Exact compact control Lie algebra of finite scalar coordinates]
\label{p4a:thm:finite-coordinate-lie-so}
Let $F_K=\bigoplus_{k=0}^{K}E_k^x$ be the complete real scalar-harmonic
band of the round reciprocal-position three-sphere, $K\ge1$, and put
$D_K=(K+1)(K+2)(2K+3)/6$. On its complexification let
$T_a=P_KM_{u_a}P_K$, $a=1,\dots,4$, where $u_a=x_a/a$
are normalized round coordinates. Choose a real harmonic basis and
$W|_{E_k^x}=1$ for even $k$ and $W|_{E_k^x}=i$ for odd $k$.
Then $B_a=W(iT_a)W^*$ are real skew-symmetric and
\[
 \boxed{\operatorname{Lie}_{\mathbb R}(B_1,B_2,B_3,B_4)
          =\mathfrak{so}(D_K).}
\]
This is a mathematical kinematic-control assertion. Actual coherent
implementation of $e^{-itT_a}$ requires a same-host, action-owned
closed-dynamics $\mathsf{DYN}$ certificate and does not follow from the
algebraic inclusion $T_a\in C^*(T_1,\ldots,T_4)$.
\end{theorem}
\begin{proof}
Coordinate multiplication changes spherical harmonic degree by $\pm1$;
$W$ therefore makes $iT_a$ real skew-symmetric. For $K=1$, use the
basis $(1,2u_1,\ldots,2u_4)$. Each $T_a$ has only the entries
$(T_a)_{0a}=(T_a)_{a0}=1/2$, so their skew phase transforms give all
$0\wedge a$ and their brackets all $a\wedge b$, a basis of
$\mathfrak{so}(5)$.

Proceed by induction. Let $V=F_{K-1}$ and $U=E_K^x$.
The exact coordinate commutator identity on $F_K$ is
\[
 [B_a,B_b]=\mp\frac{1}{2(K+1)}
  (0_V\oplus L_{ab}|_U),\qquad
 L_{ab}=u_a\partial_b-u_b\partial_a,
\]
where the overall sign depends only on the chosen parity-phase convention.
These operators generate the compact $\operatorname{Spin}(4)$ representation
on $U$ and are zero on $V$. Since $U$ is a nontrivial irreducible
$\operatorname{Spin}(4)$ module, it has no invariant vector.
Haar averaging $\operatorname{Ad}_g(B_a)$ over this already generated
subgroup therefore removes the $U$--$V$ block and retains the lower block
$B_a^{(K-1)}\oplus0$. The Lie algebra, being a finite-dimensional real
vector space invariant under the subgroup's adjoint action, contains this
average. Induction gives the embedded $\mathfrak{so}(V)$.

Subtract the lower block from each $B_a$ to obtain pure cross elements
$C_a\in U\otimes V$. If $R_a:E_{K-1}^x\to E_K^x$ is the harmonic-raising
block, the exact boundary-defect identity gives
\[
 \sum_{a=1}^{4}R_aR_a^*=\frac{K}{2(K+1)}I_U>0.
\]
Thus the images of the $C_a$ jointly span $U$. For $K\ge2$ the
standard defining $\mathfrak{so}(V)$ module $V$ has dimension at least
five and is absolutely irreducible. Consequently the
$\mathfrak{so}(V)$ submodule of $U\otimes V$ generated by these $C_a$
is all of $U\otimes V$. Commutators of its cross wedges give
$\mathfrak{so}(U)$. Therefore all three blocks
$\mathfrak{so}(V)$, $U\otimes V$, and $\mathfrak{so}(U)$ occur, yielding
$\mathfrak{so}(V\oplus U)$.\end{proof}

\begin{corollary}[One oriented effect and coordinate-controlled tomography]
\label{p4a:cor:finite-coordinate-controlled-tomography}
For each $K\ge1$, let $E_K(du)=|\kappa_{K,u}\rangle
\langle\kappa_{K,u}|\,d\mu(u)$ be the canonical positive
band-evaluation POVM. For any $u\in S^3$, the real span of the conjugates
of its pointwise density by the connected coordinate-control group
is $\operatorname{Herm}(\mathbb C^{D_K})$. A single sufficiently small,
non-antipodally-identified finite-aperture effect therefore has a full
conjugation span as well. Finitely many such conjugated effects are
informationally complete on the fixed finite host.
\end{corollary}
\begin{proof}
In the parity-phase basis the evaluation kernel is $a_u+i b_u$,
with nonzero orthogonal real even/odd parts for every $u$ (from
$E_0^x$ and $E_1^x$). Its rank-one projector has nonzero real symmetric
and imaginary antisymmetric components. Rotating orthogonal real
2-frames by $SO(D_K)$ and comparing determinant-one sign-flip pairs
spans all of $\operatorname{Sym}^2\mathbb R^{D_K}$ and
$i\Lambda^2\mathbb R^{D_K}$, hence all Hermitian matrices.
Finite aperture follows by continuity of a finite nonzero determinant
witness under averaging over a shrinking neighborhood of $u$.
\end{proof}

\begin{corollary}[Exact four-setting primitive readout and antipodal no-go]
\label{p4a:cor:primitive-four-setting-coordinate-tomography}
At $K=1$, the canonical evaluation POVM preceded respectively by
$U_0=I$ and $U_j=\exp(-i\pi T_j)$, $j=1,2,3$, has cumulative real
effect-span ranks $14,21,24,25$. Four coherent coordinate-preprocessed
settings are necessary and sufficient for full five-dimensional
state tomography when no other preprocessing Hamiltonians or detectors
are allowed. Conversely, if every evaluation effect is identified
with its antipodal mate $u\sim-u$, then even arbitrary coordinate-control
preprocessing leaves the entire $i\Lambda^2\mathbb R^{D_K}$ invisible.
\end{corollary}
\begin{proof}
At $K=1$, $W\exp(-i\pi T_j)W^*$ rotates the distinguished even unit
vector to the $j$th odd unit vector. With that direction denoted $n$,
the real effect span is
$\mathbb RI\oplus\operatorname{Sym}_0(n^\perp)
\oplus i(n\wedge n^\perp)$.
For $m$ independent distinguished directions, its skew part has rank
$\binom52-\binom{5-m}{2}$, so three settings leave a one-dimensional
invisible wedge. The chosen four axes generate all ten skew directions;
three of them already generate all fifteen real symmetric directions.
This gives the asserted ranks and optimality.
In the same parity-phase basis, $u\mapsto-u$ replaces
$a_u+ib_u$ by $a_u-ib_u$. Their combined effect is real symmetric;
conjugation by the available real orthogonal controls cannot generate
any imaginary antisymmetric effect.
\end{proof}

% OCCURRENCE FIREWALL: The theorem is on the mathematical scalar comparison
% band F_K, not the native momentum carrier H_1.  It neither actualizes a
% position measurement apparatus nor lifts the independent SF_{Lambda_F}
% occurrence.  It shows that ONE oriented source-owned effect AND a coherent
% source-owned coordinate-control DYN would suffice on any already
% constituted complete finite scalar band.  A test algebra does not imply
% an actual time-controllable Hamiltonian.

% END FRAMEWORK v3.067: CONDITIONAL SCALAR COORDINATE CONTROL

\printbibliography
\end{document}
